\documentclass{unina_doc_class}
% ============================================================
% AGGIUNTE ALLA CLASSE: spiegazioni, diagrammi, marcatori
% ============================================================
%
% UTILIZZO:
%   \begin{spiegazione}[Titolo opzionale] ... \end{spiegazione}
%       riquadro verde che segue ogni dimostrazione: l'idea in parole semplici,
%       con un diagramma o un esempio passo passo (non si spezza tra due pagine)
%   \begin{passi} \item ... \item ... \end{passi}
%       elenco di passi numerati con un cerchietto (per gli esempi passo passo)
%   \senzadim
%       targhetta «senza dimostrazione»: il programma non chiede la dimostrazione
%   \facoltativa
%       targhetta per le dimostrazioni che il programma non chiede ma che aiutano a capire
%
%   Stili TikZ:
%     Hasse:      [hasse] sull'ambiente; nodi \node[hv] (etichette) o \node[hpunto,label=...] (pallini);
%                 linee \draw[hlato]; evidenziazioni hv evid / hv evid2
%     Relazioni:  nodi \node[rv]; frecce \draw[rfreccia]; cappi \draw[rcappio] (n) to (n)
%     Funzioni:   insiemi \node[insieme, fit=...]; elementi \node[elem]; frecce \draw[fun]
% ============================================================

\RequirePackage{tikz-cd}

% Indice: i numeri di pagina a tre cifre non entrano nella casella standard di 1,55em e la
% classe (\overfullrule) segnerebbe ogni riga con un quadratino nero; la casella si allarga
% (a una misura fissa, così tutti i numeri finiscono sullo stesso margine destro).
\makeatletter
\renewcommand{\@pnumwidth}{22pt}
\makeatother
% Nelle righe delle parti la casella del numero, composta in corpo 14, sporgerebbe oltre il
% bordo del foglio: il numero si allinea al margine destro senza casella.
\titlecontents{part}
    [0pt]
    {\addvspace{1cm}\centering}
    {\color{primary!60!black}\fontsize{16}{0}\selectfont\Gobold Parte \thecontentslabel\quad}
    {\color{primary!60!black}\fontsize{16}{0}\selectfont\Gobold}
    {\hfill\color{primary!80!black}\bfseries\Scientifica\fontsize{14}{0}\selectfont\thecontentspage}
    [\addvspace{0.3cm}\hrule height 1.5pt\addvspace{0.5cm}]

% ── Colori per le classi di equivalenza e le parti di una partizione ──
\colorlet{classeA}{primary!16}
\colorlet{classeB}{orange!22}
\colorlet{classeC}{seagreen!22}
\colorlet{classeD}{pink!30}
\colorlet{classeE}{naplesyellow!45}
\colorlet{classeF}{violet!12}

% ── Riquadro SPIEGAZIONE ──────────────────────────────────────
\newtcolorbox{spiegazionebox}[1]{
  enhanced, breakable=false,
  colback=seagreen!5!white, colframe=seagreen!65!black,
  boxrule=0pt, leftrule=3pt, arc=0pt, outer arc=0pt,
  left=3mm, right=2mm, top=1.2mm, bottom=1.4mm,
  before skip=0mm, after skip=6mm,
  fontupper=\normalsize,
}
\NewDocumentEnvironment{spiegazione}{O{Spiegazione}}{%
  \begin{concetto}\vspace{-3mm}%
  \begin{spiegazionebox}{#1}%
  {\color{seagreen!55!black}{\mdseries\iidea}\ \bfseries #1}\par\nopagebreak\vspace{0.6mm}%
  \setlength{\parskip}{0.6em}%
}{%
  \end{spiegazionebox}\end{concetto}%
}

% ── Passi numerati ────────────────────────────────────────────
\NewDocumentCommand{\numcerchio}{m}{\tikz[baseline=(n.base)]\node[circle, fill=seagreen!65!black, text=white,
  font=\scriptsize\bfseries, inner sep=0pt, minimum size=4mm] (n) {#1};}
\newenvironment{passi}{%
  \begin{enumerate}[label=\numcerchio{\arabic*}, leftmargin=7.5mm, labelsep=2mm,
                    itemsep=0.8mm, topsep=0.6mm, parsep=0.3em]%
}{\end{enumerate}}

% ── Targhette ─────────────────────────────────────────────────
\newcommand{\targhetta}[2]{%
  \tikz[baseline=(t.base)]\node[rounded corners=2pt, fill=#1!12, draw=#1!60!black, line width=0.4pt,
    font=\scriptsize\bfseries, text=#1!60!black, inner xsep=3pt, inner ysep=1.5pt] (t) {#2};}
\newcommand{\senzadim}{\targhetta{neutral}{senza dimostrazione}\par\smallskip}
\newcommand{\facoltativa}{\targhetta{neutral}{dimostrazione non richiesta dal programma}\par\smallskip}

% ── Blocchi colorati (classi di equivalenza, parti di una partizione) ──
% \blocco[larghezza]{colore}{percorso}: «macchia» colorata attorno a un gruppo di elementi.
% Il percorso unisce i centri dei nodi:
%   due elementi:  \blocco{classeA}{(n1.center) -- (n2.center)}
%   tre elementi:  \blocco{classeB}{(n1.center) -- (n2.center) -- (n3.center) -- cycle}
%   un elemento:   \blocco{classeC}{(n1.center) -- ++(0.001,0)}
\newcommand{\blocco}[3][6.4mm]{%
  \begin{scope}[on background layer]
    \draw[line width=#1+0.8pt, line cap=round, line join=round, draw=#2!55!black!45] #3;
    \draw[line width=#1, line cap=round, line join=round, draw=#2, fill=#2] #3;
  \end{scope}}

% ── Stili TikZ ────────────────────────────────────────────────
\tikzset{
  % diagrammi di Hasse
  hasse/.style={every node/.append style={font=\small}},
  hv/.style={draw=primary!70!black, fill=primary!7, rounded corners=2pt, inner sep=2pt,
             minimum height=4.6mm, minimum width=4.6mm},
  hv evid/.style={hv, fill=orange!25, draw=orange!85!black},
  hv evid2/.style={hv, fill=seagreen!25, draw=seagreen!75!black},
  hpunto/.style={circle, fill=primary!80!black, inner sep=0pt, minimum size=5pt},
  hlato/.style={draw=black!60, line width=0.75pt},
  hlato evid/.style={draw=orange!85!black, line width=1.4pt},
  % grafi di relazioni binarie
  rv/.style={circle, draw=primary!75!black, fill=primary!7, inner sep=1pt, minimum size=6mm, font=\small},
  rfreccia/.style={-{Stealth[length=2.3mm]}, line width=0.7pt, draw=black!70},
  % cappi: \draw[rfreccia] (n) to[rcappio] (n);
  rcappio/.style={out=120, in=60, min distance=8mm},
  rcappio sotto/.style={out=-120, in=-60, min distance=8mm},
  rcappio sx/.style={out=150, in=210, min distance=8mm},
  rcappio dx/.style={out=30, in=-30, min distance=8mm},
  % diagrammi a frecce delle funzioni
  insieme/.style={draw=black!60, rounded corners=9pt, line width=0.6pt, inner sep=6pt},
  elem/.style={font=\small, inner sep=1.5pt},
  fun/.style={-{Stealth[length=2.2mm]}, line width=0.7pt, draw=primary!80!black},
  fun2/.style={-{Stealth[length=2.2mm]}, line width=0.7pt, draw=cherryred!85!black},
  fun3/.style={-{Stealth[length=2.2mm]}, line width=0.7pt, draw=seagreen!70!black},
}


\gdef\documentname{Algebra per Informatica}
\gdef\documentyear{A.A. \ 2025-2026}
\gdef\author{Fabio Rosiello}
\pagestyle{plain}


\gdef\xauthors{\hi{\iuser Fabio Rosiello}}

\setcounter{secnumdepth}{0}

\widowpenalty10000
\clubpenalty10000
\raggedbottom

\begin{document}
  \coverpage{26}{\documentyear}{30}{\documentname}
  \tableofcontents

  \chapter{Prefazione}

\introduction{Mattia Brescia}{
    \textit{Il corso fornisce il linguaggio e gli strumenti di base della matematica per l'informatica:
    logica, teoria degli insiemi, strutture algebriche, relazioni, calcolo combinatorio, aritmetica, insiemi ordinati,
    reticoli e algebre di Boole.}
  }
\pagestyle{plain}

\section{Il programma}

La dispensa segue il programma del corso da \textbf{6 CFU} (canale 3), cioè le prime sei parti del programma ufficiale da 9 CFU:
logica e insiemi, strutture algebriche, relazioni binarie, calcolo combinatorio, insiemi ordinati (con reticoli e algebre di Boole)
e divisibilità. \textbf{Polinomi e grafi non fanno parte del corso} e non sono trattati.
Gli argomenti sono divisi in sette parti: Logica, Teoria degli insiemi, Strutture algebriche, Relazioni binarie, Numeri naturali
e calcolo combinatorio, Aritmetica, Insiemi ordinati, reticoli e algebre di Boole.

All'orale si portano le dimostrazioni di tutti i risultati del programma, tranne quelli indicati «senza dimostrazione». Nella
dispensa questi risultati hanno la targhetta \targhetta{neutral}{senza dimostrazione}; alcune dimostrazioni in più, utili per capire ma
non richieste, hanno la targhetta \targhetta{neutral}{dimostrazione non richiesta dal programma}.

\section{Come leggere questa dispensa}

\textbf{L'ordine degli argomenti.} Gli argomenti non seguono necessariamente l'ordine delle lezioni, ma un
\textbf{ordine topologico}: ogni definizione, teorema o proposizione compare \textbf{dopo} tutto ciò da cui dipende.
Leggendo dall'inizio alla fine non si incontra mai qualcosa che usa un concetto non ancora introdotto.

\textbf{La numerazione.} Definizioni, teoremi, proposizioni, lemmi e corollari condividono un'unica numerazione per
capitolo (ad esempio Definizione 1.3, Proposizione 1.4, Teorema 1.5), così è facile ritrovarli.

\textbf{Il box «Prerequisiti».} All'inizio di ogni teorema (o proposizione, lemma, corollario) c'è un elenco di
\textbf{cosa serve sapere} per capirlo al meglio: definizioni, altri risultati o sezioni. Ogni voce è
\textbf{cliccabile} e riporta la pagina in cui si trova; se una voce non ti è chiara, conviene ripassarla prima di
leggere la dimostrazione.

\begin{concetto}
\begin{Proposizione}{Esempio di enunciato}{esempio-legenda}
    \begin{prerequisiti}
        \dipvoce{Qui trovi l'elenco di cosa serve sapere, con numero, titolo e pagina di ogni voce.}
    \end{prerequisiti}
    Qui trovi l'enunciato.
\end{Proposizione}
\begin{dimostrazione}
    Qui trovi la dimostrazione, che termina con il simbolo a destra.
\end{dimostrazione}
\end{concetto}

\begin{spiegazione}
    Dopo ogni dimostrazione trovi una \textbf{spiegazione}: la stessa idea detta in modo informale, con un diagramma (di Venn, di Hasse,
    a frecce, \dots) oppure con un esempio svolto \textbf{passo per passo}, come qui:
    \begin{passi}
        \item il primo passo;
        \item il secondo passo.
    \end{passi}
    Conviene leggerla subito dopo la dimostrazione, oppure prima, se la dimostrazione sembra difficile.
\end{spiegazione}

\textbf{I box.}
\begin{itemize}
    \item \textcolor{red!75!black}{\textbf{Rosso}}: definizioni.
    \item \textcolor{accent!80!black}{\textbf{Viola}}: assiomi.
    \item \textcolor{orange!85!black}{\textbf{Arancione}}: teoremi, seguiti dalla dimostrazione.
    \item \textcolor{primary!80!black}{\textbf{Blu}}: proposizioni e lemmi, seguiti dalla dimostrazione.
    \item \textcolor{seagreen!65!black}{\textbf{Verde scuro}}: corollari.
    \item \textcolor{seagreen!65!black}{\textbf{Verde con la lampadina}}: spiegazioni, dopo ogni dimostrazione.
    \item \textcolor{green-gradient-6!70!black}{\textbf{Verde chiaro}}: note e osservazioni.
    \item \textcolor{yellow-gradient-6!70!black}{\textbf{Giallo}}: attenzione, errori comuni.
    \item Gli esempi sono racchiusi tra due linee orizzontali.
\end{itemize}

\Needspace{8\baselineskip}
\textbf{Le tavole di verità.} Nelle tavole i valori sono $\vero$ (vero) e $\falso$ (falso). Le colonne evidenziate
sono quelle da guardare: per una tautologia, la colonna deve contenere solo $\vero$; per un'equivalenza, le colonne
con lo \textbf{stesso colore} devono coincidere riga per riga.

\textbf{I diagrammi.} Nei diagrammi di Hasse gli elementi in \colorbox{orange!25}{arancione} sono i «dati» di un esempio e quelli
in \colorbox{seagreen!25}{verde} i risultati; nei diagrammi a frecce le frecce blu sono la funzione principale e quelle rosse
tratteggiate una sua inversa (sezione, retrazione).


  % Parte 1: Logica (logica proposizionale e logica dei predicati)

  \part{Logica}

  % Logica proposizionale: alfabeto, sintassi, semantica, tautologie, connettivi
  \chapter{Logica proposizionale}

Prima di parlare di insiemi, funzioni e strutture algebriche ci serve un linguaggio preciso con cui scrivere
(e verificare) affermazioni matematiche: la \textbf{logica}. Cominciamo dalla \textbf{logica proposizionale},
in cui le affermazioni più semplici sono indicate da lettere ($p$, $q$, $r$, \dots) e vengono combinate
tramite i connettivi (\textit{non}, \textit{e}, \textit{oppure}, \dots).

Il percorso del capitolo è sempre lo stesso che seguiremo anche per la logica dei predicati:
\begin{enumerate}
    \item \textbf{alfabeto}: quali simboli possiamo usare;
    \item \textbf{sintassi}: quali sequenze di simboli sono «scritte correttamente» (le formule ben formate);
    \item \textbf{semantica}: quando una formula ben formata è vera o falsa;
    \item \textbf{tautologie}: le formule sempre vere, che useremo come «regole di calcolo» per tutto il corso.
\end{enumerate}

% ──────────────────────────────────────────────────────────────
\section{L'alfabeto}\label{sec:alfabeto-prop}

\dfn[alfabeto-prop]{Alfabeto della logica proposizionale}{
    L'\textbf{alfabeto} della logica proposizionale è la terna
    \[ \mcA = (P, C, B) \]
    dove:
    \begin{itemize}
        \item $P = \{p, q, r, \dots\}$ è l'insieme delle \textbf{variabili proposizionali};
        \item $C = \{\neg, \wedge, \vee, \rightarrow, \leftrightarrow\}$ è l'insieme dei \textbf{connettivi};
        \item $B = \{(,\ )\}$ è l'insieme delle \textbf{parentesi}.
    \end{itemize}
}

I connettivi hanno vari nomi; i più comuni sono i seguenti.

\begin{table}[H]
\centering
\renewcommand{\arraystretch}{1.4}
\begin{tabular}{c|l|l|l}
\intestazione \textbf{Simbolo} & \textbf{Nome} & \textbf{Si legge} & \textbf{Tipo} \\ \hline
$\neg$ & NOT & «non» & unario \\
$\wedge$ & AND & «e» & binario \\
$\vee$ & OR & «oppure» & binario \\
$\rightarrow$ & IMPLIES & «se \dots\ allora \dots», «implica» & binario \\
$\leftrightarrow$ & IFF & «se e solo se» & binario \\
\end{tabular}
\caption{I connettivi della logica proposizionale.}
\end{table}

NOT si dice \textbf{connettivo unario} perché agisce su una sola formula, mentre gli altri si dicono
\textbf{connettivi binari} perché agiscono su due formule.

\dfn[formula]{Formula}{
    Una qualunque sequenza (finita) di elementi dell'alfabeto $\mcA$ si dice \textbf{formula}.
}

Non tutte le formule hanno senso: ad esempio $)\,p \wedge \vee\, ($ è una formula, ma è chiaramente
«scritta male». Quali sono allora le formule ammesse?

% ──────────────────────────────────────────────────────────────
\section{Sintassi: le formule ben formate}\label{sec:sintassi-prop}

La \textbf{sintassi} è l'insieme di regole che permettono di capire quali sequenze di simboli sono
ammesse e quali no. Le formule ammesse si dicono \textbf{formule ben formate} (in breve \textbf{FBF}).

\dfn[fbf-prop]{Formule ben formate (FBF)}{
    Le \textbf{formule ben formate} della logica proposizionale sono le formule costruite con le seguenti regole:
    \begin{enumerate}
        \item ogni variabile proposizionale è una FBF;
        \item se $f$ è una FBF, allora $\neg f$ è una FBF;
        \item se $f$ e $g$ sono FBF e $\sim$ è uno dei connettivi $\wedge$, $\vee$, $\rightarrow$, $\leftrightarrow$,
              allora $f \sim g$ è una FBF;
        \item se $f$ è una FBF, allora anche $(f)$ è una FBF.
    \end{enumerate}
    Indichiamo con $\mcF$ l'\textbf{insieme di tutte le FBF}.
}

In altre parole, le FBF sono le stringhe \textbf{sintatticamente corrette}, cioè quelle che si ottengono
applicando (un numero finito di volte) le regole grammaticali qui sopra.

\begin{esempio}{Costruire una FBF passo per passo}
    Verifichiamo che $(p \wedge q) \rightarrow \neg r$ è una FBF:
    \begin{itemize}
        \item $p$, $q$, $r$ sono FBF (regola 1);
        \item $p \wedge q$ è una FBF (regola 3) e quindi anche $(p \wedge q)$ lo è (regola 4);
        \item $\neg r$ è una FBF (regola 2);
        \item $(p \wedge q) \rightarrow \neg r$ è una FBF (regola 3, con $f = (p \wedge q)$ e $g = \neg r$).
    \end{itemize}
    Invece $p \wedge \vee q$ \textbf{non} è una FBF: nessuna regola permette di scrivere due connettivi binari uno dopo l'altro.
\end{esempio}

\info{Parentesi e ordine di precedenza}{
    Le parentesi non sono realmente necessarie, ma servono per gestire l'ordine di precedenza dei connettivi.
    Ad esempio, $\neg(p \wedge q)$ sarà diverso da $(\neg p) \wedge q$.
    In teoria potremmo anche non utilizzare mai le parentesi, concordando il seguente ordine gerarchico:
    \begin{center}
        NOT, \ AND, \ OR, \ IMPLIES, \ IFF.
    \end{center}
    Tuttavia useremo le parentesi ogni volta che potremo, utilizzando implicitamente la sola convenzione
    che \textbf{NOT ha la precedenza su tutto} (quindi $\neg p \wedge q$ significa $(\neg p) \wedge q$).
}

% ──────────────────────────────────────────────────────────────
\section{Semantica: la funzione di valutazione}\label{sec:semantica-prop}

Passiamo ora alla \textbf{semantica}, ovvero al «senso» delle FBF: vogliamo stabilire quando una FBF è vera
e quando è falsa. Per farlo assegniamo a ogni FBF un valore di verità, $\vero$ (vero) oppure $\falso$ (falso),
in modo coerente con il significato dei connettivi.

\dfn[valutazione]{Funzione di valutazione}{
    Sia $\mcF$ l'insieme delle FBF. Una funzione $v \colon \mcF \to \{\vero, \falso\}$ si dice
    \textbf{(funzione di) valutazione} se per ogni coppia di FBF $f$, $g$ valgono le seguenti regole:
    \begin{enumerate}
        \item se $v(f) = \vero$, allora $v(\neg f) = \falso$;
        \item se $v(f) = \falso$, allora $v(\neg f) = \vero$;
        \item $v(f \wedge g) = \vero$ \ se e solo se \ $v(f) = \vero$ \textbf{e} $v(g) = \vero$;
        \item $v(f \vee g) = \vero$ \ se e solo se \ $v(f) = \vero$ \textbf{oppure} $v(g) = \vero$;
        \item $v(f \rightarrow g) = \vero$ \ se e solo se \ $v(f) = \falso$ \textbf{oppure} $v(g) = \vero$;
        \item $v(f \leftrightarrow g) = \vero$ \ se e solo se \ $v(f) = v(g)$.
    \end{enumerate}
    Diremo che una FBF $f$ è \textbf{vera} (nella valutazione $v$) se $v(f) = \vero$, altrimenti diremo che è \textbf{falsa}.
}

\warning{Perché nelle regole 3--6 serve il «se e solo se»}{
    Se nella regola 3 ci fosse solo «se $v(f \wedge g) = \vero$ allora $v(f) = \vero$ e $v(g) = \vero$», nulla
    vieterebbe di avere $v(f \wedge g) = \falso$ anche quando $f$ e $g$ sono entrambe vere.
    Il «se e solo se» dice invece che il valore di $f \wedge g$ è \textbf{completamente determinato} dai valori
    di $f$ e di $g$ (e lo stesso per $\vee$, $\rightarrow$, $\leftrightarrow$).
}

\begin{esempio}{Calcolare una valutazione}
    Sia $v$ una valutazione con $v(p) = \vero$ e $v(q) = \falso$. Calcoliamo $v\big((p \wedge q) \rightarrow \neg q\big)$:
    \begin{itemize}
        \item $v(p \wedge q) = \falso$ per la regola 3, perché $v(q) = \falso$;
        \item $v(\neg q) = \vero$ per la regola 2;
        \item $v\big((p \wedge q) \rightarrow \neg q\big) = \vero$ per la regola 5, perché $v(p \wedge q) = \falso$.
    \end{itemize}
\end{esempio}

Come si vede dall'esempio, per calcolare $v$ su una FBF basta conoscere i valori che $v$ assegna alle
\textbf{variabili proposizionali} che vi compaiono: tutto il resto si ricava applicando le regole
(le parentesi non cambiano il valore: $v((f)) = v(f)$). Quindi, se in una FBF compaiono $n$ variabili,
ci sono esattamente $2^n$ casi da controllare. È proprio l'idea delle tavole di verità.

\subsection{Tavole di verità}\label{sec:tavole}

Le \textbf{tavole di verità} servono a controllare tutti i possibili valori di verità di una FBF al variare
dei valori di verità delle sue «componenti»: ogni riga corrisponde a una possibile scelta dei valori delle variabili.

\begin{concetto}
\begin{Proposizione}{Tavole di verità dei connettivi}{tavole-connettivi}
    \begin{prerequisiti}
        \dip{Definizione}{def:valutazione}
    \end{prerequisiti}
    Per ogni valutazione $v$, i valori di $\neg p$, $p \wedge q$, $p \vee q$, $p \rightarrow q$ e $p \leftrightarrow q$
    in funzione dei valori di $p$ e $q$ sono dati dalle seguenti tavole di verità.

    \begin{tavola}{c|c}
  \intestazione $p$ & $\neg p$ \\ \hline
  $\vero$ & $\falso$ \\
  $\falso$ & $\vero$ \\
\end{tavola}

    \begin{tavola}{cc|cccc}
  \intestazione $p$ & $q$ & $p \wedge q$ & $p \vee q$ & $p \rightarrow q$ & $p \leftrightarrow q$ \\ \hline
  $\vero$ & $\vero$ & $\vero$ & $\vero$ & $\vero$ & $\vero$ \\
  $\vero$ & $\falso$ & $\falso$ & $\vero$ & $\falso$ & $\falso$ \\
  $\falso$ & $\vero$ & $\falso$ & $\vero$ & $\vero$ & $\falso$ \\
  $\falso$ & $\falso$ & $\falso$ & $\falso$ & $\vero$ & $\vero$ \\
\end{tavola}
\end{Proposizione}

\begin{dimostrazione}
    Ogni riga si ricava direttamente da una regola della definizione di valutazione.
    Ad esempio, nella riga $v(p) = \vero$, $v(q) = \falso$ della colonna $p \rightarrow q$: per la regola 5,
    $v(p \rightarrow q) = \vero$ se e solo se $v(p) = \falso$ oppure $v(q) = \vero$; qui non succede né l'una
    né l'altra cosa, quindi $v(p \rightarrow q) \neq \vero$, cioè $v(p \rightarrow q) = \falso$.
    Tutte le altre righe si ottengono allo stesso modo.
\end{dimostrazione}
\end{concetto}

\begin{spiegazione}
    Pensiamo a $p$ = «piove» e $q$ = «prendo l'ombrello». L'implicazione $p \rightarrow q$ è una \textbf{promessa}: «se piove, prendo l'ombrello».
    \begin{center}
    \renewcommand{\arraystretch}{1.2}
    \begin{tabular}{cc|c|l}
        $p$ & $q$ & $p \rightarrow q$ & la promessa \\ \hline
        $\vero$ & $\vero$ & $\vero$ & piove e prendo l'ombrello: mantenuta \\
        $\vero$ & $\falso$ & $\falso$ & piove e non lo prendo: \textbf{violata} \\
        $\falso$ & $\vero$ & $\vero$ & non piove e lo prendo lo stesso: non violata \\
        $\falso$ & $\falso$ & $\vero$ & non piove e non lo prendo: non violata \\
    \end{tabular}
    \end{center}
    L'unico modo di smentire la promessa è la seconda riga. Allo stesso modo $p \wedge q$ è vera solo se sono vere entrambe, $p \vee q$ è falsa solo
    se sono false entrambe (l'«oppure» della logica è \textbf{inclusivo}: vale anche quando sono vere tutte e due) e $p \leftrightarrow q$ è vera quando
    $p$ e $q$ hanno lo stesso valore.
\end{spiegazione}


È importante notare come abbiamo \textbf{ricavato} queste tavole dalle definizioni semantiche dei connettivi stessi.

\warning{L'implicazione con premessa falsa}{
    Dalla tavola, $p \rightarrow q$ è falsa \textbf{solo} quando $p$ è vera e $q$ è falsa.
    In particolare, se la premessa $p$ è falsa, allora $p \rightarrow q$ è vera qualunque sia $q$:
    «se $2 + 2 = 5$, allora la luna è di formaggio» è un'implicazione vera!
    Questo fatto tornerà utile con i quantificatori ristretti (Capitolo~\ref{cap:predicati}).
}

% ──────────────────────────────────────────────────────────────
\section{Tautologie, contraddizioni ed equivalenza logica}\label{sec:tautologie}

\dfn[tautologia]{Tautologia e contraddizione}{
    Una FBF $f$ che risulta \textbf{sempre vera}, cioè tale che $v(f) = \vero$ per ogni valutazione $v$,
    si dice \textbf{tautologia}.

    Una FBF $f$ che risulta \textbf{sempre falsa}, cioè tale che $v(f) = \falso$ per ogni valutazione $v$,
    si dice \textbf{contraddizione}.
}

In termini di tavole di verità: una FBF è una tautologia se la sua colonna contiene solo $\vero$,
ed è una contraddizione se la sua colonna contiene solo $\falso$.

\dfn[equivalenza]{Formule logicamente equivalenti}{
    Siano $f$ e $g$ due FBF. Se $f \leftrightarrow g$ è una tautologia, allora $f$ e $g$ si dicono
    \textbf{logicamente equivalenti}.
}

\begin{concetto}
\begin{Proposizione}{Equivalenza logica e tavole di verità}{criterio-colonne}
    \begin{prerequisiti}
        \dip{Definizione}{def:valutazione}
        \dip{Definizione}{def:tautologia}
        \dip{Definizione}{def:equivalenza}
    \end{prerequisiti}
    Due FBF $f$ e $g$ sono logicamente equivalenti se e solo se $v(f) = v(g)$ per ogni valutazione $v$,
    cioè se e solo se nella tavola di verità \textbf{le colonne di $f$ e di $g$ coincidono riga per riga}.
\end{Proposizione}

\begin{dimostrazione}
    \[
    \begin{aligned}
        f \text{ e } g \text{ logicamente equivalenti}
        &\sse f \leftrightarrow g \text{ è una tautologia} && \text{(Definizione~\ref{def:equivalenza})}\\
        &\sse v(f \leftrightarrow g) = \vero \text{ per ogni valutazione } v && \text{(Definizione~\ref{def:tautologia})}\\
        &\sse v(f) = v(g) \text{ per ogni valutazione } v && \text{(regola 6 della Definizione~\ref{def:valutazione})}.
    \end{aligned}
    \]
\end{dimostrazione}
\end{concetto}

\begin{spiegazione}
    Per decidere se due formule «dicono la stessa cosa» non serve calcolare la colonna di $f \leftrightarrow g$: si confrontano direttamente le
    colonne di $f$ e di $g$. Ad esempio, con $f = p \wedge q$ e $g = q \wedge p$:
    \begin{passi}
        \item si scrivono le $2^2 = 4$ righe con i valori di $p$ e $q$;
        \item si calcolano la colonna di $f$ e quella di $g$: in entrambe c'è $\vero$ solo nella riga in cui $p$ e $q$ sono vere;
        \item le colonne coincidono riga per riga, quindi $f$ e $g$ sono logicamente equivalenti.
    \end{passi}
    Basta invece \textbf{una sola} riga diversa per concludere che due formule non sono equivalenti.
\end{spiegazione}


È questo il metodo che useremo per dimostrare quasi tutte le tautologie notevoli: invece di calcolare la
colonna di $f \leftrightarrow g$, basta confrontare la colonna di $f$ con quella di $g$.

\begin{concetto}
\begin{Corollario}{Catene di equivalenze}{catene}
    \begin{prerequisiti}
        \dip{Proposizione}{prop:criterio-colonne}
    \end{prerequisiti}
    Siano $f$, $g$, $h$ FBF.
    \begin{enumerate}
        \item Se $f$ è logicamente equivalente a $g$, allora $g$ è logicamente equivalente a $f$.
        \item Se $f$ è logicamente equivalente a $g$ e $g$ è logicamente equivalente a $h$, allora
              $f$ è logicamente equivalente a $h$.
    \end{enumerate}
\end{Corollario}

\begin{dimostrazione}
    Per la Proposizione~\ref{prop:criterio-colonne} basta ragionare sui valori delle valutazioni.
    \begin{caso}{Punto 1: se $f$ è equivalente a $g$, allora $g$ è equivalente a $f$}
        Per ogni valutazione $v$ abbiamo $v(f) = v(g)$, quindi anche $v(g) = v(f)$.
    \end{caso}
    \begin{caso}{Punto 2: se $f$ è equivalente a $g$ e $g$ a $h$, allora $f$ è equivalente a $h$}
        Per ogni valutazione $v$ abbiamo $v(f) = v(g)$ e $v(g) = v(h)$, quindi $v(f) = v(h)$.
    \end{caso}
\end{dimostrazione}
\end{concetto}

\begin{spiegazione}
    L'equivalenza logica si comporta come l'uguaglianza tra numeri: se $a = b$ allora $b = a$, e se $a = b$ e $b = c$ allora $a = c$. Per questo con le
    formule si possono fare «conti» come in algebra, passando da una formula all'altra un'equivalenza alla volta:
    \[ \underbrace{p \rightarrow q}_{f_1} \ \leftrightarrow\ \underbrace{\neg p \vee q}_{f_2} \ \leftrightarrow\ \underbrace{q \vee \neg p}_{f_3}
       \qquad\text{quindi}\qquad f_1 \text{ e } f_3 \text{ sono equivalenti.} \]
\end{spiegazione}


Grazie al corollario possiamo scrivere \textbf{catene} del tipo $f_1 \leftrightarrow f_2 \leftrightarrow \dots \leftrightarrow f_n$,
in cui ogni passaggio è un'equivalenza logica: alla fine il primo e l'ultimo membro sono logicamente equivalenti.

Le due proposizioni che seguono sono usate di continuo, spesso senza dirlo esplicitamente, ogni volta che
«applichiamo una tautologia» all'interno di una formula più grande: conviene quindi averle enunciate una volta per tutte.

\begin{concetto}
\begin{Proposizione}{Tautologie e sostituzione delle variabili}{sost-variabili}
    \begin{prerequisiti}
        \dip{Definizione}{def:fbf-prop}
        \dip{Definizione}{def:valutazione}
        \dip{Definizione}{def:tautologia}
    \end{prerequisiti}
    Sia $\tau$ una tautologia in cui compare la variabile proposizionale $p$ e sia $f$ una FBF qualsiasi.
    Allora la FBF $\tau'$ ottenuta da $\tau$ sostituendo \textbf{ogni} occorrenza di $p$ con $f$ è ancora una tautologia.
\end{Proposizione}

\begin{dimostrazione}[Idea della dimostrazione]
    Fissiamo una valutazione $v$ e consideriamo la valutazione $w$ che assegna alle variabili gli stessi valori di $v$,
    tranne che a $p$, a cui assegna il valore $w(p) = v(f)$. Calcolare $v(\tau')$ con le regole della valutazione
    è esattamente lo stesso calcolo che si fa per $w(\tau)$: dove in $\tau$ compare $p$, in $\tau'$ compare $f$,
    che ha lo stesso valore. Quindi $v(\tau') = w(\tau) = \vero$, perché $\tau$ è una tautologia.
    Poiché $v$ era arbitraria, $\tau'$ è una tautologia.
\end{dimostrazione}
\end{concetto}

\begin{spiegazione}
    In una tautologia le variabili sono \textbf{segnaposto}: la formula è vera qualunque valore abbiano, quindi resta vera se al loro posto
    mettiamo formule qualsiasi. Ad esempio:
    \begin{passi}
        \item $p \vee \neg p$ è una tautologia (il terzo escluso, Proposizione~\ref{prop:terzoescluso});
        \item sostituendo \textbf{ogni} $p$ con $q \wedge r$ si ottiene $(q \wedge r) \vee \neg(q \wedge r)$, che è ancora una tautologia;
        \item attenzione: sostituendo solo \textbf{una} delle due $p$, ad esempio con $q$, si ottiene $q \vee \neg p$, che è falsa quando $q$ è falsa e $p$ è vera.
    \end{passi}
\end{spiegazione}


In pratica: tutte le tautologie che dimostreremo con le variabili $p, q, r$ restano vere se al posto delle
variabili mettiamo \textbf{FBF qualsiasi}. Ad esempio, dalla doppia negazione $\neg(\neg p) \leftrightarrow p$
otterremo che $\neg(\neg f) \leftrightarrow f$ per ogni FBF $f$.

\begin{concetto}
\begin{Proposizione}{Sostituzione di sottoformule equivalenti}{sostituzione}
    \begin{prerequisiti}
        \dip{Definizione}{def:fbf-prop}
        \dip{Definizione}{def:valutazione}
        \dip{Proposizione}{prop:criterio-colonne}
    \end{prerequisiti}
    Siano $f$ e $g$ due FBF logicamente equivalenti e sia $h$ una FBF in cui compare $f$.
    Se $h'$ è la FBF ottenuta da $h$ sostituendo $f$ con $g$, allora $h$ e $h'$ sono logicamente equivalenti.
\end{Proposizione}

\begin{dimostrazione}[Idea della dimostrazione]
    Fissata una valutazione $v$, il valore $v(h)$ si calcola a partire dai valori delle sottoformule di $h$,
    applicando le regole della Definizione~\ref{def:valutazione}. In questo calcolo $f$ interviene solo tramite
    il suo valore $v(f)$. Poiché $f$ e $g$ sono equivalenti, $v(f) = v(g)$ (Proposizione~\ref{prop:criterio-colonne}),
    quindi sostituendo $f$ con $g$ il calcolo procede in modo identico e $v(h) = v(h')$.
    Questo vale per ogni $v$, quindi $h$ e $h'$ sono logicamente equivalenti (di nuovo per la Proposizione~\ref{prop:criterio-colonne}).
\end{dimostrazione}
\end{concetto}

\begin{spiegazione}
    È come in algebra: in $2 \cdot (3 + 4)$ si può sostituire $3 + 4$ con $7$ senza cambiare il risultato. Sapendo che $p \rightarrow q$ equivale a $\neg p \vee q$
    (Proposizione~\ref{prop:impor}), nella formula $h = \neg(p \rightarrow q)$ possiamo sostituire la sottoformula $p \rightarrow q$ e ottenere $h' = \neg(\neg p \vee q)$, equivalente
    ad $h$. La differenza con la proposizione precedente: lì si sostituisce una \textbf{variabile} con una formula \textbf{qualsiasi}, in tutte le sue
    occorrenze; qui si sostituisce una \textbf{sottoformula} con una formula \textbf{equivalente}.
\end{spiegazione}


% ──────────────────────────────────────────────────────────────
\section{Tautologie notevoli}\label{sec:tautologie-notevoli}

Vediamo ora le tautologie più importanti. Tutte si dimostrano con lo stesso metodo: si costruisce la tavola
di verità e si controlla che la colonna della formula contenga solo $\vero$ oppure, per le equivalenze,
che le \textbf{colonne evidenziate} coincidano riga per riga (Proposizione~\ref{prop:criterio-colonne}).
Quando nella stessa tavola ci sono più equivalenze da controllare, le colonne da confrontare hanno lo
\textbf{stesso colore}.

Le proprietà commutativa, associativa e distributiva \textbf{fanno parte della teoria}, perché verranno usate più avanti.

% --- Doppia negazione
\begin{concetto}
\begin{Proposizione}{Doppia negazione}{doppianeg}
    \begin{prerequisiti}
        \dip{Proposizione}{prop:tavole-connettivi}
        \dip{Definizione}{def:tautologia}
    \end{prerequisiti}
    La FBF
    \[ \neg(\neg p) \leftrightarrow p \]
    è una tautologia.
\end{Proposizione}

\begin{dimostrazione}
    \begin{tavola}{c|ccT}
  \intestazione $p$ & $\neg p$ & $\neg(\neg p)$ & $\neg(\neg p) \leftrightarrow p$ \\ \hline
  $\vero$ & $\falso$ & $\vero$ & $\vero$ \\
  $\falso$ & $\vero$ & $\falso$ & $\vero$ \\
\end{tavola}
    L'ultima colonna contiene solo $\vero$. A parole: negare due volte riporta al valore di partenza.
\end{dimostrazione}
\end{concetto}

\begin{spiegazione}
    «Non è vero che non piove» vuol dire «piove»: due negazioni si annullano, come due cambi di segno, $-(-3) = 3$. Nel linguaggio di tutti i giorni a
    volte la doppia negazione è solo un rafforzativo («non ho visto nessuno»); in logica, invece, $\neg(\neg p)$ è sempre $p$.
\end{spiegazione}


% --- Idempotenza
\begin{concetto}
\begin{Proposizione}{Idempotenza di AND e OR}{idempotenza}
    \begin{prerequisiti}
        \dip{Proposizione}{prop:tavole-connettivi}
        \dip{Definizione}{def:tautologia}
    \end{prerequisiti}
    Le FBF
    \[ (p \wedge p) \leftrightarrow p \qquad\qquad (p \vee p) \leftrightarrow p \]
    sono tautologie.
\end{Proposizione}

\begin{dimostrazione}
    \begin{tavola}{c|ccTT}
  \intestazione $p$ & $p \wedge p$ & $p \vee p$ & $(p \wedge p) \leftrightarrow p$ & $(p \vee p) \leftrightarrow p$ \\ \hline
  $\vero$ & $\vero$ & $\vero$ & $\vero$ & $\vero$ \\
  $\falso$ & $\falso$ & $\falso$ & $\vero$ & $\vero$ \\
\end{tavola}
    Le ultime due colonne contengono solo $\vero$.
\end{dimostrazione}
\end{concetto}

\begin{spiegazione}
    Dire due volte la stessa cosa non aggiunge nulla: «piove e piove» e «piove oppure piove» significano semplicemente «piove». La stessa proprietà
    tornerà per gli insiemi ($x \cap x = x$, $x \cup x = x$) e per i reticoli ($x \wedge x = x$).
\end{spiegazione}


% --- Terzo escluso
\begin{concetto}
\begin{Proposizione}{Principio del terzo escluso}{terzoescluso}
    \begin{prerequisiti}
        \dip{Proposizione}{prop:tavole-connettivi}
        \dip{Definizione}{def:tautologia}
    \end{prerequisiti}
    La FBF
    \[ p \vee \neg p \]
    è una tautologia.
\end{Proposizione}

\begin{dimostrazione}
    \begin{tavola}{c|cT}
  \intestazione $p$ & $\neg p$ & $p \vee \neg p$ \\ \hline
  $\vero$ & $\falso$ & $\vero$ \\
  $\falso$ & $\vero$ & $\vero$ \\
\end{tavola}
    L'ultima colonna contiene solo $\vero$. A parole: una proposizione è vera oppure è falsa, non c'è una terza possibilità.
\end{dimostrazione}
\end{concetto}

\begin{spiegazione}
    «Oggi è lunedì oppure oggi non è lunedì» è vera qualunque giorno sia. Il terzo escluso giustifica le \textbf{dimostrazioni per casi}: per dimostrare
    un'affermazione si può trattare separatamente il caso in cui $p$ è vera e quello in cui $p$ è falsa, perché uno dei due si verifica sempre.
\end{spiegazione}


% --- Non contraddizione
\begin{concetto}
\begin{Proposizione}{Principio di non contraddizione}{noncontraddizione}
    \begin{prerequisiti}
        \dip{Proposizione}{prop:tavole-connettivi}
        \dip{Definizione}{def:tautologia}
    \end{prerequisiti}
    La FBF
    \[ \neg(p \wedge \neg p) \]
    è una tautologia.
\end{Proposizione}

\begin{dimostrazione}
    \begin{tavola}{c|ccT}
  \intestazione $p$ & $\neg p$ & $p \wedge \neg p$ & $\neg(p \wedge \neg p)$ \\ \hline
  $\vero$ & $\falso$ & $\falso$ & $\vero$ \\
  $\falso$ & $\vero$ & $\falso$ & $\vero$ \\
\end{tavola}
    L'ultima colonna contiene solo $\vero$. Equivalentemente, $p \wedge \neg p$ è una contraddizione:
    una proposizione non può essere contemporaneamente vera e falsa.
\end{dimostrazione}
\end{concetto}

\begin{spiegazione}
    «Piove e non piove» non può mai essere vera. È il principio su cui si basano le \textbf{dimostrazioni per assurdo}: se supponendo che $p$ sia falsa si
    arriva a una cosa e al suo contrario, allora $p$ non può essere falsa.
\end{spiegazione}


% --- Commutatività
\begin{concetto}
\begin{Proposizione}{Commutatività di AND, OR e IFF}{commutativita}
    \begin{prerequisiti}
        \dip{Proposizione}{prop:tavole-connettivi}
        \dip{Proposizione}{prop:criterio-colonne}
    \end{prerequisiti}
    Le FBF
    \[ (p \wedge q) \leftrightarrow (q \wedge p) \qquad (p \vee q) \leftrightarrow (q \vee p) \qquad (p \leftrightarrow q) \leftrightarrow (q \leftrightarrow p) \]
    sono tautologie.
\end{Proposizione}

\begin{dimostrazione}
    \begin{tavola}{cc|TTUUKK}
  \intestazione $p$ & $q$ & $p \wedge q$ & $q \wedge p$ & $p \vee q$ & $q \vee p$ & $p \leftrightarrow q$ & $q \leftrightarrow p$ \\ \hline
  $\vero$ & $\vero$ & $\vero$ & $\vero$ & $\vero$ & $\vero$ & $\vero$ & $\vero$ \\
  $\vero$ & $\falso$ & $\falso$ & $\falso$ & $\vero$ & $\vero$ & $\falso$ & $\falso$ \\
  $\falso$ & $\vero$ & $\falso$ & $\falso$ & $\vero$ & $\vero$ & $\falso$ & $\falso$ \\
  $\falso$ & $\falso$ & $\falso$ & $\falso$ & $\falso$ & $\falso$ & $\vero$ & $\vero$ \\
\end{tavola}
    Le colonne dello stesso colore coincidono ($p \wedge q$ con $q \wedge p$, $p \vee q$ con $q \vee p$,
    $p \leftrightarrow q$ con $q \leftrightarrow p$), quindi per la Proposizione~\ref{prop:criterio-colonne}
    le tre doppie implicazioni sono tautologie.
\end{dimostrazione}
\end{concetto}

\begin{spiegazione}
    «Piove e fa freddo» e «fa freddo e piove» dicono la stessa cosa: per $\wedge$, $\vee$ e $\leftrightarrow$ l'ordine non conta. Attenzione al linguaggio di tutti
    i giorni, dove a volte «e» indica anche un ordine nel tempo («si alzò e uscì»): in logica conta solo se le due affermazioni sono vere.
\end{spiegazione}


% --- Non commutatività di ->
\begin{concetto}
\begin{Proposizione}{L'implicazione non è commutativa}{noncommimp}
    \begin{prerequisiti}
        \dip{Proposizione}{prop:tavole-connettivi}
        \dip{Definizione}{def:tautologia}
    \end{prerequisiti}
    La FBF $(p \rightarrow q) \leftrightarrow (q \rightarrow p)$ \textbf{non} è una tautologia; quindi, in generale,
    $p \rightarrow q$ e $q \rightarrow p$ non sono logicamente equivalenti.
\end{Proposizione}

\begin{dimostrazione}
    \begin{tavola}{cc|ccT}
  \intestazione $p$ & $q$ & $p \rightarrow q$ & $q \rightarrow p$ & $(p \rightarrow q) \leftrightarrow (q \rightarrow p)$ \\ \hline
  $\vero$ & $\vero$ & $\vero$ & $\vero$ & $\vero$ \\
  $\vero$ & $\falso$ & $\falso$ & $\vero$ & $\falso$ \\
  $\falso$ & $\vero$ & $\vero$ & $\falso$ & $\falso$ \\
  $\falso$ & $\falso$ & $\vero$ & $\vero$ & $\vero$ \\
\end{tavola}
    Nelle righe $v(p) = \vero, v(q) = \falso$ e $v(p) = \falso, v(q) = \vero$ l'ultima colonna vale $\falso$,
    quindi la formula non è sempre vera. Ad esempio, «se piove allora prendo l'ombrello» non dice la stessa
    cosa di «se prendo l'ombrello allora piove».
\end{dimostrazione}
\end{concetto}

\begin{spiegazione}
    $p \rightarrow q$ e $q \rightarrow p$ sono affermazioni diverse. Con $p$ = «la figura è un quadrato» e $q$ = «la figura è un rettangolo»:
    \begin{center}
    \begin{tikzpicture}[font=\small]
        \draw[primary!70!black, fill=primary!10, line width=0.8pt] (0,0) rectangle (1.2,1.2);
        \node[align=center] at (0.6,-0.5) {quadrato:\\ $p$ e $q$ veri};
        \draw[orange!85!black, fill=orange!12, line width=0.8pt] (2.6,0.2) rectangle (4.8,1.0);
        \node[align=center] at (3.7,-0.5) {rettangolo non quadrato:\\ $q$ vero, $p$ falso};
        \node[align=left, anchor=west] at (5.6,0.55) {«se è un quadrato, è un rettangolo» ($p \rightarrow q$): vera;\\
            «se è un rettangolo, è un quadrato» ($q \rightarrow p$): falsa,\\ \hspace*{3mm}perché la figura arancione la smentisce.};
    \end{tikzpicture}
    \end{center}
\end{spiegazione}


% --- Associatività
\begin{concetto}
\begin{Proposizione}{Associatività di AND e OR}{associativita}
    \begin{prerequisiti}
        \dip{Proposizione}{prop:tavole-connettivi}
        \dip{Proposizione}{prop:criterio-colonne}
    \end{prerequisiti}
    Le FBF
    \[ \big((p \wedge q) \wedge r\big) \leftrightarrow \big(p \wedge (q \wedge r)\big) \qquad\qquad
       \big((p \vee q) \vee r\big) \leftrightarrow \big(p \vee (q \vee r)\big) \]
    sono tautologie.
\end{Proposizione}

\begin{dimostrazione}
    Con tre variabili le righe sono $2^3 = 8$. Per $\wedge$:
    \begin{tavola}{ccc|cTcT}
  \intestazione $p$ & $q$ & $r$ & $p \wedge q$ & $(p \wedge q) \wedge r$ & $q \wedge r$ & $p \wedge (q \wedge r)$ \\ \hline
  $\vero$ & $\vero$ & $\vero$ & $\vero$ & $\vero$ & $\vero$ & $\vero$ \\
  $\vero$ & $\vero$ & $\falso$ & $\vero$ & $\falso$ & $\falso$ & $\falso$ \\
  $\vero$ & $\falso$ & $\vero$ & $\falso$ & $\falso$ & $\falso$ & $\falso$ \\
  $\vero$ & $\falso$ & $\falso$ & $\falso$ & $\falso$ & $\falso$ & $\falso$ \\
  $\falso$ & $\vero$ & $\vero$ & $\falso$ & $\falso$ & $\vero$ & $\falso$ \\
  $\falso$ & $\vero$ & $\falso$ & $\falso$ & $\falso$ & $\falso$ & $\falso$ \\
  $\falso$ & $\falso$ & $\vero$ & $\falso$ & $\falso$ & $\falso$ & $\falso$ \\
  $\falso$ & $\falso$ & $\falso$ & $\falso$ & $\falso$ & $\falso$ & $\falso$ \\
\end{tavola}
    e per $\vee$:
    \begin{tavola}{ccc|cTcT}
  \intestazione $p$ & $q$ & $r$ & $p \vee q$ & $(p \vee q) \vee r$ & $q \vee r$ & $p \vee (q \vee r)$ \\ \hline
  $\vero$ & $\vero$ & $\vero$ & $\vero$ & $\vero$ & $\vero$ & $\vero$ \\
  $\vero$ & $\vero$ & $\falso$ & $\vero$ & $\vero$ & $\vero$ & $\vero$ \\
  $\vero$ & $\falso$ & $\vero$ & $\vero$ & $\vero$ & $\vero$ & $\vero$ \\
  $\vero$ & $\falso$ & $\falso$ & $\vero$ & $\vero$ & $\falso$ & $\vero$ \\
  $\falso$ & $\vero$ & $\vero$ & $\vero$ & $\vero$ & $\vero$ & $\vero$ \\
  $\falso$ & $\vero$ & $\falso$ & $\vero$ & $\vero$ & $\vero$ & $\vero$ \\
  $\falso$ & $\falso$ & $\vero$ & $\falso$ & $\vero$ & $\vero$ & $\vero$ \\
  $\falso$ & $\falso$ & $\falso$ & $\falso$ & $\falso$ & $\falso$ & $\falso$ \\
\end{tavola}
    In entrambe le tavole le colonne evidenziate coincidono, quindi si conclude con la Proposizione~\ref{prop:criterio-colonne}.
\end{dimostrazione}
\end{concetto}

\begin{spiegazione}
    Pensiamo a $p$, $q$, $r$ come a \textbf{interruttori}: vero vuol dire chiuso, cioè la corrente passa. Interruttori \textbf{in serie} realizzano $\wedge$ (la corrente
    va da $A$ a $B$ solo se sono tutti chiusi); interruttori \textbf{in parallelo} realizzano $\vee$ (basta che uno sia chiuso). Come li si raggruppa non conta:
    \begin{center}
    \begin{tikzpicture}[font=\small]
        \node (A) at (0,0) {$A$}; \node[interr] (p) at (1.1,0) {$p$}; \node[interr] (q) at (2.3,0) {$q$}; \node[interr] (r) at (3.5,0) {$r$};
        \node (B) at (4.6,0) {$B$};
        \draw[filo] (A) -- (p) -- (q) -- (r) -- (B);
        \node at (2.3,-1.25) {in serie: $p \wedge q \wedge r$};
        \begin{scope}[xshift=7.2cm]
            \node (A2) at (0,0) {$A$}; \node (B2) at (3.6,0) {$B$};
            \draw[filo] (A2) -- (0.6,0); \draw[filo] (3.0,0) -- (B2);
            \draw[filo] (0.6,0.7) -- (0.6,-0.7); \draw[filo] (3.0,0.7) -- (3.0,-0.7);
            \foreach \y/\l in {0.7/p, 0/q, -0.7/r} { \node[interr] (s\l) at (1.8,\y) {$\l$}; \draw[filo] (0.6,\y) -- (s\l) -- (3.0,\y); }
            \node at (1.8,-1.25) {in parallelo: $p \vee q \vee r$};
        \end{scope}
    \end{tikzpicture}
    \end{center}
\end{spiegazione}


Grazie all'associatività possiamo scrivere $p \wedge q \wedge r$ e $p \vee q \vee r$ \textbf{senza parentesi}:
comunque le mettiamo otteniamo formule logicamente equivalenti. Si noti che $p \wedge q \wedge r$ è vera solo
quando $p$, $q$, $r$ sono tutte vere, mentre $p \vee q \vee r$ è falsa solo quando sono tutte false.

% --- Distributività
\begin{concetto}
\begin{Proposizione}{Distributività reciproca di AND e OR}{distributivita}
    \begin{prerequisiti}
        \dip{Proposizione}{prop:tavole-connettivi}
        \dip{Proposizione}{prop:criterio-colonne}
    \end{prerequisiti}
    $\wedge$ e $\vee$ sono reciprocamente distributivi, cioè le FBF
    \[ \big(p \wedge (q \vee r)\big) \leftrightarrow \big((p \wedge q) \vee (p \wedge r)\big) \qquad\qquad
       \big(p \vee (q \wedge r)\big) \leftrightarrow \big((p \vee q) \wedge (p \vee r)\big) \]
    sono tautologie.
\end{Proposizione}

\begin{dimostrazione}
    Distributività di $\wedge$ rispetto a $\vee$:
    \begin{tavola}{ccc|cTccT}
  \intestazione $p$ & $q$ & $r$ & $q \vee r$ & $p \wedge (q \vee r)$ & $p \wedge q$ & $p \wedge r$ & $(p \wedge q) \vee (p \wedge r)$ \\ \hline
  $\vero$ & $\vero$ & $\vero$ & $\vero$ & $\vero$ & $\vero$ & $\vero$ & $\vero$ \\
  $\vero$ & $\vero$ & $\falso$ & $\vero$ & $\vero$ & $\vero$ & $\falso$ & $\vero$ \\
  $\vero$ & $\falso$ & $\vero$ & $\vero$ & $\vero$ & $\falso$ & $\vero$ & $\vero$ \\
  $\vero$ & $\falso$ & $\falso$ & $\falso$ & $\falso$ & $\falso$ & $\falso$ & $\falso$ \\
  $\falso$ & $\vero$ & $\vero$ & $\vero$ & $\falso$ & $\falso$ & $\falso$ & $\falso$ \\
  $\falso$ & $\vero$ & $\falso$ & $\vero$ & $\falso$ & $\falso$ & $\falso$ & $\falso$ \\
  $\falso$ & $\falso$ & $\vero$ & $\vero$ & $\falso$ & $\falso$ & $\falso$ & $\falso$ \\
  $\falso$ & $\falso$ & $\falso$ & $\falso$ & $\falso$ & $\falso$ & $\falso$ & $\falso$ \\
\end{tavola}
    Distributività di $\vee$ rispetto a $\wedge$:
    \begin{tavola}{ccc|cTccT}
  \intestazione $p$ & $q$ & $r$ & $q \wedge r$ & $p \vee (q \wedge r)$ & $p \vee q$ & $p \vee r$ & $(p \vee q) \wedge (p \vee r)$ \\ \hline
  $\vero$ & $\vero$ & $\vero$ & $\vero$ & $\vero$ & $\vero$ & $\vero$ & $\vero$ \\
  $\vero$ & $\vero$ & $\falso$ & $\falso$ & $\vero$ & $\vero$ & $\vero$ & $\vero$ \\
  $\vero$ & $\falso$ & $\vero$ & $\falso$ & $\vero$ & $\vero$ & $\vero$ & $\vero$ \\
  $\vero$ & $\falso$ & $\falso$ & $\falso$ & $\vero$ & $\vero$ & $\vero$ & $\vero$ \\
  $\falso$ & $\vero$ & $\vero$ & $\vero$ & $\vero$ & $\vero$ & $\vero$ & $\vero$ \\
  $\falso$ & $\vero$ & $\falso$ & $\falso$ & $\falso$ & $\vero$ & $\falso$ & $\falso$ \\
  $\falso$ & $\falso$ & $\vero$ & $\falso$ & $\falso$ & $\falso$ & $\vero$ & $\falso$ \\
  $\falso$ & $\falso$ & $\falso$ & $\falso$ & $\falso$ & $\falso$ & $\falso$ & $\falso$ \\
\end{tavola}
    In entrambe le tavole le colonne evidenziate coincidono.
\end{dimostrazione}
\end{concetto}

\begin{spiegazione}
    Con gli interruttori: $p \wedge (q \vee r)$ è $p$ in serie con il parallelo di $q$ e $r$; $(p \wedge q) \vee (p \wedge r)$ è il parallelo di due rami, ciascuno con
    una copia di $p$. In entrambi i circuiti la corrente passa se e solo se $p$ è chiuso e almeno uno tra $q$ e $r$ è chiuso:
    \begin{center}
    \begin{tikzpicture}[font=\small]
        \node (A) at (0,0) {$A$}; \node[interr] (p) at (1.0,0) {$p$};
        \node[interr] (q) at (2.8,0.5) {$q$}; \node[interr] (r) at (2.8,-0.5) {$r$}; \node (B) at (4.6,0) {$B$};
        \draw[filo] (A) -- (p) -- (1.9,0); \draw[filo] (1.9,0.5) -- (1.9,-0.5); \draw[filo] (3.7,0.5) -- (3.7,-0.5);
        \draw[filo] (1.9,0.5) -- (q) -- (3.7,0.5); \draw[filo] (1.9,-0.5) -- (r) -- (3.7,-0.5); \draw[filo] (3.7,0) -- (B);
        \node at (2.3,-1.2) {$p \wedge (q \vee r)$};
        \begin{scope}[xshift=7.2cm]
            \node (A2) at (0,0) {$A$}; \node (B2) at (4.0,0) {$B$};
            \draw[filo] (A2) -- (0.5,0); \draw[filo] (0.5,0.5) -- (0.5,-0.5); \draw[filo] (3.5,0.5) -- (3.5,-0.5); \draw[filo] (3.5,0) -- (B2);
            \node[interr] (p1) at (1.4,0.5) {$p$}; \node[interr] (q1) at (2.6,0.5) {$q$};
            \node[interr] (p2) at (1.4,-0.5) {$p$}; \node[interr] (r2) at (2.6,-0.5) {$r$};
            \draw[filo] (0.5,0.5) -- (p1) -- (q1) -- (3.5,0.5); \draw[filo] (0.5,-0.5) -- (p2) -- (r2) -- (3.5,-0.5);
            \node at (2.0,-1.2) {$(p \wedge q) \vee (p \wedge r)$};
        \end{scope}
    \end{tikzpicture}
    \end{center}
    La prima distributività ricorda quella dei numeri, $2 \cdot (3 + 4) = 2 \cdot 3 + 2 \cdot 4$. La seconda invece no: $2 + 3 \cdot 4 \neq (2 + 3)(2 + 4)$. È una
    particolarità della logica, che ritroveremo negli insiemi.
\end{spiegazione}


% --- De Morgan
\begin{concetto}
\begin{Proposizione}{Leggi di De Morgan}{demorgan}
    \begin{prerequisiti}
        \dip{Proposizione}{prop:tavole-connettivi}
        \dip{Proposizione}{prop:criterio-colonne}
    \end{prerequisiti}
    Le FBF
    \[ \neg(p \vee q) \leftrightarrow (\neg p \wedge \neg q) \qquad\qquad \neg(p \wedge q) \leftrightarrow (\neg p \vee \neg q) \]
    sono tautologie.
\end{Proposizione}

\begin{dimostrazione}
    \begin{tavola}{cc|cTccT}
  \intestazione $p$ & $q$ & $p \vee q$ & $\neg(p \vee q)$ & $\neg p$ & $\neg q$ & $\neg p \wedge \neg q$ \\ \hline
  $\vero$ & $\vero$ & $\vero$ & $\falso$ & $\falso$ & $\falso$ & $\falso$ \\
  $\vero$ & $\falso$ & $\vero$ & $\falso$ & $\falso$ & $\vero$ & $\falso$ \\
  $\falso$ & $\vero$ & $\vero$ & $\falso$ & $\vero$ & $\falso$ & $\falso$ \\
  $\falso$ & $\falso$ & $\falso$ & $\vero$ & $\vero$ & $\vero$ & $\vero$ \\
\end{tavola}
    \begin{tavola}{cc|cTccT}
  \intestazione $p$ & $q$ & $p \wedge q$ & $\neg(p \wedge q)$ & $\neg p$ & $\neg q$ & $\neg p \vee \neg q$ \\ \hline
  $\vero$ & $\vero$ & $\vero$ & $\falso$ & $\falso$ & $\falso$ & $\falso$ \\
  $\vero$ & $\falso$ & $\falso$ & $\vero$ & $\falso$ & $\vero$ & $\vero$ \\
  $\falso$ & $\vero$ & $\falso$ & $\vero$ & $\vero$ & $\falso$ & $\vero$ \\
  $\falso$ & $\falso$ & $\falso$ & $\vero$ & $\vero$ & $\vero$ & $\vero$ \\
\end{tavola}
    In entrambe le tavole le colonne evidenziate coincidono. A parole: «non è vero che (A oppure B)»
    significa «non A e non B»; «non è vero che (A e B)» significa «non A oppure non B».
\end{dimostrazione}
\end{concetto}

\begin{spiegazione}
    Negare un «oppure» dà un «e», e negare un «e» dà un «oppure»:
    \begin{itemize}
        \item «non è vero che (studio oppure lavoro)» vuol dire «non studio \textbf{e} non lavoro»;
        \item «non è vero che (sono alto e biondo)» vuol dire «non sono alto \textbf{oppure} non sono biondo».
    \end{itemize}
    Con gli interruttori: un circuito in serie è interrotto, cioè $\neg(p \wedge q)$, se e solo se \textbf{almeno uno} degli interruttori è aperto, cioè $\neg p \vee \neg q$.
\end{spiegazione}


% --- Relazione -> e v
\begin{concetto}
\begin{Proposizione}{Relazione tra IMPLIES e OR}{impor}
    \begin{prerequisiti}
        \dip{Proposizione}{prop:tavole-connettivi}
        \dip{Proposizione}{prop:criterio-colonne}
    \end{prerequisiti}
    La FBF
    \[ (p \rightarrow q) \leftrightarrow (\neg p \vee q) \]
    è una tautologia.
\end{Proposizione}

\begin{dimostrazione}
    \begin{tavola}{cc|TcT}
  \intestazione $p$ & $q$ & $p \rightarrow q$ & $\neg p$ & $\neg p \vee q$ \\ \hline
  $\vero$ & $\vero$ & $\vero$ & $\falso$ & $\vero$ \\
  $\vero$ & $\falso$ & $\falso$ & $\falso$ & $\falso$ \\
  $\falso$ & $\vero$ & $\vero$ & $\vero$ & $\vero$ \\
  $\falso$ & $\falso$ & $\vero$ & $\vero$ & $\vero$ \\
\end{tavola}
    Le colonne evidenziate coincidono.
\end{dimostrazione}
\end{concetto}

\begin{spiegazione}
    La promessa «se piove, prendo l'ombrello» è mantenuta esattamente quando «non piove, oppure prendo l'ombrello»: o non c'è nulla da mantenere,
    o la si mantiene. Questa equivalenza permette di eliminare $\rightarrow$ e di ragionare solo con $\neg$, $\wedge$ e $\vee$.
\end{spiegazione}


% --- Contrapposizione
\begin{concetto}
\begin{Proposizione}{Contrapposizione}{contrapposizione}
    \begin{prerequisiti}
        \dip{Proposizione}{prop:tavole-connettivi}
        \dip{Proposizione}{prop:criterio-colonne}
        \dipgruppo{Solo per la seconda dimostrazione:}
        \dip{Proposizione}{prop:impor}
        \dip{Proposizione}{prop:commutativita}
        \dip{Proposizione}{prop:doppianeg}
        \dip{Proposizione}{prop:sost-variabili}
        \dip{Proposizione}{prop:sostituzione}
        \dip{Corollario}{cor:catene}
    \end{prerequisiti}
    La FBF
    \[ (p \rightarrow q) \leftrightarrow (\neg q \rightarrow \neg p) \]
    è una tautologia.
\end{Proposizione}

\begin{dimostrazione}
    \begin{tavola}{cc|TccT}
  \intestazione $p$ & $q$ & $p \rightarrow q$ & $\neg q$ & $\neg p$ & $\neg q \rightarrow \neg p$ \\ \hline
  $\vero$ & $\vero$ & $\vero$ & $\falso$ & $\falso$ & $\vero$ \\
  $\vero$ & $\falso$ & $\falso$ & $\vero$ & $\falso$ & $\falso$ \\
  $\falso$ & $\vero$ & $\vero$ & $\falso$ & $\vero$ & $\vero$ \\
  $\falso$ & $\falso$ & $\vero$ & $\vero$ & $\vero$ & $\vero$ \\
\end{tavola}
    Le colonne evidenziate coincidono.
\end{dimostrazione}

\begin{dimostrazione}[Seconda dimostrazione, senza tavole]
    Usiamo le tautologie già dimostrate come regole di calcolo, concatenando equivalenze (Corollario~\ref{cor:catene}):
    \[
    \begin{aligned}
        p \rightarrow q &\leftrightarrow \neg p \vee q && \text{(relazione tra $\rightarrow$ e $\vee$)}\\
        &\leftrightarrow q \vee \neg p && \text{(commutatività di $\vee$)}\\
        &\leftrightarrow \neg(\neg q) \vee \neg p && \text{(doppia negazione, sostituita dentro la formula)}\\
        &\leftrightarrow \neg q \rightarrow \neg p && \text{(relazione tra $\rightarrow$ e $\vee$, con $\neg q$ al posto di $p$ e $\neg p$ al posto di $q$)}.
    \end{aligned}
    \]
\end{dimostrazione}
\end{concetto}

\begin{spiegazione}
    Da un'implicazione se ne ottengono altre tre scambiando e negando premessa e conclusione; solo la \textbf{contronominale} le è equivalente.
    Con $p$ = «è un quadrato» e $q$ = «è un rettangolo» (stesso colore = equivalenti):
    \begin{center}
    \begin{tikzpicture}[font=\small, sc/.style={rounded corners=3pt, inner sep=4pt, align=center, minimum width=6.4cm}]
        \node[sc, draw=seagreen!70!black, fill=seagreen!10] at (0,1.25) {$p \rightarrow q$ (l'implicazione)\\ se è un quadrato, è un rettangolo: $\vero$};
        \node[sc, draw=orange!85!black, fill=orange!10] at (7.2,1.25) {$q \rightarrow p$ (l'inversa)\\ se è un rettangolo, è un quadrato: $\falso$};
        \node[sc, draw=orange!85!black, fill=orange!10] at (0,0) {$\neg p \rightarrow \neg q$ (la contraria)\\ se non è un quadrato, non è un rettangolo: $\falso$};
        \node[sc, draw=seagreen!70!black, fill=seagreen!10] at (7.2,0) {$\neg q \rightarrow \neg p$ (la contronominale)\\ se non è un rettangolo, non è un quadrato: $\vero$};
    \end{tikzpicture}
    \end{center}
    Per dimostrare $p \rightarrow q$ si può quindi dimostrare $\neg q \rightarrow \neg p$: è la dimostrazione \textbf{per contrapposizione}.
\end{spiegazione}


% --- Negazione dell'implicazione
\begin{concetto}
\begin{Proposizione}{Negazione dell'implicazione}{negimp}
    \begin{prerequisiti}
        \dip{Proposizione}{prop:tavole-connettivi}
        \dip{Proposizione}{prop:criterio-colonne}
        \dipgruppo{Solo per la seconda dimostrazione:}
        \dip{Proposizione}{prop:impor}
        \dip{Proposizione}{prop:demorgan}
        \dip{Proposizione}{prop:doppianeg}
        \dip{Proposizione}{prop:sost-variabili}
        \dip{Proposizione}{prop:sostituzione}
        \dip{Corollario}{cor:catene}
    \end{prerequisiti}
    La FBF
    \[ \neg(p \rightarrow q) \leftrightarrow (p \wedge \neg q) \]
    è una tautologia.
\end{Proposizione}

\begin{dimostrazione}
    \begin{tavola}{cc|cTcT}
  \intestazione $p$ & $q$ & $p \rightarrow q$ & $\neg(p \rightarrow q)$ & $\neg q$ & $p \wedge \neg q$ \\ \hline
  $\vero$ & $\vero$ & $\vero$ & $\falso$ & $\falso$ & $\falso$ \\
  $\vero$ & $\falso$ & $\falso$ & $\vero$ & $\vero$ & $\vero$ \\
  $\falso$ & $\vero$ & $\vero$ & $\falso$ & $\falso$ & $\falso$ \\
  $\falso$ & $\falso$ & $\vero$ & $\falso$ & $\vero$ & $\falso$ \\
\end{tavola}
    Le colonne evidenziate coincidono.
\end{dimostrazione}

\begin{dimostrazione}[Seconda dimostrazione, senza tavole]
    \[
    \begin{aligned}
        \neg(p \rightarrow q) &\leftrightarrow \neg(\neg p \vee q) && \text{(relazione tra $\rightarrow$ e $\vee$)}\\
        &\leftrightarrow \neg(\neg p) \wedge \neg q && \text{(De Morgan, con $\neg p$ al posto di $p$)}\\
        &\leftrightarrow p \wedge \neg q && \text{(doppia negazione)}.
    \end{aligned}
    \]
    A parole: «se A allora B» è falsa esattamente quando A è vera ma B no.
\end{dimostrazione}
\end{concetto}

\begin{spiegazione}
    Per \textbf{smentire} un'implicazione basta un controesempio: un caso in cui la premessa è vera e la conclusione è falsa. «Non è vero che se studi
    passi l'esame» vuol dire «hai studiato \textbf{e} non hai passato l'esame». Ad esempio «se $n$ è dispari, allora $n$ è primo» è falsa: $n = 9$ è dispari
    e non è primo.
\end{spiegazione}


% --- Doppia implicazione
\begin{concetto}
\begin{Proposizione}{Doppia implicazione}{doppiaimp}
    \begin{prerequisiti}
        \dip{Proposizione}{prop:tavole-connettivi}
        \dip{Proposizione}{prop:criterio-colonne}
    \end{prerequisiti}
    La FBF
    \[ (p \leftrightarrow q) \leftrightarrow \big((p \rightarrow q) \wedge (q \rightarrow p)\big) \]
    è una tautologia.
\end{Proposizione}

\begin{dimostrazione}
    \begin{tavola}{cc|TccT}
  \intestazione $p$ & $q$ & $p \leftrightarrow q$ & $p \rightarrow q$ & $q \rightarrow p$ & $(p \rightarrow q) \wedge (q \rightarrow p)$ \\ \hline
  $\vero$ & $\vero$ & $\vero$ & $\vero$ & $\vero$ & $\vero$ \\
  $\vero$ & $\falso$ & $\falso$ & $\falso$ & $\vero$ & $\falso$ \\
  $\falso$ & $\vero$ & $\falso$ & $\vero$ & $\falso$ & $\falso$ \\
  $\falso$ & $\falso$ & $\vero$ & $\vero$ & $\vero$ & $\vero$ \\
\end{tavola}
    Le colonne evidenziate coincidono. È il motivo per cui, per dimostrare un «se e solo se», si dimostrano
    separatamente le due implicazioni $(\Rightarrow)$ e $(\Leftarrow)$.
\end{dimostrazione}
\end{concetto}

\begin{spiegazione}
    «$p$ se e solo se $q$» sono due promesse in una: «se $p$ allora $q$» \textbf{e} «se $q$ allora $p$». Per questo le dimostrazioni di un «se e solo se»
    hanno sempre due parti. Ad esempio, «un intero $n$ è pari se e solo se $n + 1$ è dispari» si dimostra provando che se $n$ è pari allora $n + 1$ è dispari
    $(\rightarrow)$ e che se $n + 1$ è dispari allora $n$ è pari $(\leftarrow)$.
\end{spiegazione}


% --- Negazione della doppia implicazione
\begin{concetto}
\begin{Proposizione}{Negazione della doppia implicazione}{negdoppiaimp}
    \begin{prerequisiti}
        \dip{Proposizione}{prop:tavole-connettivi}
        \dip{Proposizione}{prop:criterio-colonne}
    \end{prerequisiti}
    La FBF
    \[ (p \leftrightarrow q) \leftrightarrow (\neg p \leftrightarrow \neg q) \]
    è una tautologia.
\end{Proposizione}

\begin{dimostrazione}
    \begin{tavola}{cc|TccT}
  \intestazione $p$ & $q$ & $p \leftrightarrow q$ & $\neg p$ & $\neg q$ & $\neg p \leftrightarrow \neg q$ \\ \hline
  $\vero$ & $\vero$ & $\vero$ & $\falso$ & $\falso$ & $\vero$ \\
  $\vero$ & $\falso$ & $\falso$ & $\falso$ & $\vero$ & $\falso$ \\
  $\falso$ & $\vero$ & $\falso$ & $\vero$ & $\falso$ & $\falso$ \\
  $\falso$ & $\falso$ & $\vero$ & $\vero$ & $\vero$ & $\vero$ \\
\end{tavola}
    Le colonne evidenziate coincidono.
\end{dimostrazione}
\end{concetto}

\begin{spiegazione}
    Se due affermazioni sono vere esattamente negli stessi casi, lo sono anche le loro negazioni. Da «un triangolo è equilatero se e solo se ha i tre
    angoli uguali» segue subito «un triangolo non è equilatero se e solo se non ha i tre angoli uguali».
\end{spiegazione}


\begin{concetto}
\begin{Proposizione}{Tautologie e contraddizioni dentro una formula}{neutri}
    \begin{prerequisiti}
        \dip{Definizione}{def:valutazione}
        \dip{Definizione}{def:tautologia}
        \dip{Proposizione}{prop:criterio-colonne}
    \end{prerequisiti}
    Sia $t$ una tautologia, $c$ una contraddizione e $p$ una FBF. Allora sono logicamente equivalenti:
    \[ t \wedge p \leftrightarrow p \qquad\quad c \vee p \leftrightarrow p \qquad\quad c \wedge p \leftrightarrow c \qquad\quad t \vee p \leftrightarrow t. \]
    Negli appunti si scrive spesso $\vero$ al posto di una tautologia e $\falso$ al posto di una contraddizione,
    ad esempio $(\falso \vee p) \leftrightarrow p$ e $(p \wedge \vero) \leftrightarrow p$.
\end{Proposizione}
\begin{dimostrazione}
    Fissiamo una valutazione $v$: per definizione $v(t) = \vero$ e $v(c) = \falso$. In ogni caso mostriamo che i due
    membri hanno lo stesso valore; poiché $v$ è arbitraria, si conclude con la Proposizione~\ref{prop:criterio-colonne}.
    \begin{caso}{$t \wedge p \leftrightarrow p$}
        Per la regola 3, $v(t \wedge p) = \vero$ se e solo se $v(t) = \vero$ e $v(p) = \vero$, cioè se e solo se $v(p) = \vero$.
        Quindi $v(t \wedge p) = v(p)$.
    \end{caso}
    \begin{caso}{$c \vee p \leftrightarrow p$}
        Per la regola 4, $v(c \vee p) = \vero$ se e solo se $v(c) = \vero$ oppure $v(p) = \vero$, cioè se e solo se $v(p) = \vero$.
        Quindi $v(c \vee p) = v(p)$.
    \end{caso}
    \begin{caso}{$c \wedge p \leftrightarrow c$}
        Poiché $v(c) = \falso$, per la regola 3 $v(c \wedge p) = \falso = v(c)$.
    \end{caso}
    \begin{caso}{$t \vee p \leftrightarrow t$}
        Poiché $v(t) = \vero$, per la regola 4 $v(t \vee p) = \vero = v(t)$.
    \end{caso}
\end{dimostrazione}
\end{concetto}

\begin{spiegazione}
    Se pensiamo a $\vero$ come al numero $1$ e a $\falso$ come al numero $0$, $\wedge$ si comporta come il prodotto e $\vee$ come il massimo (con $x \in \{0, 1\}$):
    \begin{center}
    \renewcommand{\arraystretch}{1.2}
    \begin{tabular}{l|l}
        logica & numeri \\ \hline
        $\vero \wedge p \leftrightarrow p$ & $1 \cdot x = x$ \\
        $\falso \wedge p \leftrightarrow \falso$ & $0 \cdot x = 0$ \\
        $\falso \vee p \leftrightarrow p$ & $\max(0, x) = x$ \\
        $\vero \vee p \leftrightarrow \vero$ & $\max(1, x) = 1$ \\
    \end{tabular}
    \end{center}
    Una tautologia è un $\vero$ «travestito»: in un AND non conta, in un OR assorbe tutto.
\end{spiegazione}


\info{Come si usa la negazione della doppia implicazione}{
    Il nome viene dal fatto che si \textbf{negano entrambi i membri} della doppia implicazione.
    Insieme alla Proposizione~\ref{prop:sost-variabili}, dice che per ogni coppia di FBF $f$, $g$:
    \begin{center}
        se $f \leftrightarrow g$ è vera, allora anche $\neg f \leftrightarrow \neg g$ è vera (e viceversa).
    \end{center}
    Negli appunti del prof.\ è usata proprio così, per passare da un'equivalenza tra due affermazioni
    all'equivalenza tra le loro negazioni (ad esempio nelle caratterizzazioni delle funzioni iniettive e
    negli insiemi ordinati).
}

% ──────────────────────────────────────────────────────────────
\section{Altri connettivi: XOR, NAND, NOR}\label{sec:altri-connettivi}

\dfn[xor-nand-nor]{I connettivi XOR, NAND e NOR}{
    Definiamo, attraverso le loro tavole di verità, tre nuovi connettivi binari:
    \begin{itemize}
        \item \textbf{XOR} (anche detto \textit{Aut}, simbolo $\veebar$): «o l'uno o l'altro, ma non entrambi»;
        \item \textbf{NAND} (\textit{not-and}, simbolo $\mid$);
        \item \textbf{NOR} (\textit{not-or}, simbolo $\downarrow$).
    \end{itemize}
    \begin{tavola}{cc|ccc}
  \intestazione $p$ & $q$ & $p \XOR q$ & $p \NAND q$ & $p \NOR q$ \\ \hline
  $\vero$ & $\vero$ & $\falso$ & $\falso$ & $\falso$ \\
  $\vero$ & $\falso$ & $\vero$ & $\vero$ & $\falso$ \\
  $\falso$ & $\vero$ & $\vero$ & $\vero$ & $\falso$ \\
  $\falso$ & $\falso$ & $\falso$ & $\vero$ & $\vero$ \\
\end{tavola}
    Li aggiungiamo all'insieme $C$ dei connettivi: la regola 3 della Definizione~\ref{def:fbf-prop} vale anche per loro
    e il valore di una valutazione su $p \XOR q$, $p \NAND q$, $p \NOR q$ è dato dalla tavola.
}

A parole: $p \XOR q$ è vera quando \textbf{esattamente una} tra $p$ e $q$ è vera; $p \NAND q$ è falsa solo quando
$p$ e $q$ sono entrambe vere; $p \NOR q$ è vera solo quando $p$ e $q$ sono entrambe false.

\begin{concetto}
\begin{Proposizione}{XOR come negazione di IFF}{xoriff}
    \begin{prerequisiti}
        \dip{Definizione}{def:xor-nand-nor}
        \dip{Proposizione}{prop:tavole-connettivi}
        \dip{Proposizione}{prop:criterio-colonne}
    \end{prerequisiti}
    La FBF
    \[ (p \XOR q) \leftrightarrow \neg(p \leftrightarrow q) \]
    è una tautologia.
\end{Proposizione}

\begin{dimostrazione}
    \begin{tavola}{cc|TcT}
  \intestazione $p$ & $q$ & $p \XOR q$ & $p \leftrightarrow q$ & $\neg(p \leftrightarrow q)$ \\ \hline
  $\vero$ & $\vero$ & $\falso$ & $\vero$ & $\falso$ \\
  $\vero$ & $\falso$ & $\vero$ & $\falso$ & $\vero$ \\
  $\falso$ & $\vero$ & $\vero$ & $\falso$ & $\vero$ \\
  $\falso$ & $\falso$ & $\falso$ & $\vero$ & $\falso$ \\
\end{tavola}
    Le colonne evidenziate coincidono: $p \XOR q$ è vera esattamente quando $p$ e $q$ hanno valori \textbf{diversi},
    cioè quando $p \leftrightarrow q$ è falsa.
\end{dimostrazione}
\end{concetto}

\begin{spiegazione}
    XOR è il «diverso da»: $p \XOR q$ è vera quando $p$ e $q$ hanno valori \textbf{diversi}, mentre $p \leftrightarrow q$ è vera quando hanno valori \textbf{uguali}. È l'«o»
    esclusivo di «o mangi la minestra o salti dalla finestra»: una delle due cose, ma non entrambe.
\end{spiegazione}


\begin{concetto}
\begin{Proposizione}{Esplicitazioni di XOR}{xorespl}
    \begin{prerequisiti}
        \dip{Definizione}{def:xor-nand-nor}
        \dip{Proposizione}{prop:tavole-connettivi}
        \dip{Proposizione}{prop:criterio-colonne}
        \dip{Corollario}{cor:catene}
    \end{prerequisiti}
    Le seguenti FBF sono logicamente equivalenti:
    \[ p \XOR q \qquad\qquad (p \wedge \neg q) \vee (q \wedge \neg p) \qquad\qquad (p \vee q) \wedge \neg(p \wedge q). \]
\end{Proposizione}

\begin{dimostrazione}
    Confrontiamo la prima formula con la seconda
    \begin{tavola}{cc|TccT}
  \intestazione $p$ & $q$ & $p \XOR q$ & $p \wedge \neg q$ & $q \wedge \neg p$ & $(p \wedge \neg q) \vee (q \wedge \neg p)$ \\ \hline
  $\vero$ & $\vero$ & $\falso$ & $\falso$ & $\falso$ & $\falso$ \\
  $\vero$ & $\falso$ & $\vero$ & $\vero$ & $\falso$ & $\vero$ \\
  $\falso$ & $\vero$ & $\vero$ & $\falso$ & $\vero$ & $\vero$ \\
  $\falso$ & $\falso$ & $\falso$ & $\falso$ & $\falso$ & $\falso$ \\
\end{tavola}
    e la prima con la terza:
    \begin{tavola}{cc|TccT}
  \intestazione $p$ & $q$ & $p \XOR q$ & $p \vee q$ & $\neg(p \wedge q)$ & $(p \vee q) \wedge \neg(p \wedge q)$ \\ \hline
  $\vero$ & $\vero$ & $\falso$ & $\vero$ & $\falso$ & $\falso$ \\
  $\vero$ & $\falso$ & $\vero$ & $\vero$ & $\vero$ & $\vero$ \\
  $\falso$ & $\vero$ & $\vero$ & $\vero$ & $\vero$ & $\vero$ \\
  $\falso$ & $\falso$ & $\falso$ & $\falso$ & $\vero$ & $\falso$ \\
\end{tavola}
    In entrambi i casi le colonne evidenziate coincidono; per il Corollario~\ref{cor:catene} anche la seconda
    e la terza formula sono equivalenti tra loro. A parole: «esattamente una tra $p$ e $q$ è vera» equivale a
    «$p$ vera e $q$ falsa, oppure $q$ vera e $p$ falsa», e a «almeno una è vera, ma non entrambe».
\end{dimostrazione}
\end{concetto}

\begin{spiegazione}
    Le tre formule sono tre modi di dire «esattamente una delle due»:
    \begin{passi}
        \item $p \XOR q$: esattamente una tra $p$ e $q$ è vera;
        \item $(p \wedge \neg q) \vee (q \wedge \neg p)$: si elencano i due casi buoni, «$p$ sì e $q$ no» oppure «$q$ sì e $p$ no»;
        \item $(p \vee q) \wedge \neg(p \wedge q)$: «almeno una», ma togliendo il caso cattivo «tutte e due».
    \end{passi}
    Per gli insiemi diventeranno le due espressioni della differenza simmetrica (Proposizione~\ref{prop:esplicitazioni-delta}).
\end{spiegazione}


\begin{concetto}
\begin{Proposizione}{Commutatività e associatività di XOR}{xorcommassoc}
    \begin{prerequisiti}
        \dip{Definizione}{def:xor-nand-nor}
        \dip{Proposizione}{prop:criterio-colonne}
    \end{prerequisiti}
    Le FBF
    \[ (p \XOR q) \leftrightarrow (q \XOR p) \qquad\qquad \big((p \XOR q) \XOR r\big) \leftrightarrow \big(p \XOR (q \XOR r)\big) \]
    sono tautologie.
\end{Proposizione}

\begin{dimostrazione}
    Commutatività:
    \begin{tavola}{cc|TT}
  \intestazione $p$ & $q$ & $p \XOR q$ & $q \XOR p$ \\ \hline
  $\vero$ & $\vero$ & $\falso$ & $\falso$ \\
  $\vero$ & $\falso$ & $\vero$ & $\vero$ \\
  $\falso$ & $\vero$ & $\vero$ & $\vero$ \\
  $\falso$ & $\falso$ & $\falso$ & $\falso$ \\
\end{tavola}
    Associatività:
    \begin{tavola}{ccc|cTcT}
  \intestazione $p$ & $q$ & $r$ & $p \XOR q$ & $(p \XOR q) \XOR r$ & $q \XOR r$ & $p \XOR (q \XOR r)$ \\ \hline
  $\vero$ & $\vero$ & $\vero$ & $\falso$ & $\vero$ & $\falso$ & $\vero$ \\
  $\vero$ & $\vero$ & $\falso$ & $\falso$ & $\falso$ & $\vero$ & $\falso$ \\
  $\vero$ & $\falso$ & $\vero$ & $\vero$ & $\falso$ & $\vero$ & $\falso$ \\
  $\vero$ & $\falso$ & $\falso$ & $\vero$ & $\vero$ & $\falso$ & $\vero$ \\
  $\falso$ & $\vero$ & $\vero$ & $\vero$ & $\falso$ & $\falso$ & $\falso$ \\
  $\falso$ & $\vero$ & $\falso$ & $\vero$ & $\vero$ & $\vero$ & $\vero$ \\
  $\falso$ & $\falso$ & $\vero$ & $\falso$ & $\vero$ & $\vero$ & $\vero$ \\
  $\falso$ & $\falso$ & $\falso$ & $\falso$ & $\falso$ & $\falso$ & $\falso$ \\
\end{tavola}
    In entrambe le tavole le colonne evidenziate coincidono.
\end{dimostrazione}
\end{concetto}

\begin{spiegazione}
    Pensiamo alla luce delle scale comandata da tre interruttori: ogni volta che se ne aziona uno, la luce cambia stato. Se all'inizio è spenta, è accesa
    esattamente quando è stato azionato un numero \textbf{dispari} di interruttori, cioè quando $p \XOR q \XOR r$ è vera:
    \begin{center}
    \renewcommand{\arraystretch}{1.15}
    \begin{tabular}{ccc|c|c}
        $p$ & $q$ & $r$ & quanti sono veri & $p \XOR q \XOR r$ \\ \hline
        $\vero$ & $\vero$ & $\vero$ & $3$ & $\vero$ \\
        $\vero$ & $\vero$ & $\falso$ & $2$ & $\falso$ \\
        $\vero$ & $\falso$ & $\falso$ & $1$ & $\vero$ \\
        $\falso$ & $\falso$ & $\falso$ & $0$ & $\falso$ \\
    \end{tabular}
    \end{center}
    Che si contino prima $p$ e $q$ e poi $r$, o prima $q$ e $r$ e poi $p$, conta solo se il totale è dispari: ecco perché XOR è associativo. In informatica
    è il \textbf{bit di parità}.
\end{spiegazione}


% ──────────────────────────────────────────────────────────────
\section{Riduzione dei connettivi}\label{sec:riduzione}

Attraverso le tautologie possiamo definire alcuni connettivi mediante altri. Ad esempio, per la
Proposizione~\ref{prop:impor}
\[ (p \rightarrow q) \leftrightarrow (\neg p \vee q), \]
dunque possiamo \textbf{eliminare} $\rightarrow$: ogni volta che compare, lo sostituiamo con $\neg$ e $\vee$
(Proposizione~\ref{prop:sostituzione}). Allo stesso modo, per la Proposizione~\ref{prop:doppiaimp}, possiamo eliminare
$\leftrightarrow$ e per la Proposizione~\ref{prop:xoriff} possiamo eliminare XOR. Quindi ci si può ridurre a casi più semplici.

\begin{concetto}
\begin{Proposizione}{Riduzioni a NAND e NOR}{riduzioni}
    \begin{prerequisiti}
        \dip{Definizione}{def:xor-nand-nor}
        \dip{Proposizione}{prop:tavole-connettivi}
        \dip{Proposizione}{prop:criterio-colonne}
        \dip{Corollario}{cor:catene}
    \end{prerequisiti}
    Valgono le seguenti equivalenze logiche:
    \begin{enumerate}
        \item $(p \NAND p) \leftrightarrow \neg p \leftrightarrow (p \NOR p)$;
        \item $(p \wedge q) \leftrightarrow \neg(p \NAND q)$ \quad e \quad $(p \NAND q) \leftrightarrow \neg(p \wedge q)$;
        \item $(p \vee q) \leftrightarrow \neg(p \NOR q)$ \quad e \quad $(p \NOR q) \leftrightarrow \neg(p \vee q)$.
    \end{enumerate}
\end{Proposizione}

\begin{dimostrazione}
    \begin{caso}{Punto 1: $(p \NAND p) \leftrightarrow \neg p \leftrightarrow (p \NOR p)$}
        \begin{tavola}{c|TTT}
  \intestazione $p$ & $\neg p$ & $p \NAND p$ & $p \NOR p$ \\ \hline
  $\vero$ & $\falso$ & $\falso$ & $\falso$ \\
  $\falso$ & $\vero$ & $\vero$ & $\vero$ \\
\end{tavola}
        Le tre colonne evidenziate coincidono; si conclude con il Corollario~\ref{cor:catene}.
    \end{caso}
    \begin{caso}{Punto 2: $(p \wedge q) \leftrightarrow \neg(p \NAND q)$ e $(p \NAND q) \leftrightarrow \neg(p \wedge q)$}
        \begin{tavola}{cc|TTUU}
  \intestazione $p$ & $q$ & $p \wedge q$ & $\neg(p \NAND q)$ & $p \NAND q$ & $\neg(p \wedge q)$ \\ \hline
  $\vero$ & $\vero$ & $\vero$ & $\vero$ & $\falso$ & $\falso$ \\
  $\vero$ & $\falso$ & $\falso$ & $\falso$ & $\vero$ & $\vero$ \\
  $\falso$ & $\vero$ & $\falso$ & $\falso$ & $\vero$ & $\vero$ \\
  $\falso$ & $\falso$ & $\falso$ & $\falso$ & $\vero$ & $\vero$ \\
\end{tavola}
        Le colonne dello stesso colore coincidono.
    \end{caso}
    \begin{caso}{Punto 3: $(p \vee q) \leftrightarrow \neg(p \NOR q)$ e $(p \NOR q) \leftrightarrow \neg(p \vee q)$}
        \begin{tavola}{cc|TTUU}
  \intestazione $p$ & $q$ & $p \vee q$ & $\neg(p \NOR q)$ & $p \NOR q$ & $\neg(p \vee q)$ \\ \hline
  $\vero$ & $\vero$ & $\vero$ & $\vero$ & $\falso$ & $\falso$ \\
  $\vero$ & $\falso$ & $\vero$ & $\vero$ & $\falso$ & $\falso$ \\
  $\falso$ & $\vero$ & $\vero$ & $\vero$ & $\falso$ & $\falso$ \\
  $\falso$ & $\falso$ & $\falso$ & $\falso$ & $\vero$ & $\vero$ \\
\end{tavola}
        Le colonne dello stesso colore coincidono.
    \end{caso}
    Del resto i nomi lo dicono già: NAND è «not-and» e NOR è «not-or».
\end{dimostrazione}
\end{concetto}

\begin{spiegazione}
    I nomi dicono già tutto: NAND è «not-and» e NOR è «not-or». E $p \NAND p$ vuol dire «non è vero che ($p$ e $p$)», cioè, per l'idempotenza, «non $p$»;
    allo stesso modo $p \NOR p$ è «non ($p$ oppure $p$)», cioè ancora «non $p$».
\end{spiegazione}


Insomma, non è difficile dimostrare che \textbf{tutti} i connettivi si possono definire a partire dal solo NAND,
o anche dal solo NOR.

\begin{concetto}
\begin{Teorema}{Completezza di NAND e di NOR}{completezza-nand}
    \begin{prerequisiti}
        \dip{Definizione}{def:fbf-prop}
        \dip{Definizione}{def:xor-nand-nor}
        \dip{Proposizione}{prop:riduzioni}
        \dip{Proposizione}{prop:doppianeg}
        \dip{Proposizione}{prop:demorgan}
        \dip{Proposizione}{prop:impor}
        \dip{Proposizione}{prop:doppiaimp}
        \dip{Proposizione}{prop:xoriff}
        \dip{Proposizione}{prop:sost-variabili}
        \dip{Proposizione}{prop:sostituzione}
        \dip{Corollario}{cor:catene}
    \end{prerequisiti}
    Ogni FBF è logicamente equivalente a una FBF in cui compare come unico connettivo NAND.
    Lo stesso vale per NOR.
\end{Teorema}

\begin{dimostrazione}
    Basta esprimere con il solo NAND ciascuno dei connettivi $\neg$, $\wedge$, $\vee$, $\rightarrow$, $\leftrightarrow$
    (e anche XOR e NOR): poi, in una FBF qualsiasi, si sostituisce un connettivo alla volta partendo dalle sottoformule
    più interne, e ogni sostituzione produce una formula equivalente (Proposizione~\ref{prop:sostituzione}).
    Nei passaggi seguenti le tautologie sono applicate anche con FBF al posto delle variabili (Proposizione~\ref{prop:sost-variabili})
    e le equivalenze sono concatenate (Corollario~\ref{cor:catene}).
    \begin{caso}{Tutto con il solo NAND}
    \begin{itemize}
        \item \textbf{NOT}: $\neg p \leftrightarrow p \NAND p$ (Proposizione~\ref{prop:riduzioni}, punto 1).
        \item \textbf{AND}: $p \wedge q \leftrightarrow \neg(p \NAND q) \leftrightarrow (p \NAND q) \NAND (p \NAND q)$,
              per il punto 2 e poi per il punto 1 applicato alla FBF $p \NAND q$.
        \item \textbf{OR}: per la doppia negazione e De Morgan
              \[ p \vee q \leftrightarrow \neg(\neg p) \vee \neg(\neg q) \leftrightarrow \neg(\neg p \wedge \neg q). \]
              Per il punto 2, $\neg(f \wedge g) \leftrightarrow f \NAND g$ per ogni coppia di FBF $f$, $g$; quindi
              \[ p \vee q \leftrightarrow (\neg p) \NAND (\neg q) \leftrightarrow (p \NAND p) \NAND (q \NAND q). \]
        \item \textbf{IMPLIES}: per la relazione tra $\rightarrow$ e $\vee$, la doppia negazione e De Morgan
              \[ p \rightarrow q \leftrightarrow \neg p \vee q \leftrightarrow \neg p \vee \neg(\neg q) \leftrightarrow \neg(p \wedge \neg q)
                 \leftrightarrow p \NAND (\neg q) \leftrightarrow p \NAND (q \NAND q). \]
        \item \textbf{IFF}: per la doppia implicazione $p \leftrightarrow q \leftrightarrow (p \rightarrow q) \wedge (q \rightarrow p)$,
              e $\rightarrow$, $\wedge$ sono già stati espressi con il solo NAND.
        \item \textbf{XOR} e \textbf{NOR}: $p \XOR q \leftrightarrow \neg(p \leftrightarrow q)$ (Proposizione~\ref{prop:xoriff}) e
              $p \NOR q \leftrightarrow \neg(p \vee q)$ (Proposizione~\ref{prop:riduzioni}, punto 3), e $\neg$, $\leftrightarrow$, $\vee$ sono già stati espressi.
    \end{itemize}
    \end{caso}
    \begin{caso}{Tutto con il solo NOR}
        Si ragiona in modo analogo scambiando i ruoli di $\wedge$ e $\vee$:
        \begin{itemize}
            \item \textbf{NOT}: $\neg p \leftrightarrow p \NOR p$;
            \item \textbf{OR}: $p \vee q \leftrightarrow \neg(p \NOR q) \leftrightarrow (p \NOR q) \NOR (p \NOR q)$;
            \item \textbf{AND}: $p \wedge q \leftrightarrow \neg(\neg p \vee \neg q) \leftrightarrow (\neg p) \NOR (\neg q) \leftrightarrow (p \NOR p) \NOR (q \NOR q)$;
            \item \textbf{IMPLIES}, \textbf{IFF}, \textbf{XOR}, \textbf{NAND}: si esprimono tramite $\neg$, $\vee$, $\wedge$ come nel caso precedente.
        \end{itemize}
    \end{caso}
\end{dimostrazione}
\end{concetto}

\begin{spiegazione}
    Il teorema ha una conseguenza pratica importante: ogni circuito digitale si può costruire con \textbf{un solo tipo} di porta logica, la NAND (ed è
    così che sono fatti molti chip). Ecco NOT, AND e OR costruiti solo con porte NAND:
    \begin{center}
    \begin{tikzpicture}[circuit logic US, font=\small, every circuit symbol/.style={draw=primary!70!black, fill=primary!7, line width=0.8pt}]
        \node[nand gate, inputs=nn] (a) at (0,0) {};
        \draw (a.input 1) -- ++(-0.35,0) coordinate (a1); \draw (a.input 2) -- ++(-0.35,0) coordinate (a2); \draw (a1) -- (a2);
        \draw ($(a1)!0.5!(a2)$) -- ++(-0.45,0) node[left] {$p$};
        \draw (a.output) -- ++(0.45,0) node[right] {$\neg p$};
        \node at (0,-1.3) {NOT};
        \node[nand gate, inputs=nn] (b1) at (4.4,0) {};
        \draw (b1.input 1) -- ++(-0.6,0) node[left] {$p$}; \draw (b1.input 2) -- ++(-0.6,0) node[left] {$q$};
        \node[nand gate, inputs=nn] (b2) at (6.2,0) {};
        \draw (b1.output) -- ++(0.3,0) coordinate (bo); \draw (bo) |- (b2.input 1); \draw (bo) |- (b2.input 2);
        \draw (b2.output) -- ++(0.45,0) node[right] {$p \wedge q$};
        \node at (5.3,-1.3) {AND};
        \node[nand gate, inputs=nn] (c1) at (10.6,0.65) {};
        \node[nand gate, inputs=nn] (c2) at (10.6,-0.65) {};
        \foreach \g/\v in {c1/p, c2/q} {
            \draw (\g.input 1) -- ++(-0.35,0) coordinate (x1); \draw (\g.input 2) -- ++(-0.35,0) coordinate (x2); \draw (x1) -- (x2);
            \draw ($(x1)!0.5!(x2)$) -- ++(-0.4,0) node[left] {$\v$};
        }
        \node[nand gate, inputs=nn] (c3) at (12.6,0) {};
        \draw (c1.output) -- ++(0.35,0) |- (c3.input 1); \draw (c2.output) -- ++(0.35,0) |- (c3.input 2);
        \draw (c3.output) -- ++(0.45,0) node[right] {$p \vee q$};
        \node at (11.5,-1.45) {OR};
    \end{tikzpicture}
    \end{center}
    Sono esattamente le formule della dimostrazione: $\neg p \leftrightarrow p \NAND p$, $p \wedge q \leftrightarrow (p \NAND q) \NAND (p \NAND q)$ e $p \vee q \leftrightarrow (p \NAND p) \NAND (q \NAND q)$.
\end{spiegazione}



  % Logica dei predicati: termini, FBF, quantificatori, variabili libere, semantica
  \chapter{Logica dei predicati}\label{cap:predicati}

La logica proposizionale non basta per fare matematica: non sa parlare di \textbf{oggetti} (numeri, insiemi, \dots)
né dire «\textbf{per ogni}» o «\textbf{esiste}». La \textbf{logica dei predicati} (o \textbf{logica del primo ordine})
aggiunge proprio questi ingredienti. Seguiremo lo stesso percorso del capitolo precedente: alfabeto, sintassi
(termini e formule ben formate) e infine semantica.

% ──────────────────────────────────────────────────────────────
\section{Il linguaggio dell'aritmetica}\label{sec:aritmetica}

Informalmente, il modello che abbiamo per obiettivo è l'\textbf{aritmetica}, cioè il linguaggio del primo ordine
\[ (\mbN, +, \cdot, 0, s, \leq, =) \]
dove:
\begin{itemize}
    \item $\mbN$ è l'\textbf{insieme} dei numeri naturali, cioè gli «oggetti» di cui vogliamo parlare;
    \item $+$ e $\cdot$ sono \textbf{operazioni binarie};
    \item $0$ è una \textbf{costante};
    \item $s$ è una \textbf{funzione unaria}, il \textit{successore}: $s \colon n \in \mbN \mapsto n + 1 \in \mbN$;
    \item $\leq$ è un \textbf{predicato} (binario);
    \item $=$ è l'\textbf{uguale}.
\end{itemize}
Vogliamo poter scrivere cose come $x \leq y$ oppure $(x = y) \rightarrow \big(\forall z\,(z + x = 0)\big)$.

\dfn[arieta]{Arietà}{
    L'\textbf{arietà} (o $i$-arietà) di una funzione o di un predicato è il \textbf{numero dei suoi argomenti}.
    Una funzione o un predicato con $n$ argomenti si dice \textbf{$n$-ario} (\textit{unario} se $n = 1$,
    \textit{binario} se $n = 2$).
}

Nell'aritmetica, $+$ è una funzione binaria (o «2-aria», di arietà 2), la funzione $s$ è una funzione unaria,
mentre $\leq$ è un predicato binario, e così via.

% ──────────────────────────────────────────────────────────────
\section{L'alfabeto}\label{sec:alfabeto-pred}

\dfn[alfabeto-pred]{Alfabeto della logica dei predicati}{
    L'\textbf{alfabeto} della logica dei predicati è
    \[ \mcA = (V,\ \mathrm{Cost},\ P,\ F,\ B,\ \mathrm{Conn},\ Q,\ \{=\}) \]
    dove:
    \begin{itemize}
        \item $V = \{x_1, \dots, x_n\}$ è l'insieme delle \textbf{variabili};
        \item $\mathrm{Cost} = \{c_1, \dots, c_n\}$ è l'insieme delle \textbf{costanti} (individuali);
        \item $P = \{p_1^{i_1}, \dots, p_n^{i_n}\}$ è l'insieme delle \textbf{lettere predicative} (o \textbf{relazionali}) con $i$-arietà;
        \item $F = \{f_1^{i_1}, \dots, f_n^{i_n}\}$ è l'insieme delle \textbf{lettere funzionali} con $i$-arietà;
        \item $B = \{(,\ )\} \cup \{,\}$ è l'insieme degli \textbf{accessori} (parentesi e virgola);
        \item $\mathrm{Conn} = \{\neg, \wedge, \vee, \rightarrow, \leftrightarrow\}$ è l'insieme dei \textbf{connettivi logici};
        \item $Q = \{\forall\}$ è l'insieme dei \textbf{quantificatori};
        \item $\{=\}$ contiene il simbolo di \textbf{uguaglianza}.
    \end{itemize}
}

L'apice di una lettera ne indica l'arietà: ad esempio $p_1^{2}$ è una lettera predicativa binaria e $f_3^{1}$ è
una lettera funzionale unaria. Nell'aritmetica: $0$ è una costante, $s$, $+$ e $\cdot$ sono lettere funzionali
(di arietà 1, 2 e 2) e $\leq$ è una lettera predicativa binaria.

\warning{Attenzione ai nomi}{
    Qui $F$ è l'insieme delle \textbf{lettere funzionali}: non va confuso con il valore di verità $\falso$
    né con l'insieme $\mcF$ delle FBF del Capitolo~1. Allo stesso modo, $P$ qui è l'insieme delle lettere predicative
    e non quello delle variabili proposizionali.
}

% ──────────────────────────────────────────────────────────────
\section{Sintassi: i termini}\label{sec:termini}

Sintassi, ovvero: quali formule ha senso scrivere? Partiamo dai \textbf{termini}, che rappresentano gli
«oggetti» di cui parliamo (nell'aritmetica: i numeri).

\dfn[termini]{Termini}{
    I \textbf{termini} sono definiti dalle seguenti regole:
    \begin{enumerate}
        \item ogni variabile è un termine;
        \item ogni costante è un termine;
        \item se $f$ è una lettera funzionale $n$-aria e $t_1, \dots, t_n$ sono termini, allora $f(t_1, \dots, t_n)$ è un termine.
    \end{enumerate}
}

\begin{esempio}{Termini dell'aritmetica}
    Sono termini: $0$, \ $x$, \ $s(0)$, \ $s(s(0))$, \ $+(x, s(0))$, che scriveremo più comodamente $x + s(0)$.

    Non sono termini: $s(0, x)$ (perché $s$ è unaria e le si danno due argomenti) e $x \leq 0$
    (perché $\leq$ è una lettera \textbf{predicativa}: $x \leq 0$ non è un oggetto, ma un'affermazione, cioè una formula).
\end{esempio}

% ──────────────────────────────────────────────────────────────
\section{Sintassi: le formule ben formate}\label{sec:fbf-pred}

\dfn[fbf-pred]{Formule ben formate (FBF) del primo ordine}{
    Le \textbf{formule ben formate} della logica dei predicati sono definite dalle seguenti regole:
    \begin{enumerate}
        \item se $p$ è una lettera predicativa $n$-aria e $t_1, \dots, t_n$ sono termini, allora $p(t_1, \dots, t_n)$
              è una FBF, detta \textbf{atomica} (ad esempio $0 \leq s(0)$);
        \item se $t_1$ e $t_2$ sono termini, allora $t_1 = t_2$ è una FBF;
        \item se $a$ e $b$ sono FBF, allora $\neg a$, \ $a \wedge b$, \ $a \vee b$, \ $a \rightarrow b$, \ $a \leftrightarrow b$ sono FBF;
        \item se $a$ è una FBF e $x$ è una variabile, allora $(\forall x)(a)$ è una FBF.
              La FBF $a$ si dice \textbf{scopo} del quantificatore $\forall x$.
    \end{enumerate}
}

\info{Lettere predicative come proprietà}{
    Se $p$ è una lettera predicativa e $x$ una variabile, scriveremo $p(x)$ stando a significare che
    «\textbf{$x$ ha la proprietà $p$}».
}

\textbf{Significato di $\forall$.} La FBF $(\forall x)(a)$ vuol dire che $a$ (vista come un predicato su $x$) vale
\textbf{per ogni} termine $t$ messo al posto di $x$. Ad esempio $(\forall x)(0 \leq x)$ dice che ogni numero naturale
è maggiore o uguale a $0$.

Notiamo che $a$ può anche \textbf{non contenere} $x$: in quel caso $(\forall x)(a)$ equivale ad $a$.
Esempio: $\forall x\big((s(0) < s(s(0))) \wedge (2 = s(1))\big)$ dice semplicemente che $s(0) < s(s(0))$ e $2 = s(1)$.

\info{Abbreviazioni usate negli esempi}{
    Per comodità scriveremo:
    \begin{itemize}
        \item $t_1 \neq t_2$ al posto di $\neg(t_1 = t_2)$;
        \item $t_1 < t_2$ al posto di $(t_1 \leq t_2) \wedge (t_1 \neq t_2)$, e analogamente $\geq$ e $>$;
        \item $1$ al posto di $s(0)$, $2$ al posto di $s(s(0))$, e così via;
        \item $xy$ al posto di $x \cdot y$.
    \end{itemize}
}

% ──────────────────────────────────────────────────────────────
\section{Quantificatore esistenziale ed esistenziale unico}\label{sec:esiste}

L'unico quantificatore dell'alfabeto è $\forall$. A questa lista aggiungiamo altri due quantificatori,
\textbf{definiti} a partire da quello che abbiamo.

\dfn[esiste]{Quantificatore esistenziale}{
    Sia $a$ una FBF e $x$ una variabile. Definiamo
    \[ (\exists x)(a) \ \defsse\ \neg\big((\forall x)(\neg a)\big). \]
    Si legge «\textbf{esiste} $x$ tale che $a$»: non è vero che, per ogni $x$, vale $\neg a$.
}

\dfn[esiste-unico]{Quantificatore esistenziale unico}{
    Sia $p(x)$ una FBF. Definiamo
    \[ (\exists! x)\big(p(x)\big) \ \defsse\ (\exists x)\Big((\forall y)\big(p(y) \leftrightarrow x = y\big)\Big) \]
    dove $x$ e $y$ sono due variabili tali che $x \neq y$ e $y$ non compare in $p(x)$.
    Si legge «\textbf{esiste un unico} $x$ tale che $p(x)$».
}

Perché la definizione di $\exists!$ funziona? Dice che esiste un $x$ tale che, per ogni $y$, $p(y)$ vale
\textbf{esattamente quando} $y = x$. Prendendo $y = x$ si ottiene che $x$ ha la proprietà $p$; e se un qualunque
$y$ ha la proprietà $p$, allora $y = x$. Cioè: $x$ ha la proprietà $p$ ed è l'unico ad averla.

\info{Notazione}{
    Invece di $(\forall x)(f)$ scriveremo anche $\forall x(f)$, e così anche per $\exists$ ed $\exists!$.
}

% ──────────────────────────────────────────────────────────────
\section{Variabili libere e vincolate}\label{sec:libere}

\dfn[libere-vincolate]{Variabili libere e vincolate}{
    Sia $f$ una FBF. Un'occorrenza di una variabile $x$ si dice \textbf{vincolata} in $f$ se appare nella forma
    «$\forall x$» (e quindi anche «$\exists x$») oppure se si trova nello \textbf{scopo} di un quantificatore
    $\forall x$ (o $\exists x$) all'interno di $f$.

    Un'occorrenza che non è nello scopo di nessun quantificatore sulla stessa variabile si dice \textbf{libera}.
}

\begin{esempio}{Occorrenze libere e vincolate}
    \begin{itemize}
        \item In $\forall x\,(x < 0)$ sia la prima sia la seconda occorrenza di $x$ sono vincolate.
        \item In $\exists x\,\big(p(x) \rightarrow p(y)\big)$, $x$ è sempre vincolata, mentre $y$ è libera.
    \end{itemize}
\end{esempio}

\dfn[chiusa]{Formula chiusa}{
    Una FBF si dice \textbf{chiusa} quando non ha variabili libere (cioè è priva di occorrenze libere di variabili).
}

% ──────────────────────────────────────────────────────────────
\section{Cenni di semantica: formule valide}\label{sec:valide}

Per quanto riguarda la semantica della logica del primo ordine, diversamente da quanto accadeva nella logica
proposizionale, ci troviamo davanti a FBF che possono \textbf{non avere alcun valore di verità}.
È il caso di formule come $x < 0$, in cui la verità dipenderebbe dal «valore assunto da $x$», oltre che
dall'interpretazione dei simboli $<$ e $0$.

D'altra parte, se $p$ è un predicato unario, allora $p(x) \vee \neg p(x)$ sarà \textbf{sempre vera},
qualunque sia il valore di $x$ e qualunque sia il significato di $p$.

\dfn[valida]{Formula valida}{
    Una FBF si dice \textbf{valida} se è vera in tutti i modelli, cioè qualunque sia l'interpretazione dei simboli
    (l'insieme su cui variano le variabili e il significato di costanti, lettere funzionali e lettere predicative)
    e qualunque sia il valore assegnato alle sue variabili libere.
}

Un esempio di \textit{modello} è proprio l'aritmetica $(\mbN, +, \cdot, 0, s, \leq, =)$ della Sezione~«\nameref{sec:aritmetica}».
Le formule chiuse (Definizione~\ref{def:chiusa}) sono quelle che, fissato un modello, sono vere o false senza bisogno
di dare un valore alle variabili.

\info{Le regole del Capitolo 1 continuano a valere}{
    Le formule valide giocano il ruolo che nella logica proposizionale avevano le tautologie. In particolare
    (e sarà usato di continuo):
    \begin{itemize}
        \item se in una tautologia sostituiamo le variabili proposizionali con FBF del primo ordine, otteniamo una
              formula valida (come nella Proposizione~\ref{prop:sost-variabili}); ad esempio $\neg(\neg a) \leftrightarrow a$
              è valida per ogni FBF $a$;
        \item diremo che due FBF $a$ e $b$ sono \textbf{logicamente equivalenti} se $a \leftrightarrow b$ è valida; come nella
              Proposizione~\ref{prop:sostituzione}, sostituendo in una FBF una sottoformula con una equivalente
              (anche dentro lo scopo di un quantificatore) si ottiene una FBF equivalente;
        \item le equivalenze si possono concatenare come nel Corollario~\ref{cor:catene}.
    \end{itemize}
}

% ──────────────────────────────────────────────────────────────
\section{Negazione dei quantificatori}\label{sec:negazione-quant}

Come nego «Ogni gatto è bianco»? Dicendo che «\textbf{Esiste} un gatto che \textbf{non} è bianco».
Quindi negare un $\forall$ produce un $\exists$ (e viceversa), e la negazione «entra» nello scopo.

\begin{concetto}
\begin{Teorema}{Negazione dei quantificatori}{negazione-quantificatori}
    \begin{prerequisiti}
        \dip{Definizione}{def:esiste}
        \dip{Definizione}{def:valida}
        \dip{Proposizione}{prop:doppianeg}
        \dip{Proposizione}{prop:negdoppiaimp}
        \dip{Proposizione}{prop:sost-variabili}
        \dip{Proposizione}{prop:sostituzione}
        \dip{Corollario}{cor:catene}
        \dipsez{sec:valide}
    \end{prerequisiti}
    Per ogni FBF $a$ e ogni variabile $x$ le seguenti formule sono valide:
    \begin{enumerate}
        \item $(\forall x)(\neg a) \leftrightarrow \neg\big((\exists x)(a)\big)$;
        \item $\neg\big((\forall x)(a)\big) \leftrightarrow (\exists x)(\neg a)$;
        \item $\neg\big((\exists x)(\neg a)\big) \leftrightarrow (\forall x)(a)$.
    \end{enumerate}
\end{Teorema}

\begin{dimostrazione}
    \begin{caso}{Punto 1: $(\forall x)(\neg a) \leftrightarrow \neg\big((\exists x)(a)\big)$}
        Per la Definizione~\ref{def:esiste}, $(\exists x)(a)$ è per definizione $\neg\big((\forall x)(\neg a)\big)$.
    Quindi $\neg\big((\exists x)(a)\big)$ è la formula $\neg\neg\big((\forall x)(\neg a)\big)$, che per la doppia negazione
    (applicata alla FBF $(\forall x)(\neg a)$) è equivalente a $(\forall x)(\neg a)$.
    \end{caso}
    \begin{caso}{Punto 2: $\neg\big((\forall x)(a)\big) \leftrightarrow (\exists x)(\neg a)$}
        Applichiamo la Definizione~\ref{def:esiste} alla FBF $\neg a$:
    \[ (\exists x)(\neg a) \ \text{ è per definizione } \ \neg\big((\forall x)(\neg\neg a)\big). \]
    Per la doppia negazione $\neg\neg a \leftrightarrow a$; sostituendo dentro lo scopo del quantificatore
    (Proposizione~\ref{prop:sostituzione}) otteniamo $(\forall x)(\neg\neg a) \leftrightarrow (\forall x)(a)$ e quindi,
    per la negazione della doppia implicazione, $\neg\big((\forall x)(\neg\neg a)\big) \leftrightarrow \neg\big((\forall x)(a)\big)$.
    Dunque $(\exists x)(\neg a) \leftrightarrow \neg\big((\forall x)(a)\big)$, che è il punto 2 letto da destra a sinistra.
    \end{caso}
    \begin{caso}{Punto 3: $\neg\big((\exists x)(\neg a)\big) \leftrightarrow (\forall x)(a)$}
        Applicando la negazione della doppia implicazione al punto 2 otteniamo
    \[ \neg\neg\big((\forall x)(a)\big) \leftrightarrow \neg\big((\exists x)(\neg a)\big), \]
    e per la doppia negazione $\neg\neg\big((\forall x)(a)\big) \leftrightarrow (\forall x)(a)$.
    Concatenando le due equivalenze (Corollario~\ref{cor:catene}) si ottiene il punto 3.
    \end{caso}
\end{dimostrazione}
\end{concetto}

\begin{spiegazione}
    Pensiamo a una classe di studenti e a $P(x)$ = «$x$ è stato promosso». Per smentire «tutti sono stati promossi» basta \textbf{uno} studente bocciato;
    per smentire «qualcuno è stato promosso» bisogna che siano stati bocciati \textbf{tutti}:
    \begin{center}
    \begin{tikzpicture}[font=\small, st/.style={circle, minimum size=4.5mm, inner sep=0pt, line width=0.6pt}]
        \foreach \i in {0,...,5} \node[st, draw=seagreen!70!black, fill=seagreen!25] at (\i*0.6,0) {};
        \node[align=center] at (1.5,-0.75) {$\forall x\,P(x)$: tutti promossi};
        \begin{scope}[xshift=5.2cm]
            \foreach \i in {0,...,4} \node[st, draw=seagreen!70!black, fill=seagreen!25] at (\i*0.6,0) {};
            \node[st, draw=cherryred!80!black, fill=cherryred!25] at (3.0,0) {};
            \node[align=center] at (1.5,-0.75) {$\neg\forall x\,P(x) \leftrightarrow \exists x\,\neg P(x)$:\\ c'è almeno un bocciato};
        \end{scope}
        \begin{scope}[xshift=10.4cm]
            \foreach \i in {0,...,5} \node[st, draw=cherryred!80!black, fill=cherryred!25] at (\i*0.6,0) {};
            \node[align=center] at (1.5,-0.75) {$\neg\exists x\,P(x) \leftrightarrow \forall x\,\neg P(x)$:\\ nessun promosso};
        \end{scope}
    \end{tikzpicture}
    \end{center}
    Attenzione a non confondere «non tutti sono promossi» (almeno uno è bocciato) con «tutti sono bocciati».
\end{spiegazione}


In pratica: \textbf{per negare un quantificatore si scambia $\forall$ con $\exists$ e si nega lo scopo}.

\begin{esempio}{Ogni gatto è bianco}
    Siano $G(x)$ = «$x$ è un gatto» e $B(x)$ = «$x$ è bianco». «Ogni gatto è bianco» si scrive
    $\forall x\,\big(G(x) \rightarrow B(x)\big)$. Negando:
    \[
    \begin{aligned}
        \neg\forall x\,\big(G(x) \rightarrow B(x)\big)
        &\leftrightarrow \exists x\,\neg\big(G(x) \rightarrow B(x)\big) && \text{(Teorema~\ref{teo:negazione-quantificatori}, punto 2)}\\
        &\leftrightarrow \exists x\,\big(G(x) \wedge \neg B(x)\big) && \text{(negazione dell'implicazione)}
    \end{aligned}
    \]
    cioè «esiste un gatto che non è bianco».
\end{esempio}

Il prossimo esempio fa parte della teoria, perché verrà usato più avanti.

\begin{esempio}{Negare $(\exists x)\big(f(x) \wedge \forall y(\neg g(x,y))\big)$}
    \begin{prerequisiti}
        \dip{Teorema}{teo:negazione-quantificatori}
        \dip{Definizione}{def:esiste}
        \dip{Proposizione}{prop:demorgan}
        \dip{Proposizione}{prop:impor}
        \dip{Proposizione}{prop:sostituzione}
        \dip{Corollario}{cor:catene}
    \end{prerequisiti}
    Qui $f$ e $g$ sono lettere predicative (unaria e binaria).
    \[
    \begin{aligned}
        \neg(\exists x)\big(f(x) \wedge \forall y(\neg g(x,y))\big)
        &\leftrightarrow (\forall x)\,\neg\big(f(x) \wedge \forall y(\neg g(x,y))\big) && \text{(Teorema~\ref{teo:negazione-quantificatori}, punto 1)}\\
        &\leftrightarrow (\forall x)\big(\neg f(x) \vee \neg\forall y(\neg g(x,y))\big) && \text{(De Morgan)}\\
        &\leftrightarrow (\forall x)\big(\neg f(x) \vee \exists y\,(g(x,y))\big) && \text{(definizione di $\exists$)}\\
        &\leftrightarrow (\forall x)\big(f(x) \rightarrow \exists y\,(g(x,y))\big) && \text{(relazione tra $\rightarrow$ e $\vee$)}.
    \end{aligned}
    \]
    A parole: «per ogni $x$, se $x$ ha la proprietà $f$, allora esiste $y$ tale che $g(x,y)$».
\end{esempio}

% ──────────────────────────────────────────────────────────────
\section{Quantificatori e congiunzione}\label{sec:quant-congiunzione}

\begin{concetto}
\begin{Proposizione}{Il quantificatore universale e la congiunzione}{forall-and}
    \begin{prerequisiti}
        \dip{Definizione}{def:fbf-pred}
        \dip{Definizione}{def:valida}
    \end{prerequisiti}
    Siano $p$ e $q$ FBF e $x$ una variabile. La formula
    \[ \big(\forall x\,(p \wedge q)\big) \leftrightarrow \big(\forall x\,(p) \wedge \forall x\,(q)\big) \]
    è valida.
\end{Proposizione}
\begin{dimostrazione}
    \begin{caso}{$(\rightarrow)$ Se $\forall x\,(p \wedge q)$ è vera, allora $\forall x\,(p) \wedge \forall x\,(q)$ è vera}
        Per ogni valore di $x$ è vera $p \wedge q$, quindi sono vere sia $p$ sia $q$. Allora $p$ è vera per ogni valore di $x$,
        cioè $\forall x\,(p)$, e lo stesso per $q$.
    \end{caso}
    \begin{caso}{$(\leftarrow)$ Se $\forall x\,(p) \wedge \forall x\,(q)$ è vera, allora $\forall x\,(p \wedge q)$ è vera}
        Fissato un qualunque valore di $x$, $p$ è vera (perché vale $\forall x\,(p)$) e $q$ è vera (perché vale $\forall x\,(q)$),
        quindi $p \wedge q$ è vera. Poiché il valore di $x$ era arbitrario, vale $\forall x\,(p \wedge q)$.
    \end{caso}
\end{dimostrazione}
\end{concetto}

\begin{spiegazione}
    «Tutti gli studenti hanno studiato e hanno superato l'esame» dice la stessa cosa di «tutti hanno studiato e tutti hanno superato l'esame»:
    $\forall$ si distribuisce su $\wedge$. Non si distribuisce invece su $\vee$: «ogni intero è pari o dispari» è vera, mentre «ogni intero è pari, oppure
    ogni intero è dispari» è falsa. E, come dice l'avviso qui sotto, $\exists$ non si distribuisce su $\wedge$.
\end{spiegazione}


\warning{Con $\exists$ non funziona}{
    L'analoga formula $\big(\exists x\,(p \wedge q)\big) \leftrightarrow \big(\exists x\,(p) \wedge \exists x\,(q)\big)$
    \textbf{non} è valida. Nell'aritmetica, con $p$ = «$x$ è pari» e $q$ = «$x$ è dispari», $\exists x\,(p)$ e
    $\exists x\,(q)$ sono vere, ma $\exists x\,(p \wedge q)$ è falsa: nessun numero è sia pari sia dispari.
}

% ──────────────────────────────────────────────────────────────
\section{Sostituire termini alle variabili libere}\label{sec:sostituzione-termini}

\dfn[notazione-sostituzione]{Sostituzione di termini}{
    Se la formula $p$ ha come variabili libere $x_1, \dots, x_n$, possiamo scrivere $p(x_1, \dots, x_n)$.
    Se $t_1, \dots, t_n$ sono termini, scriviamo $p(t_1, \dots, t_n)$ per la formula $p$ con $t_1$ al posto di $x_1$,
    \dots, $t_n$ al posto di $x_n$.
}

Attenzione: scrivere $p(x_1, \dots, x_n)$ \textbf{non} vuol dire che queste siano \textbf{tutte} le variabili libere di $p$.
L'effetto della sostituzione è quello di mettere, al posto delle occorrenze libere delle variabili, dei termini «concreti».

\begin{esempio}{Sostituzione}
    Sia $p$ la formula
    \[ (x = y \wedge x \neq 0) \vee \exists z\,(x + z = y). \]
    Possiamo scrivere $p(x)$, come anche $p(x, y)$, ma \textbf{non} $p(x, z)$, perché ogni occorrenza di $z$ è vincolata.

    Prendiamo i termini $0$ e $1$ e sostituiamoli al posto di $x$ e $y$: $p(0, 1)$ diventa
    \[ (0 = 1 \wedge 0 \neq 0) \vee \exists z\,(0 + z = 1) \]
    ed è \textbf{vera} nell'aritmetica: la prima parte è falsa, ma la seconda è vera (basta prendere $z = 1$).
\end{esempio}

% ──────────────────────────────────────────────────────────────
\section{Convenzioni di scrittura}\label{sec:convenzioni}

\subsection{Quantificatori multipli}

Per brevità si usa raggruppare i quantificatori dello stesso tipo:
\[ \forall x_1 \Big(\forall x_2 \big(\cdots \forall x_n (p) \cdots\big)\Big) \quad \text{si scrive} \quad \forall x_1, x_2, \dots, x_n\,(p), \]
e lo stesso vale per $\exists$.

\subsection{Notazione infissa}

\dfn[infissa]{Notazione infissa}{
    Sia $p$ una lettera predicativa binaria e siano $t_1$, $t_2$ termini. Scriveremo
    \[ t_1 \, p \, t_2 \quad \text{invece di} \quad p(t_1, t_2), \]
    proprio come si fa per $\leq$ (scriviamo $x \leq y$ e non $\leq(x, y)$).
}

\subsection{Quantificatori ristretti}\label{sec:ristretti}

\dfn[ristretti]{Quantificatori ristretti}{
    Sia $p$ una lettera predicativa binaria, $x$ una variabile, $y$ un termine e $f$ una FBF. Scriveremo
    \[ (\forall x\, p\, y)(f) \ \defsse\ (\forall x)\big(p(x, y) \rightarrow f\big) \]
    \[ (\exists x\, p\, y)(f) \ \defsse\ (\exists x)\big(p(x, y) \wedge f\big) \]
    Al posto di $y$ può esserci un qualunque termine $t$.
}

A parole: $(\forall x\, p\, y)(f)$ vuol dire «per ogni $x$ che sta nella relazione $p$ con $y$, vale $f$» e
$(\exists x\, p\, y)(f)$ vuol dire «esiste un $x$ che sta nella relazione $p$ con $y$ e per cui vale $f$».

\begin{esempio}{Quantificatori ristretti nell'aritmetica}
    \begin{itemize}
        \item $(\forall x < s(0))(f)$ significa $\forall x\,\big(x < s(0) \rightarrow f\big)$: «per ogni $x$ minore di $1$ vale $f$».
        \item La formula $(\forall x < 0)(x \neq x)$ è \textbf{vera}. Perché? Significa $\forall x\,(x < 0 \rightarrow x \neq x)$
              e in $\mbN$ nessun $x$ è minore di $0$: la premessa $x < 0$ è sempre falsa, quindi l'implicazione è sempre vera
              (Proposizione~\ref{prop:tavole-connettivi}), per ogni $x$. Si dice che è vera «a vuoto»: non ci sono
              controesempi perché non c'è nessun $x$ da controllare.
        \item La formula $(\exists x < 0)(f)$, interpretata nell'aritmetica su $\mbN$, è invece \textbf{falsa} qualunque sia $f$:
              significa $\exists x\,(x < 0 \wedge f)$ e non esiste alcun $x$ con $x < 0$.
    \end{itemize}
\end{esempio}

\warning{{Con $\forall$ si usa $\rightarrow$, con $\exists$ si usa $\wedge$}}{
    È un errore comune scambiarli.
    \begin{itemize}
        \item $(\forall x)\big(p(x,y) \wedge f\big)$ direbbe che \textbf{ogni} $x$ sta nella relazione $p$ con $y$ e soddisfa $f$:
              molto più forte di quello che vogliamo.
        \item $(\exists x)\big(p(x,y) \rightarrow f\big)$ sarebbe vera appena esiste un $x$ che \textbf{non} sta nella
              relazione $p$ con $y$ (premessa falsa, implicazione vera): molto più debole di quello che vogliamo.
    \end{itemize}
    Ad esempio «alcuni psichiatri sono baristi» si scrive $\exists x\,\big(P(x) \wedge B(x)\big)$, non $\exists x\,\big(P(x) \rightarrow B(x)\big)$.
}

\begin{concetto}
\begin{Proposizione}{Negazione dei quantificatori ristretti}{neg-ristretti}
    \begin{prerequisiti}
        \dip{Definizione}{def:ristretti}
        \dip{Teorema}{teo:negazione-quantificatori}
        \dip{Proposizione}{prop:negimp}
        \dip{Proposizione}{prop:demorgan}
        \dip{Proposizione}{prop:impor}
        \dip{Proposizione}{prop:sostituzione}
        \dip{Corollario}{cor:catene}
    \end{prerequisiti}
    Con le notazioni della Definizione~\ref{def:ristretti}, le seguenti formule sono valide:
    \begin{enumerate}
        \item $\neg(\forall x\, p\, y)(f) \leftrightarrow (\exists x\, p\, y)(\neg f)$;
        \item $\neg(\exists x\, p\, y)(f) \leftrightarrow (\forall x\, p\, y)(\neg f)$.
    \end{enumerate}
\end{Proposizione}

\begin{dimostrazione}
    \begin{caso}{Punto 1: $\neg(\forall x\, p\, y)(f) \leftrightarrow (\exists x\, p\, y)(\neg f)$}
    \[
    \begin{aligned}
        \neg(\forall x\, p\, y)(f) &\leftrightarrow \neg(\forall x)\big(p(x,y) \rightarrow f\big) && \text{(definizione)}\\
        &\leftrightarrow (\exists x)\,\neg\big(p(x,y) \rightarrow f\big) && \text{(Teorema~\ref{teo:negazione-quantificatori}, punto 2)}\\
        &\leftrightarrow (\exists x)\big(p(x,y) \wedge \neg f\big) && \text{(negazione dell'implicazione)}\\
        &\leftrightarrow (\exists x\, p\, y)(\neg f) && \text{(definizione)}.
    \end{aligned}
    \]
    \end{caso}
    \begin{caso}{Punto 2: $\neg(\exists x\, p\, y)(f) \leftrightarrow (\forall x\, p\, y)(\neg f)$}
    \[
    \begin{aligned}
        \neg(\exists x\, p\, y)(f) &\leftrightarrow \neg(\exists x)\big(p(x,y) \wedge f\big) && \text{(definizione)}\\
        &\leftrightarrow (\forall x)\,\neg\big(p(x,y) \wedge f\big) && \text{(Teorema~\ref{teo:negazione-quantificatori}, punto 1)}\\
        &\leftrightarrow (\forall x)\big(\neg p(x,y) \vee \neg f\big) && \text{(De Morgan)}\\
        &\leftrightarrow (\forall x)\big(p(x,y) \rightarrow \neg f\big) && \text{(relazione tra $\rightarrow$ e $\vee$)}\\
        &\leftrightarrow (\forall x\, p\, y)(\neg f) && \text{(definizione)}.
    \end{aligned}
    \]
    \end{caso}
\end{dimostrazione}
\end{concetto}

\begin{spiegazione}
    La restrizione dice \textbf{dove} si cerca, e negando non cambia: si nega solo quello che si afferma sugli elementi. Con $x$ tra gli studenti
    di questo corso:
    \begin{passi}
        \item «tutti gli studenti del corso sono promossi»: $(\forall x \text{ del corso})(\text{promosso})$;
        \item la negazione è «c'è uno studente \textbf{del corso} che non è promosso»: $(\exists x \text{ del corso})(\neg\,\text{promosso})$;
        \item non è «c'è uno studente \textbf{non del corso} che \dots»: chi non frequenta il corso non c'entra.
    \end{passi}
    Allo stesso modo $\neg(\forall n \leq 12)(n \text{ è primo})$ è $(\exists n \leq 12)(n \text{ non è primo})$, testimoniata ad esempio da $n = 4$.
\end{spiegazione}


In pratica: per negare un quantificatore ristretto si scambia $\forall$ con $\exists$, \textbf{si lascia invariata la restrizione}
e si nega lo scopo. Ad esempio $\neg(\forall x \leq 12)(f) \leftrightarrow (\exists x \leq 12)(\neg f)$.


  % Parte 2: Teoria degli insiemi (assiomi, operazioni, corrispondenze, funzioni)

  \part{Teoria degli insiemi}

  % Linguaggio, assiomi, operazioni, traduzioni delle tautologie, diagrammi di Venn
  \chapter{Insiemi}\label{cap:insiemi}

In questo capitolo costruiamo la \textbf{teoria degli insiemi} come un linguaggio del primo ordine
(Capitolo~\ref{cap:predicati}): prima in modo informale, per capire perché servono delle regole, poi in modo
formale attraverso gli \textbf{assiomi}. Vedremo poi che le operazioni tra insiemi sono la «traduzione» dei
connettivi logici, e che le tautologie del Capitolo~1 diventano proprietà degli insiemi.

% ──────────────────────────────────────────────────────────────
\section{Il linguaggio degli insiemi}\label{sec:linguaggio-insiemi}

Costruiamo un linguaggio del primo ordine in cui, variabili, accessori, connettivi e quantificatori a parte,
non ci sono costanti né lettere funzionali e c'è \textbf{un solo predicato binario}: «$\in$».

\dfn[linguaggio-insiemi]{Linguaggio della teoria degli insiemi}{
    Il linguaggio della teoria degli insiemi è $(I, \in)$, dove $\in$ è una lettera predicativa binaria.
    Se $a$ è in relazione $\in$ con $b$, cioè $a \in b$, diremo che $a$ è un \textbf{elemento} di $b$, o che
    $a$ \textbf{appartiene} a $b$. Scriveremo $a \notin b$ al posto di $\neg(a \in b)$.
}

\textbf{Informalmente}, gli insiemi saranno descritti da parentesi graffe che «contengono» qualcosa: se $a \in b$
scriveremo $b = \{a, \dots\}$, dove i puntini indicano la presenza eventuale di altri elementi.
Ad esempio, gli unici elementi dell'insieme $\{i, j\}$ sono $i$ e $j$.

Se $f(x)$ è una FBF, scrivere
\[ a = \{x \mid f(x)\} \quad \text{equivale a dire} \quad \forall x\,\big(x \in a \leftrightarrow f(x)\big), \]
il che è un modo per descrivere \textbf{tutti e soli} gli elementi di $a$.

\warning{Il paradosso di Russell}{
    Se restiamo informali giungiamo a paradossi. Supponiamo che $r = \{x \mid \neg(x \in x)\}$ sia un insieme,
    cioè che $\forall x\,\big(x \in r \leftrightarrow \neg(x \in x)\big)$. $r$ appartiene a se stesso?
    \begin{itemize}
        \item Se $r \in r$, allora $\neg(r \in r)$, perché questa è la proprietà che caratterizza gli elementi di $r$.
        \item Se $\neg(r \in r)$, allora $r$ soddisfa la proprietà dei propri elementi, quindi $r \in r$.
    \end{itemize}
    Assurdo! Il problema è aver ammesso una teoria \textbf{senza regole}, in cui troppe cose sono «insiemi».
    Servono quindi delle regole, gli \textbf{assiomi}, che dicano quali insiemi esistono.
}

% ──────────────────────────────────────────────────────────────
\section{Inclusione}\label{sec:inclusione}

Prima degli assiomi diamo due definizioni, che servono per enunciarli.

\dfn[inclusione]{Inclusione e inclusione propria}{
    Siano $x$ e $y$ insiemi. Definiamo
    \[ x \subseteq y \ \defsse\ \forall z\,(z \in x \rightarrow z \in y) \]
    \[ x \subset y \ \defsse\ (x \subseteq y) \wedge (x \neq y). \]
    Se $x \subseteq y$ diremo che $x$ è una \textbf{parte} (o \textbf{sottoinsieme}) di $y$;
    se $x \subset y$ diremo che $x$ è una \textbf{parte propria} di $y$.
}

% ──────────────────────────────────────────────────────────────
\section{Gli assiomi}\label{sec:assiomi}

\ass[estensionalita]{Estensionalità}{
    \[ \forall x, y\,\big(x = y \leftrightarrow \forall z\,(z \in x \leftrightarrow z \in y)\big) \]
    Ovvero: «due insiemi con gli stessi elementi sono lo stesso insieme». Ad esempio $\{a, b\} = \{b, a\}$ e $\{x, x\} = \{x\}$.
}

\ass[vuoto]{Insieme vuoto}{
    \[ \exists x\,\big(\forall y\ \neg(y \in x)\big) \]
    Ovvero: «esiste un insieme senza elementi».
}

\ass[coppia]{Coppia}{
    \[ \forall x, y\ \exists z\ \forall u\,\big(u \in z \leftrightarrow (u = x \vee u = y)\big) \]
    Ovvero: «dati due insiemi $x$ e $y$, esiste un insieme i cui elementi sono esattamente $x$ e $y$».
}

\ass[parti]{Parti}{
    \[ \forall x\ \exists y\ \forall z\,(z \in y \leftrightarrow z \subseteq x) \]
    Ovvero: «dato un insieme $x$, esiste l'insieme i cui elementi sono esattamente le parti di $x$».
}

\ass[unione]{Unione (unaria)}{
    \[ \forall f\ \exists u\ \forall z\,\big(z \in u \leftrightarrow \exists y\,(z \in y \wedge y \in f)\big) \]
    Ovvero: «dato un insieme $f$, esiste l'insieme i cui elementi sono esattamente gli elementi degli elementi di $f$».
}

\ass[separazione]{Separazione (o Comprensione)}{
    Sia $f$ un predicato unario. Allora
    \[ \forall x\ \exists y\ \forall z\,\big(z \in y \leftrightarrow (z \in x \wedge f(z))\big) \]
    Ovvero: «dato un insieme $x$, esiste l'insieme di tutti gli elementi di $x$ che soddisfano $f$».
}

Gli assiomi dicono che certi insiemi \textbf{esistono}; l'Estensionalità dice che sono anche \textbf{unici}.

\begin{concetto}
\begin{Proposizione}{Unicità degli insiemi dati dagli assiomi}{unicita-assiomi}
    \begin{prerequisiti}
        \dip{Assioma}{ass:estensionalita}
        \dip{Corollario}{cor:catene}
    \end{prerequisiti}
    Sia $\varphi(z)$ una FBF. Se $y$ e $y'$ sono insiemi tali che
    $\forall z\,(z \in y \leftrightarrow \varphi(z))$ e $\forall z\,(z \in y' \leftrightarrow \varphi(z))$, allora $y = y'$.
    In particolare, gli insiemi dati dagli assiomi dell'Insieme vuoto, della Coppia, delle Parti, dell'Unione e
    della Separazione sono \textbf{unici}.
\end{Proposizione}
\begin{dimostrazione}
    Per ogni $z$ abbiamo $z \in y \leftrightarrow \varphi(z) \leftrightarrow z \in y'$, quindi (concatenando)
    $\forall z\,(z \in y \leftrightarrow z \in y')$ e, per Estensionalità, $y = y'$.
    Per gli assiomi basta scegliere $\varphi$: per il vuoto $\varphi(z) = \neg(z = z)$ (sempre falsa, così
    «$z \in y \leftrightarrow \varphi(z)$» dice $\neg(z \in y)$), per la coppia $\varphi(z) = (z = x \vee z = y)$, e così via.
\end{dimostrazione}
\end{concetto}

\begin{spiegazione}
    Un insieme è determinato \textbf{solo dai suoi elementi}, non da come lo si descrive. «L'insieme dei numeri primi pari» e «l'insieme delle soluzioni
    di $x - 2 = 0$» sono descritti in modi diversi, ma hanno gli stessi elementi: sono lo stesso insieme $\{2\}$. Per questo possiamo dire \textbf{il}
    vuoto, \textbf{la} coppia, \textbf{l'}insieme delle parti: l'insieme dato da un assioma è uno solo.
\end{spiegazione}


\dfn[notazioni-assiomi]{Notazioni}{
    Grazie all'unicità possiamo dare un nome agli insiemi costruiti dagli assiomi:
    \begin{itemize}
        \item $\varnothing$ è l'\textbf{insieme vuoto};
        \item $\{x, y\}$ è la \textbf{coppia} di $x$ e $y$; se $x = y$ scriviamo $\{x\} = \{x, x\}$, detto \textbf{singleton} di $x$;
        \item $\mcP(x)$ è l'\textbf{insieme delle parti} di $x$;
        \item $\bigcup f$ è l'\textbf{unione} di $f$;
        \item $\{z \in x \mid f(z)\}$ è l'insieme degli elementi di $x$ che soddisfano $f$ (Separazione).
    \end{itemize}
}

\begin{concetto}
\begin{Lemma}{Caratterizzazione dell'insieme vuoto}{vuoto}
    \begin{prerequisiti}
        \dip{Assioma}{ass:estensionalita}
        \dip{Assioma}{ass:vuoto}
        \dip{Definizione}{def:notazioni-assiomi}
    \end{prerequisiti}
    Per ogni insieme $a$: \quad $a = \varnothing \leftrightarrow \forall z\,\neg(z \in a)$.
\end{Lemma}
\begin{dimostrazione}
    \begin{caso}{$(\rightarrow)$ Se $a = \varnothing$, allora $a$ non ha elementi}
        È l'assioma dell'Insieme vuoto: $\varnothing$ non ha elementi.
    \end{caso}
    \begin{caso}{$(\leftarrow)$ Se $a$ non ha elementi, allora $a = \varnothing$}
        Per ogni $z$ le formule $z \in a$ e $z \in \varnothing$ sono entrambe false, quindi
        $\forall z\,(z \in a \leftrightarrow z \in \varnothing)$ e, per Estensionalità, $a = \varnothing$.
    \end{caso}
\end{dimostrazione}
\end{concetto}

\begin{spiegazione}
    Per riconoscere il vuoto basta controllare che non ci siano elementi. Così $\{n \in \mbN \mid n + 1 = 0\}$, $\{x \in \mbN \mid x < 0\}$ e $\{n \in \mbN \mid n \text{ è pari e dispari}\}$
    sono tutti uguali a $\varnothing$: di insiemi vuoti ce n'è \textbf{uno solo}, comunque lo si descriva.
\end{spiegazione}


% ──────────────────────────────────────────────────────────────
\section{La collezione di tutti gli insiemi}\label{sec:russell}

Grazie agli assiomi non è più possibile costruire l'insieme di Russell.

\begin{concetto}
\begin{Teorema}{La collezione di tutti gli insiemi non è un insieme}{russell}
    \begin{prerequisiti}
        \dip{Assioma}{ass:separazione}
        \dip{Proposizione}{prop:terzoescluso}
        \dip{Proposizione}{prop:noncontraddizione}
    \end{prerequisiti}
    Non esiste un insieme $I$ tale che $\forall x\,(x \in I)$.
\end{Teorema}
\begin{dimostrazione}
    Per assurdo, sia $I$ un insieme che contiene tutti gli insiemi. Per l'assioma di \textbf{Separazione} esiste l'insieme
    \[ R = \{x \mid x \in I \wedge \neg(x \in x)\}. \]
    Poiché $R$ è un insieme, $R \in I$. Per il terzo escluso $R \in R$ oppure $R \notin R$: vediamo i due casi.
    \begin{caso}{Caso $R \in R$}
        Gli elementi di $R$ soddisfano $\neg(x \in x)$, quindi $R \notin R$.
    \end{caso}
    \begin{caso}{Caso $R \notin R$}
        Allora $R$ soddisfa $R \in I \wedge \neg(R \in R)$, cioè la proprietà che definisce $R$, quindi $R \in R$.
    \end{caso}
    In entrambi i casi si ottiene $R \in R \wedge R \notin R$, contro il principio di non contraddizione. Assurdo.
\end{dimostrazione}
\end{concetto}

\begin{spiegazione}
    È il \textbf{paradosso del barbiere}: in un villaggio il barbiere rade tutti e soli gli uomini che non si radono da soli. Chi rade il barbiere?
    \begin{passi}
        \item se si rade da solo, allora è uno di quelli che si radono da soli, e il barbiere non lo rade: non si rade;
        \item se non si rade da solo, allora è uno di quelli che il barbiere rade: si rade.
    \end{passi}
    In entrambi i casi si arriva al contrario: un barbiere così non può esistere. Allo stesso modo non può esistere l'insieme $R$ degli insiemi che
    non appartengono a sé stessi, e quindi nemmeno un insieme $I$ di \textbf{tutti} gli insiemi, da cui $R$ si otterrebbe per Separazione.
\end{spiegazione}


% ──────────────────────────────────────────────────────────────
\section{Operazioni tra insiemi}\label{sec:operazioni}

\dfn[unione-binaria]{Unione binaria}{
    Siano $x$ e $y$ insiemi. Definiamo l'\textbf{unione} di $x$ e $y$ come
    \[ x \cup y := \bigcup\{x, y\}. \]
}

\dfn[intersezione]{Intersezione unaria e binaria}{
    Sia $f$ un insieme. Definiamo l'\textbf{intersezione} di $f$ come
    \[ \bigcap f := \{z \in \textstyle\bigcup f \mid \forall y\,(y \in f \rightarrow z \in y)\}. \]
    Siano $x$ e $y$ insiemi. Definiamo l'\textbf{intersezione} di $x$ e $y$ come $x \cap y := \bigcap\{x, y\}$.
}

\dfn[differenza]{Differenza insiemistica}{
    Siano $x$ e $y$ insiemi. Definiamo la \textbf{differenza} tra $x$ e $y$ come
    \[ x \setminus y := \{z \in x \mid \neg(z \in y)\}, \]
    cioè «l'insieme di tutti gli elementi di $x$ che non sono in $y$».
}

\dfn[diff-simmetrica]{Differenza simmetrica}{
    Siano $x$ e $y$ insiemi. Definiamo la \textbf{differenza simmetrica} tra $x$ e $y$ come
    \[ x \,\Delta\, y := \{z \in x \cup y \mid z \in x \XOR z \in y\}. \]
}

Tutte queste definizioni hanno senso perché i relativi insiemi esistono grazie agli assiomi di Coppia, Unione e
Separazione. La proposizione seguente è quella che useremo in tutte le dimostrazioni: dice quando un elemento
appartiene al risultato di un'operazione.

\begin{concetto}
\begin{Proposizione}{Appartenenza alle operazioni}{appartenenza}
    \begin{prerequisiti}
        \dip{Assioma}{ass:coppia}
        \dip{Assioma}{ass:unione}
        \dip{Assioma}{ass:separazione}
        \dip{Definizione}{def:unione-binaria}
        \dip{Definizione}{def:intersezione}
        \dip{Definizione}{def:differenza}
        \dip{Definizione}{def:diff-simmetrica}
    \end{prerequisiti}
    Per ogni insieme $x$, $y$ e per ogni $z$:
    \begin{enumerate}
        \item $z \in x \cup y \leftrightarrow z \in x \vee z \in y$;
        \item $z \in x \cap y \leftrightarrow z \in x \wedge z \in y$;
        \item $z \in x \setminus y \leftrightarrow z \in x \wedge \neg(z \in y)$;
        \item $z \in x \,\Delta\, y \leftrightarrow z \in x \XOR z \in y$.
    \end{enumerate}
\end{Proposizione}
\begin{dimostrazione}
    \begin{caso}{Punto 1: unione}
        Per l'assioma dell'Unione, $z \in \bigcup\{x, y\}$ se e solo se esiste $w$ con $z \in w$ e $w \in \{x, y\}$;
        per l'assioma della Coppia $w \in \{x,y\}$ vuol dire $w = x$ oppure $w = y$. Quindi $z \in x \cup y$ se e solo se
        $z \in x$ oppure $z \in y$.
    \end{caso}
    \begin{caso}{Punto 2: intersezione}
        Per Separazione, $z \in \bigcap\{x, y\}$ se e solo se $z \in x \cup y$ e $z$ appartiene a ogni elemento di
        $\{x, y\}$, cioè $z \in x$ e $z \in y$. La condizione $z \in x \cup y$ è allora automaticamente vera.
    \end{caso}
    \begin{caso}{Punto 3: differenza}
        È la definizione, per Separazione.
    \end{caso}
    \begin{caso}{Punto 4: differenza simmetrica}
        Per Separazione, $z \in x \,\Delta\, y$ se e solo se $z \in x \cup y$ e $z \in x \XOR z \in y$; ma se
        $z \in x \XOR z \in y$ allora $z$ sta in almeno uno dei due, quindi $z \in x \cup y$ è automaticamente vera.
    \end{caso}
\end{dimostrazione}
\end{concetto}

\begin{spiegazione}
    Ogni operazione corrisponde a una regione del diagramma di Venn (Sezione~\ref{sec:venn}) e a un connettivo, con $p$ = «$z \in x$» e $q$ = «$z \in y$»:
    \begin{center}
    \begin{tikzpicture}
        \begin{scope}
            \fill[venn] \vennUno; \fill[venn] \vennDue; \vennduecerchi{x}{y}
            \node[font=\small] at (0.5,-1.15) {$x \cup y$: \ $p \vee q$};
        \end{scope}
        \begin{scope}[xshift=3.7cm]
            \begin{scope} \clip \vennUno; \fill[venn] \vennDue; \end{scope} \vennduecerchi{x}{y}
            \node[font=\small] at (0.5,-1.15) {$x \cap y$: \ $p \wedge q$};
        \end{scope}
        \begin{scope}[xshift=7.4cm]
            \begin{scope}[even odd rule] \clip \vennDue \vennFuori; \fill[venn] \vennUno; \end{scope} \vennduecerchi{x}{y}
            \node[font=\small] at (0.5,-1.15) {$x \setminus y$: \ $p \wedge \neg q$};
        \end{scope}
        \begin{scope}[xshift=11.1cm]
            \fill[venn, even odd rule] \vennUno \vennDue; \vennduecerchi{x}{y}
            \node[font=\small] at (0.5,-1.15) {$x \,\Delta\, y$: \ $p \XOR q$};
        \end{scope}
    \end{tikzpicture}
    \end{center}
\end{spiegazione}


\textbf{Operazioni insiemistiche e connettivi logici.} La proposizione mostra che ogni operazione
\textbf{traduce} un connettivo:

\begin{table}[H]
\centering
\renewcommand{\arraystretch}{1.4}
\begin{tabular}{c|c|l}
\intestazione \textbf{Insiemi} & \textbf{Logica} & \textbf{Perché} \\ \hline
$=$ & $\leftrightarrow$ & Estensionalità: $x = y \leftrightarrow \forall z\,(z \in x \leftrightarrow z \in y)$ \\
$\subseteq$ & $\rightarrow$ & $x \subseteq y \leftrightarrow \forall z\,(z \in x \rightarrow z \in y)$ \\
$\cup$ & $\vee$ & $z \in x \cup y \leftrightarrow z \in x \vee z \in y$ \\
$\cap$ & $\wedge$ & $z \in x \cap y \leftrightarrow z \in x \wedge z \in y$ \\
$\Delta$ & XOR & $z \in x \,\Delta\, y \leftrightarrow z \in x \XOR z \in y$ \\
$\setminus$ & $\neg$ & $z \in x \setminus y \leftrightarrow z \in x \wedge \neg(z \in y)$ \\
\end{tabular}
\caption{Traduzione tra operazioni insiemistiche e connettivi.}
\end{table}

Perché per $\neg$ serve la differenza, e non il «complementare» $\{x \mid \neg(x \in a)\}$? Perché quest'ultimo non è un insieme.

\begin{concetto}
\begin{Corollario}{Il complementare assoluto non è un insieme}{complementare}
    \begin{prerequisiti}
        \dip{Teorema}{teo:russell}
        \dip{Proposizione}{prop:appartenenza}
        \dip{Proposizione}{prop:terzoescluso}
    \end{prerequisiti}
    Sia $a$ un insieme. Non esiste un insieme $c$ tale che $\forall x\,\big(x \in c \leftrightarrow \neg(x \in a)\big)$.
\end{Corollario}
\begin{dimostrazione}
    Se $c$ esistesse, allora per ogni $x$ avremmo $x \in c \cup a \leftrightarrow \neg(x \in a) \vee x \in a$, che è sempre
    vera per il terzo escluso. Quindi $c \cup a$ conterrebbe tutti gli insiemi, contro il Teorema~\ref{teo:russell}.
\end{dimostrazione}
\end{concetto}

\begin{spiegazione}
    «Tutto ciò che non sta in $a$» è una collezione troppo grande: conterrebbe tutti gli insiemi tranne pochi. Il complementare ha senso solo
    \textbf{dentro un insieme ambiente} $s$: è la differenza $s \setminus a$, che esiste per Separazione.
    \begin{center}
    \begin{tikzpicture}
        \fill[venn] (-1.4,-0.9) rectangle (1.9,0.9); \fill[white] (0.25,0) circle (0.6);
        \draw[black!60] (-1.4,-0.9) rectangle (1.9,0.9); \draw[black!60] (0.25,0) circle (0.6);
        \node[font=\small] at (0.25,0) {$a$}; \node[font=\small] at (-1.1,0.62) {$s$};
        \node[font=\small, align=left, anchor=west] at (2.3,0) {in colore $s \setminus a$:\\ il complementare di $a$ \textbf{in} $s$};
    \end{tikzpicture}
    \end{center}
\end{spiegazione}


% ──────────────────────────────────────────────────────────────
\section{Tautologie tradotte in proprietà degli insiemi}\label{sec:traduzioni}

Adesso possiamo usare le tautologie per ottenere formule valide nel linguaggio della teoria degli insiemi.
Il metodo è sempre lo stesso:
\begin{enumerate}
    \item per dimostrare che due insiemi $A$ e $B$ sono uguali, per \textbf{Estensionalità} basta dimostrare
          che $\forall z\,(z \in A \leftrightarrow z \in B)$;
    \item si scrive $z \in A$ con i connettivi usando la Proposizione~\ref{prop:appartenenza};
    \item si applica la tautologia corrispondente, e si torna indietro.
\end{enumerate}

\begin{concetto}
\begin{Proposizione}{Regola della doppia inclusione}{doppia-inclusione}
    \begin{prerequisiti}
        \dip{Assioma}{ass:estensionalita}
        \dip{Definizione}{def:inclusione}
        \dip{Proposizione}{prop:doppiaimp}
        \dip{Proposizione}{prop:forall-and}
        \dip{Proposizione}{prop:sostituzione}
        \dip{Corollario}{cor:catene}
    \end{prerequisiti}
    \[ \forall x, y\,\big(x = y \leftrightarrow (x \subseteq y \wedge y \subseteq x)\big) \]
\end{Proposizione}
\begin{dimostrazione}
    Prendo $x$ e $y$:
    \[
    \begin{aligned}
        x = y &\leftrightarrow \forall z\,(z \in x \leftrightarrow z \in y) && \text{(Estensionalità)}\\
        &\leftrightarrow \forall z\,\big((z \in x \rightarrow z \in y) \wedge (z \in y \rightarrow z \in x)\big) && \text{(doppia implicazione)}\\
        &\leftrightarrow \forall z\,(z \in x \rightarrow z \in y) \wedge \forall z\,(z \in y \rightarrow z \in x) && \text{(Proposizione~\ref{prop:forall-and})}\\
        &\leftrightarrow x \subseteq y \wedge y \subseteq x && \text{(definizione di $\subseteq$)}.
    \end{aligned}
    \]
\end{dimostrazione}
\end{concetto}

\begin{spiegazione}
    Per dimostrare che due insiemi sono uguali si fanno due verifiche separate. Con $x = \{n \in \mbN \mid n \text{ è pari e } n < 6\}$ e $y = \{0, 2, 4\}$:
    \begin{passi}
        \item $x \subseteq y$: un naturale pari minore di $6$ è $0$, $2$ oppure $4$;
        \item $y \subseteq x$: $0$, $2$ e $4$ sono pari e minori di $6$;
        \item quindi $x = y$.
    \end{passi}
    Una sola inclusione non basta: $\{0, 2\} \subseteq \{0, 2, 4\}$, ma i due insiemi sono diversi.
\end{spiegazione}


È il metodo più usato per dimostrare un'uguaglianza tra insiemi: si dimostrano le due inclusioni.

\begin{concetto}
\begin{Proposizione}{Transitività dell'inclusione}{transitivita-inclusione}
    \begin{prerequisiti}
        \dip{Definizione}{def:inclusione}
        \dip{Proposizione}{prop:tavole-connettivi}
    \end{prerequisiti}
    Siano $x$, $y$, $z$ insiemi. Se $x \subseteq y$ e $y \subseteq z$, allora $x \subseteq z$.
\end{Proposizione}
\begin{dimostrazione}
    Sia $w \in x$. Poiché $x \subseteq y$, abbiamo $w \in y$; poiché $y \subseteq z$, abbiamo $w \in z$.
    Quindi $\forall w\,(w \in x \rightarrow w \in z)$, cioè $x \subseteq z$. È la traduzione della transitività di
    $\rightarrow$: la formula $\big((p \rightarrow q) \wedge (q \rightarrow r)\big) \rightarrow (p \rightarrow r)$ è una tautologia.
\end{dimostrazione}
\end{concetto}

\begin{spiegazione}
    Le inclusioni si «concatenano», come scatole una dentro l'altra:
    \begin{center}
    \begin{tikzpicture}[font=\small]
        \draw[black!60, fill=primary!8] (0,0) ellipse (2.6 and 1.15);
        \draw[black!60, fill=primary!18] (-0.45,0) ellipse (1.75 and 0.8);
        \draw[black!60, fill=primary!32] (-0.9,0) ellipse (0.85 and 0.45);
        \node at (-0.9,0) {$x$}; \node at (0.75,0) {$y$}; \node at (1.95,0) {$z$};
        \node[align=left, anchor=west] at (3.1,0) {quadrati $\subseteq$ rettangoli $\subseteq$ parallelogrammi,\\ quindi quadrati $\subseteq$ parallelogrammi.};
    \end{tikzpicture}
    \end{center}
\end{spiegazione}


\begin{concetto}
\begin{Lemma}{Tautologie per l'inclusione}{taut-inclusione}
    \begin{prerequisiti}
        \dip{Proposizione}{prop:tavole-connettivi}
        \dip{Definizione}{def:xor-nand-nor}
        \dip{Proposizione}{prop:criterio-colonne}
        \dip{Corollario}{cor:catene}
    \end{prerequisiti}
    Le seguenti FBF sono tutte logicamente equivalenti tra loro:
    \[ p \rightarrow q \qquad (p \wedge q) \leftrightarrow p \qquad (p \vee q) \leftrightarrow q \qquad \neg(p \wedge \neg q) \qquad (p \XOR q) \leftrightarrow (q \wedge \neg p). \]
\end{Lemma}
\begin{dimostrazione}
    \begin{tavola}{cc|TTTTT}
  \intestazione $p$ & $q$ & $p \rightarrow q$ & $(p \wedge q) \leftrightarrow p$ & $(p \vee q) \leftrightarrow q$ & $\neg(p \wedge \neg q)$ & $(p \XOR q) \leftrightarrow (q \wedge \neg p)$ \\ \hline
  $\vero$ & $\vero$ & $\vero$ & $\vero$ & $\vero$ & $\vero$ & $\vero$ \\
  $\vero$ & $\falso$ & $\falso$ & $\falso$ & $\falso$ & $\falso$ & $\falso$ \\
  $\falso$ & $\vero$ & $\vero$ & $\vero$ & $\vero$ & $\vero$ & $\vero$ \\
  $\falso$ & $\falso$ & $\vero$ & $\vero$ & $\vero$ & $\vero$ & $\vero$ \\
\end{tavola}
    Tutte le colonne evidenziate coincidono con quella di $p \rightarrow q$.
\end{dimostrazione}
\end{concetto}

\begin{spiegazione}
    Tutte e cinque le formule dicono la stessa cosa: «non succede che $p$ sia vera e $q$ falsa». Nella tavola sono false solo nella riga $p = \vero$, $q = \falso$.
    Con $p$ = «$z \in x$» e $q$ = «$z \in y$», quella riga corrisponde agli elementi di $x$ che non stanno in $y$: dire $x \subseteq y$ è dire che non ce ne sono,
    cioè che nel diagramma di Venn la lunetta di $x$ fuori da $y$ è vuota.
\end{spiegazione}


\begin{concetto}
\begin{Proposizione}{Caratterizzazioni dell'inclusione}{caratt-inclusione}
    \begin{prerequisiti}
        \dip{Lemma}{lem:taut-inclusione}
        \dip{Proposizione}{prop:appartenenza}
        \dip{Assioma}{ass:estensionalita}
        \dip{Lemma}{lem:vuoto}
        \dip{Definizione}{def:inclusione}
        \dip{Proposizione}{prop:sostituzione}
    \end{prerequisiti}
    Per ogni coppia di insiemi $x$, $y$ sono equivalenti:
    \[ x \subseteq y \ \leftrightarrow\ x \cap y = x \ \leftrightarrow\ x \cup y = y \ \leftrightarrow\ x \setminus y = \varnothing \ \leftrightarrow\ x \,\Delta\, y = y \setminus x. \]
\end{Proposizione}
\begin{dimostrazione}
    Poniamo $p$ = «$z \in x$» e $q$ = «$z \in y$». Ricordiamo che $x \subseteq y \leftrightarrow \forall z\,(p \rightarrow q)$.
    Confrontiamo ogni condizione con $x \subseteq y$, usando Estensionalità, la Proposizione~\ref{prop:appartenenza}
    e il Lemma~\ref{lem:taut-inclusione} (dentro $\forall z$, per la Proposizione~\ref{prop:sostituzione}).
    \begin{caso}{$x \subseteq y \leftrightarrow x \cap y = x$}
        $x \cap y = x \leftrightarrow \forall z\,\big((p \wedge q) \leftrightarrow p\big) \leftrightarrow \forall z\,(p \rightarrow q) \leftrightarrow x \subseteq y$.
    \end{caso}
    \begin{caso}{$x \subseteq y \leftrightarrow x \cup y = y$}
        $x \cup y = y \leftrightarrow \forall z\,\big((p \vee q) \leftrightarrow q\big) \leftrightarrow \forall z\,(p \rightarrow q) \leftrightarrow x \subseteq y$.
    \end{caso}
    \begin{caso}{$x \subseteq y \leftrightarrow x \setminus y = \varnothing$}
        Per il Lemma~\ref{lem:vuoto},
        $x \setminus y = \varnothing \leftrightarrow \forall z\,\neg(p \wedge \neg q) \leftrightarrow \forall z\,(p \rightarrow q) \leftrightarrow x \subseteq y$.
    \end{caso}
    \begin{caso}{$x \subseteq y \leftrightarrow x \,\Delta\, y = y \setminus x$}
        $x \,\Delta\, y = y \setminus x \leftrightarrow \forall z\,\big((p \XOR q) \leftrightarrow (q \wedge \neg p)\big) \leftrightarrow \forall z\,(p \rightarrow q) \leftrightarrow x \subseteq y$.
    \end{caso}
\end{dimostrazione}
\end{concetto}

\begin{spiegazione}
    Se $x \subseteq y$, nel diagramma di Venn il cerchio di $x$ sta dentro quello di $y$, e le quattro uguaglianze si leggono a colpo d'occhio:
    \begin{center}
    \begin{tikzpicture}[font=\small]
        \foreach \k/\t in {0/{$x \cap y = x$}, 1/{$x \cup y = y$}, 2/{$x \setminus y = \varnothing$}, 3/{$x \,\Delta\, y = y \setminus x$}} {
            \begin{scope}[xshift=\k*3.6cm]
                \ifnum\k=0 \fill[venn] (0.25,0) circle (0.45); \fi
                \ifnum\k=1 \fill[venn] (0.5,0) circle (0.95); \fi
                \ifnum\k=3 \fill[venn, even odd rule] (0.5,0) circle (0.95) (0.25,0) circle (0.45); \fi
                \draw[black!60] (0.5,0) circle (0.95); \draw[black!60] (0.25,0) circle (0.45);
                \node at (0.25,0) {$x$}; \node at (1.1,0.35) {$y$};
                \node at (0.5,-1.25) {\t};
            \end{scope}
        }
    \end{tikzpicture}
    \end{center}
\end{spiegazione}


\begin{concetto}
\begin{Proposizione}{Proprietà di unione, intersezione e differenza simmetrica}{prop-operazioni}
    \begin{prerequisiti}
        \dip{Proposizione}{prop:appartenenza}
        \dip{Assioma}{ass:estensionalita}
        \dip{Proposizione}{prop:commutativita}
        \dip{Proposizione}{prop:associativita}
        \dip{Proposizione}{prop:idempotenza}
        \dip{Proposizione}{prop:xorcommassoc}
    \end{prerequisiti}
    Per ogni insieme $x$, $y$, $z$:
    \begin{enumerate}
        \item \textbf{commutatività}: $x \cup y = y \cup x$, \ $x \cap y = y \cap x$, \ $x \,\Delta\, y = y \,\Delta\, x$;
        \item \textbf{associatività}: $(x \cup y) \cup z = x \cup (y \cup z)$, \ $(x \cap y) \cap z = x \cap (y \cap z)$, \
              $(x \,\Delta\, y) \,\Delta\, z = x \,\Delta\, (y \,\Delta\, z)$;
        \item \textbf{idempotenza}: $x \cup x = x$, \ $x \cap x = x$.
    \end{enumerate}
\end{Proposizione}
\begin{dimostrazione}
    Mostriamo come esempio la commutatività di $\cup$; tutte le altre sono identiche, cambiando la tautologia usata.
    Per ogni $w$:
    \[ w \in x \cup y \leftrightarrow w \in x \vee w \in y \leftrightarrow w \in y \vee w \in x \leftrightarrow w \in y \cup x, \]
    dove il passaggio centrale è la commutatività di $\vee$. Per Estensionalità $x \cup y = y \cup x$.
    Per $\cap$ si usano le proprietà di $\wedge$, per $\Delta$ quelle di XOR (Proposizione~\ref{prop:xorcommassoc});
    per l'idempotenza, $(p \vee p) \leftrightarrow p$ e $(p \wedge p) \leftrightarrow p$.
\end{dimostrazione}
\end{concetto}

\begin{spiegazione}
    Ogni proprietà degli insiemi è la «traduzione» di una tautologia. Verifichiamo l'associatività di $\Delta$ su un esempio, con $x = \{1, 2\}$, $y = \{2, 3\}$, $z = \{3, 4\}$:
    \begin{passi}
        \item $(x \,\Delta\, y) \,\Delta\, z = \{1, 3\} \,\Delta\, \{3, 4\} = \{1, 4\}$;
        \item $x \,\Delta\, (y \,\Delta\, z) = \{1, 2\} \,\Delta\, \{2, 4\} = \{1, 4\}$;
        \item lo stesso risultato: restano gli elementi che stanno in un numero dispari dei tre insiemi ($1$ e $4$ in uno solo, $2$ e $3$ in due).
    \end{passi}
\end{spiegazione}


\begin{concetto}
\begin{Proposizione}{Doppia negazione insiemistica}{doppia-neg-ins}
    \begin{prerequisiti}
        \dip{Proposizione}{prop:appartenenza}
        \dip{Assioma}{ass:estensionalita}
        \dip{Proposizione}{prop:demorgan}
        \dip{Proposizione}{prop:doppianeg}
        \dip{Proposizione}{prop:distributivita}
        \dip{Proposizione}{prop:neutri}
        \dip{Proposizione}{prop:caratt-inclusione}
    \end{prerequisiti}
    Per ogni coppia di insiemi $x$, $y$ vale $x \setminus (x \setminus y) = x \cap y$. In particolare
    \[ \forall x, y\,\big(y \subseteq x \rightarrow x \setminus (x \setminus y) = y\big). \]
\end{Proposizione}
\begin{dimostrazione}
    \begin{caso}{Per ogni $x$, $y$: $x \setminus (x \setminus y) = x \cap y$}
    \[
    \begin{aligned}
        x \setminus (x \setminus y) &= \{z \in x \mid \neg(z \in x \setminus y)\} = \{z \mid z \in x \wedge \neg(z \in x \wedge \neg(z \in y))\} && \text{(def.)}\\
        &= \{z \mid z \in x \wedge (\neg(z \in x) \vee z \in y)\} && \text{(De Morgan + doppia negazione)}\\
        &= \{z \mid (z \in x \wedge \neg(z \in x)) \vee (z \in x \wedge z \in y)\} && \text{(distributività di $\wedge$ rispetto a $\vee$)}\\
        &= \{z \mid z \in x \wedge z \in y\} = x \cap y && \text{($\falso \vee p \leftrightarrow p$)}.
    \end{aligned}
    \]
    \end{caso}
    \begin{caso}{Se $y \subseteq x$: $x \setminus (x \setminus y) = y$}
        Per il caso precedente $x \setminus (x \setminus y) = x \cap y$, e se $y \subseteq x$ allora
        $x \cap y = y \cap x = y$ (Proposizione~\ref{prop:caratt-inclusione}).
    \end{caso}
\end{dimostrazione}
\end{concetto}

\begin{spiegazione}
    Togliendo da $x$ la parte di $x$ fuori da $y$, resta la parte di $x$ dentro $y$:
    \begin{center}
    \begin{tikzpicture}[font=\small]
        \begin{scope}
            \begin{scope}[even odd rule] \clip \vennDue \vennFuori; \fill[venn2] \vennUno; \end{scope} \vennduecerchi{x}{y}
            \node at (0.5,-1.15) {$x \setminus y$};
        \end{scope}
        \node at (2.75,0) {$\longrightarrow$};
        \begin{scope}[xshift=4.5cm]
            \begin{scope} \clip \vennUno; \fill[venn] \vennDue; \end{scope} \vennduecerchi{x}{y}
            \node at (0.5,-1.15) {$x \setminus (x \setminus y) = x \cap y$};
        \end{scope}
        \node[align=left, anchor=west] at (7.1,0) {se poi $y \subseteq x$, la parte di $x$ dentro $y$\\ è tutto $y$: \ $x \setminus (x \setminus y) = y$.};
    \end{tikzpicture}
    \end{center}
    È la doppia negazione: «stare in $x$ e non (stare in $x$ ma non in $y$)» equivale a «stare in $x$ e in $y$».
\end{spiegazione}


\begin{concetto}
\begin{Proposizione}{Terzo escluso insiemistico}{terzo-escluso-ins}
    \begin{prerequisiti}
        \dip{Proposizione}{prop:appartenenza}
        \dip{Assioma}{ass:estensionalita}
        \dip{Proposizione}{prop:distributivita}
        \dip{Proposizione}{prop:terzoescluso}
        \dip{Proposizione}{prop:neutri}
        \dip{Proposizione}{prop:caratt-inclusione}
    \end{prerequisiti}
    Per ogni coppia di insiemi $x$, $y$ vale $y \cup (x \setminus y) = x \cup y$. In particolare
    \[ \forall x, y\,\big(y \subseteq x \rightarrow y \cup (x \setminus y) = x\big). \]
\end{Proposizione}
\begin{dimostrazione}
    \begin{caso}{Per ogni $x$, $y$: $y \cup (x \setminus y) = x \cup y$}
    \[
    \begin{aligned}
        y \cup (x \setminus y) &= \{z \mid z \in y \vee (z \in x \setminus y)\} = \{z \mid z \in y \vee (z \in x \wedge \neg(z \in y))\} && \text{(def.)}\\
        &= \{z \mid (z \in y \vee z \in x) \wedge (z \in y \vee \neg(z \in y))\} && \text{(distributività di $\vee$ rispetto a $\wedge$)}\\
        &= \{z \mid z \in y \vee z \in x\} = x \cup y && \text{($p \wedge \vero \leftrightarrow p$, terzo escluso)}.
    \end{aligned}
    \]
    \end{caso}
    \begin{caso}{Se $y \subseteq x$: $y \cup (x \setminus y) = x$}
        Per il caso precedente $y \cup (x \setminus y) = x \cup y$, e se $y \subseteq x$ allora $x \cup y = y \cup x = x$
        (Proposizione~\ref{prop:caratt-inclusione}).
    \end{caso}
\end{dimostrazione}
\end{concetto}

\begin{spiegazione}
    Ogni elemento di $x$ o sta in $y$ o non ci sta (terzo escluso): rimettendo insieme $y$ e la parte di $x$ fuori da $y$ si ottiene tutto $x \cup y$.
    \begin{center}
    \begin{tikzpicture}[font=\small]
        \fill[venn] \vennDue;
        \begin{scope}[even odd rule] \clip \vennDue \vennFuori; \fill[venn2] \vennUno; \end{scope}
        \vennduecerchi{x}{y}
        \node[align=left, anchor=west] at (2.2,0) {in blu $y$, in arancione $x \setminus y$:\\ insieme formano $x \cup y$.};
    \end{tikzpicture}
    \end{center}
\end{spiegazione}


\begin{concetto}
\begin{Proposizione}{Non contraddizione insiemistica}{non-contr-ins}
    \begin{prerequisiti}
        \dip{Proposizione}{prop:appartenenza}
        \dip{Proposizione}{prop:associativita}
        \dip{Proposizione}{prop:commutativita}
        \dip{Proposizione}{prop:noncontraddizione}
        \dip{Proposizione}{prop:neutri}
        \dip{Lemma}{lem:vuoto}
    \end{prerequisiti}
    Per ogni coppia di insiemi $x$, $y$ vale $y \cap (x \setminus y) = \varnothing$.
\end{Proposizione}
\begin{dimostrazione}
    \[
    \begin{aligned}
        y \cap (x \setminus y) &= \{z \mid z \in y \wedge (z \in x \wedge \neg(z \in y))\} && \text{(def.)}\\
        &= \{z \mid (z \in y \wedge \neg(z \in y)) \wedge z \in x\} && \text{(associatività e commutatività di $\wedge$)}\\
        &= \varnothing && \text{($\falso \wedge p \leftrightarrow \falso$, Lemma~\ref{lem:vuoto})}.
    \end{aligned}
    \]
\end{dimostrazione}
\end{concetto}

\begin{spiegazione}
    $y$ e $x \setminus y$ non hanno elementi in comune: un elemento non può stare in $y$ e, insieme, non starci (non contraddizione). Nel diagramma
    precedente le due regioni colorate, $y$ in blu e $x \setminus y$ in arancione, non si sovrappongono. Insieme alla proposizione precedente dice che
    $y$ e $x \setminus y$ «spezzano» $x \cup y$ in due parti disgiunte.
\end{spiegazione}


\begin{concetto}
\begin{Proposizione}{Leggi di De Morgan insiemistiche}{demorgan-ins}
    \begin{prerequisiti}
        \dip{Proposizione}{prop:appartenenza}
        \dip{Assioma}{ass:estensionalita}
        \dip{Proposizione}{prop:demorgan}
        \dip{Proposizione}{prop:distributivita}
        \dip{Proposizione}{prop:sostituzione}
    \end{prerequisiti}
    Per ogni insieme $x$, $y$, $z$:
    \[ x \setminus (y \cap z) = (x \setminus y) \cup (x \setminus z) \qquad\qquad x \setminus (y \cup z) = (x \setminus y) \cap (x \setminus z). \]
\end{Proposizione}
\begin{dimostrazione}
    Siano $p$ = «$w \in x$», $q$ = «$w \in y$», $r$ = «$w \in z$».
    \begin{caso}{Prima legge: $x \setminus (y \cap z) = (x \setminus y) \cup (x \setminus z)$}
    \[
    \begin{aligned}
        x \setminus (y \cap z) &= \{w \mid p \wedge \neg(q \wedge r)\} = \{w \mid p \wedge (\neg q \vee \neg r)\} && \text{(De Morgan)}\\
        &= \{w \mid (p \wedge \neg q) \vee (p \wedge \neg r)\} && \text{(distributività)}\\
        &= \{w \mid w \in x \setminus y \vee w \in x \setminus z\} = (x \setminus y) \cup (x \setminus z).
    \end{aligned}
    \]
    \end{caso}
    \begin{caso}{Seconda legge: $x \setminus (y \cup z) = (x \setminus y) \cap (x \setminus z)$}
    È identica, usando l'altra legge di De Morgan: $p \wedge \neg(q \vee r) \leftrightarrow p \wedge (\neg q \wedge \neg r)
    \leftrightarrow (p \wedge \neg q) \wedge (p \wedge \neg r)$ (idempotenza, associatività e commutatività di $\wedge$).
    \end{caso}
\end{dimostrazione}
\end{concetto}

\begin{spiegazione}
    Verifichiamo su un esempio, con $x = \{1, 2, 3, 4, 5, 6\}$, $y$ = i pari di $x$ = $\{2, 4, 6\}$ e $z$ = i multipli di $3$ di $x$ = $\{3, 6\}$:
    \begin{passi}
        \item $x \setminus (y \cap z) = x \setminus \{6\} = \{1, 2, 3, 4, 5\}$ e $(x \setminus y) \cup (x \setminus z) = \{1, 3, 5\} \cup \{1, 2, 4, 5\} = \{1, 2, 3, 4, 5\}$;
        \item $x \setminus (y \cup z) = x \setminus \{2, 3, 4, 6\} = \{1, 5\}$ e $(x \setminus y) \cap (x \setminus z) = \{1, 3, 5\} \cap \{1, 2, 4, 5\} = \{1, 5\}$.
    \end{passi}
    A parole: non essere «pari e multiplo di $3$» vuol dire non essere pari \textbf{oppure} non essere multiplo di $3$.
\end{spiegazione}


\begin{concetto}
\begin{Proposizione}{Distributività insiemistiche}{distributivita-ins}
    \begin{prerequisiti}
        \dip{Proposizione}{prop:appartenenza}
        \dip{Assioma}{ass:estensionalita}
        \dip{Proposizione}{prop:distributivita}
    \end{prerequisiti}
    $\cup$ e $\cap$ sono reciprocamente distributive: per ogni $x$, $y$, $z$
    \[ x \cup (y \cap z) = (x \cup y) \cap (x \cup z) \qquad\qquad x \cap (y \cup z) = (x \cap y) \cup (x \cap z). \]
\end{Proposizione}
\begin{dimostrazione}
    Poniamo $p$ = «$w \in x$», $q$ = «$w \in y$», $r$ = «$w \in z$» e traduciamo con la Proposizione~\ref{prop:appartenenza}.
    \begin{caso}{$x \cup (y \cap z) = (x \cup y) \cap (x \cup z)$}
        $w \in x \cup (y \cap z) \leftrightarrow p \vee (q \wedge r)$ e $w \in (x \cup y) \cap (x \cup z) \leftrightarrow (p \vee q) \wedge (p \vee r)$,
        equivalenti per la distributività di $\vee$ rispetto a $\wedge$. Si conclude per Estensionalità.
    \end{caso}
    \begin{caso}{$x \cap (y \cup z) = (x \cap y) \cup (x \cap z)$}
        $w \in x \cap (y \cup z) \leftrightarrow p \wedge (q \vee r)$ e $w \in (x \cap y) \cup (x \cap z) \leftrightarrow (p \wedge q) \vee (p \wedge r)$,
        equivalenti per la distributività di $\wedge$ rispetto a $\vee$. Si conclude per Estensionalità.
    \end{caso}
\end{dimostrazione}
\end{concetto}

\begin{spiegazione}
    Le due distributività nel diagramma di Venn con tre insiemi:
    \begin{center}
    \begin{tikzpicture}[font=\small]
        \begin{scope}
            \begin{scope} \clip \vennA; \fill[venn] \vennB; \fill[venn] \vennC; \end{scope}
            \venntrecerchi{x}{y}{z}
            \node[align=center] at (0,-1.55) {$x \cap (y \cup z) = (x \cap y) \cup (x \cap z)$:\\ la parte di $x$ che sta in $y$ o in $z$};
        \end{scope}
        \begin{scope}[xshift=6.5cm]
            \fill[venn] \vennA; \begin{scope} \clip \vennB; \fill[venn] \vennC; \end{scope}
            \venntrecerchi{x}{y}{z}
            \node[align=center] at (0,-1.55) {$x \cup (y \cap z) = (x \cup y) \cap (x \cup z)$:\\ tutto $x$, più la parte comune a $y$ e $z$};
        \end{scope}
    \end{tikzpicture}
    \end{center}
\end{spiegazione}


\begin{concetto}
\begin{Proposizione}{Distributività di $\cap$ rispetto a $\Delta$}{distr-delta}
    \begin{prerequisiti}
        \dip{Proposizione}{prop:appartenenza}
        \dip{Assioma}{ass:estensionalita}
        \dip{Proposizione}{prop:criterio-colonne}
    \end{prerequisiti}
    Per ogni $x$, $y$, $z$: \quad $x \cap (y \,\Delta\, z) = (x \cap y) \,\Delta\, (x \cap z)$.
\end{Proposizione}
\begin{dimostrazione}
    Traducendo, basta che $p \wedge (q \XOR r)$ e $(p \wedge q) \XOR (p \wedge r)$ siano equivalenti:
    \begin{tavola}{ccc|cTccT}
  \intestazione $p$ & $q$ & $r$ & $q \XOR r$ & $p \wedge (q \XOR r)$ & $p \wedge q$ & $p \wedge r$ & $(p \wedge q) \XOR (p \wedge r)$ \\ \hline
  $\vero$ & $\vero$ & $\vero$ & $\falso$ & $\falso$ & $\vero$ & $\vero$ & $\falso$ \\
  $\vero$ & $\vero$ & $\falso$ & $\vero$ & $\vero$ & $\vero$ & $\falso$ & $\vero$ \\
  $\vero$ & $\falso$ & $\vero$ & $\vero$ & $\vero$ & $\falso$ & $\vero$ & $\vero$ \\
  $\vero$ & $\falso$ & $\falso$ & $\falso$ & $\falso$ & $\falso$ & $\falso$ & $\falso$ \\
  $\falso$ & $\vero$ & $\vero$ & $\falso$ & $\falso$ & $\falso$ & $\falso$ & $\falso$ \\
  $\falso$ & $\vero$ & $\falso$ & $\vero$ & $\falso$ & $\falso$ & $\falso$ & $\falso$ \\
  $\falso$ & $\falso$ & $\vero$ & $\vero$ & $\falso$ & $\falso$ & $\falso$ & $\falso$ \\
  $\falso$ & $\falso$ & $\falso$ & $\falso$ & $\falso$ & $\falso$ & $\falso$ & $\falso$ \\
\end{tavola}
    Le colonne evidenziate coincidono; si conclude per Estensionalità.
\end{dimostrazione}
\end{concetto}

\begin{spiegazione}
    Su un esempio, con $x = \{1, 2, 3\}$, $y = \{2, 3, 4\}$ e $z = \{3, 5\}$:
    \begin{passi}
        \item $y \,\Delta\, z = \{2, 4, 5\}$, quindi $x \cap (y \,\Delta\, z) = \{2\}$;
        \item $x \cap y = \{2, 3\}$ e $x \cap z = \{3\}$, quindi $(x \cap y) \,\Delta\, (x \cap z) = \{2\}$.
    \end{passi}
    Invece $\cup$ \textbf{non} è distributiva rispetto a $\Delta$: lo mostra l'esempio con i diagrammi di Venn alla fine del capitolo. Questa proprietà renderà
    $(\mcP(s), \Delta, \cap)$ un anello (Capitolo~\ref{cap:anelli}).
\end{spiegazione}


\begin{concetto}
\begin{Proposizione}{Esplicitazioni della differenza simmetrica}{esplicitazioni-delta}
    \begin{prerequisiti}
        \dip{Proposizione}{prop:appartenenza}
        \dip{Assioma}{ass:estensionalita}
        \dip{Proposizione}{prop:xorespl}
    \end{prerequisiti}
    Per ogni $x$, $y$: \quad $x \,\Delta\, y = (x \setminus y) \cup (y \setminus x) = (x \cup y) \setminus (y \cap x)$.
\end{Proposizione}
\begin{dimostrazione}
    Traducendo, i tre insiemi corrispondono a $p \XOR q$, \ $(p \wedge \neg q) \vee (q \wedge \neg p)$ \ e
    $(p \vee q) \wedge \neg(q \wedge p)$, che sono equivalenti per le esplicitazioni di XOR
    (e la commutatività di $\wedge$). Si conclude per Estensionalità.
\end{dimostrazione}
\end{concetto}

\begin{spiegazione}
    Le due espressioni descrivono la stessa regione in due modi: «le due lunette» oppure «tutto, tranne la parte comune».
    \begin{center}
    \begin{tikzpicture}[font=\small]
        \begin{scope}
            \begin{scope}[even odd rule] \clip \vennDue \vennFuori; \fill[venn] \vennUno; \end{scope}
            \begin{scope}[even odd rule] \clip \vennUno \vennFuori; \fill[venn2] \vennDue; \end{scope}
            \vennduecerchi{x}{y}
            \node at (0.5,-1.15) {$(x \setminus y) \cup (y \setminus x)$};
        \end{scope}
        \node at (3.0,0) {$=$};
        \begin{scope}[xshift=4.8cm]
            \fill[venn, even odd rule] \vennUno \vennDue; \vennduecerchi{x}{y}
            \node at (0.5,-1.15) {$(x \cup y) \setminus (x \cap y)$};
        \end{scope}
    \end{tikzpicture}
    \end{center}
\end{spiegazione}


% ──────────────────────────────────────────────────────────────
\section{Diagrammi di Venn}\label{sec:venn}

I \textbf{diagrammi di Venn} rappresentano tutte le possibili intersezioni tra un numero finito di insiemi,
disegnati come regioni di un rettangolo $U$. Ecco i diagrammi generali per 2, 3 e 4 insiemi.

\begin{concetto}
\begin{center}
\begin{tikzpicture}[scale=0.8]
    \draw (-2.2,-1.6) rectangle (2.2,1.6); \node at (-1.9,1.3) {$U$};
    \draw (-0.6,0) circle (1.1); \draw (0.6,0) circle (1.1);
    \node at (-1.1,0) {$a$}; \node at (1.1,0) {$b$};
\end{tikzpicture}
\hspace{0.8cm}
\begin{tikzpicture}[scale=0.8]
    \draw (-2.2,-2.0) rectangle (2.2,1.9); \node at (-1.9,1.6) {$U$};
    \draw (0,0.55) circle (1.1); \draw (-0.6,-0.45) circle (1.1); \draw (0.6,-0.45) circle (1.1);
    \node at (0,1.15) {$a$}; \node at (-1.15,-0.8) {$b$}; \node at (1.15,-0.8) {$c$};
\end{tikzpicture}
\hspace{0.8cm}
\begin{tikzpicture}[scale=0.8]
    \draw (-2.9,-2.3) rectangle (2.9,2.3); \node at (-2.6,2.0) {$U$};
    \draw (-0.55,-0.35) ellipse [x radius=2.0, y radius=0.95, rotate=-50];
    \draw (0.55,-0.35) ellipse [x radius=2.0, y radius=0.95, rotate=50];
    \draw (-0.05,0.35) ellipse [x radius=1.9, y radius=0.85, rotate=-50];
    \draw (0.05,0.35) ellipse [x radius=1.9, y radius=0.85, rotate=50];
    \node at (-2.1,1.3) {$a$}; \node at (-0.75,1.95) {$b$}; \node at (0.75,1.95) {$c$}; \node at (2.1,1.3) {$d$};
\end{tikzpicture}
\end{center}
\end{concetto}

Possono servire a visualizzare le operazioni, ad esempio l'intersezione e la differenza simmetrica:

\begin{concetto}
\begin{center}
\begin{tikzpicture}[scale=0.8]
    \draw (-2.2,-1.6) rectangle (2.2,1.6); \node at (-1.9,1.3) {$U$};
    \begin{scope} \clip (-0.6,0) circle (1.1); \fill[primary!35] (0.6,0) circle (1.1); \end{scope}
    \draw (-0.6,0) circle (1.1); \draw (0.6,0) circle (1.1);
    \node at (-1.1,0) {$a$}; \node at (1.1,0) {$b$}; \node at (0,-1.35) {$a \cap b$};
\end{tikzpicture}
\hspace{1cm}
\begin{tikzpicture}[scale=0.8]
    \draw (-2.2,-1.6) rectangle (2.2,1.6); \node at (-1.9,1.3) {$U$};
    \fill[primary!35,even odd rule] (-0.6,0) circle (1.1) (0.6,0) circle (1.1);
    \draw (-0.6,0) circle (1.1); \draw (0.6,0) circle (1.1);
    \node at (-1.1,0) {$a$}; \node at (1.1,0) {$b$}; \node at (0,-1.35) {$a \,\Delta\, b$};
\end{tikzpicture}
\end{center}
\end{concetto}

L'utilizzo più vantaggioso dei diagrammi di Venn, però, è quello di trovare \textbf{controesempi} ad alcune
relazioni insiemistiche: se i diagrammi dei due membri sono diversi, si scelgono insiemi che differiscano proprio
nelle aree in cui i diagrammi differiscono.

\begin{esempio}{$a \cup (b \,\Delta\, c) = (a \cup b) \,\Delta\, (a \cup c)$ non vale sempre}
    Disegniamo i due membri:
    \begin{center}
    \begin{tikzpicture}[scale=0.75]
        \draw (-2.2,-2.1) rectangle (2.2,1.9);
        \fill[primary!35] (0,0.55) circle (1.1);
        \fill[primary!35,even odd rule] (-0.6,-0.45) circle (1.1) (0.6,-0.45) circle (1.1);
        \draw (0,0.55) circle (1.1); \draw (-0.6,-0.45) circle (1.1); \draw (0.6,-0.45) circle (1.1);
        \node at (0,1.15) {$a$}; \node at (-1.15,-0.8) {$b$}; \node at (1.15,-0.8) {$c$};
        \node at (0,-1.85) {$a \cup (b \,\Delta\, c)$};
    \end{tikzpicture}
    \hspace{1cm}
    \begin{tikzpicture}[scale=0.75]
        \draw (-2.2,-2.1) rectangle (2.2,1.9);
        % (a ∪ b) Δ (a ∪ c) = (b Δ c) \ a
        \begin{scope}[even odd rule]
            \clip (-2.2,-2.1) rectangle (2.2,1.9) (0,0.55) circle (1.1);
            \fill[primary!35,even odd rule] (-0.6,-0.45) circle (1.1) (0.6,-0.45) circle (1.1);
        \end{scope}
        \draw (0,0.55) circle (1.1); \draw (-0.6,-0.45) circle (1.1); \draw (0.6,-0.45) circle (1.1);
        \node at (0,1.15) {$a$}; \node at (-1.15,-0.8) {$b$}; \node at (1.15,-0.8) {$c$};
        \node at (0,-1.85) {$(a \cup b) \,\Delta\, (a \cup c)$};
    \end{tikzpicture}
    \end{center}
    Sono diversi: nel secondo manca tutta la regione di $a$. Basta allora prendere $a$ con almeno un elemento
    che non sta né in $b$ né in $c$. La scelta più facile è $a = \{x\}$, $b = c = \varnothing$: allora
    \[ a \cup (b \,\Delta\, c) = a \neq \varnothing = (a \cup b) \,\Delta\, (a \cup c), \]
    e ciò dimostra che l'uguaglianza non vale sempre.
\end{esempio}

\begin{esempio}{La differenza non è associativa}
    Confrontiamo $a \setminus (b \setminus c)$ e $(a \setminus b) \setminus c$:
    \begin{center}
    \begin{tikzpicture}[scale=0.75]
        \draw (-2.2,-2.1) rectangle (2.2,1.9);
        % a \ (b \ c) = (a \ b) ∪ (a ∩ c)
        \begin{scope}
            \clip (0,0.55) circle (1.1);
            \fill[primary!35] (-2.2,-2.1) rectangle (2.2,1.9);
            \begin{scope}[even odd rule]
                \clip (-2.2,-2.1) rectangle (2.2,1.9) (0.6,-0.45) circle (1.1);
                \fill[white] (-0.6,-0.45) circle (1.1);
            \end{scope}
        \end{scope}
        \draw (0,0.55) circle (1.1); \draw (-0.6,-0.45) circle (1.1); \draw (0.6,-0.45) circle (1.1);
        \node at (0,1.15) {$a$}; \node at (-1.15,-0.8) {$b$}; \node at (1.15,-0.8) {$c$};
        \node at (0,-1.85) {$a \setminus (b \setminus c)$};
    \end{tikzpicture}
    \hspace{1cm}
    \begin{tikzpicture}[scale=0.75]
        \draw (-2.2,-2.1) rectangle (2.2,1.9);
        % (a \ b) \ c
        \begin{scope}
            \clip (0,0.55) circle (1.1);
            \fill[primary!35] (-2.2,-2.1) rectangle (2.2,1.9);
            \fill[white] (-0.6,-0.45) circle (1.1);
            \fill[white] (0.6,-0.45) circle (1.1);
        \end{scope}
        \draw (0,0.55) circle (1.1); \draw (-0.6,-0.45) circle (1.1); \draw (0.6,-0.45) circle (1.1);
        \node at (0,1.15) {$a$}; \node at (-1.15,-0.8) {$b$}; \node at (1.15,-0.8) {$c$};
        \node at (0,-1.85) {$(a \setminus b) \setminus c$};
    \end{tikzpicture}
    \end{center}
    Differiscono nella regione degli elementi che stanno in $a$ e in $c$ ma non in $b$. Prendiamo quindi
    $a = c = \{x\}$ e $b = \varnothing$: allora $a \setminus (b \setminus c) = \{x\} \setminus \varnothing = \{x\}$, mentre
    $(a \setminus b) \setminus c = \{x\} \setminus \{x\} = \varnothing$.
\end{esempio}

\begin{esempio}{L'associatività di $\Delta$ vista con i diagrammi}
    Sia $(a \,\Delta\, b) \,\Delta\, c$ sia $a \,\Delta\, (b \,\Delta\, c)$ contengono esattamente gli elementi che stanno in
    \textbf{uno solo} oppure in \textbf{tutti e tre} gli insiemi: il diagramma è lo stesso, come dimostrato nella
    Proposizione~\ref{prop:prop-operazioni}.
    \begin{center}
    \begin{tikzpicture}[scale=0.75]
        \draw (-2.2,-2.1) rectangle (2.2,1.9);
        \fill[primary!35,even odd rule] (0,0.55) circle (1.1) (-0.6,-0.45) circle (1.1) (0.6,-0.45) circle (1.1);
        \draw (0,0.55) circle (1.1); \draw (-0.6,-0.45) circle (1.1); \draw (0.6,-0.45) circle (1.1);
        \node at (0,1.15) {$a$}; \node at (-1.15,-0.8) {$b$}; \node at (1.15,-0.8) {$c$};
        \node at (0,-1.85) {$a \,\Delta\, b \,\Delta\, c$};
    \end{tikzpicture}
    \end{center}
\end{esempio}


  % Coppie ordinate, prodotto cartesiano, corrispondenze
  \chapter{Coppie ordinate e corrispondenze}\label{cap:corrispondenze}

Gli insiemi non hanno un ordine: $\{x, y\} = \{y, x\}$. Per parlare di funzioni e relazioni ci servono invece
coppie in cui conta \textbf{chi viene prima}. In questo capitolo costruiamo le coppie ordinate usando solo gli
assiomi, poi il prodotto cartesiano e infine le corrispondenze tra insiemi.

% ──────────────────────────────────────────────────────────────
\section{Coppie ordinate}\label{sec:coppie}

\dfn[coppia-ordinata]{Coppia ordinata}{
    Siano $x$ e $y$ insiemi. L'insieme $\{\{x\}, \{x, y\}\}$ si chiama \textbf{coppia ordinata} di
    \textbf{prima coordinata} $x$ e \textbf{seconda coordinata} $y$, e si scrive $(x, y)$.
}

L'insieme esiste per l'assioma della Coppia, applicato prima a $x, x$ e $x, y$ e poi a $\{x\}, \{x, y\}$.
La proprietà che ci interessa è che due coppie ordinate sono uguali solo se hanno le stesse coordinate
\textbf{nello stesso ordine}.

\begin{concetto}
\begin{Teorema}{Proprietà caratteristica della coppia ordinata}{coppia-ordinata}
    \begin{prerequisiti}
        \dip{Definizione}{def:coppia-ordinata}
        \dip{Assioma}{ass:estensionalita}
        \dip{Assioma}{ass:coppia}
        \dip{Definizione}{def:notazioni-assiomi}
        \dip{Proposizione}{prop:doppiaimp}
    \end{prerequisiti}
    \[ \forall x, y, w, z\,\big((x, y) = (w, z) \leftrightarrow (x = w \wedge y = z)\big) \]
\end{Teorema}
\begin{dimostrazione}
    \begin{caso}{$(\leftarrow)$ Se $x = w$ e $y = z$, allora $(x, y) = (w, z)$}
        Ovvio: sostituendo $w$ con $x$ e $z$ con $y$ si ottiene lo stesso insieme $\{\{x\}, \{x, y\}\}$.
    \end{caso}
    \begin{caso}{$(\rightarrow)$ Se $(x, y) = (w, z)$, allora $x = w$ e $y = z$}
        Per ipotesi $\{\{x\}, \{x, y\}\} = \{\{w\}, \{w, z\}\}$. Distinguiamo due casi.
        \begin{caso}{Caso $x = y$}
            Allora $\{\{x\}, \{x, x\}\} = \{\{x\}\}$ ha un solo elemento, quindi anche $\{\{w\}, \{w, z\}\}$ ne ha uno solo:
            $\{w\} = \{w, z\} = \{x\}$. Quindi $w = x$ e $z = w$, cioè $x = y = w = z$.
        \end{caso}
        \begin{caso}{Caso $x \neq y$}
            Allora $\{x\} \neq \{x, y\}$ e $\{x, y\}$ ha due elementi. Per Estensionalità $\{x, y\}$ deve essere uguale a
            $\{w\}$ oppure a $\{w, z\}$; non può essere $\{w\}$ (altrimenti $x = w = y$), quindi $\{x, y\} = \{w, z\}$.
            Allora l'altro elemento, $\{x\}$, deve essere uguale a $\{w\}$, cioè $x = w$. Ma allora $\{x, y\} = \{x, z\}$
            e $y$, non potendo essere uguale a $x$, deve essere uguale a $z$.
        \end{caso}
        In entrambi i casi $x = w$ e $y = z$.
    \end{caso}
\end{dimostrazione}
\end{concetto}

\begin{spiegazione}
    Il trucco della definizione è il \textbf{singleton}, che «segna» la prima coordinata. Confrontiamo $\{1, 2\}$ con le coppie $(1, 2)$ e $(2, 1)$:
    \begin{passi}
        \item $\{1, 2\} = \{2, 1\}$: in un insieme l'ordine non conta;
        \item $(1, 2) = \{\{1\}, \{1, 2\}\}$ e $(2, 1) = \{\{2\}, \{1, 2\}\}$: il secondo elemento è lo stesso, ma i singleton $\{1\}$ e $\{2\}$ sono diversi,
              quindi $(1, 2) \neq (2, 1)$;
        \item se le coordinate sono uguali la coppia «collassa»: $(1, 1) = \{\{1\}, \{1, 1\}\} = \{\{1\}\}$, ed è questo il caso $x = y$ della dimostrazione.
    \end{passi}
\end{spiegazione}


\dfn[terna-ordinata]{Terna ordinata}{
    Siano $x$, $y$, $z$ insiemi. La \textbf{terna ordinata} di coordinate $x$, $y$, $z$ è
    \[ (x, y, z) := \big((x, y), z\big). \]
}

\begin{concetto}
\begin{Corollario}{Proprietà caratteristica della terna ordinata}{terna}
    \begin{prerequisiti}
        \dip{Teorema}{teo:coppia-ordinata}
        \dip{Definizione}{def:terna-ordinata}
    \end{prerequisiti}
    $(x, y, z) = (x', y', z') \leftrightarrow (x = x' \wedge y = y' \wedge z = z')$.
\end{Corollario}
\begin{dimostrazione}
    Per il Teorema~\ref{teo:coppia-ordinata} applicato due volte:
    $\big((x,y),z\big) = \big((x',y'),z'\big) \leftrightarrow (x,y) = (x',y') \wedge z = z' \leftrightarrow x = x' \wedge y = y' \wedge z = z'$.
\end{dimostrazione}
\end{concetto}

\begin{spiegazione}
    Una terna è una coppia «a scatole cinesi»: la prima coordinata è a sua volta una coppia. Per confrontare $\big((1, 2), 3\big)$ e $\big((1, 2), 4\big)$
    si confrontano prima le scatole esterne (la coppia interna e l'ultima coordinata) e poi quella interna: qui le coppie interne coincidono, ma $3 \neq 4$,
    quindi le terne sono diverse. Allo stesso modo si costruiscono quadruple, quintuple e così via.
\end{spiegazione}


% ──────────────────────────────────────────────────────────────
\section{Prodotto cartesiano}\label{sec:prodotto-cartesiano}

Vogliamo l'insieme di \textbf{tutte} le coppie ordinate con prima coordinata in $a$ e seconda coordinata in $b$.
Per costruirlo con la Separazione ci serve un insieme che le contenga già tutte.

\begin{concetto}
\begin{Lemma}{Dove vivono le coppie ordinate}{coppia-parti}
    \begin{prerequisiti}
        \dip{Definizione}{def:coppia-ordinata}
        \dip{Assioma}{ass:parti}
        \dip{Definizione}{def:inclusione}
        \dip{Proposizione}{prop:appartenenza}
    \end{prerequisiti}
    Siano $a$, $b$ insiemi, $x \in a$ e $y \in b$. Allora $(x, y) \in \mcP(\mcP(a \cup b))$.
\end{Lemma}
\begin{dimostrazione}
    \begin{caso}{Passo 1: $\{x\}$ e $\{x, y\}$ stanno in $\mcP(a \cup b)$}
        Dato che $x \in a \cup b$, il singleton $\{x\}$ è un sottoinsieme di $a \cup b$, quindi $\{x\} \in \mcP(a \cup b)$.
        Allo stesso modo, poiché $x, y \in a \cup b$, anche $\{x, y\} \in \mcP(a \cup b)$.
    \end{caso}
    \begin{caso}{Passo 2: $(x, y)$ sta in $\mcP(\mcP(a \cup b))$}
        Poiché sia $\{x\}$ sia $\{x, y\}$ appartengono a $\mcP(a \cup b)$, l'insieme che li racchiude entrambi,
        $\{\{x\}, \{x, y\}\} = (x, y)$, è un sottoinsieme di $\mcP(a \cup b)$. Quindi $(x, y) \in \mcP(\mcP(a \cup b))$.
    \end{caso}
\end{dimostrazione}
\end{concetto}

\begin{spiegazione}
    Seguiamo il lemma con $a = \{1\}$ e $b = \{2\}$:
    \begin{passi}
        \item $a \cup b = \{1, 2\}$ e $\mcP(a \cup b) = \{\varnothing, \{1\}, \{2\}, \{1, 2\}\}$;
        \item $(1, 2) = \{\{1\}, \{1, 2\}\}$ è formata da due elementi di $\mcP(a \cup b)$, quindi è una sua parte;
        \item dunque $(1, 2) \in \mcP(\mcP(a \cup b))$: le coppie si trovano «due piani sopra» gli elementi di $a$ e $b$.
    \end{passi}
\end{spiegazione}


\dfn[prodotto-cartesiano]{Prodotto cartesiano}{
    Siano $a$ e $b$ insiemi. Il \textbf{prodotto cartesiano} di $a$ e $b$ è
    \[ a \times b := \{z \in \mcP(\mcP(a \cup b)) \mid \exists x, y\,(x \in a \wedge y \in b \wedge z = (x, y))\}. \]
    Per il Lemma~\ref{lem:coppia-parti}, $a \times b$ contiene \textbf{tutte} le coppie $(x, y)$ con $x \in a$ e $y \in b$.
}

\info{Perché le parti delle parti?}{
    Facciamo le parti delle parti di $a \cup b$ perché, partendo dalla definizione di coppia
    $\{\{x\}, \{x, y\}\}$, sappiamo solo che $x \in a$ e $y \in b$: di conseguenza sia $x$ sia $y$ appartengono
    all'insieme unione $a \cup b$. È il più piccolo insieme «già costruito» in cui cercare le coppie con la Separazione.
}

% ──────────────────────────────────────────────────────────────
\section{Corrispondenze}\label{sec:corrispondenze}

\dfn[corrispondenza]{Corrispondenza}{
    Siano $a$, $b$, $g$ insiemi con $g \subseteq a \times b$. Allora la terna $\rho = (a, b, g)$ si dice
    \textbf{corrispondenza} tra $a$ e $b$, e $g$ si dice \textbf{grafico} della corrispondenza.
    L'insieme delle corrispondenze tra $a$ e $b$ lo chiamiamo $\mathrm{Corr}(a, b)$.

    Se $\rho = (a, b, g) \in \mathrm{Corr}(a, b)$, $x \in a$, $y \in b$ e $(x, y) \in g$, scriveremo $x \rho y$.
}

\dfn[relazione-binaria]{Relazione binaria}{
    Una corrispondenza $(a, a, g)$ di un insieme con se stesso si dice \textbf{relazione binaria} su $a$.
    Definiamo $\mathrm{Rel}(a) := \mathrm{Corr}(a, a)$ e la \textbf{diagonale} $\mathrm{diag}(a) := \{(x, x) \mid x \in a\}$.
    \begin{itemize}
        \item Se $g = a \times a$, la relazione si dice \textbf{totale}.
        \item Se $g = \mathrm{diag}(a)$, la relazione si dice \textbf{identica}.
    \end{itemize}
}

\dfn[prodotto-corrispondenze]{Prodotto tra corrispondenze}{
    Siano $\rho = (a, b, g_1)$ e $\sigma = (c, d, g_2)$ corrispondenze. Definisco la corrispondenza
    \[ \sigma\rho = (a, d, g_3) \]
    e la chiamo \textbf{prodotto tra $\rho$ e $\sigma$}, dove
    \[ \forall x, y\,\Big((x, y) \in g_3 \leftrightarrow \exists z\,\big((x, z) \in g_1 \wedge (z, y) \in g_2\big)\Big). \]
}

Attenzione all'ordine: in $\sigma\rho$ si applica \textbf{prima} $\rho$ e \textbf{poi} $\sigma$. Le coppie di $g_3$ sono quelle
«collegate attraverso elementi in comune tra $b$ e $c$».

\begin{esempio}{Prodotto tra corrispondenze}
    Sia $\rho$ da $a = \{x_1, x_2\}$ a $b = \{z_1, z_2\}$ e $\sigma$ da $c = b$ a $d = \{y_1, y_2\}$:
    \begin{center}
    \begin{tikzpicture}[>=Stealth, every node/.style={font=\small}]
        \node (x1) at (0,0.6) {$x_1$}; \node (x2) at (0,-0.6) {$x_2$};
        \node (z1) at (3,0.6) {$z_1$}; \node (z2) at (3,-0.6) {$z_2$};
        \node (y1) at (6,0.6) {$y_1$}; \node (y2) at (6,-0.6) {$y_2$};
        \draw[rounded corners] (-0.5,-1.1) rectangle (0.5,1.1); \node at (0,1.4) {$a$};
        \draw[rounded corners] (2.5,-1.1) rectangle (3.5,1.1); \node at (3,1.4) {$b = c$};
        \draw[rounded corners] (5.5,-1.1) rectangle (6.5,1.1); \node at (6,1.4) {$d$};
        \draw[->,primary] (x1) -- (z1); \draw[->,primary] (x1) -- (z2);
        \draw[->,cherryred] (z1) -- (y2);
        \node[primary] at (1.5,-1.4) {$\rho$}; \node[cherryred] at (4.5,-1.4) {$\sigma$};
    \end{tikzpicture}
    \end{center}
    Qui $x_1 \rho z_1$ e $z_1 \sigma y_2$, quindi $x_1 \,\sigma\rho\, y_2$. Nessun'altra coppia è collegata:
    il grafico di $\sigma\rho$ è $\{(x_1, y_2)\}$.
\end{esempio}

\begin{esempio}{Altri prodotti}
    \begin{itemize}
        \item Siano $\rho$ e $\sigma$ le relazioni binarie tra gli esseri umani «essere padre di» ed «essere fratello di».
              Allora $x \,\rho\sigma\, y$ se e solo se «$x$ è padre del fratello di $y$», $x \,\sigma\rho\, y$ se e solo se
              «$x$ è fratello del padre (cioè è uno zio) di $y$», mentre $x \rho^2 y$ (con $\rho^2 = \rho\rho$) se e solo se
              «$x$ è il nonno paterno di $y$».
        \item Consideriamo l'inclusione su $\mcP(a)$. Allora $\subseteq^2 = \subseteq$, grazie alla transitività dell'inclusione
              (Proposizione~\ref{prop:transitivita-inclusione}).
    \end{itemize}
\end{esempio}

\begin{concetto}
\begin{Proposizione}{Associatività del prodotto tra corrispondenze}{assoc-prodotto}
    \begin{prerequisiti}
        \dip{Definizione}{def:prodotto-corrispondenze}
        \dip{Corollario}{cor:terna}
        \dip{Assioma}{ass:estensionalita}
        \dip{Proposizione}{prop:associativita}
    \end{prerequisiti}
    Il prodotto di corrispondenze è \textbf{associativo}: per ogni terna di corrispondenze $\rho$, $\sigma$, $\tau$
    \[ \tau(\sigma\rho) = (\tau\sigma)\rho. \]
\end{Proposizione}
\begin{dimostrazione}
    Siano $\rho = (a, b, g_1)$, $\sigma = (c, d, g_2)$, $\tau = (e, h, g_3)$. Entrambi i prodotti hanno $a$ come primo
    insieme e $h$ come secondo, quindi (Corollario~\ref{cor:terna}) basta mostrare che i grafici coincidono.
    Per definizione di prodotto, applicata due volte:
    \[
    \begin{aligned}
        (x, y) \in \text{grafico di } \tau(\sigma\rho) &\leftrightarrow \exists w\,\Big(\exists z\,\big((x, z) \in g_1 \wedge (z, w) \in g_2\big) \wedge (w, y) \in g_3\Big)\\
        &\leftrightarrow \exists z\,\exists w\,\big((x, z) \in g_1 \wedge (z, w) \in g_2 \wedge (w, y) \in g_3\big),\\
        (x, y) \in \text{grafico di } (\tau\sigma)\rho &\leftrightarrow \exists z\,\Big((x, z) \in g_1 \wedge \exists w\,\big((z, w) \in g_2 \wedge (w, y) \in g_3\big)\Big)\\
        &\leftrightarrow \exists z\,\exists w\,\big((x, z) \in g_1 \wedge (z, w) \in g_2 \wedge (w, y) \in g_3\big).
    \end{aligned}
    \]
    Nei passaggi abbiamo portato fuori il quantificatore $\exists$ (la variabile quantificata non compare nell'altro
    membro della congiunzione) e usato l'associatività di $\wedge$. Le due condizioni coincidono, quindi per
    Estensionalità i grafici sono uguali.
\end{dimostrazione}
\end{concetto}

\begin{spiegazione}
    In entrambi i casi $x$ è collegato a $y$ se c'è un \textbf{cammino in tre passi} $x \to z \to w \to y$, fatto con un passo di $\rho$, uno di $\sigma$ e uno di
    $\tau$; raggruppare i primi due passi o gli ultimi due non cambia il cammino:
    \begin{center}
    \begin{tikzpicture}[font=\small]
        \foreach \x/\n in {0/x, 2.4/z, 4.8/w, 7.2/y} \node[circle, fill=primary!75!black, inner sep=1.6pt, label=below:{$\n$}] (\n) at (\x,0) {};
        \draw[fun] (x) -- node[above] {$\rho$} (z); \draw[fun2] (z) -- node[above] {$\sigma$} (w); \draw[fun3] (w) -- node[above] {$\tau$} (y);
        \draw[decorate, decoration={brace, amplitude=4pt, mirror}, black!60] (0,-0.55) -- node[below=3pt] {$\sigma\rho$} (4.8,-0.55);
        \draw[decorate, decoration={brace, amplitude=4pt}, black!60] (2.4,0.45) -- node[above=3pt] {$\tau\sigma$} (7.2,0.45);
    \end{tikzpicture}
    \end{center}
    È la stessa ragione per cui la composizione di funzioni è associativa.
\end{spiegazione}



  % Funzioni: composizione, iniettive/suriettive/biettive, immagine, sezioni e retrazioni
  \chapter{Funzioni}\label{cap:funzioni}

Le funzioni sono particolari corrispondenze: quelle che associano a \textbf{ogni} elemento del primo insieme
\textbf{uno e un solo} elemento del secondo. In questo capitolo le definiamo, le componiamo e studiamo le tre
proprietà fondamentali (iniettività, suriettività, biettività), caratterizzandole in più modi.

% ──────────────────────────────────────────────────────────────
\section{Definizione di funzione}\label{sec:def-funzione}

\dfn[funzione]{Funzione}{
    Siano $a$ e $b$ insiemi. Una corrispondenza $f = (a, b, g)$ si dice \textbf{applicazione} o \textbf{funzione} se
    \[ \forall x \in a\ \exists!\, y \in b\ (x f y). \]
    Visto che $y$ è unico, lo denoto $f(x)$. Inoltre:
    \begin{itemize}
        \item $f(x)$ si dice \textbf{immagine di $x$ mediante $f$};
        \item $a$ e $b$ si dicono rispettivamente \textbf{dominio} e \textbf{codominio} di $f$;
        \item $\mathrm{Im} f := \{y \in b \mid \exists x \in a\,(f(x) = y)\} = \{f(x) \mid x \in a\}$ si dice \textbf{immagine di $f$};
        \item scriviamo brevemente $f \colon a \to b$.
    \end{itemize}
}

Nella scrittura $\{f(x) \mid x \in a\}$ non si indica l'insieme di appartenenza perché è implicitamente il codominio di $f$.
È a volte utile dare una \textbf{descrizione esplicita} della funzione, che evidenzi per ogni $x$ qual è la sua immagine:
\[ f \colon x \in a \mapsto f(x) \in b. \]

% ──────────────────────────────────────────────────────────────
\section{Composizione e funzioni notevoli}\label{sec:composizione}

\begin{concetto}
\begin{Proposizione}{Funzione composta}{composta}
    \begin{prerequisiti}
        \dip{Definizione}{def:funzione}
        \dip{Definizione}{def:prodotto-corrispondenze}
    \end{prerequisiti}
    Siano $f \colon a \to b$ e $g \colon b \to c$ funzioni. Il prodotto $gf$ (come corrispondenze) è una funzione da $a$ a $c$,
    detta \textbf{funzione composta} e indicata con $g \circ f = gf$, con descrizione esplicita
    \[ g \circ f \colon x \in a \mapsto g(f(x)) \in c. \]
\end{Proposizione}
\begin{dimostrazione}
    Sia $x \in a$. Dobbiamo mostrare che esiste uno e un solo $z \in c$ con $x \,(gf)\, z$.
    \begin{caso}{Esistenza}
        Posto $y = f(x) \in b$ e $z = g(y) \in c$, si ha $x f y$ e $y g z$, quindi $x \,(gf)\, z$ per definizione di prodotto.
    \end{caso}
    \begin{caso}{Unicità}
        Se $x \,(gf)\, z'$, esiste $w$ con $x f w$ e $w g z'$; per l'unicità in $f$ si ha $w = f(x)$, e per l'unicità in $g$
        si ha $z' = g(w) = g(f(x))$.
    \end{caso}
    Quindi ogni $x \in a$ ha un'unica immagine, che è $g(f(x))$.
\end{dimostrazione}
\end{concetto}

\begin{spiegazione}
    $g \circ f$ si legge «$g$ dopo $f$»: prima si applica $f$, poi $g$. Con $f(x) = x + 1$ e $g(y) = 2y$:
    \begin{center}
    \begin{tikzpicture}[font=\small]
        \foreach \e/\y in {1/0.3, 2/-0.3} \node[elem] (a\e) at (0,\y) {$\e$};
        \foreach \e/\y in {2/0.3, 3/-0.3} \node[elem] (b\e) at (2.6,\y) {$\e$};
        \foreach \e/\y in {4/0.3, 6/-0.3} \node[elem] (c\e) at (5.2,\y) {$\e$};
        \foreach \x/\n in {0/a, 2.6/b, 5.2/c} { \draw[black!55, rounded corners=8pt, line width=0.6pt] (\x-0.42,-0.65) rectangle (\x+0.42,0.65); \node at (\x,0.9) {$\n$}; }
        \draw[fun] (a1) -- (b2); \draw[fun] (a2) -- (b3); \draw[fun2] (b2) -- (c4); \draw[fun2] (b3) -- (c6);
        \node[primary!75!black] at (1.3,0.95) {$f$}; \node[cherryred!80!black] at (3.9,0.95) {$g$};
        \draw[fun3, dashed] (a2) to[bend right=28] node[below] {$g \circ f$} (c6);
        \node[align=left, anchor=west] at (6.2,0) {$g \circ f(x) = g(x + 1) = 2x + 2$;\\ invece $f \circ g(x) = f(2x) = 2x + 1$:\\ l'ordine conta.};
    \end{tikzpicture}
    \end{center}
\end{spiegazione}


\dfn[funzioni-notevoli]{Funzioni notevoli}{
    \begin{itemize}
        \item Se $c \in b$, la funzione $f \colon x \in a \mapsto c \in b$ la dico \textbf{funzione costante} (di valore $c$).
        \item Se $a \subseteq b$, la funzione $f \colon x \in a \mapsto x \in b$ la dico \textbf{immersione} di $a$ in $b$.
        \item Se $s \subseteq a$, la funzione $f_{|s} \colon x \in s \mapsto f(x) \in b$ si chiama \textbf{restrizione} di $f$ a $s$,
              e $f$ si dice \textbf{prolungamento} di $f_{|s}$ ad $a$.
        \item La funzione $x \in a \mapsto x \in a$ la dico \textbf{identità} (o funzione identica) su $a$ e la indico con $\mathrm{id}_a$.
    \end{itemize}
}

% ──────────────────────────────────────────────────────────────
\section{Iniettività, suriettività, biettività}\label{sec:iniettive}

\dfn[iniettiva]{Funzione iniettiva}{
    $f \colon a \to b$ si dice \textbf{iniettiva} se
    \[ \forall x_1, x_2 \in a\,\big(x_1 \neq x_2 \rightarrow f(x_1) \neq f(x_2)\big). \]
}

\begin{concetto}
\begin{Proposizione}{Altre forme dell'iniettività}{forme-iniettivita}
    \begin{prerequisiti}
        \dip{Definizione}{def:iniettiva}
        \dip{Proposizione}{prop:contrapposizione}
        \dip{Teorema}{teo:negazione-quantificatori}
        \dip{Proposizione}{prop:negimp}
        \dip{Proposizione}{prop:doppianeg}
    \end{prerequisiti}
    Sia $f \colon a \to b$.
    \begin{enumerate}
        \item $f$ è iniettiva se e solo se $\forall x_1, x_2 \in a\,\big(f(x_1) = f(x_2) \rightarrow x_1 = x_2\big)$.
        \item $f$ \textbf{non} è iniettiva se e solo se $\exists x_1, x_2 \in a\,\big(x_1 \neq x_2 \wedge f(x_1) = f(x_2)\big)$.
    \end{enumerate}
\end{Proposizione}
\begin{dimostrazione}
    \begin{caso}{Punto 1: forma contrapposta}
        Per contrapposizione $(x_1 \neq x_2 \rightarrow f(x_1) \neq f(x_2)) \leftrightarrow (\neg(f(x_1) \neq f(x_2)) \rightarrow \neg(x_1 \neq x_2))$,
        e per la doppia negazione questo è $f(x_1) = f(x_2) \rightarrow x_1 = x_2$.
    \end{caso}
    \begin{caso}{Punto 2: negazione}
        Negando la definizione: $\neg\forall x_1, x_2\,(\dots)$ diventa $\exists x_1, x_2\,\neg(x_1 \neq x_2 \rightarrow f(x_1) \neq f(x_2))$
        (negazione dei quantificatori), cioè $\exists x_1, x_2\,(x_1 \neq x_2 \wedge f(x_1) = f(x_2))$ (negazione dell'implicazione e doppia negazione).
    \end{caso}
\end{dimostrazione}
\end{concetto}

\begin{spiegazione}
    Nel diagramma a frecce, $f$ è iniettiva se in ogni elemento del codominio arriva \textbf{al più una} freccia. Per dimostrare che non lo è basta
    trovare due elementi diversi con la stessa immagine (in arancione):
    \begin{center}
    \begin{tikzpicture}[font=\small]
        \diagfreccie[name prefix=i-]{a}{1,2,3}{b}{x,y,z,w}{1/1,2/2,3/4}
        \node at (1.2,-1.65) {iniettiva};
        \diagfreccie[xshift=5.2cm, name prefix=n-]{a}{1,2,3}{b}{x,y,z}{1/1,2/1,3/3}
        \node[evid] at (n-s1) {}; \node[evid] at (n-s2) {};
        \node at (6.4,-1.65) {non iniettiva: $1 \neq 2$, ma $f(1) = f(2) = x$};
    \end{tikzpicture}
    \end{center}
\end{spiegazione}


La forma 1 è quella più comoda nelle dimostrazioni: \textbf{per dimostrare che $f$ è iniettiva si prendono $x_1, x_2$ con
$f(x_1) = f(x_2)$ e si mostra che $x_1 = x_2$}.

\dfn[suriettiva]{Funzione suriettiva}{
    $f \colon a \to b$ si dice \textbf{suriettiva} se
    \[ \forall y \in b\ \exists x \in a\,\big(f(x) = y\big). \]
    Se $\mathrm{Im} f \subseteq t \subseteq b$, la funzione $x \in a \mapsto f(x) \in t$ si chiama \textbf{ridotta} di $f$ a $t$.
}

\begin{concetto}
\begin{Proposizione}{La ridotta all'immagine è suriettiva}{ridotta}
    \begin{prerequisiti}
        \dip{Definizione}{def:suriettiva}
        \dip{Definizione}{def:funzione}
    \end{prerequisiti}
    Sia $f \colon a \to b$. La ridotta di $f$ a $\mathrm{Im} f$, cioè $x \in a \mapsto f(x) \in \mathrm{Im} f$, è suriettiva.
\end{Proposizione}
\begin{dimostrazione}
    Sia $y \in \mathrm{Im} f$. Per definizione di immagine esiste $x \in a$ con $f(x) = y$.
\end{dimostrazione}
\end{concetto}

\begin{spiegazione}
    Una funzione può non essere suriettiva solo perché il codominio è «troppo grande»: restringendolo all'immagine diventa suriettiva.
    \begin{center}
    \begin{tikzpicture}[font=\small]
        \diagfreccie[name prefix=f-]{a}{1,2,3}{b}{x,y,z,w}{1/1,2/2,3/2}
        \node[evid2] at (f-d1) {}; \node[evid2] at (f-d2) {};
        \node[align=center] at (1.2,-1.75) {$f$: in verde $\mathrm{Im} f = \{x, y\}$;\\ $z$ e $w$ non sono raggiunti};
        \diagfreccie[xshift=6.0cm, name prefix=r-]{a}{1,2,3}{\mathrm{Im} f}{x,y}{1/1,2/2,3/2}
        \node at (7.2,-1.75) {la ridotta: è suriettiva};
    \end{tikzpicture}
    \end{center}
    Ad esempio $x \in \mbZ \mapsto x^2 \in \mbZ$ non è suriettiva, ma la sua ridotta a $\{0, 1, 4, 9, \dots\}$ lo è.
\end{spiegazione}


\dfn[biettiva]{Funzione biettiva}{
    Se $f$ è suriettiva e iniettiva, $f$ si dice \textbf{biettiva} (o \textbf{biunivoca}).
}

\begin{concetto}
\begin{Teorema}{Composta di funzioni iniettive}{comp-iniettive}
    \begin{prerequisiti}
        \dip{Definizione}{def:iniettiva}
        \dip{Proposizione}{prop:forme-iniettivita}
        \dip{Proposizione}{prop:composta}
    \end{prerequisiti}
    Siano $f \colon a \to b$ e $g \colon b \to c$ iniettive. Allora $g \circ f$ è iniettiva.
\end{Teorema}
\begin{dimostrazione}
    Siano $x_1, x_2 \in a$ tali che $g \circ f(x_1) = g \circ f(x_2)$, cioè $g(f(x_1)) = g(f(x_2))$.
    Ma $g$ è iniettiva, allora $f(x_1) = f(x_2)$; ma $f$ è iniettiva, allora $x_1 = x_2$.
\end{dimostrazione}
\end{concetto}

\begin{spiegazione}
    Se nessuno dei due passaggi manda elementi diversi nello stesso elemento, non lo fa neanche il percorso completo. Con $f(x) = x + 1$ e $g(y) = 2y$ su $\mbN$:
    \begin{passi}
        \item supponiamo $g(f(x_1)) = g(f(x_2))$, cioè $2(x_1 + 1) = 2(x_2 + 1)$;
        \item $g$ è iniettiva, quindi $x_1 + 1 = x_2 + 1$, cioè $f(x_1) = f(x_2)$;
        \item $f$ è iniettiva, quindi $x_1 = x_2$.
    \end{passi}
    Si «risale» la composizione un passo alla volta, dall'ultima funzione alla prima.
\end{spiegazione}


\begin{concetto}
\begin{Teorema}{Composta di funzioni suriettive}{comp-suriettive}
    \begin{prerequisiti}
        \dip{Definizione}{def:suriettiva}
        \dip{Proposizione}{prop:composta}
    \end{prerequisiti}
    Siano $f \colon a \to b$ e $g \colon b \to c$ suriettive. Allora $g \circ f$ è suriettiva.
\end{Teorema}
\begin{dimostrazione}
    Sia $y \in c$. Poiché $g$ è suriettiva, trovo $z \in b$ con $g(z) = y$; ma $f$ è suriettiva, allora trovo
    $x \in a$ con $f(x) = z$. Quindi $g(f(x)) = g(z) = y$.
\end{dimostrazione}
\end{concetto}

\begin{spiegazione}
    Per raggiungere un elemento $y$ dell'ultimo insieme si cammina \textbf{all'indietro}: prima si trova chi va in $y$ con $g$, poi chi va lì con $f$. Con
    $f(x) = x + 1$ e $g(z) = -z$ su $\mbZ$, per ottenere $g(f(x)) = 7$ serve $f(x) = -7$, cioè $x = -8$: infatti $g(f(-8)) = -(-8 + 1) = 7$.
\end{spiegazione}


\begin{concetto}
\begin{Corollario}{Composta di funzioni biettive}{comp-biettive}
    \begin{prerequisiti}
        \dip{Teorema}{teo:comp-iniettive}
        \dip{Teorema}{teo:comp-suriettive}
        \dip{Definizione}{def:biettiva}
    \end{prerequisiti}
    La composta di funzioni biettive è biettiva.
\end{Corollario}
\begin{dimostrazione}
    Segue immediatamente dai due teoremi precedenti.
\end{dimostrazione}
\end{concetto}

\begin{spiegazione}
    Una biezione di un insieme in sé è un «rimescolamento» dei suoi elementi: mescolando due volte un mazzo di carte si ottiene ancora un ordine delle
    stesse carte, senza perderne né duplicarne nessuna.
\end{spiegazione}


\begin{concetto}
\begin{Teorema}{Funzioni composte e caratteristiche}{composte-caratteristiche}
    \begin{prerequisiti}
        \dip{Proposizione}{prop:composta}
        \dip{Proposizione}{prop:forme-iniettivita}
        \dip{Definizione}{def:suriettiva}
    \end{prerequisiti}
    Siano $f \colon a \to b$ e $g \colon b \to c$ due funzioni.
    \begin{enumerate}
        \item Se $g \circ f$ è iniettiva, allora $f$ è iniettiva.
        \item Se $g \circ f$ è suriettiva, allora $g$ è suriettiva.
    \end{enumerate}
    Cioè, se $g \circ f$ è biettiva, allora $f$ è iniettiva e $g$ è suriettiva.
\end{Teorema}
\begin{dimostrazione}
    \begin{caso}{Punto 1: $g \circ f$ iniettiva $\Rightarrow$ $f$ iniettiva}
        Prendo $x_1, x_2 \in a$ con $f(x_1) = f(x_2)$; quindi $g \circ f(x_1) = g(f(x_1)) = g(f(x_2)) = g \circ f(x_2)$,
        ma $g \circ f$ è iniettiva, quindi $x_1 = x_2$. Dunque $f$ è iniettiva.
    \end{caso}
    \begin{caso}{Punto 2: $g \circ f$ suriettiva $\Rightarrow$ $g$ suriettiva}
        Sia $y \in c$. Visto che $g \circ f$ è suriettiva e il suo codominio è $c$, trovo $x \in a$ con
        $g \circ f(x) = y$, cioè $g(f(x)) = y$: l'elemento $f(x) \in b$ ha immagine $y$ mediante $g$. Questo vale per ogni $y$,
        quindi $g$ è suriettiva.
    \end{caso}
\end{dimostrazione}
\end{concetto}

\begin{spiegazione}
    Se $f$ manda due elementi diversi nello stesso elemento, $g$ non può più separarli: quindi, se $g \circ f$ è iniettiva, già $f$ doveva esserlo. Se
    $g \circ f$ raggiunge tutto $c$, lo raggiunge con l'ultimo passo, quindi $g$ è suriettiva.
    \begin{center}
    \begin{tikzpicture}[font=\small]
        \node[elem] (x1) at (0,0.3) {$x_1$}; \node[elem] (x2) at (0,-0.3) {$x_2$};
        \node[elem] (y) at (2.4,0) {$y$}; \node[elem] (z) at (4.8,0) {$g(y)$};
        \foreach \x/\n in {0/a, 2.4/b, 4.8/c} { \draw[black!55, rounded corners=8pt, line width=0.6pt] (\x-0.5,-0.65) rectangle (\x+0.5,0.65); \node at (\x,0.9) {$\n$}; }
        \draw[fun] (x1) -- (y); \draw[fun] (x2) -- (y); \draw[fun2] (y) -- (z);
        \node[primary!75!black] at (1.2,0.55) {$f$}; \node[cherryred!80!black] at (3.6,0.3) {$g$};
        \node[align=left, anchor=west] at (5.8,0) {$f(x_1) = f(x_2)$, quindi anche $g(f(x_1)) = g(f(x_2))$:\\ se $f$ non è iniettiva, non lo è nemmeno $g \circ f$.};
    \end{tikzpicture}
    \end{center}
    Il viceversa non vale: nel terzo dei «Tre esempi importanti» qui sotto $g \circ f$ è biettiva, ma $f$ non è suriettiva e $g$ non è iniettiva.
\end{spiegazione}


\begin{esempio}{Tre esempi importanti}
    \begin{itemize}
        \item Sia $a$ un insieme. La funzione $x \in a \mapsto \{x\} \in \mcP(a)$ è \textbf{iniettiva} (se $\{x\} = \{y\}$, per
              Estensionalità $x = y$) ma \textbf{non suriettiva}: $\varnothing \in \mcP(a)$ non è un singleton, quindi non è immagine di nessun $x$.
        \item La funzione $x \in \mcP(a) \mapsto a \setminus x \in \mcP(a)$ è una importantissima funzione \textbf{biettiva}.
              Infatti per la doppia negazione insiemistica (Proposizione~\ref{prop:doppia-neg-ins}) $a \setminus (a \setminus x) = x$
              per ogni $x \subseteq a$: se $a \setminus x = a \setminus y$, applicando $a \setminus$ a entrambi si ottiene $x = y$
              (iniettiva); e ogni $y \subseteq a$ è l'immagine di $a \setminus y$ (suriettiva).
        \item Siano $a = \{1\}$, $b = \{1, 2\}$, $f$ l'immersione di $a$ in $b$ e $g$ la funzione costante da $b$ in $a$.
              Allora $g \circ f = \mathrm{id}_a$, ma $f$ non è suriettiva e $g$ non è iniettiva: nel Teorema~\ref{teo:composte-caratteristiche}
              non si può dire di più.
    \end{itemize}
    \begin{center}
    \begin{tikzpicture}[>=Stealth, every node/.style={font=\small}]
        \node (a1) at (0,0) {$1$};
        \node (b1) at (3,0.5) {$1$}; \node (b2) at (3,-0.5) {$2$};
        \draw[rounded corners] (-0.5,-1) rectangle (0.5,1); \node at (0,1.3) {$a = \{1\}$};
        \draw[rounded corners] (2.5,-1) rectangle (3.5,1); \node at (3,1.3) {$b = \{1,2\}$};
        \draw[->,primary] (a1) to[bend left=15] node[above,sloped] {$f$} (b1);
        \draw[->,cherryred] (b1) to[bend left=15] node[below,sloped] {$g$} (a1);
        \draw[->,cherryred] (b2) to[bend left=10] (a1);
    \end{tikzpicture}
    \end{center}
\end{esempio}

% ──────────────────────────────────────────────────────────────
\section{Immagine e antiimmagine}\label{sec:immagine}

\dfn[immagine-antiimmagine]{Applicazioni immagine e antiimmagine}{
    Sia $f \colon a \to b$. Definisco l'applicazione \textbf{immagine}
    \[ \overrightarrow{f} \colon x \in \mcP(a) \mapsto \overrightarrow{f}(x) = \{f(z) \mid z \in x\} \in \mcP(b) \]
    e l'applicazione \textbf{antiimmagine}
    \[ \overleftarrow{f} \colon y \in \mcP(b) \mapsto \overleftarrow{f}(y) = \{x \in a \mid f(x) \in y\} \in \mcP(a). \]
}

\begin{concetto}
\begin{Proposizione}{Proprietà di immagine e antiimmagine}{prop-immagine}
    \begin{prerequisiti}
        \dip{Definizione}{def:immagine-antiimmagine}
        \dip{Definizione}{def:funzione}
        \dip{Definizione}{def:suriettiva}
        \dip{Lemma}{lem:vuoto}
        \dip{Proposizione}{prop:doppia-inclusione}
    \end{prerequisiti}
    Sia $f \colon a \to b$ e sia $x \in a$. Allora:
    \begin{enumerate}
        \item $\overrightarrow{f}(\varnothing) = \varnothing = \overleftarrow{f}(\varnothing)$;
        \item $\overrightarrow{f}(a) = \mathrm{Im} f$;
        \item $f$ è suriettiva $\leftrightarrow$ $\mathrm{Im} f = b$ $\leftrightarrow$ $\overrightarrow{f}(a) = b$;
        \item $\overrightarrow{f}(\{x\}) = \{f(x)\}$;
        \item $\overleftarrow{f}(b) = a = \overleftarrow{f}(\mathrm{Im} f)$.
    \end{enumerate}
\end{Proposizione}
\begin{dimostrazione}
    \begin{caso}{Punto 1: $\overrightarrow{f}(\varnothing) = \varnothing = \overleftarrow{f}(\varnothing)$}
        Non esiste $z \in \varnothing$, quindi $\{f(z) \mid z \in \varnothing\}$ non ha elementi; allo stesso modo nessun
        $x$ soddisfa $f(x) \in \varnothing$. Per il Lemma~\ref{lem:vuoto} entrambi sono $\varnothing$.
    \end{caso}
    \begin{caso}{Punto 2: $\overrightarrow{f}(a) = \mathrm{Im} f$}
        $\overrightarrow{f}(a) = \{f(z) \mid z \in a\}$, che è la definizione di $\mathrm{Im} f$.
    \end{caso}
    \begin{caso}{Punto 3: $f$ suriettiva $\leftrightarrow \mathrm{Im} f = b \leftrightarrow \overrightarrow{f}(a) = b$}
        Sempre $\mathrm{Im} f \subseteq b$. $(\rightarrow)$ Se $f$ è suriettiva, ogni $y \in b$ sta in $\mathrm{Im} f$, cioè
        $b \subseteq \mathrm{Im} f$, e per la doppia inclusione $\mathrm{Im} f = b$. $(\leftarrow)$ Se $\mathrm{Im} f = b$, ogni $y \in b$ sta in
        $\mathrm{Im} f$, cioè è immagine di qualche $x$: $f$ è suriettiva. L'ultima equivalenza segue dal punto 2.
    \end{caso}
    \begin{caso}{Punto 4: $\overrightarrow{f}(\{x\}) = \{f(x)\}$}
        $\overrightarrow{f}(\{x\}) = \{f(z) \mid z \in \{x\}\} = \{f(z) \mid z = x\} = \{f(x)\}$.
    \end{caso}
    \begin{caso}{Punto 5: $\overleftarrow{f}(b) = a = \overleftarrow{f}(\mathrm{Im} f)$}
        $(\supseteq)$ Per ogni $x \in a$ si ha $f(x) \in \mathrm{Im} f \subseteq b$, quindi $x$ sta sia in $\overleftarrow{f}(b)$ sia in
        $\overleftarrow{f}(\mathrm{Im} f)$. $(\subseteq)$ Entrambi sono per definizione sottoinsiemi di $a$.
        Per doppia inclusione sono uguali ad $a$.
    \end{caso}
\end{dimostrazione}
\end{concetto}

\begin{spiegazione}
    L'immagine $\overrightarrow{f}$ trasporta una \textbf{parte} del dominio in avanti; l'antiimmagine $\overleftarrow{f}$ raccoglie tutti gli elementi che finiscono in una parte
    del codominio. Con la funzione $f$ qui sotto:
    \begin{center}
    \begin{tikzpicture}[font=\small]
        \diagfreccie[name prefix=f-]{a}{1,2,3,4}{b}{0,1,2}{1/2,2/3,3/1,4/2}
        \node[evid] at (f-s1) {}; \node[evid] at (f-s4) {}; \node[evid] at (f-d2) {};
        \node[align=left, anchor=west] at (3.4,0) {$\overrightarrow{f}(\{1, 4\}) = \{1\}$ \ e \ $\overleftarrow{f}(\{1\}) = \{1, 4\}$ (in arancione);\\
            $\overleftarrow{f}(\{0, 2\}) = \{2, 3\}$; \ $\overrightarrow{f}(\{2\}) = \{f(2)\} = \{2\}$;\\
            $\overrightarrow{f}(a) = \mathrm{Im} f = \{0, 1, 2\} = b$: $f$ è suriettiva; \ $\overleftarrow{f}(b) = a$.};
    \end{tikzpicture}
    \end{center}
\end{spiegazione}


% ──────────────────────────────────────────────────────────────
\section{Caratterizzazioni con l'antiimmagine}\label{sec:caratt-antiimmagine}

\dfn[map-in-su-bi]{Insiemi di funzioni}{
    Siano $a$, $b$ due insiemi. Dico
    \begin{itemize}
        \item $\mathrm{Map}(a, b)$ l'insieme delle funzioni da $a$ a $b$;
        \item $\mathrm{In}(a, b)$ l'insieme delle funzioni iniettive da $a$ a $b$;
        \item $\mathrm{Su}(a, b)$ l'insieme delle funzioni suriettive da $a$ a $b$;
        \item $\mathrm{Bi}(a, b)$ l'insieme delle funzioni biettive da $a$ a $b$.
    \end{itemize}
}

\begin{concetto}
\begin{Teorema}{Caratterizzazione delle funzioni iniettive}{caratt-iniettive}
    \begin{prerequisiti}
        \dip{Definizione}{def:map-in-su-bi}
        \dip{Definizione}{def:immagine-antiimmagine}
        \dip{Proposizione}{prop:forme-iniettivita}
        \dip{Proposizione}{prop:negdoppiaimp}
        \dip{Teorema}{teo:negazione-quantificatori}
        \dip{Definizione}{def:esiste-unico}
    \end{prerequisiti}
    $f \colon a \to b$ è iniettiva se e solo se, per ogni $y \in b$, $\overleftarrow{f}(\{y\})$ ha al più un elemento, cioè:
    \[ \forall a, b\ \forall f\,\Big(f \in \mathrm{In}(a, b) \leftrightarrow \forall y \in b\,\big(\overleftarrow{f}(\{y\}) = \varnothing \vee \exists!\, x \in a\,(f(x) = y)\big)\Big). \]
\end{Teorema}
\begin{dimostrazione}
    Lo dimostro usando la \textbf{negazione della doppia implicazione}: basta dimostrare che le \textbf{negazioni} dei due
    membri sono equivalenti. Negandoli:
    \begin{itemize}
        \item $f$ \textbf{non} è iniettiva se e solo se esistono $x_1 \neq x_2$ in $a$ con $f(x_1) = f(x_2)$
              (Proposizione~\ref{prop:forme-iniettivita}, punto 2);
        \item il secondo membro è falso se e solo se esiste $y \in b$ tale che $\overleftarrow{f}(\{y\})$ non è vuoto e non ha
              un solo elemento, cioè $\overleftarrow{f}(\{y\})$ ha \textbf{almeno due} elementi.
    \end{itemize}
    \begin{caso}{$(\rightarrow)$ Se $f$ non è iniettiva, allora esiste $y$ con $\overleftarrow{f}(\{y\})$ di almeno due elementi}
        Siano $x_1 \neq x_2$ con $f(x_1) = f(x_2)$. Detto $y = f(x_1)$, l'insieme $\overleftarrow{f}(\{y\})$ contiene almeno $x_1$ e $x_2$.
    \end{caso}
    \begin{caso}{$(\leftarrow)$ Se esiste $y$ con $\overleftarrow{f}(\{y\})$ di almeno due elementi, allora $f$ non è iniettiva}
        Se $\overleftarrow{f}(\{y\})$ contiene due elementi distinti $x_1$, $x_2$, allora $f(x_1) = y = f(x_2)$ con $x_1 \neq x_2$.
    \end{caso}
\end{dimostrazione}
\end{concetto}

\begin{spiegazione}
    $\overleftarrow{f}(\{y\})$ è l'insieme degli elementi da cui parte una freccia verso $y$ (la «fibra» di $y$). $f$ è iniettiva quando ogni fibra ha al più un elemento:
    \begin{center}
    \begin{tikzpicture}[font=\small]
        \diagfreccie[name prefix=f-]{a}{1,2,3}{b}{x,y,z}{1/1,2/1,3/2}
        \draw[orange!85!black, line width=0.9pt, rounded corners=5pt] (-0.3,-0.12) rectangle (0.3,0.72);
        \node[align=left, anchor=west] at (3.4,0) {la fibra di $x$ è $\{1, 2\}$ (in arancione): ha due elementi,\\
            quindi $f$ non è iniettiva; le fibre di $y$ e di $z$ sono $\{3\}$ e $\varnothing$.};
    \end{tikzpicture}
    \end{center}
\end{spiegazione}


\begin{concetto}
\begin{Teorema}{Caratterizzazione delle funzioni suriettive}{caratt-suriettive}
    \begin{prerequisiti}
        \dip{Definizione}{def:map-in-su-bi}
        \dip{Definizione}{def:immagine-antiimmagine}
        \dip{Definizione}{def:suriettiva}
        \dip{Proposizione}{prop:negdoppiaimp}
    \end{prerequisiti}
    \[ \forall a, b\ \forall f\,\Big(f \in \mathrm{Su}(a, b) \leftrightarrow \forall y \in b\,\big(\overleftarrow{f}(\{y\}) \neq \varnothing\big)\Big). \]
\end{Teorema}
\begin{dimostrazione}
    Di nuovo con la negazione della doppia implicazione: basta dimostrare che le negazioni dei due membri sono equivalenti.
    Ricordiamo che $\overleftarrow{f}(\{y\}) = \{x \in a \mid f(x) = y\}$.
    \begin{caso}{$(\rightarrow)$ Se $f$ non è suriettiva, allora esiste $y \in b$ con $\overleftarrow{f}(\{y\}) = \varnothing$}
        Esiste $y \in b$ che non è immagine di nessun elemento di $a$: nessun $x \in a$ ha $f(x) = y$, quindi $\overleftarrow{f}(\{y\}) = \varnothing$.
    \end{caso}
    \begin{caso}{$(\leftarrow)$ Se esiste $y \in b$ con $\overleftarrow{f}(\{y\}) = \varnothing$, allora $f$ non è suriettiva}
        Nessun $x \in a$ ha $f(x) = y$, cioè $y$ non è immagine di nessun elemento: $f$ non è suriettiva.
    \end{caso}
\end{dimostrazione}
\end{concetto}

\begin{spiegazione}
    $f$ è suriettiva quando nessuna fibra è vuota, cioè quando in ogni elemento del codominio arriva \textbf{almeno una} freccia. Nell'esempio precedente
    la fibra di $z$ è vuota: nessuna freccia arriva in $z$, e infatti $f$ non è suriettiva.
\end{spiegazione}


\begin{concetto}
\begin{Corollario}{Caratterizzazione delle funzioni biettive con l'antiimmagine}{caratt-biettive-anti}
    \begin{prerequisiti}
        \dip{Teorema}{teo:caratt-iniettive}
        \dip{Teorema}{teo:caratt-suriettive}
        \dip{Definizione}{def:biettiva}
    \end{prerequisiti}
    $f \colon a \to b$ è biettiva se e solo se, per ogni $y \in b$, $\overleftarrow{f}(\{y\})$ ha esattamente un elemento, cioè
    \[ \forall y \in b\ \exists!\, x \in a\,\big(f(x) = y\big). \]
\end{Corollario}
\begin{dimostrazione}
    Si combinano i due teoremi precedenti: «al più un elemento» e «non vuoto» insieme danno «esattamente un elemento».
\end{dimostrazione}
\end{concetto}

\begin{spiegazione}
    Biettiva vuol dire che in ogni elemento del codominio arriva \textbf{esattamente una} freccia: le frecce accoppiano gli elementi dei due insiemi uno a uno.
    \begin{center}
    \begin{tikzpicture}[font=\small]
        \diagfreccie{a}{1,2,3}{b}{x,y,z}{1/2,2/3,3/1}
        \node[align=left, anchor=west] at (3.4,0) {ogni fibra ha un solo elemento: $\overleftarrow{f}(\{x\}) = \{3\}$,\\ $\overleftarrow{f}(\{y\}) = \{1\}$, $\overleftarrow{f}(\{z\}) = \{2\}$.};
    \end{tikzpicture}
    \end{center}
\end{spiegazione}


% ──────────────────────────────────────────────────────────────
\section{Sezioni, retrazioni e inversa}\label{sec:sezioni}

\dfn[sezione-retrazione]{Sezione, retrazione, inversa}{
    Siano $f \colon a \to b$ e $g \colon b \to a$.
    \begin{itemize}
        \item Se $g \circ f = \mathrm{Id}_a$, allora $g$ si dice \textbf{retrazione} di $f$.
        \item Se $f \circ g = \mathrm{Id}_b$, allora $g$ si dice \textbf{sezione} di $f$.
        \item Se $g$ è sia sezione che retrazione, allora $g$ è l'\textbf{inversa} di $f$, e $f$ si dice \textbf{invertibile}.
    \end{itemize}
}

\begin{concetto}
\begin{Teorema}{Funzione iniettiva e retrazione}{iniettiva-retrazione}
    \begin{prerequisiti}
        \dip{Definizione}{def:sezione-retrazione}
        \dip{Definizione}{def:map-in-su-bi}
        \dip{Teorema}{teo:composte-caratteristiche}
        \dip{Definizione}{def:funzioni-notevoli}
        \dip{Definizione}{def:funzione}
    \end{prerequisiti}
    $f \colon a \to b$ è iniettiva $\leftrightarrow$ $a = \varnothing \ \vee\ \exists g \in \mathrm{Map}(b, a)\,(g \text{ è una retrazione di } f)$.
\end{Teorema}
\begin{dimostrazione}
    \begin{caso}{$(\rightarrow)$ Se $f$ è iniettiva, allora $a = \varnothing$ oppure $f$ ha una retrazione}
        \begin{caso}{Caso $a = \varnothing$}
            Non c'è niente da dimostrare.
        \end{caso}
        \begin{caso}{Caso $a \neq \varnothing$}
            Prendo $\bar{x} \in a$. Per ogni $y \in \mathrm{Im} f$ esiste un unico $x \in a$ con $f(x) = y$ (perché $f$ è iniettiva). Definisco
            \[ g \colon y \in b \mapsto \begin{cases} x & \text{se } y \in \mathrm{Im} f \wedge f(x) = y \\ \bar{x} & \text{se } y \in b \setminus \mathrm{Im} f \end{cases} \in a. \]
            Allora per ogni $x \in a$ si ha $g \circ f(x) = g(f(x)) = x$, cioè $g \circ f = \mathrm{Id}_a$: $g$ è una retrazione di $f$.
        \end{caso}
    \end{caso}
    \begin{caso}{$(\leftarrow)$ Se $a = \varnothing$ oppure $f$ ha una retrazione, allora $f$ è iniettiva}
        \begin{caso}{Caso $a = \varnothing$}
            $f = (\varnothing, b, \varnothing)$ è iniettiva: la condizione di iniettività è un'implicazione con antecedente sempre falso.
        \end{caso}
        \begin{caso}{Caso $f$ ha una retrazione $g$}
            $g \circ f = \mathrm{Id}_a$ e l'identità è iniettiva, quindi $f$ è iniettiva per il Teorema~\ref{teo:composte-caratteristiche}, punto 1.
        \end{caso}
    \end{caso}
\end{dimostrazione}
\end{concetto}

\begin{spiegazione}
    Una retrazione «disfa» $f$: riporta ogni immagine al suo unico punto di partenza, e manda gli elementi non raggiunti dove si vuole (in $\bar{x}$).
    \begin{center}
    \begin{tikzpicture}[font=\small]
        \diagfreccie[name prefix=f-]{a}{1,2}{b}{x,y,z}{1/1,2/2}
        \draw[fun2, dashed] (f-d1) to[bend right=25] (f-s1); \draw[fun2, dashed] (f-d2) to[bend left=10] (f-s2); \draw[fun2, dashed] (f-d3) to[bend left=20] (f-s1);
        \node[align=left, anchor=west] at (3.4,0) {in blu $f$, in rosso una retrazione $g$ con $\bar{x} = 1$:\\ $g(x) = 1$, $g(y) = 2$, $g(z) = 1$; \ $g(f(1)) = 1$ e $g(f(2)) = 2$.};
    \end{tikzpicture}
    \end{center}
    Se $f$ non fosse iniettiva, con $f(1) = f(2)$, la retrazione dovrebbe riportare quell'immagine sia in $1$ sia in $2$: impossibile.
\end{spiegazione}


\begin{concetto}
\begin{Teorema}{Funzione suriettiva e sezione}{suriettiva-sezione}
    \begin{prerequisiti}
        \dip{Definizione}{def:sezione-retrazione}
        \dip{Teorema}{teo:composte-caratteristiche}
    \end{prerequisiti}
    $f \colon a \to b$ è suriettiva $\leftrightarrow$ $f$ ha sezioni.
\end{Teorema}
\begin{dimostrazione}[Dimostrazione parziale]
    \begin{caso}{$(\rightarrow)$ Se $f$ è suriettiva, allora $f$ ha sezioni}
        Di questa implicazione non si fornisce dimostrazione, perché richiederebbe l'introduzione di un ulteriore
        assioma: l'\textbf{Assioma della Scelta} (per ogni $y \in b$ bisogna «scegliere» un elemento di $\overleftarrow{f}(\{y\})$).
    \end{caso}
    \begin{caso}{$(\leftarrow)$ Se $f$ ha una sezione, allora $f$ è suriettiva}
        Se $g$ è una sezione, $f \circ g = \mathrm{Id}_b$ è suriettiva, quindi $f$ è suriettiva per il
        Teorema~\ref{teo:composte-caratteristiche}, punto 2.
    \end{caso}
\end{dimostrazione}
\end{concetto}

\begin{spiegazione}
    Una sezione \textbf{sceglie}, per ogni elemento del codominio, uno dei suoi punti di partenza:
    \begin{center}
    \begin{tikzpicture}[font=\small]
        \diagfreccie[name prefix=f-]{a}{1,2,3}{b}{x,y}{1/1,2/1,3/2}
        \draw[fun2, dashed] (f-d1) to[bend right=20] (f-s1); \draw[fun2, dashed] (f-d2) to[bend left=20] (f-s3);
        \node[align=left, anchor=west] at (3.4,0) {in rosso la sezione $s$ con $s(x) = 1$ e $s(y) = 3$: $f(s(x)) = x$, $f(s(y)) = y$;\\
            un'altra sezione è $s'(x) = 2$, $s'(y) = 3$.};
    \end{tikzpicture}
    \end{center}
    Con infinite scelte da fare tutte insieme serve l'assioma della scelta: per questo una delle due implicazioni resta senza dimostrazione.
\end{spiegazione}


\begin{concetto}
\begin{Teorema}{Sezioni e retrazioni}{unicita-inversa}
    \begin{prerequisiti}
        \dip{Definizione}{def:sezione-retrazione}
        \dip{Proposizione}{prop:assoc-prodotto}
        \dip{Definizione}{def:funzioni-notevoli}
    \end{prerequisiti}
    Sia $f \colon a \to b$. Sia $s$ una sezione di $f$ ed $r$ una retrazione di $f$. Allora $s = r$.
    In particolare, l'inversa, se esiste, è \textbf{unica}.
\end{Teorema}
\begin{dimostrazione}
    \begin{caso}{Ogni sezione è uguale a ogni retrazione}
        Usando l'associatività del prodotto (composizione):
        \[ r = r \circ \mathrm{Id}_b = r \circ (f \circ s) = (r \circ f) \circ s = \mathrm{Id}_a \circ s = s. \]
    \end{caso}
    \begin{caso}{L'inversa è unica}
        Se $g$ e $g'$ sono due inverse, $g$ è in particolare una sezione e $g'$ una retrazione, quindi $g = g'$.
    \end{caso}
\end{dimostrazione}
\end{concetto}

\begin{spiegazione}
    La catena di uguaglianze si legge così:
    \begin{passi}
        \item $r = r \circ \mathrm{Id}_b$: comporre con l'identità non cambia nulla;
        \item $\mathrm{Id}_b = f \circ s$, perché $s$ è una sezione; quindi $r = r \circ (f \circ s)$;
        \item per l'associatività si spostano le parentesi: $(r \circ f) \circ s$;
        \item $r \circ f = \mathrm{Id}_a$, perché $r$ è una retrazione; quindi $r = \mathrm{Id}_a \circ s = s$.
    \end{passi}
    È lo stesso ragionamento che darà l'unicità dell'inverso in un monoide (Teorema~\ref{teo:unicita-inverso}).
\end{spiegazione}


Data la sua unicità, l'inversa di $f$, se esiste, la denoto con $f^{-1}$.

\begin{concetto}
\begin{Teorema}{Funzione biunivoca e inversa}{biettiva-inversa}
    \begin{prerequisiti}
        \dip{Teorema}{teo:iniettiva-retrazione}
        \dip{Teorema}{teo:suriettiva-sezione}
        \dip{Teorema}{teo:unicita-inversa}
        \dip{Definizione}{def:biettiva}
    \end{prerequisiti}
    $f \colon a \to b$ è biettiva $\leftrightarrow$ $f$ ha un'inversa.
\end{Teorema}
\begin{dimostrazione}
    È conseguenza dei teoremi precedenti.
    \begin{caso}{$(\rightarrow)$ Se $f$ è biettiva, allora $f$ ha un'inversa}
        $f$ è suriettiva, quindi ha una sezione $s$ (Teorema~\ref{teo:suriettiva-sezione}).
        \begin{caso}{Caso $a \neq \varnothing$}
            $f$ è iniettiva, quindi ha una retrazione $r$ (Teorema~\ref{teo:iniettiva-retrazione}); per il
            Teorema~\ref{teo:unicita-inversa} $s = r$, che è quindi sia sezione sia retrazione: l'inversa.
        \end{caso}
        \begin{caso}{Caso $a = \varnothing$}
            La suriettività obbliga $b = \varnothing$, e l'inversa è la funzione vuota.
        \end{caso}
    \end{caso}
    \begin{caso}{$(\leftarrow)$ Se $f$ ha un'inversa, allora $f$ è biettiva}
        L'inversa è una retrazione, quindi $f$ è iniettiva (Teorema~\ref{teo:iniettiva-retrazione}); è anche una
        sezione, quindi $f$ è suriettiva (Teorema~\ref{teo:suriettiva-sezione}).
    \end{caso}
\end{dimostrazione}
\end{concetto}

\begin{spiegazione}
    $x \in \mbZ \mapsto x + 3 \in \mbZ$ è biettiva e la sua inversa è $y \mapsto y - 3$: aggiungere $3$ e togliere $3$ si annullano. Invece $x \in \mbN \mapsto x + 3 \in \mbN$
    è iniettiva ma non suriettiva ($0$, $1$, $2$ non sono raggiunti): ha retrazioni, ad esempio $y \mapsto y - 3$ per $y \geq 3$ e $y \mapsto 0$ altrimenti, ma nessuna
    sezione, e quindi nessuna inversa.
\end{spiegazione}


\begin{concetto}
\begin{Teorema}{Funzione biettiva}{caratt-biettive}
    \begin{prerequisiti}
        \dip{Teorema}{teo:biettiva-inversa}
        \dip{Teorema}{teo:unicita-inversa}
        \dip{Corollario}{cor:caratt-biettive-anti}
        \dip{Teorema}{teo:suriettiva-sezione}
        \dip{Proposizione}{prop:forme-iniettivita}
        \dip{Definizione}{def:sezione-retrazione}
        \dip{Definizione}{def:immagine-antiimmagine}
    \end{prerequisiti}
    Sia $f \colon a \to b$. Le seguenti condizioni sono equivalenti:
    \begin{enumerate}
        \item $f$ è biettiva;
        \item $f$ ha sezioni e retrazioni;
        \item $f$ ha inversa;
        \item $f$ ha una sola sezione e una sola retrazione;
        \item $\forall y \in b\ \exists!\, x \in a\,\big(f(x) = y\big)$;
        \item $f$ ha una ed una sola sezione.
    \end{enumerate}
\end{Teorema}
\begin{dimostrazione}
    Dimostriamo che ogni condizione è equivalente a (1) o a (3), che sono equivalenti tra loro.
    \begin{caso}{(1) $\leftrightarrow$ (3)}
        È il Teorema~\ref{teo:biettiva-inversa}.
    \end{caso}
    \begin{caso}{(1) $\leftrightarrow$ (5)}
        È il Corollario~\ref{cor:caratt-biettive-anti}.
    \end{caso}
    \begin{caso}{(3) $\rightarrow$ (2)}
        L'inversa è sia sezione sia retrazione.
    \end{caso}
    \begin{caso}{(2) $\rightarrow$ (3)}
        Se $s$ è una sezione e $r$ una retrazione, per il Teorema~\ref{teo:unicita-inversa} $s = r$, che è quindi sia sezione
        sia retrazione, cioè l'inversa.
    \end{caso}
    \begin{caso}{(3) $\rightarrow$ (4)}
        Sia $f^{-1}$ l'inversa. Ogni sezione $s$ è uguale alla retrazione $f^{-1}$, e ogni retrazione $r$ è uguale alla sezione
        $f^{-1}$ (Teorema~\ref{teo:unicita-inversa}): quindi c'è una sola sezione e una sola retrazione.
    \end{caso}
    \begin{caso}{(4) $\rightarrow$ (2)}
        Ovvio: se ha una sola sezione e una sola retrazione, in particolare ha sezioni e retrazioni.
    \end{caso}
    \begin{caso}{(3) $\rightarrow$ (6)}
        Come sopra, ogni sezione è uguale a $f^{-1}$, quindi la sezione è una sola.
    \end{caso}
    \begin{caso}{(6) $\rightarrow$ (1)}
    $f$ ha una sezione, quindi è suriettiva (Teorema~\ref{teo:suriettiva-sezione}).
    Per assurdo, $f$ non sia iniettiva: siano $x_1 \neq x_2$ elementi di $a$ con $f(x_1) = f(x_2) =: z$, e sia $s$ l'unica sezione di $f$.
    Dalla definizione di sezione $f(s(z)) = z$, cioè $s(z) \in \overleftarrow{f}(\{z\})$, che contiene almeno due elementi distinti.
    Posso quindi supporre, senza ledere la generalità, che $s(z) \neq x_1$ (altrimenti scambio $x_1$ e $x_2$). Definisco
    \[ g \colon y \in b \mapsto \begin{cases} s(y) & \text{se } y \neq z \\ x_1 & \text{se } y = z \end{cases} \in a. \]
    Per ogni $y \neq z$ si ha $f(g(y)) = f(s(y)) = y$, e $f(g(z)) = f(x_1) = z$: quindi $g$ è una sezione di $f$, diversa da $s$
    perché $g(z) = x_1 \neq s(z)$. Assurdo, perché la sezione era unica. Quindi $f$ è iniettiva, dunque biettiva.
    \end{caso}
\end{dimostrazione}
\end{concetto}

\begin{spiegazione}
    La mappa della dimostrazione: ogni freccia è un passaggio dimostrato, e da ogni condizione si arriva a ogni altra seguendo le frecce.
    \begin{center}
    \begin{tikzpicture}[font=\small, cond/.style={circle, draw=primary!70!black, fill=primary!8, minimum size=7mm, inner sep=0pt}]
        \node[cond] (1) at (0,0) {$1$}; \node[cond] (3) at (2.6,0) {$3$}; \node[cond] (5) at (-2.6,0) {$5$};
        \node[cond] (2) at (4.6,1.0) {$2$}; \node[cond] (4) at (6.6,0) {$4$}; \node[cond] (6) at (2.6,-1.5) {$6$};
        \draw[{Stealth}-{Stealth}, line width=0.8pt] (1) -- (3); \draw[{Stealth}-{Stealth}, line width=0.8pt] (1) -- (5);
        \draw[{Stealth}-{Stealth}, line width=0.8pt] (3) -- (2);
        \draw[-{Stealth}, line width=0.8pt] (3) -- (4); \draw[-{Stealth}, line width=0.8pt] (4) -- (2);
        \draw[-{Stealth}, line width=0.8pt] (3) -- (6); \draw[-{Stealth}, line width=0.8pt] (6) -- (1);
        \node[align=left, anchor=west] at (7.6,0) {1: biettiva \quad 2: sezioni e retrazioni\\ 3: inversa \quad 4: una sola sezione e retrazione\\
            5: ogni fibra ha un elemento \quad 6: una sola sezione};
    \end{tikzpicture}
    \end{center}
\end{spiegazione}



  % Parte 3: Strutture algebriche (semigruppi, monoidi, gruppi, anelli, campi)

  \part{Strutture algebriche}

  % Operazioni, semigruppi, monoidi, gruppi, sottostrutture, cancellabilità, Cayley, omomorfismi
  \chapter{Semigruppi, monoidi e gruppi}\label{cap:gruppi}

Una \textbf{struttura algebrica} è un insieme dotato di una o più operazioni. In questo capitolo studiamo le
strutture con \textbf{una sola operazione binaria}: semigruppi, monoidi e gruppi, le loro sottostrutture e le funzioni
che ne rispettano l'operazione (gli omomorfismi).

% ──────────────────────────────────────────────────────────────
\section{Operazioni e strutture algebriche}\label{sec:operazioni-interne}

\dfn[operazione-interna]{Operazione interna}{
    Sia $s$ un insieme non vuoto. Una funzione
    \[ * \colon \underbrace{s \times \dots \times s}_{n \text{ volte}} \to s \]
    si dice \textbf{operazione $n$-aria interna} di $s$. Per $n = 2$ si parla di \textbf{operazione binaria interna}
    e, se $x, y \in s$, scriveremo $x * y := *(x, y)$, abbandonando la notazione funzionale.
}

\dfn[commutativa-associativa]{Operazione commutativa e associativa}{
    Un'operazione binaria interna $*$ di $s$ si dice
    \begin{enumerate}
        \item \textbf{commutativa} se $\forall a, b \in s\,(a * b = b * a)$;
        \item \textbf{associativa} se $\forall a, b, c \in s\,\big(a * (b * c) = (a * b) * c\big)$.
    \end{enumerate}
}

\dfn[struttura-algebrica]{Struttura algebrica}{
    Siano $*_1, \dots, *_n$ operazioni interne di $s$. La $(n+1)$-upla $(s, *_1, \dots, *_n)$ si dice
    \textbf{struttura algebrica} ad $n$ operazioni. In particolare $(s, *)$, con $*$ binaria, è una struttura algebrica
    ad un'operazione binaria interna.
}

\dfn[semigruppo]{Semigruppo}{
    Sia $(s, *)$ una struttura algebrica con un'operazione binaria interna. $(s, *)$ si dice \textbf{semigruppo} se $*$ è
    associativa; si dice \textbf{semigruppo commutativo} se $*$ è sia associativa che commutativa.
}

% ──────────────────────────────────────────────────────────────
\section{Elemento neutro e monoidi}\label{sec:neutro}

\dfn[neutro]{Elemento neutro}{
    Sia $(s, *)$ una struttura algebrica con un'operazione binaria interna.
    \begin{itemize}
        \item $l \in s$ tale che $\forall x \in s\,(l * x = x)$ si dice \textbf{elemento neutro sinistro};
        \item $r \in s$ tale che $\forall x \in s\,(x * r = x)$ si dice \textbf{elemento neutro destro};
        \item $u \in s$ che sia neutro sia a sinistra sia a destra si dice \textbf{elemento neutro}.
    \end{itemize}
}

\begin{concetto}
\begin{Teorema}{Unicità dell'elemento neutro}{unicita-neutro}
    \begin{prerequisiti}
        \dip{Definizione}{def:neutro}
        \dip{Definizione}{def:semigruppo}
    \end{prerequisiti}
    Sia $(s, *)$ un semigruppo. Siano $l$ neutro a sinistra e $r$ neutro a destra. Allora $l = r$.
    In particolare, il neutro, se esiste, è unico.
\end{Teorema}
\begin{dimostrazione}
    \begin{caso}{$l = r$}
        $l = l * r = r$: la prima uguaglianza perché $r$ è neutro a destra, la seconda perché $l$ è neutro a sinistra.
    \end{caso}
    \begin{caso}{Il neutro è unico}
        Se $u$ e $u'$ sono due elementi neutri, $u$ è in particolare neutro a sinistra e $u'$ neutro a destra, quindi $u = u'$.
    \end{caso}
\end{dimostrazione}
\end{concetto}

\dfn[monoide]{Monoide}{
    Un semigruppo $(m, *)$ dotato di elemento neutro si dice \textbf{monoide}; per evidenziare il neutro scriveremo
    $(m, *, 1_m)$, dove $1_m$ è l'elemento neutro. Se $*$ è anche commutativa, il monoide si dice \textbf{commutativo}.
}

\begin{esempio}{Monoidi e non}
    \begin{itemize}
        \item $(\mbN, +, 0)$, $(\mbN, \cdot, 1)$, $(\mbZ, +, 0)$, $(\mbZ, \cdot, 1)$ sono tutti monoidi commutativi.
        \item Dato un insieme $a$, $(\mcP(a), \cap, a)$, $(\mcP(a), \cup, \varnothing)$ e $(\mcP(a), \Delta, \varnothing)$ sono monoidi
              commutativi, per le proprietà dimostrate nella Proposizione~\ref{prop:prop-operazioni} (e perché per ogni $x \subseteq a$
              si ha $x \cap a = x$, $x \cup \varnothing = x$, $x \,\Delta\, \varnothing = x$).
        \item $(\mbZ, -)$ \textbf{non} è un semigruppo: $(1 - 1) - 1 = -1 \neq 1 = 1 - (1 - 1)$. Ha $0$ come elemento neutro destro
              ($x - 0 = x$), che però non è neutro sinistro ($0 - x = -x$).
        \item Per un qualunque insieme non vuoto $a$, $(\mathrm{In}(a, a), \circ, \mathrm{id}_a)$, $(\mathrm{Su}(a, a), \circ, \mathrm{id}_a)$ e
              $(\mathrm{Bi}(a, a), \circ, \mathrm{id}_a)$ sono monoidi (Teoremi~\ref{teo:comp-iniettive}, \ref{teo:comp-suriettive},
              Corollario~\ref{cor:comp-biettive} e Proposizione~\ref{prop:assoc-prodotto}).
    \end{itemize}
\end{esempio}

\dfn[parole]{Monoide delle parole}{
    Sia $A$ un insieme di simboli (ad esempio $A = \{a, b, c\}$), detto \textbf{alfabeto}. L'insieme $s$ delle \textbf{parole} su $A$
    è l'insieme delle successioni finite di elementi di $A$, compresa la parola vuota $\varnothing$:
    \[ s = \{x \mid \exists n \in \mbN \setminus \{0\}\ \exists x_1, \dots, x_n \in A\,(x = x_1 \dots x_n)\} \cup \{\varnothing\}. \]
    La \textbf{concatenazione} $* \colon (x, y) \in s \times s \mapsto xy \in s$ (la parola $x$ seguita dalla parola $y$), con la
    convenzione $\varnothing * x = x = x * \varnothing$, rende $(s, *, \varnothing)$ un monoide.
}

\begin{concetto}
\begin{Proposizione}{Quando il monoide delle parole è commutativo}{parole-commutativo}
    \begin{prerequisiti}
        \dip{Definizione}{def:parole}
        \dip{Definizione}{def:commutativa-associativa}
    \end{prerequisiti}
    Sia $A$ un alfabeto non vuoto. Il monoide delle parole su $A$ è commutativo se e solo se $A = \{a\}$ ha un solo elemento.
\end{Proposizione}
\begin{dimostrazione}
    \begin{caso}{$(\leftarrow)$ Se $A = \{a\}$, allora il monoide è commutativo}
        Le parole sono formate solo dalla lettera $a$: $x = a \dots a$ ($n$ volte) e $y = a \dots a$ ($k$ volte), quindi sia
        $xy$ sia $yx$ sono la parola formata da $n + k$ lettere $a$.
    \end{caso}
    \begin{caso}{$(\rightarrow)$ Se il monoide è commutativo, allora $A$ ha un solo elemento}
        Per contrapposizione: se $A$ contiene almeno due lettere distinte $a$ e $b$, allora $a * b = ab \neq ba = b * a$,
        quindi il monoide non è commutativo.
    \end{caso}
\end{dimostrazione}
\end{concetto}

\dfn[duale]{Operazione duale (o opposta)}{
    Sia $(s, *)$ una struttura algebrica con un'operazione binaria interna. Definisco
    \[ \bar{*} \colon (a, b) \in s \times s \mapsto b * a \in s; \]
    $\bar{*}$ si dice \textbf{operazione duale} (o \textbf{opposta}) di $*$.
}

\begin{concetto}
\begin{Proposizione}{Proprietà dell'operazione duale}{duale}
    \begin{prerequisiti}
        \dip{Definizione}{def:duale}
        \dip{Definizione}{def:commutativa-associativa}
        \dip{Definizione}{def:neutro}
        \dip{Teorema}{teo:biettiva-inversa}
    \end{prerequisiti}
    Sia $(s, *)$ una struttura algebrica con un'operazione binaria interna.
    \begin{enumerate}
        \item $*$ è commutativa se e solo se lo è $\bar{*}$; $*$ è associativa se e solo se lo è $\bar{*}$.
        \item $l$ è neutro a sinistra in $(s, *)$ se e solo se è neutro a destra in $(s, \bar{*})$ (e viceversa).
        \item $\bar{\bar{*}} = *$; quindi la funzione $* \mapsto \bar{*}$, dall'insieme $\mathrm{Map}(s \times s, s)$ delle operazioni
              binarie interne di $s$ in se stesso, è l'inversa di sé stessa ed è biettiva.
    \end{enumerate}
\end{Proposizione}
\begin{dimostrazione}
    \begin{caso}{Punto 1}
        Commutatività: $a \,\bar{*}\, b = b \,\bar{*}\, a \leftrightarrow b * a = a * b$. Associatività: $a \,\bar{*}\, (b \,\bar{*}\, c) = (c * b) * a$
        e $(a \,\bar{*}\, b) \,\bar{*}\, c = c * (b * a)$, quindi l'associatività di $\bar{*}$ per $a, b, c$ è l'associatività di $*$ per $c, b, a$.
    \end{caso}
    \begin{caso}{Punto 2}
        $\forall x\,(l * x = x) \leftrightarrow \forall x\,(x \,\bar{*}\, l = x)$, per definizione di $\bar{*}$.
    \end{caso}
    \begin{caso}{Punto 3}
        $a \,\bar{\bar{*}}\, b = b \,\bar{*}\, a = a * b$. Una funzione che è l'inversa di sé stessa è biettiva (Teorema~\ref{teo:biettiva-inversa}).
    \end{caso}
\end{dimostrazione}
\end{concetto}

% ──────────────────────────────────────────────────────────────
\section{Elementi invertibili e gruppi}\label{sec:invertibili}

\dfn[invertibile]{Elemento invertibile e inverso (simmetrico)}{
    Sia $(m, *, 1_m)$ un monoide e sia $x \in m$.
    \begin{enumerate}
        \item Se esiste $l \in m$ con $l * x = 1_m$, $x$ lo dico \textbf{invertibile a sinistra} ed $l$ è detto
              \textbf{inverso (simmetrico) sinistro} di $x$.
        \item Se esiste $r \in m$ con $x * r = 1_m$, $x$ lo dico \textbf{invertibile a destra} ed $r$ è detto
              \textbf{inverso (simmetrico) destro} di $x$.
    \end{enumerate}
    Un elemento che è inverso sia destro sia sinistro di $x$ si dice \textbf{inverso} (o \textbf{simmetrico}) di $x$, e $x$ si dice
    \textbf{invertibile} (o \textbf{simmetrizzabile}).
}

\begin{concetto}
\begin{Teorema}{Unicità dell'inverso}{unicita-inverso}
    \begin{prerequisiti}
        \dip{Definizione}{def:invertibile}
        \dip{Definizione}{def:monoide}
    \end{prerequisiti}
    Sia $(m, *, 1_m)$ un monoide e sia $x \in m$ con un inverso sinistro $l$ e un inverso destro $r$. Allora $l = r$.
    In particolare, l'inverso, se esiste, è unico.
\end{Teorema}
\begin{dimostrazione}
    \begin{caso}{$l = r$}
        Usando l'associatività:
        \[ l = l * 1_m = l * (x * r) = (l * x) * r = 1_m * r = r. \]
    \end{caso}
    \begin{caso}{L'inverso è unico}
        Se $y$ e $y'$ sono due inversi di $x$, $y$ è in particolare inverso sinistro e $y'$ inverso destro, quindi $y = y'$.
    \end{caso}
\end{dimostrazione}
\end{concetto}

\info{Notazione}{
    L'inverso di $x$ lo scrivo $x^{-1}$. Se l'operazione si indica con $+$ (notazione additiva, ad esempio in $(\mbZ, +, 0)$)
    si scrive $-x$ e si chiama \textbf{opposto} di $x$.
}

\dfn[gruppo]{Gruppo, gruppo abeliano}{
    Un monoide in cui tutti gli elementi sono invertibili si chiama \textbf{gruppo}. Un gruppo commutativo si chiama
    \textbf{gruppo abeliano}.
}

\dfn[unita]{Insieme degli elementi invertibili}{
    Sia $(m, *, 1_m)$ un monoide. Scrivo $U(m)$ l'insieme di tutti gli \textbf{elementi invertibili} di $m$.
}

\begin{concetto}
\begin{Proposizione}{Gli invertibili formano un gruppo}{unita-gruppo}
    \begin{prerequisiti}
        \dip{Definizione}{def:unita}
        \dip{Teorema}{teo:unicita-inverso}
        \dip{Definizione}{def:gruppo}
    \end{prerequisiti}
    Sia $(m, *, 1_m)$ un monoide e siano $a, b \in U(m)$. Allora $a * b \in U(m)$ e $(a * b)^{-1} = b^{-1} * a^{-1}$.
    Di conseguenza $U(m)$, con l'operazione $*$, è un gruppo.
\end{Proposizione}
\begin{dimostrazione}
    \begin{caso}{$b^{-1} * a^{-1}$ è l'inverso di $a * b$}
        $(a * b) * (b^{-1} * a^{-1}) = a * (b * b^{-1}) * a^{-1} = a * a^{-1} = 1_m$ e, allo stesso modo,
        $(b^{-1} * a^{-1}) * (a * b) = b^{-1} * (a^{-1} * a) * b = b^{-1} * b = 1_m$.
    \end{caso}
    \begin{caso}{$U(m)$ è un gruppo}
        L'operazione resta associativa, $1_m \in U(m)$ (perché $1_m * 1_m = 1_m$) e ogni $a \in U(m)$ ha inverso $a^{-1}$, che è a sua volta
        invertibile (con inverso $a$), quindi sta in $U(m)$.
    \end{caso}
\end{dimostrazione}
\end{concetto}

\begin{esempio}{Elementi invertibili}
    $U(\mbN, +, 0) = \{0\}$; \ $U(\mbZ, \cdot, 1) = \{-1, 1\}$; \ se $a$ è un insieme non vuoto, $U(\mathrm{In}(a, a), \circ, \mathrm{id}_a) = \mathrm{Bi}(a, a)$
    (Teorema~\ref{teo:biettiva-inversa}), quindi $(\mathrm{Bi}(a, a), \circ)$ è un gruppo.
\end{esempio}

% ──────────────────────────────────────────────────────────────
\section{Parti chiuse e sottostrutture}\label{sec:sottostrutture}

\dfn[parte-chiusa]{Parte chiusa}{
    Sia $(s, *)$ una struttura algebrica con un'operazione binaria interna e sia $t \subseteq s$. $t$ si dice
    \textbf{parte chiusa} (o \textbf{stabile}) di $(s, *)$ se
    \[ \forall x, y \in t\,(x * y \in t). \]
}

\dfn[sottostruttura]{Operazione indotta, sottostruttura}{
    Se $t$ è una parte chiusa di $(s, *)$, definisco
    \[ *' \colon (x, y) \in t \times t \mapsto x * y \in t, \]
    detta \textbf{operazione indotta da $*$ su $t$}. $(t, *')$ la dico \textbf{sottostruttura} di $(s, *)$ e scrivo $(t, *') \leq (s, *)$
    (spesso, con abuso di notazione, $(t, *)$ o semplicemente $t \leq s$). Inoltre:
    \begin{itemize}
        \item se $(s, *)$ è un semigruppo, una sua sottostruttura si dice \textbf{sottosemigruppo};
        \item se $(m, *, 1_m)$ è un monoide, un \textbf{sottomonoide} è una sottostruttura che contiene $1_m$;
        \item se $(g, *)$ è un gruppo, un \textbf{sottogruppo} è un sottomonoide che contiene l'inverso di ogni suo elemento.
    \end{itemize}
}

\begin{esempio}{Sottostrutture}
    \begin{itemize}
        \item $\{0\}$ è un sottosemigruppo e un sottomonoide di $(\mbN, +, 0)$.
        \item $\{1\}$ è un sottomonoide di $(\mbN, \cdot, 1)$, ma non è una parte chiusa di $(\mbN, +)$: $1 + 1 = 2 \notin \{1\}$.
        \item Per $m \in \mbN$, $m\mbN = \{n \in \mbN \mid \exists k \in \mbN\,(n = km)\}$ è un sottomonoide di $(\mbN, +, 0)$.
        \item Per $m \in \mbZ$, $m\mbZ = \{n \in \mbZ \mid \exists k \in \mbZ\,(n = km)\}$, l'insieme dei multipli interi di $m$, è un
              sottogruppo di $(\mbZ, +)$ e $(m\mbZ, \cdot)$ è un sottosemigruppo di $(\mbZ, \cdot)$. Ad esempio $0\mbZ = \{0\}$ e $1\mbZ = \mbZ$.
        \item Se $a$ è un insieme non vuoto, $(\mathrm{Bi}(a, a), \circ)$ è un sottomonoide di $(\mathrm{In}(a, a), \circ)$, ed è un gruppo.
    \end{itemize}
\end{esempio}

\begin{concetto}
\begin{Teorema}{Intersezione di parti chiuse}{intersezione-chiuse}
    \begin{prerequisiti}
        \dip{Definizione}{def:parte-chiusa}
        \dip{Definizione}{def:sottostruttura}
        \dip{Definizione}{def:intersezione}
    \end{prerequisiti}
    Sia $(s, *)$ una struttura e $f$ una famiglia di parti chiuse di $(s, *)$. Allora $\bigcap f$ è una parte chiusa di $(s, *)$.
    Inoltre, se $f$ è una famiglia (non vuota) di sottomonoidi (risp.\ sottogruppi), $\bigcap f$ è un sottomonoide (risp.\ sottogruppo).
\end{Teorema}
\begin{dimostrazione}
    \begin{caso}{$\bigcap f$ è una parte chiusa}
        Siano $x, y \in \bigcap f$; allora $\forall t \in f\,(x, y \in t)$. Ma per ipotesi ogni $t \in f$ è una parte chiusa, quindi
        $\forall t \in f\,(x * y \in t)$, cioè $x * y \in \bigcap f$.
    \end{caso}
    \begin{caso}{Sottomonoidi e sottogruppi}
        Se ogni $t \in f$ contiene $1_m$, anche $\bigcap f$ lo contiene. Se ogni $t \in f$ contiene l'inverso di ogni suo elemento e
        $x \in \bigcap f$, allora $x^{-1}$ sta in ogni $t \in f$, quindi $x^{-1} \in \bigcap f$.
    \end{caso}
\end{dimostrazione}
\end{concetto}

\dfn[generata]{Sottostruttura generata}{
    Sia $(s, *)$ una struttura algebrica e sia $t \subseteq s$. Dico \textbf{sottostruttura di $(s, *)$ generata da $t$}
    \[ \langle t \rangle := \bigcap \{c \in \mcP(s) \mid c \leq s \wedge t \subseteq c\}, \]
    l'intersezione delle sottostrutture di $s$ che contengono $t$ (dello stesso tipo: sottomonoidi se $s$ è un monoide, sottogruppi
    se $s$ è un gruppo). Per il Teorema~\ref{teo:intersezione-chiuse} è anch'essa una sottostruttura: è la \textbf{più piccola}
    sottostruttura di $s$ che contiene $t$. Si scrive $\langle x \rangle$ invece di $\langle \{x\} \rangle$, $\langle a, b \rangle$ invece di
    $\langle \{a, b\} \rangle$. Una struttura $s$ si dice \textbf{ciclica} se esiste $x \in s$ con $s = \langle x \rangle$.
}

\begin{esempio}{Sottostrutture generate}
    \begin{itemize}
        \item Se $(m, *, 1_m)$ è un monoide, $\langle \varnothing \rangle = \{1_m\}$.
        \item Il sottomonoide di $(\mbN, +, 0)$ generato da $1$ è $\mbN$; quello generato da $2$ è $2\mbN$.
        \item Il sottosemigruppo di $(\mbZ, +)$ generato da $2$ è $2\mbN \setminus \{0\}$.
        \item Il sottogruppo di $(\mbZ, +)$ generato da $2$ è $2\mbZ$; quello generato da $\{2, 3\}$ è $\mbZ$.
        \item $(\mbN, +, 0)$ è un monoide ciclico ($\mbN = \langle 1 \rangle$) e $(\mbZ, +)$ è un gruppo ciclico ($\mbZ = \langle 1 \rangle$).
    \end{itemize}
\end{esempio}

\begin{concetto}
\begin{Teorema}{Caratterizzazione dei sottomonoidi generati}{sottomonoidi-generati}
    \begin{prerequisiti}
        \dip{Definizione}{def:generata}
        \dip{Teorema}{teo:intersezione-chiuse}
        \dip{Proposizione}{prop:doppia-inclusione}
        \dip{Definizione}{def:sottostruttura}
    \end{prerequisiti}
    Sia $(m, *, 1_m)$ un monoide e sia $t \in \mcP(m) \setminus \{\varnothing\}$. Allora
    \[ \langle t \rangle = \{x \in m \mid \exists n \in \mbN \setminus \{0\}\ \exists x_1, \dots, x_n \in t\,(x = x_1 * \dots * x_n)\} \cup \{1_m\}, \]
    cioè l'insieme delle «combinazioni» (prodotti finiti) di elementi di $t$, più il neutro.
\end{Teorema}
\begin{dimostrazione}
    Chiamo $S$ l'insieme a destra e, tramite la doppia inclusione, dimostro che $\langle t \rangle = S$. Ricordiamo che
    $\langle t \rangle = \bigcap\{c \in \mcP(m) \mid t \subseteq c \wedge c \leq m\}$.
    \begin{caso}{$(\subseteq)$ $\langle t \rangle \subseteq S$}
        Basta vedere che $S$ è un sottomonoide che contiene $t$: allora è uno degli insiemi che si intersecano, e l'intersezione
        è contenuta in esso.
        Se prendo $n = 1$, ho $t \subseteq S$; inoltre $1_m \in S$. Prendo $x_1 * \dots * x_{n_1}$ e $y_1 * \dots * y_{n_2}$ in $S$: allora
        $x_1 * \dots * x_{n_1} * y_1 * \dots * y_{n_2} \in S$ (è ancora un prodotto finito di elementi di $t$; se uno dei due è $1_m$ il prodotto
        è l'altro). Quindi $S \leq m$ e $\langle t \rangle \subseteq S$.
    \end{caso}
    \begin{caso}{$(\supseteq)$ $S \subseteq \langle t \rangle$}
        Prendo un qualunque sottomonoide $e \leq m$ con $t \subseteq e$. Allora $1_m \in e$, e se $x_1, \dots, x_n \in t$ allora
        $x_1, \dots, x_n \in e$, quindi $x_1 * \dots * x_n \in e$ poiché $e$ è una parte chiusa. Quindi $S \subseteq e$; per la generalità
        di $e$, $S$ è contenuto nell'intersezione, cioè $S \subseteq \langle t \rangle$.
    \end{caso}
\end{dimostrazione}
\end{concetto}

\begin{concetto}
\begin{Teorema}{Caratterizzazione dei sottogruppi generati}{sottogruppi-generati}
    \begin{prerequisiti}
        \dip{Teorema}{teo:sottomonoidi-generati}
        \dip{Definizione}{def:gruppo}
    \end{prerequisiti}
    Sia $(g, *)$ un gruppo e sia $t \in \mcP(g) \setminus \{\varnothing\}$. Allora
    \[ \langle t \rangle = \{x \in g \mid \exists n \in \mbN \setminus \{0\}\ \exists \varepsilon_1, \dots, \varepsilon_n \in \{-1, 1\}\ \exists x_1, \dots, x_n \in t\,(x = x_1^{\varepsilon_1} * \dots * x_n^{\varepsilon_n})\} \cup \{1_g\}. \]
\end{Teorema}
\begin{dimostrazione}[Nota]
    Negli appunti questo teorema è dato senza dimostrazione. L'idea è la stessa del Teorema~\ref{teo:sottomonoidi-generati}:
    in un sottogruppo, oltre ai prodotti di elementi di $t$, ci sono anche i loro inversi ($x^{-1}$, cioè esponente $-1$).
\end{dimostrazione}
\end{concetto}

% ──────────────────────────────────────────────────────────────
\section{Elementi cancellabili}\label{sec:cancellabili}

\dfn[cancellabile]{Elemento cancellabile}{
    Sia $(s, *)$ un semigruppo. Un elemento $a \in s$ si dice
    \begin{itemize}
        \item \textbf{cancellabile a sinistra} se $\forall b, c \in s\,(a * b = a * c \rightarrow b = c)$;
        \item \textbf{cancellabile a destra} se $\forall b, c \in s\,(b * a = c * a \rightarrow b = c)$;
        \item \textbf{cancellabile} se lo è sia a sinistra che a destra.
    \end{itemize}
}

\dfn[traslazioni]{Traslazione sinistra e destra}{
    Sia $(s, *)$ un semigruppo e $a \in s$.
    \begin{enumerate}
        \item \textbf{Traslazione sinistra}: $\sigma_a \colon x \in s \mapsto a * x \in s$.
        \item \textbf{Traslazione destra}: $\delta_a \colon x \in s \mapsto x * a \in s$.
    \end{enumerate}
}

\begin{concetto}
\begin{Proposizione}{Cancellabilità e traslazioni}{traslazioni}
    \begin{prerequisiti}
        \dip{Definizione}{def:cancellabile}
        \dip{Definizione}{def:traslazioni}
        \dip{Proposizione}{prop:forme-iniettivita}
        \dip{Definizione}{def:invertibile}
    \end{prerequisiti}
    Sia $(s, *)$ un semigruppo e $a \in s$.
    \begin{enumerate}
        \item $a$ è cancellabile a sinistra se e solo se $\sigma_a$ è iniettiva.
        \item $a$ è cancellabile a destra se e solo se $\delta_a$ è iniettiva.
        \item Se $s$ è un monoide e $a$ è invertibile, $\sigma_{a^{-1}}$ è l'inversa di $\sigma_a$ e $\delta_{a^{-1}}$ è l'inversa di $\delta_a$.
    \end{enumerate}
\end{Proposizione}
\begin{dimostrazione}
    \begin{caso}{Punto 1}
        $\sigma_a$ iniettiva vuol dire $\forall b, c\,(\sigma_a(b) = \sigma_a(c) \rightarrow b = c)$, cioè $\forall b, c\,(a * b = a * c \rightarrow b = c)$,
        che è la definizione di cancellabile a sinistra.
    \end{caso}
    \begin{caso}{Punto 2}
        Identico, con $\delta_a(b) = b * a$.
    \end{caso}
    \begin{caso}{Punto 3}
        $\sigma_{a^{-1}}(\sigma_a(x)) = a^{-1} * (a * x) = x$ e $\sigma_a(\sigma_{a^{-1}}(x)) = a * (a^{-1} * x) = x$; per $\delta$ è analogo.
    \end{caso}
\end{dimostrazione}
\end{concetto}

\begin{concetto}
\begin{Teorema}{Gli invertibili sono cancellabili}{invertibile-cancellabile}
    \begin{prerequisiti}
        \dip{Definizione}{def:invertibile}
        \dip{Definizione}{def:cancellabile}
    \end{prerequisiti}
    Sia $(m, *, 1_m)$ un monoide. Se $a \in m$ è invertibile a sinistra, allora $a$ è cancellabile a sinistra (e lo stesso a destra).
    In particolare, ogni elemento invertibile è cancellabile.
\end{Teorema}
\begin{dimostrazione}
    Per ipotesi trovo $b \in m$ con $b * a = 1_m$. Se $a * c = a * d$, allora
    \[ c = 1_m * c = b * a * c = b * a * d = 1_m * d = d. \]
    A destra è identico, usando un inverso destro.
\end{dimostrazione}
\end{concetto}

Il viceversa in generale è falso: in $(\mbZ, \cdot, 1)$ ogni $n \neq 0$ è cancellabile, ma solo $-1$ e $1$ sono invertibili.
Diventa vero se il monoide è \textbf{finito}: lo dimostreremo nel Teorema~\ref{teo:monoide-finito}, dopo aver introdotto gli
insiemi finiti (Capitolo~\ref{cap:finiti}).

% ──────────────────────────────────────────────────────────────
\section{Tavole di Cayley}\label{sec:cayley}

\dfn[cayley]{Tavola di Cayley}{
    Sia $(s, *)$ una struttura con $s$ finito. La \textbf{tavola di Cayley} si ottiene scrivendo gli elementi di $s$, nello stesso ordine,
    su una riga e su una colonna, e mettendo all'incrocio tra la riga di $a$ e la colonna di $b$ il risultato $a * b$.
}

\begin{esempio}{La tavola di $(t, \cdot)$ con $t = \{-1, 0, 1\}$}
    \begin{center}
    \renewcommand{\arraystretch}{1.3}
    \begin{tabular}{c|ccc}
        $\cdot$ & $-1$ & $0$ & $1$ \\ \hline
        $-1$ & \cellcolor{pink!45}$1$ & $0$ & \cellcolor{azure-gradient-2}$-1$ \\
        $0$  & $0$ & \cellcolor{pink!45}$0$ & \cellcolor{azure-gradient-2}$0$ \\
        $1$  & \cellcolor{green-gradient-3}$-1$ & \cellcolor{green-gradient-3}$0$ & \cellcolor{pink!45}$1$ \\
    \end{tabular}
    \end{center}
    \begin{itemize}
        \item \colorbox{green-gradient-3}{Riga di $1$}: ripete la prima riga, quindi $1$ è \textbf{neutro a sinistra};
              \colorbox{azure-gradient-2}{colonna di $1$}: ripete la prima colonna, quindi $1$ è \textbf{neutro a destra}.
        \item La tavola è simmetrica rispetto alla \colorbox{pink!45}{diagonale}: l'operazione è \textbf{commutativa}.
        \item Nella riga (e nella colonna) di $-1$ compare il neutro $1$: $-1$ è \textbf{invertibile} (è l'inverso di sé stesso).
        \item Nella riga di $0$ non compare $1$ e ci sono elementi ripetuti: $0$ non è né invertibile né cancellabile.
    \end{itemize}
\end{esempio}

\begin{concetto}
\begin{Proposizione}{Come si legge una tavola di Cayley}{lettura-cayley}
    \begin{prerequisiti}
        \dip{Definizione}{def:cayley}
        \dip{Definizione}{def:neutro}
        \dip{Definizione}{def:invertibile}
        \dip{Proposizione}{prop:traslazioni}
        \dip{Definizione}{def:commutativa-associativa}
    \end{prerequisiti}
    Sia $(s, *)$ con $s$ finito e sia $a \in s$. Nella tavola di Cayley:
    \begin{enumerate}
        \item $a$ è neutro a sinistra $\leftrightarrow$ sulla riga di $a$ si presentano gli stessi elementi di $s$, nello stesso ordine;
        \item $a$ è neutro a destra $\leftrightarrow$ sulla colonna di $a$ si presentano gli stessi elementi di $s$, nello stesso ordine;
        \item $a$ è simmetrizzabile a sinistra $\leftrightarrow$ nella sua colonna trovo l'elemento neutro;
        \item $a$ è simmetrizzabile a destra $\leftrightarrow$ nella sua riga trovo l'elemento neutro;
        \item $a$ è cancellabile a sinistra $\leftrightarrow$ sulla riga di $a$ ho tutti elementi distinti;
        \item $a$ è cancellabile a destra $\leftrightarrow$ sulla colonna di $a$ ho tutti elementi distinti;
        \item $*$ è commutativa $\leftrightarrow$ la tavola è simmetrica rispetto alla diagonale.
    \end{enumerate}
    (Nei punti 3 e 4 si suppone che $(s, *)$ sia un monoide.)
\end{Proposizione}
\begin{dimostrazione}
    Ricordiamo che la riga di $a$ contiene i valori $a * x$ (cioè $\sigma_a(x)$) e la colonna di $a$ i valori $x * a$ (cioè $\delta_a(x)$), per $x$ che scorre in $s$
    nell'ordine della tavola.
    \begin{caso}{Punti 1 e 2: neutri}
        La riga di $a$ ripete l'intestazione se e solo se $a * x = x$ per ogni $x$, cioè $a$ è neutro a sinistra. Per la colonna, $x * a = x$.
    \end{caso}
    \begin{caso}{Punti 3 e 4: simmetrizzabili}
        Nella colonna di $a$ c'è $1_s$ se e solo se esiste $x$ con $x * a = 1_s$, cioè $a$ ha un inverso sinistro. Nella riga, $a * x = 1_s$: inverso destro.
    \end{caso}
    \begin{caso}{Punti 5 e 6: cancellabili}
        La riga di $a$ ha tutti elementi distinti se e solo se $\sigma_a$ è iniettiva, cioè (Proposizione~\ref{prop:traslazioni}) $a$ è cancellabile
        a sinistra. Per la colonna si usa $\delta_a$.
    \end{caso}
    \begin{caso}{Punto 7: commutatività}
        L'elemento in posizione (riga $x$, colonna $y$) è $x * y$, quello in posizione simmetrica (riga $y$, colonna $x$) è $y * x$:
        coincidono tutti se e solo se $*$ è commutativa.
    \end{caso}
\end{dimostrazione}
\end{concetto}

% ──────────────────────────────────────────────────────────────
\section{Omomorfismi}\label{sec:omomorfismi}

\dfn[omomorfismo]{Omomorfismo tra strutture algebriche}{
    Siano $(s, *)$ e $(s', *')$ strutture algebriche con un'operazione binaria interna. Una funzione $\varphi \colon s \to s'$ si dice
    \textbf{omomorfismo} da $s$ a $s'$ se
    \[ \forall a, b \in s\,\big(\varphi(a * b) = \varphi(a) *' \varphi(b)\big). \]
    Se un omomorfismo è
    \begin{enumerate}
        \item iniettivo, si dice \textbf{monomorfismo};
        \item suriettivo, si dice \textbf{epimorfismo};
        \item biettivo, si dice \textbf{isomorfismo}.
    \end{enumerate}
}

\begin{esempio}{Omomorfismi}
    \begin{itemize}
        \item $m \in \mbZ \mapsto 2m \in \mbZ$ è un monomorfismo di $(\mbZ, +)$ in sé stesso: $2(m + n) = 2m + 2n$.
        \item $m \in \mbZ \mapsto -m \in \mbZ$ è un isomorfismo di $(\mbZ, +)$ in sé stesso.
        \item $\exp \colon x \in \mbR \mapsto e^x \in \mbR^+$ è un isomorfismo da $(\mbR, +)$ a $(\mbR^+, \cdot)$: $e^{x + y} = e^x e^y$.
        \item Se $a$ è un insieme, $x \in \mcP(a) \mapsto a \setminus x \in \mcP(a)$ è un isomorfismo tra $(\mcP(a), \cup)$ e $(\mcP(a), \cap)$:
              è biettiva e, per De Morgan (Proposizione~\ref{prop:demorgan-ins}), $a \setminus (x \cup y) = (a \setminus x) \cap (a \setminus y)$.
    \end{itemize}
\end{esempio}

\begin{concetto}
\begin{Teorema}{L'inversa di un isomorfismo è un isomorfismo}{inverso-isomorfismo}
    \begin{prerequisiti}
        \dip{Definizione}{def:omomorfismo}
        \dip{Teorema}{teo:biettiva-inversa}
    \end{prerequisiti}
    Sia $\varphi \colon s \to s'$ un isomorfismo tra $(s, *)$ e $(s', *')$. Allora $\varphi^{-1}$ è un isomorfismo tra $(s', *')$ e $(s, *)$.
\end{Teorema}
\begin{dimostrazione}
    $\varphi^{-1}$ è biettiva (ha come inversa $\varphi$). Siano $a, b \in s'$:
    \[ \varphi\big(\varphi^{-1}(a) * \varphi^{-1}(b)\big) = \varphi(\varphi^{-1}(a)) *' \varphi(\varphi^{-1}(b)) = a *' b; \]
    applicando $\varphi^{-1}$ a entrambi i membri ho $\varphi^{-1}(a) * \varphi^{-1}(b) = \varphi^{-1}(a *' b)$.
\end{dimostrazione}
\end{concetto}

\begin{concetto}
\begin{Teorema}{Gli epimorfismi conservano le proprietà}{epimorfismi}
    \begin{prerequisiti}
        \dip{Definizione}{def:omomorfismo}
        \dip{Definizione}{def:commutativa-associativa}
        \dip{Definizione}{def:neutro}
        \dip{Definizione}{def:invertibile}
        \dip{Definizione}{def:cancellabile}
    \end{prerequisiti}
    Sia $\varphi \colon s \to s'$ un epimorfismo tra $(s, *)$ e $(s', *')$. Allora in $(s', *')$ vengono conservate associatività,
    commutatività, neutri e simmetrici:
    \begin{enumerate}
        \item se $*$ è commutativa, $*'$ è commutativa;
        \item se $*$ è associativa, $*'$ è associativa;
        \item se $z$ è neutro destro/sinistro di $(s, *)$, allora $\varphi(z)$ è neutro destro/sinistro di $(s', *')$;
        \item se $x$ è simmetrizzabile a destra/sinistra in $(s, *)$, allora $\varphi(x)$ lo è in $(s', *')$;
        \item se $\varphi$ è anche iniettiva (isomorfismo) e $x$ è cancellabile a destra/sinistra, allora $\varphi(x)$ lo è.
    \end{enumerate}
\end{Teorema}
\begin{dimostrazione}
    In tutti i punti, dati elementi di $s'$, li scriviamo come immagini di elementi di $s$ (possiamo perché $\varphi$ è suriettiva).
    \begin{caso}{Punto 1: commutatività}
        Siano $a, b \in s'$ e $x, y \in s$ con $\varphi(x) = a$, $\varphi(y) = b$. Allora
        $a *' b = \varphi(x) *' \varphi(y) = \varphi(x * y) = \varphi(y * x) = \varphi(y) *' \varphi(x) = b *' a$.
    \end{caso}
    \begin{caso}{Punto 2: associatività}
        Siano $a, b, c \in s'$ immagini di $x, y, z$. Allora
        $a *' (b *' c) = \varphi(x) *' \varphi(y * z) = \varphi(x * (y * z)) = \varphi((x * y) * z) = \varphi(x * y) *' \varphi(z) = (a *' b) *' c$.
    \end{caso}
    \begin{caso}{Punto 3: neutri}
        Sia $z$ neutro destro, $a \in s'$ e $x \in s$ con $\varphi(x) = a$. Allora $a *' \varphi(z) = \varphi(x) *' \varphi(z) = \varphi(x * z) = \varphi(x) = a$.
        A sinistra è identico.
    \end{caso}
    \begin{caso}{Punto 4: simmetrici}
        Sia $u$ il neutro di $s$ e $y$ con $x * y = u$. Allora $\varphi(x) *' \varphi(y) = \varphi(x * y) = \varphi(u)$, che per il punto 3 è il neutro
        di $s'$: $\varphi(x)$ è simmetrizzabile a destra. A sinistra è identico.
    \end{caso}
    \begin{caso}{Punto 5: cancellabili}
        Sia $x$ cancellabile a sinistra, $b, c \in s'$ con $\varphi(x) *' b = \varphi(x) *' c$, e $y, z$ con $\varphi(y) = b$, $\varphi(z) = c$. Allora
        $\varphi(x * y) = \varphi(x) *' b = \varphi(x) *' c = \varphi(x * z)$; poiché $\varphi$ è iniettiva, $x * y = x * z$, e poiché $x$ è cancellabile
        a sinistra, $y = z$, quindi $b = c$.
    \end{caso}
\end{dimostrazione}
\end{concetto}


  % Anelli, LAP, divisori dello zero, corpi, campi, anelli booleani
  \chapter{Anelli e campi}\label{cap:anelli}

Passiamo alle strutture algebriche con \textbf{due} operazioni, pensate sul modello di $(\mbZ, +, \cdot)$: una somma e un prodotto,
legate dalla proprietà distributiva.

% ──────────────────────────────────────────────────────────────
\section{Anelli}\label{sec:anelli}

\dfn[anello]{Anello}{
    Sia $a$ un insieme non vuoto con due operazioni binarie interne $+$ e $\cdot$. $(a, +, \cdot)$ lo dico \textbf{anello} se
    \begin{enumerate}
        \item $(a, +)$ è un gruppo commutativo (con neutro $0_a$);
        \item $(a, \cdot)$ è un semigruppo;
        \item valgono le distributività: $\forall x, y, z \in a\,\big(x(y + z) = xy + xz \ \wedge\ (y + z)x = yx + zx\big)$.
    \end{enumerate}
    Un anello:
    \begin{enumerate}
        \item si dice \textbf{commutativo} se $(a, \cdot)$ è commutativo;
        \item si dice \textbf{unitario} se $(a, \cdot)$ è un monoide.
    \end{enumerate}
    Una parte di $a$ chiusa rispetto a $+$ e a $\cdot$ che sia un anello con le operazioni indotte si dice \textbf{sottoanello} di $a$.
}

\dfn[notazioni-anello]{Elemento neutro, inversi e potenze}{
    Sia $(a, +, \cdot)$ un anello.
    \begin{itemize}
        \item Dico $0_a$ il neutro di $(a, +)$ e $1_a$ il neutro di $(a, \cdot)$, se esiste (o $0$ e $1$ se non c'è ambiguità).
        \item Gli inversi di $(a, +)$ si chiamano opposti: $x + (-x) = 0_a$; l'inverso di $x$ in $(a, \cdot)$, se esiste, lo scrivo $x^{-1}$.
        \item $x - y := x + (-y)$.
        \item $\forall x \in a\ \forall n \in \mbN \setminus \{0\}\ \big(nx = \underbrace{x + \dots + x}_{n \text{ volte}}\big)$ e, per $n \in \mbZ \setminus \mbN$,
              $nx = \underbrace{(-x) + \dots + (-x)}_{-n \text{ volte}}$.
        \item $\forall x \in a\ \forall n \in \mbN \setminus \{0\}\ \big(x^n = \underbrace{x \cdot \ldots \cdot x}_{n \text{ volte}}\big)$ e, per $n \in \mbZ \setminus \mbN$
              (se $x$ è invertibile), $x^n = \underbrace{x^{-1} \cdot \ldots \cdot x^{-1}}_{-n \text{ volte}}$.
    \end{itemize}
}

\begin{esempio}{Anelli e non}
    \begin{itemize}
        \item $(\mbZ, +, \cdot)$ è un anello commutativo unitario.
        \item Se $s$ è un insieme, $(\mcP(s), \Delta, \cap)$ è un anello commutativo unitario: lo zero è $\varnothing$ e l'unità è $s$
              (segue dalle Proposizioni~\ref{prop:prop-operazioni} e \ref{prop:distr-delta}; l'opposto di $x$ è $x$ stesso, perché $x \,\Delta\, x = \varnothing$).
        \item $(\mcP(\mbN), \cup, \cap)$ non è un anello, perché $(\mcP(\mbN), \cup)$ non è un gruppo (solo $\varnothing$ ha un opposto).
        \item $(\mbN, +, \cdot)$ non è un anello, perché $(\mbN, +)$ non è un gruppo.
    \end{itemize}
\end{esempio}

\begin{concetto}
\begin{Lemma}{Regole di calcolo negli anelli}{regole-anello}
    \begin{prerequisiti}
        \dip{Definizione}{def:anello}
        \dip{Definizione}{def:notazioni-anello}
        \dip{Teorema}{teo:invertibile-cancellabile}
        \dip{Teorema}{teo:unicita-inverso}
    \end{prerequisiti}
    Sia $(a, +, \cdot)$ un anello e siano $x, y, z \in a$:
    \begin{enumerate}
        \item $0_a x = x 0_a = 0_a$;
        \item $(-x)y = -xy = x(-y)$;
        \item $x(y - z) = xy - xz$ (il prodotto è distributivo rispetto alla differenza);
        \item se $a$ è unitario e ha almeno due elementi, allora $0_a \neq 1_a$.
    \end{enumerate}
\end{Lemma}
\begin{dimostrazione}
    \begin{caso}{Punto 1: $0_a x = x 0_a = 0_a$}
        $0_a x + 0_a x = (0_a + 0_a) x = 0_a x = 0_a x + 0_a$ e, poiché in $(a, +)$ ogni elemento è cancellabile (è un gruppo), $0_a x = 0_a$.
        Allo stesso modo $x 0_a + x 0_a = x(0_a + 0_a) = x 0_a + 0_a$, quindi $x 0_a = 0_a$.
    \end{caso}
    \begin{caso}{Punto 2: $(-x)y = -xy = x(-y)$}
        $(-x)y + xy = (-x + x)y = 0_a y = 0_a$, quindi $(-x)y$ è l'opposto di $xy$: $(-x)y = -xy$.
        Allo stesso modo $x(-y) + xy = x(-y + y) = x 0_a = 0_a$, quindi $x(-y) = -xy$.
    \end{caso}
    \begin{caso}{Punto 3: $x(y - z) = xy - xz$}
        $x(y - z) = x(y + (-z)) = xy + x(-z) = xy - xz$, usando la definizione e il punto 2.
    \end{caso}
    \begin{caso}{Punto 4: $0_a \neq 1_a$}
        Se fosse $0_a = 1_a$, per ogni $x$ avremmo $x = 1_a x = 0_a x = 0_a$, cioè $a = \{0_a\}$ con un solo elemento.
    \end{caso}
\end{dimostrazione}
\end{concetto}

\dfn[omomorfismo-anelli]{Omomorfismo di anelli}{
    Siano $(a, +, \cdot)$ e $(a', +', \cdot')$ anelli. Una funzione $f \colon a \to a'$ è un \textbf{omomorfismo di anelli} se
    \begin{enumerate}
        \item $\forall x, y \in a\,\big(f(x + y) = f(x) +' f(y)\big)$;
        \item $\forall x, y \in a\,\big(f(x \cdot y) = f(x) \cdot' f(y)\big)$;
        \item se $a$ è unitario, $f(1_a)$ è l'elemento neutro di $(a', \cdot')$.
    \end{enumerate}
    Allo stesso modo si definiscono gli omomorfismi tra strutture con più operazioni. Ad esempio, per De Morgan,
    $x \in \mcP(s) \mapsto s \setminus x \in \mcP(s)$ è un omomorfismo tra $(\mcP(s), \cup, \cap)$ e $(\mcP(s), \cap, \cup)$.
}

% ──────────────────────────────────────────────────────────────
\section{Legge di annullamento del prodotto e divisori dello zero}\label{sec:lap}

\dfn[lap]{Legge di annullamento del prodotto}{
    In $(a, +, \cdot)$ si dice che vale la \textbf{legge di annullamento del prodotto} (LAP) se
    \[ \forall x, y \in a\,\big(xy = 0_a \rightarrow x = 0_a \vee y = 0_a\big). \]
    \begin{itemize}
        \item Un anello in cui vale la LAP si dice \textbf{integro}.
        \item Un anello commutativo unitario, con $0_a \neq 1_a$, in cui vale la LAP si dice \textbf{dominio di integrità}.
    \end{itemize}
}

La LAP non vale sempre: ad esempio in $(\mcP(\mbN), \Delta, \cap)$, per ogni $x \neq \varnothing, \mbN$ si ha $x \cap (\mbN \setminus x) = \varnothing$
con entrambi i fattori diversi da $\varnothing = 0_{\mcP(\mbN)}$.

\dfn[divisore-zero]{Divisore dello zero}{
    Sia $(a, +, \cdot)$ un anello. Un $x \in a$ si dice
    \begin{itemize}
        \item \textbf{divisore sinistro dello zero} se $\exists y \in a \setminus \{0_a\}\,(xy = 0_a)$;
        \item \textbf{divisore destro dello zero} se $\exists y \in a \setminus \{0_a\}\,(yx = 0_a)$;
        \item \textbf{divisore dello zero} se è sia divisore sinistro sia divisore destro.
    \end{itemize}
}

\begin{concetto}
\begin{Teorema}{Divisori dello zero e cancellabilità}{divisori-cancellabili}
    \begin{prerequisiti}
        \dip{Definizione}{def:divisore-zero}
        \dip{Definizione}{def:cancellabile}
        \dip{Lemma}{lem:regole-anello}
    \end{prerequisiti}
    Sia $(a, +, \cdot)$ un anello e $x \in a \setminus \{0_a\}$. Allora $x$ è divisore sinistro dello zero se e solo se $x$ non è
    cancellabile a sinistra (e lo stesso a destra).
\end{Teorema}
\begin{dimostrazione}
    \begin{caso}{$(\rightarrow)$ Se $x$ è divisore sinistro dello zero, allora non è cancellabile a sinistra}
        Visto che $x$ è divisore sinistro, esiste $y \in a \setminus \{0_a\}$ con $xy = 0_a = x \cdot 0_a$ (Lemma~\ref{lem:regole-anello}).
        Se $x$ fosse anche cancellabile a sinistra, avrei $y = 0_a$: assurdo.
    \end{caso}
    \begin{caso}{$(\leftarrow)$ Se $x$ non è cancellabile a sinistra, allora è divisore sinistro dello zero}
        Esistono $y, z \in a$ con $xy = xz$ e $y \neq z$. Quindi $xy - xz = 0_a$, cioè $x(y - z) = 0_a$ (Lemma~\ref{lem:regole-anello}, punto 3),
        e $y - z \neq 0_a$ perché $y \neq z$. Quindi $x$ è divisore sinistro dello zero.
    \end{caso}
    A destra la dimostrazione è identica.
\end{dimostrazione}
\end{concetto}

\begin{concetto}
\begin{Corollario}{Domini di integrità e divisori dello zero}{dominio-divisori}
    \begin{prerequisiti}
        \dip{Definizione}{def:lap}
        \dip{Definizione}{def:divisore-zero}
        \dip{Lemma}{lem:regole-anello}
    \end{prerequisiti}
    Sia $(a, +, \cdot)$ un anello commutativo unitario con $0_a \neq 1_a$. Allora $a$ è un dominio di integrità se e solo se $0_a$ è l'unico
    divisore dello zero.
\end{Corollario}
\begin{dimostrazione}
    Intanto $0_a$ è sempre un divisore dello zero: $0_a \cdot 1_a = 0_a$ con $1_a \neq 0_a$.
    \begin{caso}{$(\rightarrow)$ Se vale la LAP, nessun $x \neq 0_a$ è divisore dello zero}
        Se $x \neq 0_a$ e $xy = 0_a$, per la LAP $y = 0_a$: non esiste $y \neq 0_a$ con $xy = 0_a$.
    \end{caso}
    \begin{caso}{$(\leftarrow)$ Se $0_a$ è l'unico divisore dello zero, vale la LAP}
        Sia $xy = 0_a$. Se $x \neq 0_a$, allora $x$ non è divisore dello zero, quindi $y = 0_a$. In ogni caso $x = 0_a \vee y = 0_a$.
    \end{caso}
\end{dimostrazione}
\end{concetto}

% ──────────────────────────────────────────────────────────────
\section{Corpi e campi}\label{sec:campi}

\dfn[corpo-campo]{Corpo, campo}{
    Sia $(a, +, \cdot)$ un anello. Se $(a \setminus \{0_a\}, \cdot)$ è un gruppo, l'anello si dice \textbf{corpo}.
    Un corpo commutativo si dice \textbf{campo}.
}

\begin{esempio}{Corpi e campi}
    $(\mbZ, +, \cdot)$ è un dominio di integrità che non è un corpo ($2$ non ha inverso in $\mbZ$); $(\mbQ, +, \cdot)$ è un campo.
\end{esempio}

\begin{concetto}
\begin{Proposizione}{Corpi, campi e LAP}{campo-dominio}
    \begin{prerequisiti}
        \dip{Definizione}{def:corpo-campo}
        \dip{Teorema}{teo:invertibile-cancellabile}
        \dip{Teorema}{teo:divisori-cancellabili}
        \dip{Definizione}{def:lap}
    \end{prerequisiti}
    Ogni corpo è un anello integro; in particolare ogni campo è un dominio di integrità. Non vale il viceversa.
\end{Proposizione}
\begin{dimostrazione}
    \begin{caso}{Ogni corpo è integro}
        In un corpo tutti gli elementi diversi da $0_a$ sono invertibili, quindi cancellabili (Teorema~\ref{teo:invertibile-cancellabile}),
        quindi non sono divisori dello zero (Teorema~\ref{teo:divisori-cancellabili}). Allora se $xy = 0_a$ con $x \neq 0_a$, deve essere $y = 0_a$: vale la LAP.
    \end{caso}
    \begin{caso}{Ogni campo è un dominio di integrità}
        Un campo è commutativo, unitario ($1_a$ è il neutro del gruppo $(a \setminus \{0_a\}, \cdot)$, quindi $1_a \neq 0_a$) e integro per il caso precedente.
    \end{caso}
    \begin{caso}{Il viceversa è falso}
        $(\mbZ, +, \cdot)$ è un dominio di integrità ma non è un campo.
    \end{caso}
\end{dimostrazione}
\end{concetto}

Per gli anelli \textbf{finiti}, invece, il viceversa vale: ogni anello unitario finito e integro, con almeno due elementi, è
un corpo, e ogni dominio di integrità finito è un campo. Lo dimostreremo nel Teorema~\ref{teo:integri-finiti}, dopo aver
introdotto gli insiemi finiti (Capitolo~\ref{cap:finiti}).

% ──────────────────────────────────────────────────────────────
\section{Anelli booleani}\label{sec:booleani}

\dfn[booleano]{Anello booleano}{
    Un anello $(a, +, \cdot)$ si dice \textbf{booleano} se
    \begin{enumerate}
        \item è unitario;
        \item $\forall x \in a\,(x^2 = x)$.
    \end{enumerate}
}

Ad esempio, per ogni insieme $s$, $(\mcP(s), \Delta, \cap)$ è un anello booleano: $x \cap x = x$.

\begin{concetto}
\begin{Teorema}{Proprietà degli anelli booleani}{booleani}
    \begin{prerequisiti}
        \dip{Definizione}{def:booleano}
        \dip{Definizione}{def:anello}
        \dip{Definizione}{def:gruppo}
        \dip{Teorema}{teo:invertibile-cancellabile}
    \end{prerequisiti}
    Sia $(a, +, \cdot)$ un anello booleano. Allora $\forall x \in a\,(x = -x)$ e l'anello è commutativo.
\end{Teorema}
\begin{dimostrazione}
    \begin{caso}{$x = -x$ per ogni $x$}
        Basta dimostrare che $x + x = 0_a$. Usando $z^2 = z$ e la distributività:
        \[ x + x + 0_a = x + x = (x + x)^2 = x^2 + x^2 + x^2 + x^2 = x + x + (x + x), \]
        e dalla cancellabilità di $x + x$ in $(a, +)$ (è un gruppo) ottengo $x + x = 0_a$.
    \end{caso}
    \begin{caso}{L'anello è commutativo}
        Siano $x, y \in a$. Allora $x + y = (x + y)^2 = x^2 + xy + yx + y^2 = x + xy + yx + y$; cancellando $x$ e $y$ in $(a, +)$ ottengo
        $xy + yx = 0_a$, cioè $xy = -yx = yx$, usando il caso precedente.
    \end{caso}
\end{dimostrazione}
\end{concetto}


  % Parte 4: Relazioni binarie (equivalenze, partizioni, ordini)

  \part{Relazioni binarie}

  % Proprietà, duale, congruenze, nucleo, classi, quoziente, omomorfismo per insiemi, partizioni
  \chapter{Relazioni binarie e relazioni di equivalenza}\label{cap:equivalenze}

Nel Capitolo~\ref{cap:corrispondenze} abbiamo definito le \textbf{relazioni binarie}: le corrispondenze di un insieme con
sé stesso. In questo capitolo ne studiamo le proprietà e ci concentriamo sulle più importanti, le \textbf{relazioni di
equivalenza}. Servono a raggruppare gli elementi di un insieme in \textbf{classi} di elementi che, per lo scopo che ci
interessa, consideriamo «uguali»: i numeri pari e i numeri dispari, gli orari che segnano la stessa ora sull'orologio,
gli elementi a cui una funzione assegna lo stesso valore.

% ──────────────────────────────────────────────────────────────
\section{Rappresentare una relazione binaria}\label{sec:rappresentare-relazioni}

Ricordiamo (Definizione~\ref{def:relazione-binaria}) che una relazione binaria su un insieme $s$ è una terna
$\rho = (s, s, g)$ con $g \subseteq s \times s$, e che scriviamo $x \rho y$ quando $(x, y) \in g$. Quando $s$ è finito
una relazione si può descrivere in tre modi equivalenti, ognuno utile in situazioni diverse.

\begin{esempio}{Tre modi di descrivere la stessa relazione}
    Siano $s = \{1, 2, 3\}$ e $\rho = (s, s, g)$ con $g = \{(1, 1), (1, 2), (2, 1), (2, 3)\}$.
    \begin{center}
    \begin{minipage}[c]{0.28\linewidth}\centering
        \textbf{Elenco delle coppie}\\[2mm]
        $g = \{(1,1), (1,2),$\\[0.5mm] $(2,1), (2,3)\}$
    \end{minipage}%
    \hfill
    \begin{minipage}[c]{0.28\linewidth}\centering
        \textbf{Tabella}\\[2mm]
        \renewcommand{\arraystretch}{1.25}
        \begin{tabular}{c|ccc}
            $\rho$ & $1$ & $2$ & $3$ \\ \hline
            $1$ & $\bullet$ & $\bullet$ & \\
            $2$ & $\bullet$ & & $\bullet$ \\
            $3$ & & & \\
        \end{tabular}
    \end{minipage}%
    \hfill
    \begin{minipage}[c]{0.38\linewidth}\centering
        \textbf{Grafo con frecce}\\[1mm]
        \begin{tikzpicture}
            \node[rv] (1) at (0,0) {$1$};
            \node[rv] (2) at (1.7,0) {$2$};
            \node[rv] (3) at (3.4,0) {$3$};
            \draw[rfreccia] (1) to[rcappio] (1);
            \draw[rfreccia] (1) to[bend left=28] (2);
            \draw[rfreccia] (2) to[bend left=28] (1);
            \draw[rfreccia] (2) to (3);
        \end{tikzpicture}
    \end{minipage}
    \end{center}
    Nella tabella il pallino nella \textbf{riga} di $x$ e nella \textbf{colonna} di $y$ dice che $x \rho y$. Nel grafo la stessa
    informazione è data da una \textbf{freccia} da $x$ a $y$; una freccia da un elemento a sé stesso si chiama \textbf{cappio}.
    Qui $1 \rho 1$, $1 \rho 2$, $2 \rho 1$ e $2 \rho 3$, mentre ad esempio $3 \rho 2$ è falso.
\end{esempio}

% ──────────────────────────────────────────────────────────────
\section{Proprietà delle relazioni binarie}\label{sec:proprieta-relazioni}

\dfn[proprieta-relazioni]{Proprietà delle relazioni binarie}{
    Sia $\rho = (s, s, g)$ una relazione binaria. $\rho$ si dice
    \begin{itemize}
        \item \textbf{riflessiva} se $\forall x \in s\,(x \rho x)$;
        \item \textbf{antiriflessiva} se $\forall x \in s\,\neg(x \rho x)$;
        \item \textbf{simmetrica} se $\forall x, y \in s\,(x \rho y \rightarrow y \rho x)$;
        \item \textbf{asimmetrica} se $\forall x, y \in s\,\big(x \rho y \rightarrow \neg(y \rho x)\big)$;
        \item \textbf{antisimmetrica} se $\forall x, y \in s\,\big((x \rho y \wedge y \rho x) \rightarrow x = y\big)$;
        \item \textbf{transitiva} se $\forall x, y, z \in s\,\big((x \rho y \wedge y \rho z) \rightarrow x \rho z\big)$.
    \end{itemize}
}

\begin{concetto}
Sul grafo con frecce le proprietà si riconoscono a colpo d'occhio:
\begin{center}
\renewcommand{\arraystretch}{1.35}
\begin{tabular}{>{\bfseries}l p{0.66\linewidth}}
    \rowcolor{primary!12} Proprietà & \textbf{Come si vede sul grafo} \\
    riflessiva & c'è un cappio su \textbf{ogni} elemento \\
    antiriflessiva & non c'è \textbf{nessun} cappio \\
    simmetrica & ogni freccia ha la sua \textbf{freccia di ritorno} \\
    asimmetrica & \textbf{nessuna} freccia ha la freccia di ritorno (in particolare non ci sono cappi) \\
    antisimmetrica & tra due elementi \textbf{distinti} non ci sono mai frecce in entrambi i versi (i cappi sono ammessi) \\
    transitiva & ogni volta che si va da $x$ a $y$ e da $y$ a $z$, c'è anche la \textbf{scorciatoia} da $x$ a $z$ \\
\end{tabular}
\end{center}
\end{concetto}

\dfn[tipi-relazioni]{Relazioni di equivalenza e relazioni d'ordine}{
    Sia $\rho$ una relazione binaria su $s$. $\rho$ si dice
    \begin{itemize}
        \item \textbf{relazione di equivalenza} se è riflessiva, simmetrica e transitiva;
        \item \textbf{relazione d'ordine largo} (o semplicemente \textbf{d'ordine}) se è riflessiva, antisimmetrica e transitiva;
        \item \textbf{relazione d'ordine stretto} se è antiriflessiva e transitiva.
    \end{itemize}
    Indichiamo con $\mathrm{Eq}(s)$ l'insieme delle relazioni di equivalenza su $s$.
    \begin{center}
    \renewcommand{\arraystretch}{1.25}
    \begin{tabular}{l|ccccc}
        & riflessiva & antiriflessiva & simmetrica & antisimmetrica & transitiva \\ \hline
        equivalenza & \checkmark & & \checkmark & & \checkmark \\
        ordine largo & \checkmark & & & \checkmark & \checkmark \\
        ordine stretto & & \checkmark & & & \checkmark \\
    \end{tabular}
    \end{center}
}

\begin{esempio}{Relazioni notevoli}
    Sia $s$ un insieme.
    \begin{itemize}
        \item La relazione \textbf{identica} $(s, s, \mathrm{diag}(s))$, cioè l'uguaglianza ($x \rho y \leftrightarrow x = y$), è sia
              una relazione di equivalenza sia una relazione d'ordine largo. Non è d'ordine stretto, a meno che $s = \varnothing$.
        \item La relazione \textbf{totale} $(s, s, s \times s)$ è una relazione di equivalenza. È antiriflessiva solo se
              $s = \varnothing$ e antisimmetrica solo se $s$ ha al più un elemento.
        \item L'inclusione $\subseteq$ è una relazione d'ordine largo su $\mcP(s)$, l'inclusione stretta $\subset$ una relazione
              d'ordine stretto.
        \item $\leq$ è una relazione d'ordine largo su $\mbZ$ e $<$ una relazione d'ordine stretto.
        \item La relazione di divisibilità su $\mbN$, $m \mid n \leftrightarrow \exists k \in \mbN\,(n = km)$, è una relazione
              d'ordine largo.
        \item \textbf{Attenzione}: la divisibilità su $\mbZ$, $m \mid n \leftrightarrow \exists k \in \mbZ\,(n = km)$,
              \textbf{non} è una relazione d'ordine, perché non è antisimmetrica: $1 \mid -1$ e $-1 \mid 1$, ma $1 \neq -1$.
    \end{itemize}
\end{esempio}

Per dimostrare che una proprietà \textbf{non} vale bisogna negarne la definizione. Le negazioni hanno sempre la forma
«esiste un controesempio».

\begin{concetto}
\begin{Proposizione}{Negazione delle proprietà}{negazioni-relazioni}
    \begin{prerequisiti}
        \dip{Definizione}{def:proprieta-relazioni}
        \dip{Teorema}{teo:negazione-quantificatori}
        \dip{Proposizione}{prop:negimp}
        \dip{Proposizione}{prop:demorgan}
    \end{prerequisiti}
    Sia $\rho$ una relazione binaria su $s$. Allora:
    \begin{enumerate}
        \item $\rho$ \textbf{non} è riflessiva $\leftrightarrow \exists x \in s\,\neg(x \rho x)$;
        \item $\rho$ \textbf{non} è simmetrica $\leftrightarrow \exists x, y \in s\,\big(x \rho y \wedge \neg(y \rho x)\big)$;
        \item $\rho$ \textbf{non} è antisimmetrica $\leftrightarrow \exists x, y \in s\,(x \rho y \wedge y \rho x \wedge x \neq y)$;
        \item $\rho$ \textbf{non} è transitiva $\leftrightarrow \exists x, y, z \in s\,\big(x \rho y \wedge y \rho z \wedge \neg(x \rho z)\big)$.
    \end{enumerate}
\end{Proposizione}
\begin{dimostrazione}
    In ogni punto si nega la definizione: il $\forall$ diventa $\exists$ e si nega ciò che segue
    (Teorema~\ref{teo:negazione-quantificatori}); poi si usa $\neg(p \rightarrow q) \leftrightarrow (p \wedge \neg q)$
    (Proposizione~\ref{prop:negimp}).
    \begin{caso}{Punto 1}
        $\neg\forall x\,(x \rho x) \leftrightarrow \exists x\,\neg(x \rho x)$.
    \end{caso}
    \begin{caso}{Punti 2, 3, 4}
        Per il punto 2: $\neg\forall x, y\,(x \rho y \rightarrow y \rho x) \leftrightarrow \exists x, y\,\neg(x \rho y \rightarrow y \rho x)
        \leftrightarrow \exists x, y\,(x \rho y \wedge \neg(y \rho x))$. Per il punto 3 si nega l'implicazione
        $(x \rho y \wedge y \rho x) \rightarrow x = y$ e si ottiene $x \rho y \wedge y \rho x \wedge \neg(x = y)$; per il punto 4
        si nega $(x \rho y \wedge y \rho z) \rightarrow x \rho z$ e si ottiene $x \rho y \wedge y \rho z \wedge \neg(x \rho z)$.
    \end{caso}
\end{dimostrazione}
\end{concetto}

\begin{spiegazione}
    Per mostrare che una proprietà \textbf{vale} bisogna controllare \textbf{tutti} gli elementi; per mostrare che \textbf{non vale}
    basta \textbf{un solo controesempio}. Sul grafo i quattro controesempi sono questi (in rosso tratteggiato ciò che manca o che
    non dovrebbe esserci):
    \begin{center}
    \begin{tikzpicture}
        \begin{scope}
            \node[rv] (x) at (0,0) {$x$};
            \draw[rfreccia, cherryred, dashed] (x) to[rcappio] (x);
            \node[font=\small, align=center] at (0,-0.95) {non riflessiva:\\manca un cappio};
        \end{scope}
        \begin{scope}[xshift=3.2cm]
            \node[rv] (x) at (0,0) {$x$}; \node[rv] (y) at (1.4,0) {$y$};
            \draw[rfreccia] (x) to[bend left=25] (y);
            \draw[rfreccia, cherryred, dashed] (y) to[bend left=25] (x);
            \node[font=\small, align=center] at (0.7,-0.95) {non simmetrica:\\manca il ritorno};
        \end{scope}
        \begin{scope}[xshift=7.4cm]
            \node[rv] (x) at (0,0) {$x$}; \node[rv] (y) at (1.4,0) {$y$};
            \draw[rfreccia, cherryred] (x) to[bend left=25] (y);
            \draw[rfreccia, cherryred] (y) to[bend left=25] (x);
            \node[font=\small, align=center] at (0.7,-0.95) {non antisimmetrica:\\andata e ritorno, $x \neq y$};
        \end{scope}
        \begin{scope}[xshift=11.6cm]
            \node[rv] (x) at (0,0) {$x$}; \node[rv] (y) at (0.8,0.95) {$y$}; \node[rv] (z) at (1.6,0) {$z$};
            \draw[rfreccia] (x) to (y); \draw[rfreccia] (y) to (z);
            \draw[rfreccia, cherryred, dashed] (x) to (z);
            \node[font=\small, align=center] at (0.8,-0.95) {non transitiva:\\manca la scorciatoia};
        \end{scope}
    \end{tikzpicture}
    \end{center}
\end{spiegazione}

\warning{Negare non vuol dire passare alla proprietà «anti»}{
    «Non riflessiva» \textbf{non} significa «antiriflessiva», e «non simmetrica» non significa «antisimmetrica».
    Su $s = \{1, 2, 3\}$:
    \begin{itemize}
        \item $g = \{(1, 1)\}$ non è riflessiva ($2$ non ha il cappio) e non è antiriflessiva ($1$ ce l'ha);
        \item $g = \{(1, 2), (2, 1), (2, 3)\}$ non è simmetrica ($2 \rho 3$ ma non $3 \rho 2$) e non è antisimmetrica
              ($1 \rho 2$ e $2 \rho 1$ con $1 \neq 2$);
        \item l'uguaglianza è \textbf{sia} simmetrica \textbf{sia} antisimmetrica.
    \end{itemize}
}

Controllare la transitività su un insieme finito richiede di guardare tutte le terne di elementi. Molte terne però si
controllano da sole.

\begin{concetto}
\begin{Proposizione}{Quali terne guardare per la transitività}{transitivita-distinti}
    \begin{prerequisiti}
        \dip{Definizione}{def:proprieta-relazioni}
        \dip{Proposizione}{prop:terzoescluso}
    \end{prerequisiti}
    Sia $\rho$ una relazione binaria su $s$. $\rho$ è transitiva se e solo se valgono entrambe le condizioni:
    \begin{enumerate}
        \item[(a)] $(x \rho y \wedge y \rho z) \rightarrow x \rho z$ per ogni terna $x, y, z$ di elementi \textbf{a due a due distinti};
        \item[(b)] $(x \rho y \wedge y \rho x) \rightarrow x \rho x$ per ogni $x, y \in s$.
    \end{enumerate}
    In particolare, se $\rho$ è \textbf{riflessiva} oppure \textbf{antisimmetrica}, la (b) vale sempre, e per la transitività
    \textbf{basta guardare le terne di elementi a due a due distinti}.
\end{Proposizione}
\begin{dimostrazione}
    \begin{caso}{$(\rightarrow)$ Se $\rho$ è transitiva, valgono (a) e (b)}
        (a) è un caso particolare della transitività; (b) è la transitività applicata alla terna $x, y, x$.
    \end{caso}
    \begin{caso}{$(\leftarrow)$ Se valgono (a) e (b), $\rho$ è transitiva}
        Siano $x \rho y$ e $y \rho z$: dobbiamo mostrare $x \rho z$.
        Se $x, y, z$ sono a due a due distinti, basta (a). Se $x = y$, la tesi $x \rho z$ è l'ipotesi $y \rho z$; se $y = z$,
        la tesi è l'ipotesi $x \rho y$. Resta il caso $x = z$ (con $x \neq y$): le ipotesi sono $x \rho y$ e $y \rho x$, e
        la tesi $x \rho x$ è proprio la (b).
    \end{caso}
    \begin{caso}{Se $\rho$ è riflessiva o antisimmetrica, la (b) vale sempre}
        Se $\rho$ è riflessiva, $x \rho x$ vale comunque. Se $\rho$ è antisimmetrica, da $x \rho y$ e $y \rho x$ segue $x = y$,
        e allora $x \rho x$ è l'ipotesi $x \rho y$.
    \end{caso}
\end{dimostrazione}
\end{concetto}

\begin{spiegazione}
    In una terna con due elementi uguali la transitività è quasi sempre automatica: se $x = y$ o $y = z$ la «scorciatoia»
    coincide con una delle due frecce che abbiamo già. L'unico caso insidioso è il \textbf{viaggio di andata e ritorno}
    $x \to y \to x$, che obbliga ad avere il cappio su $x$. Ecco perché la condizione (b) non si può togliere:
    \begin{center}
    \begin{tikzpicture}
        \node[rv] (1) at (0,0) {$1$}; \node[rv] (2) at (1.8,0) {$2$};
        \draw[rfreccia] (1) to[bend left=28] (2);
        \draw[rfreccia] (2) to[bend left=28] (1);
        \draw[rfreccia, cherryred, dashed] (1) to[rcappio] (1);
        \draw[rfreccia, cherryred, dashed] (2) to[rcappio] (2);
        \node[font=\small, align=left, anchor=west] at (2.8,0.1)
            {$s = \{1, 2\}$, $g = \{(1, 2), (2, 1)\}$: non ci sono terne di elementi distinti,\\
             quindi (a) è vera; ma $1 \rho 2$ e $2 \rho 1$ senza $1 \rho 1$: $\rho$ \textbf{non} è transitiva.};
    \end{tikzpicture}
    \end{center}
    Per le relazioni di equivalenza e d'ordine largo (riflessive) e per quelle antisimmetriche questo problema non c'è.
\end{spiegazione}

La definizione di ordine stretto chiede solo antiriflessività e transitività. Ci si aspetterebbe di dover chiedere anche
che non ci siano mai «andata e ritorno», cioè l'asimmetria: la proposizione seguente mostra che viene gratis.

\begin{concetto}
\begin{Proposizione}{Per un ordine stretto bastano antiriflessività e transitività}{stretto-asimmetrico}
    \begin{prerequisiti}
        \dip{Definizione}{def:proprieta-relazioni}
        \dip{Definizione}{def:tipi-relazioni}
    \end{prerequisiti}
    Sia $\rho$ una relazione binaria transitiva su $s$. Allora $\rho$ è asimmetrica se e solo se è antiriflessiva.
    In particolare ogni relazione d'ordine stretto è asimmetrica e antisimmetrica.
\end{Proposizione}
\begin{dimostrazione}
    \begin{caso}{$(\rightarrow)$ Asimmetrica $\Rightarrow$ antiriflessiva (anche senza transitività)}
        Se fosse $x \rho x$ per qualche $x$, l'asimmetria con $y = x$ darebbe $\neg(x \rho x)$: assurdo.
    \end{caso}
    \begin{caso}{$(\leftarrow)$ Antiriflessiva e transitiva $\Rightarrow$ asimmetrica}
        Per assurdo siano $x \rho y$ e $y \rho x$. Per transitività $x \rho x$, contro l'antiriflessività.
    \end{caso}
    \begin{caso}{Ogni ordine stretto è antisimmetrico}
        È asimmetrico per quanto appena visto, quindi $x \rho y \wedge y \rho x$ è sempre falsa e l'implicazione che definisce
        l'antisimmetria è vera (antecedente falso).
    \end{caso}
\end{dimostrazione}
\end{concetto}

\begin{spiegazione}
    Un'«andata e ritorno» $x \to y \to x$, per la transitività, produrrebbe la scorciatoia $x \to x$: un cappio. Ma un ordine
    stretto non ha cappi, quindi non può avere andate e ritorni. Pensiamo a $<$ su $\mbZ$: da $3 < 5$ e $5 < 3$ seguirebbe
    $3 < 3$, che è falso; quindi $3 < 5$ e $5 < 3$ non possono essere vere insieme.
\end{spiegazione}

% ──────────────────────────────────────────────────────────────
\section{Relazione duale e diagonale}\label{sec:duale-diagonale}

\dfn[relazione-duale]{Relazione duale}{
    Sia $\rho = (s, s, g)$ una relazione binaria. La relazione $\bar{\rho} = (s, s, \bar{g})$ con
    \[ \bar{g} := \{(y, x) \mid (x, y) \in g\}, \qquad\text{cioè}\qquad \forall x, y \in s\,(x \bar{\rho} y \leftrightarrow y \rho x), \]
    si dice \textbf{relazione duale} (o \textbf{opposta}) di $\rho$. Sul grafo si ottiene \textbf{invertendo il verso di tutte le frecce}.
}

\begin{esempio}{Relazioni duali}
    La duale di $\leq$ su $\mbZ$ è $\geq$; la duale di $\subseteq$ è $\supseteq$; la duale di «essere padre di» è «essere figlio di»
    (ristretta ai maschi). La duale di una relazione simmetrica è la relazione stessa.
\end{esempio}

Ricordiamo che $\mathrm{diag}(s) = \{(x, x) \mid x \in s\}$ (Definizione~\ref{def:relazione-binaria}): nel grafo corrisponde
ai cappi, nella tabella alla diagonale. Le proprietà della Definizione~\ref{def:proprieta-relazioni} si possono scrivere come
semplici relazioni tra insiemi.

\begin{concetto}
\begin{Proposizione}{Le proprietà guardando il grafico}{proprieta-grafico}
    \begin{prerequisiti}
        \dip{Definizione}{def:proprieta-relazioni}
        \dip{Definizione}{def:relazione-duale}
        \dip{Definizione}{def:relazione-binaria}
        \dip{Definizione}{def:prodotto-corrispondenze}
        \dip{Proposizione}{prop:doppia-inclusione}
    \end{prerequisiti}
    Sia $\rho = (s, s, g)$ una relazione binaria e sia $\bar{g}$ il grafico della duale. Allora
    \begin{enumerate}
        \item $\rho$ è riflessiva $\leftrightarrow \mathrm{diag}(s) \subseteq g$;
        \item $\rho$ è antiriflessiva $\leftrightarrow \mathrm{diag}(s) \cap g = \varnothing$;
        \item $\rho$ è simmetrica $\leftrightarrow g = \bar{g}$, cioè $\rho = \bar{\rho}$;
        \item $\rho$ è antisimmetrica $\leftrightarrow g \cap \bar{g} \subseteq \mathrm{diag}(s)$;
        \item $\rho$ è asimmetrica $\leftrightarrow g \cap \bar{g} = \varnothing$;
        \item $\rho$ è transitiva $\leftrightarrow$ il grafico di $\rho\rho$ è contenuto in $g$.
    \end{enumerate}
\end{Proposizione}
\begin{dimostrazione}
    Basta tradurre le definizioni, ricordando che $(x, y) \in \bar{g} \leftrightarrow (y, x) \in g$.
    \begin{caso}{Punti 1 e 2}
        $\mathrm{diag}(s) \subseteq g$ significa che ogni $(x, x)$ sta in $g$, cioè $x \rho x$ per ogni $x$. $\mathrm{diag}(s) \cap g = \varnothing$
        significa che nessun $(x, x)$ sta in $g$, cioè $\neg(x \rho x)$ per ogni $x$.
    \end{caso}
    \begin{caso}{Punto 3}
        $(\rightarrow)$ Se $(x, y) \in g$, per simmetria $(y, x) \in g$, cioè $(x, y) \in \bar{g}$: quindi $g \subseteq \bar{g}$. Se
        $(x, y) \in \bar{g}$, allora $(y, x) \in g$ e per simmetria $(x, y) \in g$: quindi $\bar{g} \subseteq g$, e per la doppia
        inclusione $g = \bar{g}$. $(\leftarrow)$ Se $x \rho y$, allora $(x, y) \in g = \bar{g}$, cioè $(y, x) \in g$: $y \rho x$.
    \end{caso}
    \begin{caso}{Punti 4 e 5}
        $(x, y) \in g \cap \bar{g}$ significa $x \rho y \wedge y \rho x$. L'antisimmetria dice che questo accade solo se $x = y$, cioè solo
        se $(x, y) \in \mathrm{diag}(s)$; l'asimmetria dice che non accade mai, cioè $g \cap \bar{g} = \varnothing$.
    \end{caso}
    \begin{caso}{Punto 6}
        Per la Definizione~\ref{def:prodotto-corrispondenze}, $(x, z)$ sta nel grafico di $\rho\rho$ se e solo se esiste $y$ con
        $x \rho y$ e $y \rho z$. Dire che tutte queste coppie stanno in $g$ è esattamente la transitività.
    \end{caso}
\end{dimostrazione}
\end{concetto}

\begin{spiegazione}
    Conviene pensare alla \textbf{tabella}. La diagonale $\mathrm{diag}(s)$ sono le caselle in cui riga e colonna coincidono,
    e la duale si ottiene \textbf{ribaltando la tabella attorno alla diagonale}: la casella (riga $x$, colonna $y$) va in
    (riga $y$, colonna $x$).
    \begin{center}
    \begin{tikzpicture}[cella/.style={minimum size=6mm, draw=black!25, font=\small}]
        \foreach \i in {1,2,3} {
            \node[font=\small] at (\i*0.6, 0.6) {$\i$};
            \node[font=\small] at (0, -\i*0.6+0.6) {$\i$};
        }
        \foreach \r/\c in {1/1,2/2,3/3} \node[cella, fill=naplesyellow!40] at (\c*0.6, -\r*0.6+0.6) {};
        \foreach \r/\c in {1/2,2/1,2/3} \node[cella] at (\c*0.6, -\r*0.6+0.6) {$\bullet$};
        \foreach \r/\c in {1/1} \node[cella, fill=naplesyellow!40] at (\c*0.6, -\r*0.6+0.6) {$\bullet$};
        \node[font=\small] at (1.2,-1.75) {$g$};
        \draw[-{Stealth}, line width=0.8pt, primary] (2.6,-0.6) -- node[above, font=\small] {ribalto} (3.8,-0.6);
        \begin{scope}[xshift=4.2cm]
            \foreach \i in {1,2,3} {
                \node[font=\small] at (\i*0.6, 0.6) {$\i$};
                \node[font=\small] at (0, -\i*0.6+0.6) {$\i$};
            }
            \foreach \r/\c in {1/1,2/2,3/3} \node[cella, fill=naplesyellow!40] at (\c*0.6, -\r*0.6+0.6) {};
            \foreach \r/\c in {2/1,1/2,3/2} \node[cella] at (\c*0.6, -\r*0.6+0.6) {$\bullet$};
            \foreach \r/\c in {1/1} \node[cella, fill=naplesyellow!40] at (\c*0.6, -\r*0.6+0.6) {$\bullet$};
            \node[font=\small] at (1.2,-1.75) {$\bar{g}$};
        \end{scope}
        \node[font=\small, align=left, anchor=west] at (7.0,-0.55)
            {È la relazione dell'Esempio iniziale:\\
             $g = \{(1,1), (1,2), (2,1), (2,3)\}$.\\
             La diagonale (in giallo) resta ferma.};
    \end{tikzpicture}
    \end{center}
    Allora: \textbf{riflessiva} = diagonale tutta piena; \textbf{antiriflessiva} = diagonale tutta vuota; \textbf{simmetrica} =
    la tabella non cambia ribaltandola; \textbf{antisimmetrica} = fuori dalla diagonale, una casella e la sua simmetrica non sono
    mai piene insieme. Nell'esempio $(1, 2)$ e $(2, 1)$ sono piene insieme: la relazione non è antisimmetrica; e $(2, 3)$ è
    piena mentre $(3, 2)$ no: non è nemmeno simmetrica.
\end{spiegazione}

\begin{concetto}
\begin{Corollario}{Relazioni simmetriche e antisimmetriche}{simm-antisimm}
    \begin{prerequisiti}
        \dip{Proposizione}{prop:proprieta-grafico}
    \end{prerequisiti}
    Una relazione binaria $\rho = (s, s, g)$ è sia simmetrica sia antisimmetrica se e solo se $g \subseteq \mathrm{diag}(s)$.
\end{Corollario}
\begin{dimostrazione}
    \begin{caso}{$(\rightarrow)$}
        Per i punti 3 e 4 della Proposizione~\ref{prop:proprieta-grafico}, $g = \bar{g}$, quindi $g = g \cap \bar{g} \subseteq \mathrm{diag}(s)$.
    \end{caso}
    \begin{caso}{$(\leftarrow)$}
        Se $g \subseteq \mathrm{diag}(s)$, ogni coppia di $g$ è del tipo $(x, x)$, che coincide con la sua ribaltata: quindi
        $\bar{g} = g$ e $\rho$ è simmetrica. Inoltre $g \cap \bar{g} \subseteq g \subseteq \mathrm{diag}(s)$: $\rho$ è antisimmetrica.
    \end{caso}
\end{dimostrazione}
\end{concetto}

\begin{spiegazione}
    Simmetrica vuol dire «ogni freccia ha il ritorno», antisimmetrica «nessuna freccia tra elementi distinti ha il ritorno».
    Le due cose insieme sono possibili solo se \textbf{non ci sono frecce tra elementi distinti}: restano al più i cappi.
    Ad esempio su $\{1, 2, 3\}$ la relazione con grafico $\{(1,1), (3,3)\}$ è simmetrica e antisimmetrica.
\end{spiegazione}

\begin{concetto}
\begin{Proposizione}{La duale conserva le proprietà}{duale-proprieta}
    \begin{prerequisiti}
        \dip{Definizione}{def:relazione-duale}
        \dip{Definizione}{def:proprieta-relazioni}
        \dip{Definizione}{def:tipi-relazioni}
        \dip{Corollario}{cor:terna}
    \end{prerequisiti}
    Sia $\rho$ una relazione binaria su $s$. Allora $\bar{\bar{\rho}} = \rho$ e $\bar{\rho}$ ha le stesse proprietà di $\rho$ tra quelle
    della Definizione~\ref{def:proprieta-relazioni}. In particolare la duale di una relazione di equivalenza, d'ordine largo o
    d'ordine stretto è una relazione dello stesso tipo.
\end{Proposizione}
\begin{dimostrazione}
    \begin{caso}{$\bar{\bar{\rho}} = \rho$}
        $x \bar{\bar{\rho}} y \leftrightarrow y \bar{\rho} x \leftrightarrow x \rho y$: i grafici coincidono, quindi (Corollario~\ref{cor:terna}) le terne coincidono.
    \end{caso}
    \begin{caso}{Le proprietà si conservano}
        Ogni proprietà si ottiene da quella di $\rho$ \textbf{scambiando i nomi} delle variabili. Ad esempio, per la transitività:
        se $x \bar{\rho} y$ e $y \bar{\rho} z$, allora $z \rho y$ e $y \rho x$, quindi $z \rho x$ per la transitività di $\rho$, cioè
        $x \bar{\rho} z$. Per l'antisimmetria: $x \bar{\rho} y \wedge y \bar{\rho} x$ è $y \rho x \wedge x \rho y$, che implica $x = y$.
        Riflessività e antiriflessività riguardano $x \bar{\rho} x$, che è $x \rho x$; simmetria e asimmetria si trattano allo stesso modo.
        Viceversa, se $\bar{\rho}$ ha una proprietà, ce l'ha anche $\bar{\bar{\rho}} = \rho$.
    \end{caso}
\end{dimostrazione}
\end{concetto}

\begin{spiegazione}
    Invertire tutte le frecce non cambia la «forma» del grafo: i cappi restano cappi, le andate e ritorno restano andate e
    ritorno, e un percorso $x \to y \to z$ con la sua scorciatoia $x \to z$ diventa il percorso $z \to y \to x$ con la scorciatoia
    $z \to x$. Ad esempio $\geq$ è un ordine largo come $\leq$, e $>$ un ordine stretto come $<$.
\end{spiegazione}

% ──────────────────────────────────────────────────────────────
\section{Relazioni di equivalenza}\label{sec:equivalenze}

Una relazione di equivalenza si comporta come una «uguaglianza parziale»: ogni elemento è equivalente a sé stesso, se $x$ è
equivalente a $y$ anche $y$ lo è a $x$, e due elementi equivalenti a un terzo sono equivalenti tra loro.

\begin{esempio}{Relazioni di equivalenza}
    \begin{itemize}
        \item L'uguaglianza e la relazione totale su un qualunque insieme.
        \item Sull'insieme delle persone, «avere la stessa età» o «essere nati nello stesso mese».
        \item In un insieme di strutture algebriche, «essere isomorfe»: ogni struttura è isomorfa a sé stessa tramite l'identità,
              l'inversa di un isomorfismo è un isomorfismo (Teorema~\ref{teo:inverso-isomorfismo}) e la composta di due
              isomorfismi è biettiva (Corollario~\ref{cor:comp-biettive}) e rispetta l'operazione.
    \end{itemize}
\end{esempio}

\dfn[congruenza-modulo]{Congruenza modulo $m$}{
    Sia $m \in \mbZ$. Definiamo la relazione binaria $\equiv_m$ su $\mbZ$ ponendo
    \[ \forall a, b \in \mbZ\,\big(a \equiv_m b \leftrightarrow \exists k \in \mbZ\,(b - a = km)\big). \]
    Se $a \equiv_m b$ diciamo che \textbf{$a$ è congruo a $b$ modulo $m$}: $a$ e $b$ differiscono per un multiplo di $m$.
}

\begin{concetto}
\begin{Proposizione}{La congruenza modulo $m$ è di equivalenza}{congruenza-equivalenza}
    \begin{prerequisiti}
        \dip{Definizione}{def:congruenza-modulo}
        \dip{Definizione}{def:tipi-relazioni}
    \end{prerequisiti}
    Per ogni $m \in \mbZ$, $\equiv_m$ è una relazione di equivalenza su $\mbZ$.
\end{Proposizione}
\begin{dimostrazione}
    \begin{caso}{Riflessiva}
        Per ogni $a \in \mbZ$, $a - a = 0 = 0 \cdot m$, quindi $a \equiv_m a$ (con $k = 0$).
    \end{caso}
    \begin{caso}{Simmetrica}
        Se $a \equiv_m b$, esiste $k$ con $b - a = km$. Allora $a - b = (-k)m$, quindi $b \equiv_m a$ (con $-k$).
    \end{caso}
    \begin{caso}{Transitiva}
        Se $a \equiv_m b$ e $b \equiv_m c$, esistono $h, k$ con $b - a = hm$ e $c - b = km$. Sommando,
        $c - a = (c - b) + (b - a) = km + hm = (h + k)m$, quindi $a \equiv_m c$ (con $h + k$).
    \end{caso}
\end{dimostrazione}
\end{concetto}

\begin{spiegazione}
    Pensiamo ai numeri interi come a una strada e a $\equiv_m$ come a \textbf{salti di lunghezza $m$}: $a \equiv_m b$ quando da $a$ si
    arriva a $b$ con un certo numero $k$ di salti (in avanti se $k > 0$, all'indietro se $k < 0$). Per $m = 3$, colorando ogni
    numero secondo i salti che lo raggiungono partendo da $0$, $1$ o $2$:
    \begin{center}
    \begin{tikzpicture}[x=0.92cm]
        \draw[-{Stealth}, black!50] (-4.6,0) -- (8.6,0);
        \foreach \n in {-4,...,8} {
            \pgfmathtruncatemacro{\r}{mod(\n+300,3)}
            \ifcase\r \def\col{classeA}\or \def\col{classeB}\or \def\col{classeC}\fi
            \node[circle, fill=\col, draw=black!30, inner sep=0pt, minimum size=6mm, font=\small] at (\n,0) {$\n$};
        }
        \draw[-{Stealth}, line width=0.7pt, orange!80!black] (1,0.35) to[bend left=55] node[above, font=\scriptsize] {$+3$} (4,0.35);
        \draw[-{Stealth}, line width=0.7pt, orange!80!black] (4,0.35) to[bend left=55] node[above, font=\scriptsize] {$+3$} (7,0.35);
        \draw[-{Stealth}, line width=0.7pt, orange!80!black] (1,-0.35) to[bend left=55] node[below, font=\scriptsize] {$-3$} (-2,-0.35);
    \end{tikzpicture}
    \end{center}
    \begin{passi}
        \item \textbf{Riflessiva}: da $a$ si arriva ad $a$ con zero salti.
        \item \textbf{Simmetrica}: se da $1$ si arriva a $7$ con due salti avanti, da $7$ si torna a $1$ con due salti indietro.
        \item \textbf{Transitiva}: da $1$ a $4$ con un salto e da $4$ a $7$ con un altro: da $1$ a $7$ con $1 + 1 = 2$ salti.
    \end{passi}
    I numeri dello stesso colore sono tutti congrui tra loro modulo $3$.
\end{spiegazione}

\begin{esempio}{Casi particolari}
    $\equiv_0$ è l'uguaglianza ($b - a = k \cdot 0$ vuol dire $a = b$); $\equiv_1$ è la relazione totale (ogni differenza è un
    multiplo di $1$); due interi sono congrui modulo $2$ se e solo se sono entrambi pari o entrambi dispari. Infine
    $\equiv_m$ e $\equiv_{-m}$ sono la stessa relazione, perché $km = (-k)(-m)$.
\end{esempio}

\dfn[nucleo]{Nucleo di equivalenza di una funzione}{
    Sia $f \colon a \to b$ una funzione. La relazione binaria $\sim_f$ su $a$ definita da
    \[ \forall x, y \in a\,\big(x \sim_f y \leftrightarrow f(x) = f(y)\big) \]
    si chiama \textbf{nucleo di equivalenza} di $f$ e si indica anche con $\mathrm{Ker}\, f$.
}

\begin{concetto}
\begin{Proposizione}{Proprietà del nucleo di equivalenza}{nucleo-equivalenza}
    \begin{prerequisiti}
        \dip{Definizione}{def:nucleo}
        \dip{Definizione}{def:tipi-relazioni}
        \dip{Proposizione}{prop:forme-iniettivita}
        \dip{Definizione}{def:funzioni-notevoli}
    \end{prerequisiti}
    Sia $f \colon a \to b$ una funzione. Allora:
    \begin{enumerate}
        \item $\sim_f$ è una relazione di equivalenza su $a$;
        \item $\sim_f$ è la relazione identica se e solo se $f$ è iniettiva;
        \item se $a \neq \varnothing$, $\sim_f$ è la relazione totale se e solo se $f$ è costante.
    \end{enumerate}
\end{Proposizione}
\begin{dimostrazione}
    \begin{caso}{Punto 1}
        Le tre proprietà vengono da quelle dell'uguaglianza: $f(x) = f(x)$; se $f(x) = f(y)$ allora $f(y) = f(x)$; se
        $f(x) = f(y)$ e $f(y) = f(z)$ allora $f(x) = f(z)$.
    \end{caso}
    \begin{caso}{Punto 2}
        Sempre $x = y \rightarrow f(x) = f(y)$. Quindi $\sim_f$ è l'uguaglianza se e solo se vale anche
        $f(x) = f(y) \rightarrow x = y$ per ogni $x, y$, cioè (Proposizione~\ref{prop:forme-iniettivita}) se e solo se $f$ è iniettiva.
    \end{caso}
    \begin{caso}{Punto 3}
        $(\leftarrow)$ Se $f$ è costante di valore $c$, $f(x) = c = f(y)$ per ogni $x, y$: tutti sono in relazione.
        $(\rightarrow)$ Prendo $x_0 \in a$ e pongo $c = f(x_0)$. Per ogni $x \in a$ si ha $x \sim_f x_0$, cioè $f(x) = c$: $f$ è la
        funzione costante di valore $c$.
    \end{caso}
\end{dimostrazione}
\end{concetto}

\begin{spiegazione}
    $\sim_f$ mette insieme gli elementi che \textbf{$f$ non sa distinguere}, cioè che hanno la stessa immagine. Prendiamo ad
    esempio la funzione che manda ogni intero nel suo quadrato, $f(n) = n^2$:
    \begin{center}
    \begin{tikzpicture}
        \foreach \n/\y in {0/2.4, 1/1.6, -1/1.15, 2/0.35, -2/-0.1, 3/-0.9, -3/-1.35}
            \node[elem] (a\n) at (0,\y) {$\n$};
        \foreach \v/\y in {0/2.4, 1/1.375, 4/0.125, 9/-1.125}
            \node[elem] (b\v) at (3.4,\y) {$\v$};
        \blocco{classeA}{(a0.center) -- ++(0.001,0)}
        \blocco{classeB}{(a1.center) -- (a-1.center)}
        \blocco{classeC}{(a2.center) -- (a-2.center)}
        \blocco{classeD}{(a3.center) -- (a-3.center)}
        \foreach \n/\v in {0/0, 1/1, -1/1, 2/4, -2/4, 3/9, -3/9} \draw[fun] (a\n) -- (b\v);
        \node[font=\small] at (0,3.05) {$\mbZ$ (una parte)};
        \node[font=\small] at (3.4,3.05) {$\mbZ$};
        \node[font=\small, primary!80!black] at (1.7,2.75) {$f$};
        \node[font=\small, align=left, anchor=west] at (4.4,0.6)
            {$-2 \sim_f 2$ perché $f(-2) = 4 = f(2)$;\\
             $1 \not\sim_f 2$ perché $f(1) = 1 \neq 4 = f(2)$.\\[1mm]
             Ogni macchia colorata è un gruppo di elementi\\
             con la stessa immagine.};
    \end{tikzpicture}
    \end{center}
    Il punto 2 dice che se $f$ è iniettiva non ci sono gruppi con più di un elemento; il punto 3 che se $f$ è costante c'è un
    solo gruppo, con dentro tutto $a$.
\end{spiegazione}

% ──────────────────────────────────────────────────────────────
\section{Classi di equivalenza e insieme quoziente}\label{sec:classi}

\dfn[classe-equivalenza]{Classe di equivalenza, insieme quoziente}{
    Sia $\sim$ una relazione di equivalenza su $a$ e sia $x \in a$.
    \begin{itemize}
        \item L'insieme $[x]_\sim := \{y \in a \mid x \sim y\}$ si chiama \textbf{classe di equivalenza di $x$ modulo $\sim$};
              ogni elemento di $[x]_\sim$ (ad esempio $x$ stesso) si dice \textbf{rappresentante} della classe.
        \item L'insieme di tutte le classi $a/{\sim} := \{[x]_\sim \mid x \in a\}$ si chiama \textbf{insieme quoziente} di $a$
              modulo $\sim$.
    \end{itemize}
    Quando non c'è ambiguità si scrive $[x]$ invece di $[x]_\sim$.
}

\begin{esempio}{Classi e quozienti}
    \begin{itemize}
        \item Per l'uguaglianza ogni classe è un singleton: $[x] = \{x\}$, e $a/{=} = \{\{x\} \mid x \in a\}$.
        \item Per la relazione totale su $a \neq \varnothing$ c'è una sola classe: $[x] = a$ per ogni $x$, e il quoziente è $\{a\}$.
        \item Per $\equiv_2$: $[0]_{\equiv_2}$ è l'insieme dei pari, $[1]_{\equiv_2}$ quello dei dispari, e $\mbZ/{\equiv_2}$ ha due elementi.
              Per $\equiv_3$ le classi sono i tre colori della figura nella spiegazione della
              Proposizione~\ref{prop:congruenza-equivalenza}: $[0] = \{\dots, -3, 0, 3, 6, \dots\}$, $[1] = \{\dots, -2, 1, 4, 7, \dots\}$,
              $[2] = \{\dots, -1, 2, 5, 8, \dots\}$.
        \item Per il nucleo di $n \in \mbZ \mapsto n^2 \in \mbZ$: $[n] = \{n, -n\}$; ad esempio $[0] = \{0\}$ e $[-3] = [3] = \{-3, 3\}$.
    \end{itemize}
\end{esempio}

\begin{concetto}
\begin{Teorema}{Proprietà fondamentali delle classi di equivalenza}{proprieta-classi}
    \begin{prerequisiti}
        \dip{Definizione}{def:classe-equivalenza}
        \dip{Definizione}{def:tipi-relazioni}
        \dip{Proposizione}{prop:doppia-inclusione}
        \dip{Definizione}{def:notazioni-assiomi}
        \dip{Proposizione}{prop:contrapposizione}
    \end{prerequisiti}
    Sia $\sim$ una relazione di equivalenza su $a$. Allora:
    \begin{enumerate}
        \item per ogni $x \in a$, $x \in [x]$; in particolare ogni classe è \textbf{non vuota};
        \item per ogni $x, y \in a$:\quad
              $x \sim y \ \leftrightarrow\ y \sim x \ \leftrightarrow\ y \in [x] \ \leftrightarrow\ x \in [y] \ \leftrightarrow\ [x] = [y]
              \ \leftrightarrow\ [x] \cap [y] \neq \varnothing$;
        \item per ogni $x, y \in a$, se $[x] \neq [y]$ allora $[x] \cap [y] = \varnothing$: le classi sono \textbf{a due a due disgiunte};
        \item $\bigcup (a/{\sim}) = a$.
    \end{enumerate}
\end{Teorema}
\begin{dimostrazione}
    \begin{caso}{Punto 1}
        Per la riflessività $x \sim x$, cioè $x \in [x]$.
    \end{caso}
    \begin{caso}{Punto 2}
        Le prime quattro condizioni sono equivalenti per la simmetria e per la definizione di classe
        ($y \in [x] \leftrightarrow x \sim y$, $x \in [y] \leftrightarrow y \sim x$). Chiudiamo il cerchio con le altre due.
        \begin{itemize}
            \item[] $x \sim y \rightarrow [x] = [y]$: se $w \in [x]$, cioè $x \sim w$, da $y \sim x$ (simmetria) e transitività si ha
                    $y \sim w$, cioè $w \in [y]$; quindi $[x] \subseteq [y]$. Scambiando $x$ e $y$, $[y] \subseteq [x]$: per la doppia
                    inclusione $[x] = [y]$.
            \item[] $[x] = [y] \rightarrow [x] \cap [y] \neq \varnothing$: per il punto 1 $x \in [x] = [y]$, quindi $x \in [x] \cap [y]$.
            \item[] $[x] \cap [y] \neq \varnothing \rightarrow x \sim y$: se $z \in [x] \cap [y]$, allora $x \sim z$ e $y \sim z$; per
                    simmetria $z \sim y$ e per transitività $x \sim y$.
        \end{itemize}
    \end{caso}
    \begin{caso}{Punto 3}
        Dal punto 2, $[x] \cap [y] \neq \varnothing \rightarrow [x] = [y]$; per contrapposizione $[x] \neq [y] \rightarrow [x] \cap [y] = \varnothing$.
    \end{caso}
    \begin{caso}{Punto 4}
        $(\subseteq)$ Gli elementi delle classi sono elementi di $a$. $(\supseteq)$ Se $x \in a$, per il punto 1 $x \in [x] \in a/{\sim}$,
        quindi $x \in \bigcup(a/{\sim})$.
    \end{caso}
\end{dimostrazione}
\end{concetto}

\begin{spiegazione}
    Le classi sono come le \textbf{stanze di una casa}: ogni elemento sta in una stanza (punto 1), due stanze diverse non hanno
    nessuno in comune (punto 3) e tutte insieme formano la casa intera (punto 4). Il punto 2 dice che \textbf{essere in
    relazione equivale a stare nella stessa stanza}, e che una stanza non dipende da chi scegliamo per indicarla:
    \begin{center}
    \begin{tikzpicture}
        \node[elem] (a1) at (0,1.1) {$1$}; \node[elem] (a2) at (1.0,1.35) {$2$}; \node[elem] (a3) at (0.55,0.35) {$3$};
        \node[elem] (b1) at (2.5,1.25) {$4$}; \node[elem] (b2) at (3.0,0.4) {$5$};
        \node[elem] (c1) at (4.4,0.85) {$6$};
        \blocco[7mm]{classeA}{(a1.center) -- (a2.center) -- (a3.center) -- cycle}
        \blocco[7mm]{classeB}{(b1.center) -- (b2.center)}
        \blocco[7mm]{classeC}{(c1.center) -- ++(0.001,0)}
        \draw[black!45, rounded corners=12pt] (-0.75,-0.45) rectangle (5.15,2.0);
        \node[font=\small] at (-0.45,1.75) {$a$};
        \node[font=\small, align=left, anchor=west] at (5.6,0.8)
            {$[1] = [2] = [3] = \{1, 2, 3\}$\\
             $[4] = [5] = \{4, 5\}$, \quad $[6] = \{6\}$\\
             $a/{\sim} = \{\{1, 2, 3\}, \{4, 5\}, \{6\}\}$};
    \end{tikzpicture}
    \end{center}
    Per dimostrare un'equivalenza tra tante condizioni non serve provare tutte le coppie di implicazioni: basta una catena che
    torni al punto di partenza, come abbiamo fatto con
    $x \sim y \Rightarrow [x] = [y] \Rightarrow [x] \cap [y] \neq \varnothing \Rightarrow x \sim y$.
\end{spiegazione}

Per il nucleo di equivalenza di una funzione le classi si descrivono con l'antiimmagine.

\begin{concetto}
\begin{Proposizione}{Classi del nucleo di equivalenza}{classi-nucleo}
    \begin{prerequisiti}
        \dip{Definizione}{def:nucleo}
        \dip{Definizione}{def:classe-equivalenza}
        \dip{Definizione}{def:immagine-antiimmagine}
        \dip{Assioma}{ass:estensionalita}
    \end{prerequisiti}
    Sia $f \colon a \to b$ e sia $x \in a$. Allora $[x]_{\sim_f} = \overleftarrow{f}(\{f(x)\})$.
\end{Proposizione}
\begin{dimostrazione}
    Per ogni $y \in a$:\quad $y \in [x]_{\sim_f} \leftrightarrow x \sim_f y \leftrightarrow f(y) = f(x) \leftrightarrow f(y) \in \{f(x)\}
    \leftrightarrow y \in \overleftarrow{f}(\{f(x)\})$. Per Estensionalità i due insiemi coincidono.
\end{dimostrazione}
\end{concetto}

\begin{spiegazione}
    La classe di $x$ è l'insieme di \textbf{tutti gli elementi che finiscono dove finisce $x$}: si parte da $f(x)$ e si
    seguono le frecce all'indietro. Nell'esempio di $f(n) = n^2$: $[2] = \overleftarrow{f}(\{4\}) = \{-2, 2\}$.
\end{spiegazione}

\dfn[proiezione-canonica]{Proiezione canonica}{
    Sia $\sim$ una relazione di equivalenza su $a$. La funzione
    \[ \pi \colon x \in a \mapsto [x]_\sim \in a/{\sim} \]
    si chiama \textbf{proiezione canonica} di $a$ su $a/{\sim}$.
}

\begin{concetto}
\begin{Proposizione}{Proprietà della proiezione canonica}{proiezione}
    \begin{prerequisiti}
        \dip{Definizione}{def:proiezione-canonica}
        \dip{Definizione}{def:suriettiva}
        \dip{Teorema}{teo:proprieta-classi}
        \dip{Definizione}{def:nucleo}
    \end{prerequisiti}
    Sia $\sim$ una relazione di equivalenza su $a$ e sia $\pi$ la proiezione canonica. Allora $\pi$ è suriettiva e il suo nucleo
    di equivalenza è proprio $\sim$: $\pi(x) = \pi(y) \leftrightarrow x \sim y$.
\end{Proposizione}
\begin{dimostrazione}
    \begin{caso}{Suriettiva}
        Ogni elemento di $a/{\sim}$ è una classe $[x]$ con $x \in a$, e $[x] = \pi(x)$.
    \end{caso}
    \begin{caso}{Il nucleo è $\sim$}
        $\pi(x) = \pi(y) \leftrightarrow [x] = [y] \leftrightarrow x \sim y$ per il Teorema~\ref{teo:proprieta-classi}, punto 2.
    \end{caso}
\end{dimostrazione}
\end{concetto}

\begin{spiegazione}
    $\pi$ «schiaccia» ogni stanza in un unico punto: manda ogni elemento nella sua classe. Nell'esempio delle stanze,
    $\pi(1) = \pi(2) = \pi(3) = \{1, 2, 3\}$. Il fatto che il nucleo di $\pi$ sia $\sim$ dice che \textbf{ogni relazione di
    equivalenza è il nucleo di qualche funzione}: non c'è differenza tra «dare un'equivalenza» e «dare una funzione e guardare
    quali elementi hanno la stessa immagine».
\end{spiegazione}

% ──────────────────────────────────────────────────────────────
\section{Il teorema fondamentale di omomorfismo per insiemi}\label{sec:omomorfismo-insiemi}

Ogni funzione si può spezzare in tre passi elementari. È il primo esempio di un'idea che ritroveremo spesso: quella di
\textbf{passare al quoziente}.

\begin{concetto}
\begin{Teorema}{Teorema fondamentale di omomorfismo per insiemi}{omomorfismo-insiemi}
    \begin{prerequisiti}
        \dip{Definizione}{def:nucleo}
        \dip{Definizione}{def:proiezione-canonica}
        \dip{Definizione}{def:funzioni-notevoli}
        \dip{Definizione}{def:funzione}
        \dip{Definizione}{def:biettiva}
        \dip{Teorema}{teo:proprieta-classi}
        \dip{Corollario}{cor:terna}
    \end{prerequisiti}
    Sia $f \colon a \to b$ una funzione, siano $\pi \colon a \to a/{\sim_f}$ la proiezione canonica e $\iota \colon \mathrm{Im} f \to b$
    l'immersione. Allora:
    \begin{enumerate}
        \item la corrispondenza $\tilde{f} \colon [x]_{\sim_f} \in a/{\sim_f} \mapsto f(x) \in \mathrm{Im} f$ è una funzione \textbf{biettiva};
        \item $f = \iota \circ \tilde{f} \circ \pi$, cioè il diagramma seguente è commutativo:
    \end{enumerate}
    \begin{center}
    \begin{tikzcd}[column sep=2.4cm, row sep=1.1cm]
        a \arrow[r, "f"] \arrow[d, "\pi"'] & b \\
        a/{\sim_f} \arrow[r, "\tilde{f}"'] & \mathrm{Im} f \arrow[u, hook, "\iota"']
    \end{tikzcd}
    \end{center}
\end{Teorema}
\begin{dimostrazione}
    \begin{caso}{$\tilde{f}$ è una funzione (è «ben posta»)}
        $\tilde{f}$ è la corrispondenza $(a/{\sim_f}, \mathrm{Im} f, \tilde{g})$ con $\tilde{g} = \{([x], f(x)) \mid x \in a\}$. Ogni classe
        ha un'immagine, perché ogni classe è $[x]$ per qualche $x$. L'immagine è unica: se $[x] = [y]$, allora $x \sim_f y$, cioè
        $f(x) = f(y)$. Quindi \textbf{il valore non dipende dal rappresentante} scelto.
    \end{caso}
    \begin{caso}{$\tilde{f}$ è iniettiva}
        Se $\tilde{f}([x]) = \tilde{f}([y])$, allora $f(x) = f(y)$, cioè $x \sim_f y$, e quindi $[x] = [y]$
        (Teorema~\ref{teo:proprieta-classi}, punto 2).
    \end{caso}
    \begin{caso}{$\tilde{f}$ è suriettiva}
        Se $w \in \mathrm{Im} f$, esiste $x \in a$ con $f(x) = w$, e allora $\tilde{f}([x]) = w$.
    \end{caso}
    \begin{caso}{$f = \iota \circ \tilde{f} \circ \pi$}
        Le due funzioni hanno lo stesso dominio $a$ e lo stesso codominio $b$, e per ogni $x \in a$:
        \[ \iota\big(\tilde{f}(\pi(x))\big) = \iota\big(\tilde{f}([x])\big) = \iota(f(x)) = f(x). \]
        Hanno quindi lo stesso grafico e sono la stessa terna (Corollario~\ref{cor:terna}).
    \end{caso}
\end{dimostrazione}
\end{concetto}

\begin{spiegazione}
    Prendiamo $a = \{1, 2, 3, 4, 5\}$, $b = \{p, q, r, s\}$ e $f$ con $f(1) = f(2) = p$, $f(3) = q$, $f(4) = f(5) = r$. Allora $f$
    fa tre cose, una dopo l'altra:
    \begin{center}
    \begin{tikzpicture}
        \foreach \e/\y in {1/2.0, 2/1.5, 3/0.5, 4/-0.5, 5/-1.0} \node[elem] (x\e) at (0,\y) {$\e$};
        \blocco[6mm]{classeA}{(x1.center) -- (x2.center)}
        \blocco[6mm]{classeB}{(x3.center) -- ++(0.001,0)}
        \blocco[6mm]{classeC}{(x4.center) -- (x5.center)}
        \node[elem] (c1) at (3.1,1.75) {$\{1, 2\}$};
        \node[elem] (c2) at (3.1,0.5) {$\{3\}$};
        \node[elem] (c3) at (3.1,-0.75) {$\{4, 5\}$};
        \node[elem] (i1) at (6.0,1.75) {$p$}; \node[elem] (i2) at (6.0,0.5) {$q$}; \node[elem] (i3) at (6.0,-0.75) {$r$};
        \node[elem] (b1) at (8.6,1.75) {$p$}; \node[elem] (b2) at (8.6,0.5) {$q$}; \node[elem] (b3) at (8.6,-0.75) {$r$};
        \node[elem] (b4) at (8.6,-1.7) {$s$};
        \draw[fun3] (x1) -- (c1); \draw[fun3] (x2) -- (c1); \draw[fun3] (x3) -- (c2); \draw[fun3] (x4) -- (c3); \draw[fun3] (x5) -- (c3);
        \draw[fun] (c1) -- (i1); \draw[fun] (c2) -- (i2); \draw[fun] (c3) -- (i3);
        \draw[fun2, {Hooks[right]}-{Stealth[length=2.2mm]}] (i1) -- (b1);
        \draw[fun2, {Hooks[right]}-{Stealth[length=2.2mm]}] (i2) -- (b2);
        \draw[fun2, {Hooks[right]}-{Stealth[length=2.2mm]}] (i3) -- (b3);
        \node[font=\small] at (0,2.75) {$a$};
        \node[font=\small] at (3.1,2.75) {$a/{\sim_f}$};
        \node[font=\small] at (6.0,2.75) {$\mathrm{Im} f$};
        \node[font=\small] at (8.6,2.75) {$b$};
        \node[font=\small, seagreen!60!black] at (1.55,2.75) {$\pi$};
        \node[font=\small, primary!80!black] at (4.55,2.75) {$\tilde{f}$};
        \node[font=\small, cherryred!85!black] at (7.3,2.75) {$\iota$};
    \end{tikzpicture}
    \end{center}
    \begin{passi}
        \item $\pi$ \textbf{raggruppa} gli elementi che hanno la stessa immagine: $1$ e $2$ finiscono nella stessa classe.
              È suriettiva, ma in generale non iniettiva.
        \item $\tilde{f}$ dà a ogni gruppo il suo valore: ora gruppi diversi hanno valori diversi, quindi è \textbf{iniettiva}; e
              arriva solo sui valori effettivamente presi ($\mathrm{Im} f$), quindi è anche \textbf{suriettiva}.
        \item $\iota$ \textbf{rimette} i valori nel codominio $b$, dove c'è anche $s$, che nessuno raggiunge. È iniettiva, ma
              in generale non suriettiva.
    \end{passi}
    Tutto ciò che impedisce a $f$ di essere iniettiva è finito in $\pi$, tutto ciò che le impedisce di essere suriettiva è
    finito in $\iota$; in mezzo resta una biezione.
\end{spiegazione}

\begin{concetto}
\begin{Corollario}{Ogni funzione è composta di una suriettiva e di un'iniettiva}{iniettiva-suriettiva}
    \begin{prerequisiti}
        \dip{Teorema}{teo:omomorfismo-insiemi}
        \dip{Proposizione}{prop:proiezione}
        \dip{Teorema}{teo:comp-iniettive}
    \end{prerequisiti}
    Ogni funzione $f \colon a \to b$ si scrive come $f = j \circ \pi$, con $\pi$ suriettiva e $j$ iniettiva.
\end{Corollario}
\begin{dimostrazione}
    Per il Teorema~\ref{teo:omomorfismo-insiemi}, $f = \iota \circ \tilde{f} \circ \pi = (\iota \circ \tilde{f}) \circ \pi$. La proiezione
    $\pi$ è suriettiva (Proposizione~\ref{prop:proiezione}); $j = \iota \circ \tilde{f}$ è iniettiva perché composta di due funzioni
    iniettive (Teorema~\ref{teo:comp-iniettive}): $\tilde{f}$ è biettiva e l'immersione $\iota$ è iniettiva.
\end{dimostrazione}
\end{concetto}

\begin{spiegazione}
    Con $f \colon n \in \mbZ \mapsto n^2 \in \mbZ$:
    \begin{passi}
        \item $\pi \colon n \mapsto [n] = \{n, -n\}$ raggruppa ogni intero con il suo opposto (suriettiva su
              $\mbZ/{\sim_f} = \{\{0\}, \{1, -1\}, \{2, -2\}, \dots\}$);
        \item $j \colon \{n, -n\} \mapsto n^2 \in \mbZ$ manda ogni coppia nel suo quadrato: coppie diverse hanno quadrati diversi,
              quindi è iniettiva (ma non suriettiva: $2$ non è un quadrato).
    \end{passi}
    Infatti $j(\pi(-3)) = j(\{3, -3\}) = 9 = f(-3)$.
\end{spiegazione}

% ──────────────────────────────────────────────────────────────
\section{Partizioni}\label{sec:partizioni}

Le proprietà fondamentali delle classi (Teorema~\ref{teo:proprieta-classi}) dicono che un insieme quoziente divide l'insieme
in pezzi non vuoti, disgiunti, che lo ricoprono tutto. Diamo un nome a questo tipo di suddivisione.

\dfn[partizione]{Partizione}{
    Sia $a$ un insieme. Un insieme $p$ di parti di $a$ (cioè $p \subseteq \mcP(a)$) si dice \textbf{partizione} di $a$ se
    \begin{enumerate}
        \item $\forall x \in p\,(x \neq \varnothing)$: i pezzi sono non vuoti;
        \item $\forall x, y \in p\,(x \neq y \rightarrow x \cap y = \varnothing)$: i pezzi sono a due a due disgiunti;
        \item $\bigcup p = a$: i pezzi ricoprono $a$.
    \end{enumerate}
    Gli elementi di $p$ si chiamano anche \textbf{blocchi} della partizione. Indichiamo con $\mathrm{Part}(a)$ l'insieme delle
    partizioni di $a$ (è una parte di $\mcP(\mcP(a))$).
}

\begin{esempio}{Le cinque partizioni di $\{0, 1, 2\}$}
    \begin{center}
    \begin{tikzpicture}
        \foreach \k/\px in {1/0, 2/3.0, 3/6.0, 4/9.0, 5/12.0} {
            \begin{scope}[xshift=\px cm]
                \node[elem] (z\k) at (0.55,0.95) {$0$};
                \node[elem] (u\k) at (0,0) {$1$};
                \node[elem] (d\k) at (1.1,0) {$2$};
            \end{scope}
        }
        \blocco{classeA}{(z1.center) -- ++(0.001,0)} \blocco{classeB}{(u1.center) -- ++(0.001,0)} \blocco{classeC}{(d1.center) -- ++(0.001,0)}
        \blocco{classeA}{(z2.center) -- (u2.center)} \blocco{classeC}{(d2.center) -- ++(0.001,0)}
        \blocco{classeA}{(z3.center) -- (d3.center)} \blocco{classeB}{(u3.center) -- ++(0.001,0)}
        \blocco{classeB}{(u4.center) -- (d4.center)} \blocco{classeA}{(z4.center) -- ++(0.001,0)}
        \blocco{classeD}{(z5.center) -- (u5.center) -- (d5.center) -- cycle}
        \node[font=\scriptsize, align=center] at (0.55,-0.85) {$\{\{0\}, \{1\}, \{2\}\}$};
        \node[font=\scriptsize, align=center] at (3.55,-0.85) {$\{\{0, 1\}, \{2\}\}$};
        \node[font=\scriptsize, align=center] at (6.55,-0.85) {$\{\{0, 2\}, \{1\}\}$};
        \node[font=\scriptsize, align=center] at (9.55,-0.85) {$\{\{1, 2\}, \{0\}\}$};
        \node[font=\scriptsize, align=center] at (12.55,-0.85) {$\{\{0, 1, 2\}\}$};
    \end{tikzpicture}
    \end{center}
    Invece $\{\{0, 1\}, \{1, 2\}\}$ non è una partizione (i blocchi non sono disgiunti), $\{\{0\}, \{1\}\}$ non lo è (non ricopre $2$)
    e $\{\{0, 1, 2\}, \varnothing\}$ nemmeno (contiene un blocco vuoto). Altri esempi: $\{\text{pari}, \text{dispari}\}$ è una
    partizione di $\mbZ$; l'unica partizione di $\varnothing$ è $\varnothing$ (le tre condizioni sono vere a vuoto).
\end{esempio}

\begin{concetto}
\begin{Corollario}{Ogni insieme quoziente è una partizione}{quoziente-partizione}
    \begin{prerequisiti}
        \dip{Teorema}{teo:proprieta-classi}
        \dip{Definizione}{def:partizione}
    \end{prerequisiti}
    Se $\sim$ è una relazione di equivalenza su $a$, allora $a/{\sim} \in \mathrm{Part}(a)$.
\end{Corollario}
\begin{dimostrazione}
    Ogni classe è una parte di $a$, quindi $a/{\sim} \subseteq \mcP(a)$. Le tre condizioni della Definizione~\ref{def:partizione}
    sono i punti 1, 3 e 4 del Teorema~\ref{teo:proprieta-classi}.
\end{dimostrazione}
\end{concetto}

\begin{spiegazione}
    È la casa con le stanze vista dall'alto: le stanze (le classi) sono i blocchi della partizione. Ad esempio $\mbZ/{\equiv_3}$
    è la partizione di $\mbZ$ nei tre colori della figura sulla congruenza modulo $3$.
\end{spiegazione}

Vale anche il viceversa: ogni partizione è l'insieme quoziente di una (e una sola) relazione di equivalenza. Le relazioni di
equivalenza e le partizioni sono quindi \textbf{due modi di dare la stessa informazione}.

\begin{concetto}
\begin{Teorema}{Teorema fondamentale su relazioni di equivalenza e partizioni}{equivalenze-partizioni}
    \begin{prerequisiti}
        \dip{Corollario}{cor:quoziente-partizione}
        \dip{Teorema}{teo:proprieta-classi}
        \dip{Definizione}{def:partizione}
        \dip{Definizione}{def:biettiva}
        \dip{Corollario}{cor:terna}
    \end{prerequisiti}
    Per ogni insieme $a$ la funzione
    \[ \Phi \colon {\sim} \in \mathrm{Eq}(a) \mapsto a/{\sim} \in \mathrm{Part}(a) \]
    è biettiva. La sua inversa associa alla partizione $p$ la relazione «stare nello stesso blocco»:
    $x \sim_p y \leftrightarrow \exists w \in p\,(x \in w \wedge y \in w)$.
\end{Teorema}
\begin{dimostrazione}
    $\Phi$ è una funzione a valori in $\mathrm{Part}(a)$ per il Corollario~\ref{cor:quoziente-partizione}.
    \begin{caso}{$\Phi$ è iniettiva}
        Siano $\sim_1, \sim_2 \in \mathrm{Eq}(a)$ con $a/{\sim_1} = a/{\sim_2}$, e sia $x \sim_1 y$. Allora $[x]_{\sim_1} = [y]_{\sim_1}$
        è un elemento di $a/{\sim_1} = a/{\sim_2}$, quindi esiste $z \in a$ con $[x]_{\sim_1} = [z]_{\sim_2}$. Perciò $x, y \in [z]_{\sim_2}$,
        cioè $z \sim_2 x$ e $z \sim_2 y$, e per simmetria e transitività $x \sim_2 y$. Scambiando i ruoli, $x \sim_2 y \rightarrow x \sim_1 y$.
        Le due relazioni hanno lo stesso grafico, quindi sono uguali (Corollario~\ref{cor:terna}).
    \end{caso}
    \begin{caso}{$\Phi$ è suriettiva}
        Sia $p \in \mathrm{Part}(a)$ e sia $\sim_p$ come nell'enunciato.
        \begin{itemize}
            \item[] \textbf{Riflessiva}: poiché $\bigcup p = a$, ogni $x \in a$ sta in un blocco $w$, quindi $x \sim_p x$.
            \item[] \textbf{Simmetrica}: $x \in w \wedge y \in w$ equivale a $y \in w \wedge x \in w$.
            \item[] \textbf{Transitiva}: se $x, y \in w_1$ e $y, z \in w_2$ con $w_1, w_2 \in p$, allora $y \in w_1 \cap w_2$; i blocchi
                    sono a due a due disgiunti, quindi $w_1 = w_2$ e $x, z$ stanno nello stesso blocco.
        \end{itemize}
        Resta da vedere che $a/{\sim_p} = p$. Sia $x \in a$ e sia $w_x$ il blocco che contiene $x$ (è unico, perché i blocchi sono
        disgiunti). Allora $y \in [x]_{\sim_p} \leftrightarrow$ $x$ e $y$ stanno in uno stesso blocco $\leftrightarrow y \in w_x$, cioè
        $[x]_{\sim_p} = w_x$. Quindi ogni classe è un blocco; e ogni blocco $w$, non essendo vuoto, contiene un $x$ ed è $w = w_x = [x]_{\sim_p}$.
        Dunque $\Phi(\sim_p) = a/{\sim_p} = p$.
    \end{caso}
\end{dimostrazione}
\end{concetto}

\begin{spiegazione}
    Ecco la stessa informazione nelle due forme, su $a = \{1, 2, 3, 4, 5\}$. A sinistra il grafo della relazione di
    equivalenza (le doppie frecce indicano andata e ritorno; ogni elemento ha anche il suo cappio, che omettiamo); a destra
    la partizione corrispondente:
    \begin{center}
    \begin{tikzpicture}
        \node[rv] (1) at (0,1.2) {$1$}; \node[rv] (2) at (1.3,1.2) {$2$}; \node[rv] (3) at (0.65,0) {$3$};
        \node[rv] (4) at (2.7,1.2) {$4$}; \node[rv] (5) at (2.7,0) {$5$};
        \draw[{Stealth[length=2mm]}-{Stealth[length=2mm]}, line width=0.7pt, black!70] (1) -- (2);
        \draw[{Stealth[length=2mm]}-{Stealth[length=2mm]}, line width=0.7pt, black!70] (2) -- (3);
        \draw[{Stealth[length=2mm]}-{Stealth[length=2mm]}, line width=0.7pt, black!70] (1) -- (3);
        \draw[{Stealth[length=2mm]}-{Stealth[length=2mm]}, line width=0.7pt, black!70] (4) -- (5);
        \node[font=\small] at (1.35,-0.8) {la relazione $\sim$};
        \draw[-{Stealth}, line width=1pt, primary] (3.6,0.6) -- node[above, font=\small] {$\Phi$} (5.0,0.6);
        \draw[{Stealth}-, line width=1pt, seagreen!65!black] (3.6,0.2) -- node[below, font=\small] {$\Phi^{-1}$} (5.0,0.2);
        \begin{scope}[xshift=5.9cm]
            \node[elem] (p1) at (0,1.2) {$1$}; \node[elem] (p2) at (1.3,1.2) {$2$}; \node[elem] (p3) at (0.65,0) {$3$};
            \node[elem] (p4) at (2.7,1.2) {$4$}; \node[elem] (p5) at (2.7,0) {$5$};
            \blocco[7mm]{classeA}{(p1.center) -- (p2.center) -- (p3.center) -- cycle}
            \blocco[7mm]{classeB}{(p4.center) -- (p5.center)}
            \node[font=\small] at (1.35,-0.8) {la partizione $\{\{1, 2, 3\}, \{4, 5\}\}$};
        \end{scope}
    \end{tikzpicture}
    \end{center}
    \begin{passi}
        \item Da sinistra a destra ($\Phi$): ogni «gruppo di elementi tutti collegati tra loro» diventa un blocco.
        \item Da destra a sinistra ($\Phi^{-1}$): si collegano tra loro (con cappi e frecce in entrambi i versi) gli elementi di ogni
              blocco, e solo quelli.
        \item L'iniettività dice che due relazioni diverse danno partizioni diverse; la suriettività che ogni partizione si
              ottiene così.
    \end{passi}
\end{spiegazione}

% ──────────────────────────────────────────────────────────────
\section{Trovare tutte le relazioni di equivalenza}\label{sec:trovare-equivalenze}

Grazie al Teorema~\ref{teo:equivalenze-partizioni}, \textbf{trovare tutte le relazioni di equivalenza su un insieme equivale a
trovare tutte le sue partizioni}, che è molto più facile: basta disegnare i blocchi.

\begin{esempio}{Le relazioni di equivalenza su $\{0, 1, 2\}$}
    Le partizioni di $\{0, 1, 2\}$ sono le cinque disegnate prima, quindi le relazioni di equivalenza sono cinque. I loro grafici:
    \begin{center}
    \renewcommand{\arraystretch}{1.3}
    \begin{tabular}{l|l}
        \rowcolor{primary!12} \textbf{Partizione} & \textbf{Grafico della relazione} (con $D = \mathrm{diag}(\{0, 1, 2\})$) \\
        $\{\{0\}, \{1\}, \{2\}\}$ & $D$ \quad (l'uguaglianza) \\
        $\{\{0, 1\}, \{2\}\}$ & $D \cup \{(0, 1), (1, 0)\}$ \\
        $\{\{0, 2\}, \{1\}\}$ & $D \cup \{(0, 2), (2, 0)\}$ \\
        $\{\{1, 2\}, \{0\}\}$ & $D \cup \{(1, 2), (2, 1)\}$ \\
        $\{\{0, 1, 2\}\}$ & $\{0, 1, 2\} \times \{0, 1, 2\}$ \quad (la relazione totale) \\
    \end{tabular}
    \end{center}
\end{esempio}

\begin{esempio}{Equivalenze con dei vincoli}
    Siano $a, b, c, d$ elementi a due a due distinti e $s = \{a, b, c, d\}$.
    \begin{itemize}
        \item Le equivalenze su $s$ con $[a] = [b] = [c]$ corrispondono alle partizioni in cui $a, b, c$ stanno nello stesso blocco:
              $\{\{a, b, c\}, \{d\}\}$ e $\{\{a, b, c, d\}\}$. Sono \textbf{due}.
        \item Quelle con $[a] = [b]$ corrispondono alle partizioni con $a$ e $b$ nello stesso blocco: si parte da $\{a, b\}$ e si
              decide dove mettere $c$ e $d$. Si ottengono $\{\{a, b\}, \{c\}, \{d\}\}$, $\{\{a, b\}, \{c, d\}\}$, $\{\{a, b, c\}, \{d\}\}$,
              $\{\{a, b, d\}, \{c\}\}$ e $\{\{a, b, c, d\}\}$: sono \textbf{cinque}.
    \end{itemize}
\end{esempio}

\begin{esempio}{Un esempio passo passo}
    Sia $a = \{n \in \mbN \mid n \leq 7\}$. Troviamo tutte le relazioni di equivalenza $\rho$ su $a$ tali che
    $0 \rho 7$, $1 \rho 4$, $\{3, 4, 5, 7\} \subseteq [2]_\rho$ e $0 \rho 6$. Ogni vincolo costringe certi elementi nello stesso blocco:
    \begin{passi}
        \item $0 \rho 7$: il blocco di $0$ contiene $\{0, 7\}$.
        \item $1 \rho 4$: il blocco di $1$ contiene $\{1, 4\}$.
        \item $\{3, 4, 5, 7\} \subseteq [2]$: $2, 3, 4, 5, 7$ stanno nello stesso blocco. Ma $7$ è con $0$ e $4$ è con $1$, quindi
              lo stesso blocco contiene $\{0, 1, 2, 3, 4, 5, 7\}$.
        \item $0 \rho 6$: anche $6$ entra nel blocco, che diventa tutto $a$.
    \end{passi}
    C'è quindi \textbf{una sola} relazione con queste proprietà: la relazione totale su $a$.
\end{esempio}


  % Ordini larghi e stretti, insiemi ordinati, minimi e minimali, crescenti, Hasse, minoranti
  \chapter{Relazioni d'ordine}\label{cap:ordini}

Le relazioni d'ordine formalizzano l'idea di «venire prima» o di «essere più piccolo»: $\leq$ tra i numeri, $\subseteq$ tra
gli insiemi, «divide» tra i numeri naturali. A differenza dei numeri, però, in un ordine generico due elementi possono anche
\textbf{non essere confrontabili}: $\{1\}$ e $\{2\}$ non sono uno contenuto nell'altro. In questo capitolo vediamo il legame tra
ordini larghi e stretti, gli elementi «estremi» (minimo, massimo, minimali, massimali), le funzioni che rispettano l'ordine e
un modo comodo di disegnare gli insiemi ordinati finiti: i \textbf{diagrammi di Hasse}.

% ──────────────────────────────────────────────────────────────
\section{Ordini larghi e ordini stretti}\label{sec:largo-stretto}

Ricordiamo (Definizione~\ref{def:tipi-relazioni}) che un ordine \textbf{largo} è riflessivo, antisimmetrico e transitivo, come
$\leq$, mentre un ordine \textbf{stretto} è antiriflessivo e transitivo, come $<$. Da un tipo si passa all'altro aggiungendo o
togliendo l'uguaglianza.

\dfn[ol-os]{Gli insiemi $\mathrm{OL}(s)$ e $\mathrm{OS}(s)$}{
    Sia $s$ un insieme. Indichiamo con $\mathrm{OL}(s)$ l'insieme delle relazioni d'ordine largo su $s$ e con $\mathrm{OS}(s)$
    quello delle relazioni d'ordine stretto.
    \begin{itemize}
        \item Se $\rho \in \mathrm{OL}(s)$, definiamo la relazione $\rho^\wedge$ su $s$ ponendo \ $x \rho^\wedge y \leftrightarrow (x \rho y \wedge x \neq y)$.
        \item Se $\rho \in \mathrm{OS}(s)$, definiamo la relazione $\rho^\vee$ su $s$ ponendo \ $x \rho^\vee y \leftrightarrow (x \rho y \vee x = y)$.
    \end{itemize}
}

\begin{concetto}
\begin{Teorema}{Ordini larghi e ordini stretti si corrispondono}{largo-stretto}
    \begin{prerequisiti}
        \dip{Definizione}{def:ol-os}
        \dip{Definizione}{def:tipi-relazioni}
        \dip{Proposizione}{prop:distributivita}
        \dip{Teorema}{teo:biettiva-inversa}
    \end{prerequisiti}
    Sia $s$ un insieme. Allora $\rho \in \mathrm{OL}(s) \mapsto \rho^\wedge \in \mathrm{OS}(s)$ è una funzione biettiva, e la sua inversa è
    $\rho \in \mathrm{OS}(s) \mapsto \rho^\vee \in \mathrm{OL}(s)$.
\end{Teorema}
\begin{dimostrazione}
    \facoltativa
    \begin{caso}{Se $\rho \in \mathrm{OL}(s)$, allora $\rho^\wedge \in \mathrm{OS}(s)$}
        \textbf{Antiriflessiva}: $x \rho^\wedge x$ richiederebbe $x \neq x$, che è falso. \textbf{Transitiva}: siano $x \rho^\wedge y$ e
        $y \rho^\wedge z$, cioè $x \rho y$, $x \neq y$, $y \rho z$, $y \neq z$. Per la transitività di $\rho$, $x \rho z$. Se fosse $x = z$,
        avremmo $x \rho y$ e $y \rho x$, quindi $x = y$ per l'antisimmetria: assurdo. Dunque $x \neq z$ e $x \rho^\wedge z$.
    \end{caso}
    \begin{caso}{Se $\rho \in \mathrm{OS}(s)$, allora $\rho^\vee \in \mathrm{OL}(s)$}
        \textbf{Riflessiva}: $x = x$. \textbf{Antisimmetrica}: siano $x \rho^\vee y$ e $y \rho^\vee x$; se fosse $x \neq y$ avremmo
        $x \rho y$ e $y \rho x$, quindi $x \rho x$ per transitività, contro l'antiriflessività; dunque $x = y$. \textbf{Transitiva}: siano
        $x \rho^\vee y$ e $y \rho^\vee z$. Se $x = y$ o $y = z$ la tesi è una delle ipotesi; altrimenti $x \rho y$ e $y \rho z$, quindi $x \rho z$.
    \end{caso}
    \begin{caso}{Le due funzioni sono una l'inversa dell'altra}
        Sia $\rho \in \mathrm{OL}(s)$. Usando la distributività di $\vee$ rispetto a $\wedge$:
        \[ x (\rho^\wedge)^\vee y \leftrightarrow (x \rho y \wedge x \neq y) \vee x = y \leftrightarrow (x \rho y \vee x = y) \wedge (x \neq y \vee x = y)
           \leftrightarrow x \rho y \vee x = y \leftrightarrow x \rho y, \]
        dove l'ultimo passaggio usa la riflessività di $\rho$. Sia $\rho \in \mathrm{OS}(s)$. Usando la distributività di $\wedge$ rispetto a $\vee$:
        \[ x (\rho^\vee)^\wedge y \leftrightarrow (x \rho y \vee x = y) \wedge x \neq y \leftrightarrow (x \rho y \wedge x \neq y) \vee (x = y \wedge x \neq y)
           \leftrightarrow x \rho y \wedge x \neq y \leftrightarrow x \rho y, \]
        dove l'ultimo passaggio usa l'antiriflessività di $\rho$ ($x \rho y$ implica $x \neq y$). Due funzioni una inversa dell'altra
        sono biettive (Teorema~\ref{teo:biettiva-inversa}).
    \end{caso}
\end{dimostrazione}
\end{concetto}

\begin{spiegazione}
    Sul grafo con frecce, passare da un ordine largo a uno stretto vuol dire \textbf{togliere tutti i cappi}; tornare indietro vuol
    dire \textbf{rimetterli}. Le altre frecce non cambiano. Ecco $\leq$ e $<$ su $\{1, 2, 3\}$:
    \begin{center}
    \begin{tikzpicture}
        \node[rv] (1) at (0,0) {$1$}; \node[rv] (2) at (1.6,0) {$2$}; \node[rv] (3) at (3.2,0) {$3$};
        \draw[rfreccia] (1) to[rcappio] (1); \draw[rfreccia] (2) to[rcappio] (2); \draw[rfreccia] (3) to[rcappio] (3);
        \draw[rfreccia] (1) to (2); \draw[rfreccia] (2) to (3); \draw[rfreccia] (1) to[bend right=35] (3);
        \node[font=\small] at (1.6,-1.15) {$\leq$ (ordine largo)};
        \draw[-{Stealth}, line width=1pt, primary] (4.0,0.25) -- node[above, font=\small] {$\rho \mapsto \rho^\wedge$} (5.6,0.25);
        \draw[{Stealth}-, line width=1pt, seagreen!65!black] (4.0,-0.25) -- node[below, font=\small] {$\rho^\vee \mapsfrom \rho$} (5.6,-0.25);
        \begin{scope}[xshift=6.4cm]
            \node[rv] (1) at (0,0) {$1$}; \node[rv] (2) at (1.6,0) {$2$}; \node[rv] (3) at (3.2,0) {$3$};
            \draw[rfreccia] (1) to (2); \draw[rfreccia] (2) to (3); \draw[rfreccia] (1) to[bend right=35] (3);
            \node[font=\small] at (1.6,-1.15) {$<$ (ordine stretto)};
        \end{scope}
    \end{tikzpicture}
    \end{center}
    Allo stesso modo $(\subseteq)^\wedge = {\subset}$ e $(\subset)^\vee = {\subseteq}$. Il punto delicato della dimostrazione è che, togliendo i
    cappi, la transitività non si perde: un percorso $x \to y \to z$ con $x \neq y$ e $y \neq z$ non può tornare in $x$, perché un
    ordine largo non ha andate e ritorno tra elementi distinti (antisimmetria).
\end{spiegazione}

Nel seguito lavoreremo quasi sempre con ordini larghi: un ordine stretto si può sempre ricavare togliendo l'uguaglianza.

% ──────────────────────────────────────────────────────────────
\section{Insiemi ordinati}\label{sec:insiemi-ordinati}

\dfn[insieme-ordinato]{Insieme ordinato, ordine indotto}{
    Sia $s$ un insieme e sia $\rho = (s, s, g) \in \mathrm{OL}(s)$.
    \begin{itemize}
        \item La coppia $(s, \rho)$ si dice \textbf{insieme ordinato}.
        \item Se $t \subseteq s$, la relazione $\rho_t = (t, t, g \cap (t \times t))$ si dice \textbf{ordine indotto} da $\rho$ su $t$, e
              $(t, \rho_t)$ si dice \textbf{sottoinsieme ordinato} di $(s, \rho)$. Con abuso di notazione scriveremo spesso $\rho$ invece di $\rho_t$.
    \end{itemize}
    Per un ordine generico useremo spesso il simbolo $\leq$, e scriveremo $x < y$ per $x \leq y \wedge x \neq y$ (cioè $x \leq^\wedge y$).
}

\begin{concetto}
\begin{Proposizione}{L'ordine indotto è un ordine}{ordine-indotto}
    \begin{prerequisiti}
        \dip{Definizione}{def:insieme-ordinato}
        \dip{Definizione}{def:tipi-relazioni}
    \end{prerequisiti}
    Se $(s, \rho)$ è un insieme ordinato e $t \subseteq s$, allora $\rho_t \in \mathrm{OL}(t)$. Inoltre, per $x, y \in t$, $x \rho_t y \leftrightarrow x \rho y$.
\end{Proposizione}
\begin{dimostrazione}
    Per $x, y \in t$ si ha $(x, y) \in t \times t$, quindi $(x, y) \in g \cap (t \times t) \leftrightarrow (x, y) \in g$: questa è la seconda
    affermazione. Le tre proprietà di $\rho_t$ sono allora quelle di $\rho$, applicate a elementi che stanno in $t$: ad esempio se
    $x, y, z \in t$, $x \rho_t y$ e $y \rho_t z$, allora $x \rho y$ e $y \rho z$, quindi $x \rho z$, cioè $x \rho_t z$.
\end{dimostrazione}
\end{concetto}

\begin{spiegazione}
    Restringere un ordine a una parte vuol dire \textbf{cancellare dal grafo gli elementi che non stanno in $t$}, con tutte le
    frecce che li toccano. Le frecce rimaste collegano elementi di $t$ esattamente come prima, quindi riflessività, antisimmetria
    e transitività restano vere. Ad esempio la divisibilità su $\mbN$ induce un ordine su $\{1, 2, 3, 4, 6, 12\}$.
\end{spiegazione}

\dfn[confrontabili]{Elementi confrontabili, ordine totale}{
    Sia $(s, \rho)$ un insieme ordinato.
    \begin{itemize}
        \item Due elementi $x, y \in s$ si dicono \textbf{confrontabili} se $x \rho y \vee y \rho x$.
        \item Se ogni coppia di elementi di $s$ è confrontabile, $\rho$ si dice \textbf{ordine totale} e $(s, \rho)$ si dice
              \textbf{insieme totalmente ordinato} (o \textbf{catena}).
    \end{itemize}
}

\dfn[minimo-minimale]{Minimo, massimo, minimali, massimali}{
    Sia $(s, \rho)$ un insieme ordinato e sia $m \in s$. Diciamo che $m$ è
    \begin{itemize}
        \item \textbf{minimo} di $s$ se $\forall x \in s\,(m \rho x)$: sta sotto a tutti; si scrive $m = \min(s)$;
        \item \textbf{massimo} di $s$ se $\forall x \in s\,(x \rho m)$: sta sopra a tutti; si scrive $m = \max(s)$;
        \item \textbf{minimale} se $\forall x \in s\,(x \rho m \rightarrow x = m)$: non c'è nessuno strettamente sotto di lui;
        \item \textbf{massimale} se $\forall x \in s\,(m \rho x \rightarrow x = m)$: non c'è nessuno strettamente sopra di lui.
    \end{itemize}
}

\begin{concetto}
\begin{Proposizione}{Altre forme della minimalità}{forme-minimale}
    \begin{prerequisiti}
        \dip{Definizione}{def:minimo-minimale}
        \dip{Definizione}{def:confrontabili}
        \dip{Teorema}{teo:negazione-quantificatori}
        \dip{Proposizione}{prop:negimp}
    \end{prerequisiti}
    Sia $(s, \rho)$ un insieme ordinato e $m \in s$. Sono equivalenti:
    \begin{enumerate}
        \item $m$ è minimale;
        \item $\neg\exists y \in s\,(y \rho m \wedge y \neq m)$: non esiste un elemento strettamente più piccolo di $m$;
        \item $\forall x \in s\,\big((x \rho m \vee m \rho x) \rightarrow m \rho x\big)$: ogni elemento confrontabile con $m$ sta sopra $m$.
    \end{enumerate}
    Per i massimali valgono le affermazioni con $\rho$ scambiata con la duale.
\end{Proposizione}
\begin{dimostrazione}
    \begin{caso}{(1) $\leftrightarrow$ (2)}
        Negando (2) si ottiene $\exists y\,(y \rho m \wedge y \neq m)$, che è la negazione di $\forall y\,(y \rho m \rightarrow y = m)$
        (Teorema~\ref{teo:negazione-quantificatori} e Proposizione~\ref{prop:negimp}). Quindi (2) è equivalente a (1).
    \end{caso}
    \begin{caso}{(1) $\rightarrow$ (3)}
        Sia $x$ confrontabile con $m$. Se $m \rho x$ abbiamo finito; se $x \rho m$, per la minimalità $x = m$ e quindi $m \rho x$ per riflessività.
    \end{caso}
    \begin{caso}{(3) $\rightarrow$ (1)}
        Sia $x \rho m$. Allora $x$ è confrontabile con $m$, quindi per (3) $m \rho x$; per l'antisimmetria $x = m$.
    \end{caso}
\end{dimostrazione}
\end{concetto}

\begin{spiegazione}
    «Minimale» non vuol dire «sotto a tutti», ma «\textbf{sotto a tutti quelli con cui si confronta}»: un elemento minimale può
    avere accanto elementi con cui non è confrontabile. Nell'insieme ordinato qui sotto (disegnato con la convenzione dei
    diagrammi di Hasse della Sezione~\ref{sec:hasse}: «più in basso» = «più piccolo») $a$ e $b$ sono entrambi minimali, ma
    nessuno dei due è minimo: non sono confrontabili tra loro.
    \begin{center}
    \begin{tikzpicture}[hasse]
        \node[hv evid] (a) at (0,0) {$a$}; \node[hv evid] (b) at (1.6,0) {$b$};
        \node[hv] (c) at (0.8,1.1) {$c$}; \node[hv] (d) at (0.8,2.2) {$d$};
        \draw[hlato] (a) -- (c) -- (d); \draw[hlato] (b) -- (c);
        \node[font=\small, align=left, anchor=west] at (2.6,1.1)
            {$a \leq c \leq d$, \ $b \leq c$;\\ $a$ e $b$ non sono confrontabili:\\ sono minimali, ma non c'è minimo.\\ $d$ è il massimo.};
    \end{tikzpicture}
    \end{center}
\end{spiegazione}

\begin{concetto}
\begin{Teorema}{Unicità di minimo e massimo}{unicita-minimo}
    \begin{prerequisiti}
        \dip{Definizione}{def:minimo-minimale}
        \dip{Definizione}{def:tipi-relazioni}
    \end{prerequisiti}
    Sia $(s, \rho)$ un insieme ordinato. Il minimo di $s$, se esiste, è unico. Lo stesso vale per il massimo.
\end{Teorema}
\begin{dimostrazione}
    Siano $m_1$ e $m_2$ minimi di $s$. Poiché $m_1$ è minimo, $m_1 \rho m_2$; poiché $m_2$ è minimo, $m_2 \rho m_1$. Per l'antisimmetria
    $m_1 = m_2$. Per il massimo la dimostrazione è identica, con $x \rho m$ al posto di $m \rho x$.
\end{dimostrazione}
\end{concetto}

\begin{spiegazione}
    Due minimi dovrebbero stare \textbf{ciascuno sotto l'altro}; in un ordine questo succede solo se sono lo stesso elemento.
    È lo stesso ragionamento che fa dire che due numeri con $x \leq y$ e $y \leq x$ sono uguali.
\end{spiegazione}

\begin{concetto}
\begin{Teorema}{Il minimo è l'unico minimale}{minimo-minimale}
    \begin{prerequisiti}
        \dip{Definizione}{def:minimo-minimale}
        \dip{Definizione}{def:tipi-relazioni}
    \end{prerequisiti}
    Sia $(s, \rho)$ un insieme ordinato con minimo $m$. Allora $m$ è minimale ed è l'unico elemento minimale di $s$.
    Allo stesso modo il massimo, se esiste, è l'unico elemento massimale.
\end{Teorema}
\begin{dimostrazione}
    \begin{caso}{$m$ è minimale}
        Se $x \rho m$, poiché $m$ è il minimo vale anche $m \rho x$, quindi $x = m$ per l'antisimmetria.
    \end{caso}
    \begin{caso}{Ogni minimale è $m$}
        Sia $x$ minimale. Poiché $m$ è il minimo, $m \rho x$; poiché $x$ è minimale, da $m \rho x$ segue $m = x$.
    \end{caso}
    Per il massimo si ripete il ragionamento con la relazione duale.
\end{dimostrazione}
\end{concetto}

\begin{spiegazione}
    Se c'è un elemento sotto a tutti, nessun altro può essere «senza nessuno sotto»: ha sotto almeno il minimo. Attenzione al
    viceversa: avere un \textbf{unico} minimale non basta in generale per avere un minimo (lo vedremo nel
    Capitolo~\ref{cap:insiemi-ordinati}: negli insiemi finiti basta, in quelli infiniti no).
\end{spiegazione}

\dfn[ben-ordinato]{Insieme ben ordinato}{
    Un insieme ordinato $(s, \rho)$ si dice \textbf{ben ordinato} (e $\rho$ si dice \textbf{buon ordine}) se ogni parte non vuota di
    $s$ ha minimo (rispetto all'ordine indotto).
}

\begin{concetto}
\begin{Teorema}{Ogni buon ordine è totale}{ben-ordinato-totale}
    \begin{prerequisiti}
        \dip{Definizione}{def:ben-ordinato}
        \dip{Definizione}{def:confrontabili}
        \dip{Proposizione}{prop:ordine-indotto}
    \end{prerequisiti}
    Ogni insieme ben ordinato è totalmente ordinato.
\end{Teorema}
\begin{dimostrazione}
    Sia $(s, \rho)$ ben ordinato e siano $x, y \in s$. La parte $\{x, y\}$ non è vuota, quindi ha un minimo $m$, che è $x$ oppure $y$.
    Se $m = x$, allora $x \rho y$; se $m = y$, allora $y \rho x$. In ogni caso $x$ e $y$ sono confrontabili.
\end{dimostrazione}
\end{concetto}

\begin{spiegazione}
    Il buon ordine chiede il minimo di \textbf{ogni} parte non vuota, in particolare delle parti con due elementi: e il minimo di
    $\{x, y\}$ è proprio quello dei due che sta sotto l'altro. Il viceversa è falso: $(\mbZ, \leq)$ è totalmente ordinato ma la
    parte $\{\dots, -2, -1, 0\}$ non ha minimo.
\end{spiegazione}

\begin{esempio}{Insiemi ordinati notevoli}
    \begin{itemize}
        \item $(\mbN, \leq)$ è ben ordinato: ogni parte non vuota di $\mbN$ ha un elemento più piccolo di tutti.
        \item $(\mbZ, \leq)$ è totalmente ordinato ma non ben ordinato, e non ha né minimo né massimo.
        \item Per ogni insieme $s$, $(\mcP(s), \subseteq)$ ha minimo $\varnothing$ e massimo $s$. Se $s$ ha due elementi distinti $x$ e $y$,
              non è totalmente ordinato: $\{x\}$ e $\{y\}$ non sono confrontabili.
        \item $(\mbN, \mid)$ ha minimo $1$ (che divide tutti) e massimo $0$ (che è diviso da tutti: $0 = 0 \cdot n$), ma non è
              totalmente ordinato ($2 \nmid 3$ e $3 \nmid 2$).
        \item In $(\mbN \setminus \{1\}, \mid)$ non c'è minimo; gli elementi minimali sono i numeri che non hanno divisori
              diversi da $1$ e da sé stessi, cioè i numeri primi.
    \end{itemize}
\end{esempio}

\begin{concetto}
\begin{Proposizione}{Quando l'insieme delle parti è totalmente ordinato}{parti-totale}
    \begin{prerequisiti}
        \dip{Definizione}{def:confrontabili}
        \dip{Definizione}{def:inclusione}
        \dip{Proposizione}{prop:contrapposizione}
    \end{prerequisiti}
    Sia $s$ un insieme. $(\mcP(s), \subseteq)$ è totalmente ordinato se e solo se $s = \varnothing$ oppure $s$ è un singleton.
\end{Proposizione}
\begin{dimostrazione}
    \begin{caso}{$(\leftarrow)$}
        $\mcP(\varnothing) = \{\varnothing\}$ ha un solo elemento, che è confrontabile con sé stesso. $\mcP(\{x\}) = \{\varnothing, \{x\}\}$ e
        $\varnothing \subseteq \{x\}$.
    \end{caso}
    \begin{caso}{$(\rightarrow)$, per contrapposizione}
        Se $s$ non è vuoto né un singleton, contiene due elementi distinti $x \neq y$. Allora $\{x\} \not\subseteq \{y\}$ (perché
        $x \notin \{y\}$) e $\{y\} \not\subseteq \{x\}$: $\{x\}$ e $\{y\}$ non sono confrontabili.
    \end{caso}
\end{dimostrazione}
\end{concetto}

\begin{spiegazione}
    Basta che $s$ abbia due elementi perché nascano due «rami» paralleli: i diagrammi di Hasse qui sotto (Sezione~\ref{sec:hasse})
    mostrano le parti di $\varnothing$, di $\{a\}$ e di $\{a, b\}$; solo nel terzo compaiono elementi non confrontabili.
    \begin{center}
    \begin{tikzpicture}[hasse]
        \node[hv] (v0) at (0,0) {$\varnothing$};
        \node[font=\small] at (0,-0.75) {$\mcP(\varnothing)$};
        \begin{scope}[xshift=2.6cm]
            \node[hv] (v1) at (0,0) {$\varnothing$}; \node[hv] (a1) at (0,1.1) {$\{a\}$};
            \draw[hlato] (v1) -- (a1);
            \node[font=\small] at (0,-0.75) {$\mcP(\{a\})$};
        \end{scope}
        \begin{scope}[xshift=5.8cm]
            \node[hv] (v2) at (0,0) {$\varnothing$}; \node[hv evid] (a2) at (-0.9,1.1) {$\{a\}$}; \node[hv evid] (b2) at (0.9,1.1) {$\{b\}$};
            \node[hv] (ab2) at (0,2.2) {$\{a, b\}$};
            \draw[hlato] (v2) -- (a2) -- (ab2); \draw[hlato] (v2) -- (b2) -- (ab2);
            \node[font=\small] at (0,-0.75) {$\mcP(\{a, b\})$};
        \end{scope}
    \end{tikzpicture}
    \end{center}
\end{spiegazione}

% ──────────────────────────────────────────────────────────────
\section{Applicazioni crescenti e isomorfismi}\label{sec:crescenti}

Come per le strutture algebriche, anche tra insiemi ordinati ci interessano le funzioni che ne «rispettano la forma».

\dfn[crescente]{Applicazione crescente, isomorfismo di insiemi ordinati}{
    Siano $(s, \rho)$ e $(s', \rho')$ insiemi ordinati e $f \colon s \to s'$.
    \begin{itemize}
        \item $f$ si dice \textbf{crescente} se $\forall x, y \in s\,\big(x \rho y \rightarrow f(x) \rho' f(y)\big)$.
        \item $f$ si dice \textbf{isomorfismo} (di insiemi ordinati) se è biettiva e $\forall x, y \in s\,\big(x \rho y \leftrightarrow f(x) \rho' f(y)\big)$.
              In questo caso $(s, \rho)$ e $(s', \rho')$ si dicono \textbf{isomorfi}.
    \end{itemize}
}

\begin{esempio}{Isomorfismi}
    \begin{itemize}
        \item $n \in \mbN \mapsto 2n \in \{2n \mid n \in \mbN\}$ è un isomorfismo tra $(\mbN, \leq)$ e l'insieme dei pari con $\leq$:
              è biettiva e $m \leq n \leftrightarrow 2m \leq 2n$.
        \item $n \in \mbN \mapsto 2^n \in \{2^n \mid n \in \mbN\}$ è un isomorfismo tra $(\mbN, \leq)$ e le potenze di $2$ ordinate per
              divisibilità: $m \leq n \leftrightarrow 2^m \mid 2^n$.
    \end{itemize}
\end{esempio}

Ogni isomorfismo è crescente. Al contrario di quanto succede per le strutture algebriche, però, una funzione biettiva e
crescente non è detto che sia un isomorfismo.

\begin{concetto}
\begin{Proposizione}{Una biezione crescente che non è un isomorfismo}{biettiva-crescente}
    \begin{prerequisiti}
        \dip{Definizione}{def:crescente}
        \dip{Definizione}{def:funzioni-notevoli}
    \end{prerequisiti}
    Sia $s = \mbN \setminus \{0\}$. L'identità $\mathrm{id}_s$ è una funzione biettiva e crescente da $(s, \mid)$ a $(s, \leq)$, ma non è
    un isomorfismo.
\end{Proposizione}
\begin{dimostrazione}
    L'identità è biettiva. È crescente: se $m \mid n$ con $n \geq 1$, allora $n = km$ con $k \geq 1$, quindi $m \leq km = n$.
    Non è un isomorfismo: $2 \leq 3$, ma $2 \nmid 3$; cioè $\mathrm{id}(2) \leq \mathrm{id}(3)$ non implica $2 \mid 3$.
\end{dimostrazione}
\end{concetto}

\begin{spiegazione}
    Un isomorfismo deve portare le frecce in frecce \textbf{e} le «non frecce» in «non frecce». Una funzione crescente garantisce
    solo la prima metà: può \textbf{aggiungere} confronti che prima non c'erano. Ecco i diagrammi di Hasse di $\{1, 2, 3, 4, 6\}$
    con la divisibilità e con $\leq$: l'identità porta ogni segmento di sinistra in un percorso verso l'alto a destra, ma a destra
    ci sono confronti nuovi, come $2 \leq 3$ e $4 \leq 6$.
    \begin{center}
    \begin{tikzpicture}[hasse]
        \node[hv] (1) at (0,0) {$1$}; \node[hv] (2) at (-0.8,1.1) {$2$}; \node[hv] (3) at (0.8,1.1) {$3$};
        \node[hv] (4) at (-1.6,2.2) {$4$}; \node[hv] (6) at (0,2.2) {$6$};
        \draw[hlato] (1) -- (2) -- (4); \draw[hlato] (2) -- (6); \draw[hlato] (1) -- (3) -- (6);
        \node[font=\small] at (-0.4,-0.75) {$(\{1, 2, 3, 4, 6\}, \mid)$};
        \draw[-{Stealth}, line width=1pt, primary] (1.7,1.1) -- node[above, font=\small] {$\mathrm{id}$} (3.1,1.1);
        \begin{scope}[xshift=4.3cm, yshift=-0.2cm]
            \node[hv] (a) at (0,0) {$1$}; \node[hv] (b) at (0,0.6) {$2$}; \node[hv] (c) at (0,1.2) {$3$};
            \node[hv] (d) at (0,1.8) {$4$}; \node[hv] (e) at (0,2.4) {$6$};
            \draw[hlato] (a) -- (b); \draw[hlato evid] (b) -- (c); \draw[hlato] (c) -- (d); \draw[hlato evid] (d) -- (e);
            \node[font=\small] at (0,-0.55) {$(\{1, 2, 3, 4, 6\}, \leq)$};
        \end{scope}
    \end{tikzpicture}
    \end{center}
\end{spiegazione}

\begin{concetto}
\begin{Proposizione}{Inversa e composta di isomorfismi}{isomorfismi-ordinati}
    \begin{prerequisiti}
        \dip{Definizione}{def:crescente}
        \dip{Teorema}{teo:biettiva-inversa}
        \dip{Corollario}{cor:comp-biettive}
    \end{prerequisiti}
    L'inversa di un isomorfismo di insiemi ordinati è un isomorfismo, e la composta di due isomorfismi è un isomorfismo.
    Quindi «essere isomorfi» si comporta come una relazione di equivalenza tra insiemi ordinati.
\end{Proposizione}
\begin{dimostrazione}
    \begin{caso}{Inversa}
        Sia $f \colon (s, \rho) \to (s', \rho')$ un isomorfismo; $f^{-1}$ è biettiva. Per $u, v \in s'$, posto $x = f^{-1}(u)$, $y = f^{-1}(v)$:
        $f^{-1}(u) \rho f^{-1}(v) \leftrightarrow x \rho y \leftrightarrow f(x) \rho' f(y) \leftrightarrow u \rho' v$.
    \end{caso}
    \begin{caso}{Composta}
        Se $f \colon (s, \rho) \to (s', \rho')$ e $g \colon (s', \rho') \to (s'', \rho'')$ sono isomorfismi, $g \circ f$ è biettiva
        (Corollario~\ref{cor:comp-biettive}) e $x \rho y \leftrightarrow f(x) \rho' f(y) \leftrightarrow g(f(x)) \rho'' g(f(y))$.
    \end{caso}
    Infine l'identità di $(s, \rho)$ è un isomorfismo di $(s, \rho)$ in sé stesso.
\end{dimostrazione}
\end{concetto}

\begin{spiegazione}
    Un isomorfismo è un \textbf{cambio di nomi} che lascia intatto il disegno: l'inversa rimette i nomi vecchi, e due cambi di nome
    uno dopo l'altro sono ancora un cambio di nome. Per questo due insiemi ordinati isomorfi sono, per tutto ciò che riguarda
    l'ordine, «lo stesso insieme con etichette diverse».
\end{spiegazione}

% ──────────────────────────────────────────────────────────────
\section{Relazione di copertura e diagrammi di Hasse}\label{sec:hasse}

Il grafo con frecce di un ordine ha moltissime frecce: tutti i cappi e tutte le scorciatoie. Per disegnare un insieme ordinato
finito basta molto meno: le frecce tra elementi \textbf{immediatamente} uno sopra l'altro.

\dfn[copertura]{Relazione di copertura, diagramma di Hasse}{
    Sia $(s, \leq)$ un insieme ordinato e siano $x, y \in s$. Si dice che \textbf{$y$ copre $x$} se
    \[ x < y \ \wedge\ \neg\exists z \in s\,(x < z \wedge z < y), \]
    cioè $y$ sta sopra $x$ e non c'è niente in mezzo: $y$ è un \textbf{successore immediato} di $x$. La relazione binaria
    $\gamma$ su $s$ definita da $x \gamma y \leftrightarrow y \text{ copre } x$ si chiama \textbf{relazione di copertura} (o
    \textbf{diagramma di Hasse}) di $(s, \leq)$; è antiriflessiva.

    Il diagramma di Hasse di un insieme ordinato finito si \textbf{disegna} mettendo ogni elemento più in basso degli elementi che lo
    coprono e tracciando un segmento da $x$ a $y$ ogni volta che $y$ copre $x$.
}

\begin{concetto}
Per \textbf{leggere} un diagramma di Hasse: $x \leq y$ se e solo se $x = y$ oppure da $x$ si arriva a $y$ \textbf{salendo} lungo i
segmenti (per gli insiemi finiti lo dimostreremo nel Teorema~\ref{teo:hasse-isomorfismo}). Ad esempio:
\begin{center}
\begin{tikzpicture}[hasse]
    \node[hv] (v) at (0,0) {$\varnothing$}; \node[hv] (e) at (0,1.1) {$\{\varnothing\}$};
    \draw[hlato] (v) -- (e);
    \node[font=\small, align=center] at (0,-0.8) {$(\mcP(\mcP(\varnothing)), \subseteq)$};
    \begin{scope}[xshift=4.2cm]
        \node[hv] (v1) at (0,0) {$\varnothing$};
        \node[hv] (a) at (-1.5,1.1) {$\{a\}$}; \node[hv] (b) at (0,1.1) {$\{b\}$}; \node[hv] (c) at (1.5,1.1) {$\{c\}$};
        \node[hv] (ab) at (-1.5,2.2) {$\{a, b\}$}; \node[hv] (ac) at (0,2.2) {$\{a, c\}$}; \node[hv] (bc) at (1.5,2.2) {$\{b, c\}$};
        \node[hv] (abc) at (0,3.3) {$\{a, b, c\}$};
        \draw[hlato] (v1) -- (a); \draw[hlato] (v1) -- (b); \draw[hlato] (v1) -- (c);
        \draw[hlato] (a) -- (ab); \draw[hlato] (a) -- (ac); \draw[hlato] (b) -- (ab); \draw[hlato] (b) -- (bc);
        \draw[hlato] (c) -- (ac); \draw[hlato] (c) -- (bc);
        \draw[hlato] (ab) -- (abc); \draw[hlato] (ac) -- (abc); \draw[hlato] (bc) -- (abc);
        \node[font=\small] at (0,-0.8) {$(\mcP(\{a, b, c\}), \subseteq)$};
    \end{scope}
    \begin{scope}[xshift=9.4cm]
        \node[hv] (1) at (0,0) {$1$}; \node[hv] (2) at (-0.8,1.1) {$2$}; \node[hv] (3) at (0.8,1.1) {$3$};
        \node[hv] (4) at (-1.6,2.2) {$4$}; \node[hv] (6) at (0,2.2) {$6$}; \node[hv] (12) at (-0.8,3.3) {$12$};
        \draw[hlato] (1) -- (2) -- (4) -- (12); \draw[hlato] (1) -- (3) -- (6) -- (12); \draw[hlato] (2) -- (6);
        \node[font=\small] at (-0.4,-0.8) {divisori di $12$ con $\mid$};
    \end{scope}
    \begin{scope}[xshift=13.2cm]
        \node[hv] (t3) at (0,0) {$3$}; \node[hv] (t6) at (0,0.8) {$6$}; \node[hv] (t23) at (0,1.6) {$23$}; \node[hv] (t1000) at (0,2.4) {$1000$};
        \draw[hlato] (t3) -- (t6) -- (t23) -- (t1000);
        \node[font=\small] at (0,-0.8) {$(\{3, 6, 23, 1000\}, \leq)$};
    \end{scope}
\end{tikzpicture}
\end{center}
Nel secondo diagramma $\{a\} \subseteq \{a, b, c\}$ perché si sale da $\{a\}$ a $\{a, b\}$ e poi a $\{a, b, c\}$; $\{a\}$ e $\{b, c\}$ non sono
confrontabili perché nessuno dei due si raggiunge salendo dall'altro. Un insieme totalmente ordinato finito ha come diagramma
di Hasse una \textbf{catena} verticale, come l'ultimo.
\end{concetto}

\begin{esempio}{Minimi e minimali sul diagramma di Hasse}
    Sul diagramma di Hasse di un insieme ordinato finito:
    \begin{itemize}
        \item un elemento è \textbf{minimale} se non ha segmenti che scendono, \textbf{massimale} se non ha segmenti che salgono;
        \item c'è il \textbf{minimo} se c'è un elemento da cui si raggiungono tutti salendo, il \textbf{massimo} se c'è un elemento
              che si raggiunge salendo da tutti.
    \end{itemize}
    Nei divisori di $12$ il minimo è $1$ e il massimo è $12$. Togliendo $1$ e $12$ restano $\{2, 3, 4, 6\}$: i minimali sono $2$ e $3$,
    i massimali $4$ e $6$, e non ci sono né minimo né massimo.
\end{esempio}

% ──────────────────────────────────────────────────────────────
\section{Minoranti e maggioranti}\label{sec:minoranti}

\dfn[minoranti]{Minoranti e maggioranti}{
    Sia $(s, \rho)$ un insieme ordinato e sia $t \subseteq s$. Un elemento $m \in s$ si dice
    \begin{itemize}
        \item \textbf{minorante} di $t$ se $\forall z \in t\,(m \rho z)$;
        \item \textbf{maggiorante} di $t$ se $\forall z \in t\,(z \rho m)$.
    \end{itemize}
    Indichiamo con $\mathrm{Minor}_{(s, \rho)}(t)$ l'insieme dei minoranti di $t$ e con $\mathrm{Maggior}_{(s, \rho)}(t)$ quello dei
    maggioranti (omettendo il pedice quando non c'è ambiguità).
}

Attenzione: i minoranti e i maggioranti sono elementi di $s$, non necessariamente di $t$. Il minimo di $t$, se esiste, è un
minorante di $t$ che sta in $t$.

\begin{esempio}{Minoranti e maggioranti}
    \begin{itemize}
        \item In $(\mbN, \leq)$ i minoranti di una parte non vuota $t$ sono i numeri minori o uguali a tutti gli elementi di $t$.
        \item In $(\mbN, \mid)$ i minoranti di $t$ sono i \textbf{divisori comuni} di tutti gli elementi di $t$, i maggioranti i
              \textbf{multipli comuni}. Ad esempio $\mathrm{Maggior}(\{2\})$ è l'insieme dei multipli di $2$ (compreso $0$) e
              $\mathrm{Minor}(\{12, 18\}) = \{1, 2, 3, 6\}$.
        \item In $(\mcP(s), \subseteq)$, se $t$ è una parte non vuota di $\mcP(s)$, i minoranti di $t$ sono gli insiemi contenuti in tutti
              gli elementi di $t$, cioè le parti di $\bigcap t$; i maggioranti sono gli insiemi che contengono tutti gli elementi di $t$,
              cioè le parti di $s$ che contengono $\bigcup t$.
        \item Per la parte vuota: ogni elemento di $s$ è sia minorante sia maggiorante di $\varnothing$ (la condizione $\forall z \in \varnothing$
              è vera a vuoto), quindi $\mathrm{Minor}(\varnothing) = \mathrm{Maggior}(\varnothing) = s$.
    \end{itemize}
    Sul diagramma dei divisori di $12$, con $t = \{4, 6\}$ (in arancione): i minoranti (in verde) sono gli elementi da cui si
    raggiungono \textbf{sia} $4$ \textbf{sia} $6$ salendo, i maggioranti quelli che si raggiungono salendo da entrambi.
    \begin{center}
    \begin{tikzpicture}[hasse]
        \node[hv evid2] (1) at (0,0) {$1$}; \node[hv evid2] (2) at (-0.8,1.1) {$2$}; \node[hv] (3) at (0.8,1.1) {$3$};
        \node[hv evid] (4) at (-1.6,2.2) {$4$}; \node[hv evid] (6) at (0,2.2) {$6$}; \node[hv] (12) at (-0.8,3.3) {$12$};
        \draw[hlato] (1) -- (2) -- (4) -- (12); \draw[hlato] (1) -- (3) -- (6) -- (12); \draw[hlato] (2) -- (6);
        \node[font=\small, align=left, anchor=west] at (1.6,1.6)
            {$\mathrm{Minor}(\{4, 6\}) = \{1, 2\}$\\[1mm] $\mathrm{Maggior}(\{4, 6\}) = \{12\}$\\[1mm]
             ($3$ non è un minorante: $3 \nmid 4$)};
    \end{tikzpicture}
    \end{center}
\end{esempio}


  % Parte 5: Numeri naturali e calcolo combinatorio

  \part{Numeri naturali e calcolo combinatorio}

  % Equipotenza, infinito, naturali, induzione di prima e seconda forma, potenze e multipli
  \chapter{Numeri naturali e principio di induzione}\label{cap:induzione}

Finora abbiamo usato i numeri naturali negli esempi come oggetti già noti. In questo capitolo vediamo come la teoria degli
insiemi li fa nascere: da un insieme infinito si ricava un insieme ordinato con le proprietà di $\mbN$, e da queste proprietà
segue lo strumento più importante per dimostrare enunciati «per ogni $n$»: il \textbf{principio di induzione}. Lo useremo subito
per dimostrare le regole delle potenze nei semigruppi, rimaste in sospeso nella Parte~III.

Le operazioni $+$ e $\cdot$ su $\mbN$ e su $\mbZ$, con le loro proprietà, continuiamo a darle per note.

% ──────────────────────────────────────────────────────────────
\section{Insiemi infiniti}\label{sec:infiniti}

\dfn[equipotenti]{Insiemi equipotenti}{
    Due insiemi $a$ e $b$ si dicono \textbf{equipotenti} se esiste una funzione biettiva $f \colon a \to b$.
}

\begin{concetto}
\begin{Proposizione}{Proprietà dell'equipotenza}{equipotenza}
    \begin{prerequisiti}
        \dip{Definizione}{def:equipotenti}
        \dip{Teorema}{teo:biettiva-inversa}
        \dip{Corollario}{cor:comp-biettive}
    \end{prerequisiti}
    Siano $a$, $b$, $c$ insiemi. Allora $a$ è equipotente ad $a$; se $a$ è equipotente a $b$, allora $b$ è equipotente ad $a$; se $a$ è
    equipotente a $b$ e $b$ a $c$, allora $a$ è equipotente a $c$.
\end{Proposizione}
\begin{dimostrazione}
    L'identità $\mathrm{id}_a$ è biettiva. Se $f \colon a \to b$ è biettiva, ha un'inversa $f^{-1} \colon b \to a$, che è biettiva
    (Teorema~\ref{teo:biettiva-inversa}: la sua inversa è $f$). Se $f \colon a \to b$ e $g \colon b \to c$ sono biettive, $g \circ f$
    è biettiva (Corollario~\ref{cor:comp-biettive}).
\end{dimostrazione}
\end{concetto}

\begin{spiegazione}
    Due insiemi sono equipotenti quando i loro elementi si possono \textbf{accoppiare uno a uno} senza avanzi, come i posti e gli
    spettatori in un teatro tutto esaurito. Questo accoppiamento si può rovesciare (inversa) e concatenare (composta): per questo
    l'equipotenza si comporta come una relazione di equivalenza. Per gli insiemi finiti vorrà dire «avere lo stesso numero di
    elementi»; per quelli infiniti riserva qualche sorpresa.
\end{spiegazione}

\dfn[infinito]{Insieme infinito}{
    Un insieme si dice \textbf{infinito} se è equipotente a una sua parte propria.
}

La definizione sembra strana, ma coglie proprio ciò che distingue gli insiemi infiniti: un insieme finito non può essere
messo in corrispondenza biunivoca con una sua parte più piccola (lo vedremo nel Teorema~\ref{teo:in-non-infinito}), mentre
per un insieme infinito questo è possibile (Esempio~\ref{es:hilbert}). Che ne esista almeno uno non si può dimostrare dagli
assiomi visti finora: è un nuovo assioma.

\ass[cantor]{Dell'infinito (o di Cantor)}{
    Esiste un insieme infinito.
}

% ──────────────────────────────────────────────────────────────
\section{I numeri naturali}\label{sec:naturali}

\dfn[naturalmente-ordinato]{Insieme naturalmente ordinato}{
    Sia $(s, \leq)$ un insieme ordinato. Una parte $t \subseteq s$ si dice \textbf{superiormente limitata} se ha almeno un maggiorante.
    $(s, \leq)$ si dice \textbf{naturalmente ordinato} se è ben ordinato e ogni sua parte non vuota superiormente limitata ha massimo.
}

I due teoremi seguenti dicono che, a meno di isomorfismi, esiste uno e un solo insieme naturalmente ordinato senza maggioranti.

\begin{concetto}
\begin{Teorema}{Infinito e insiemi naturalmente ordinati}{infinito-naturale}
    \begin{prerequisiti}
        \dip{Definizione}{def:infinito}
        \dip{Definizione}{def:naturalmente-ordinato}
    \end{prerequisiti}
    \senzadim
    Esiste un insieme infinito se e solo se esiste un insieme naturalmente ordinato non superiormente limitato.
\end{Teorema}
\end{concetto}

\begin{concetto}
\begin{Teorema}{Unicità dei naturali}{unicita-naturali}
    \begin{prerequisiti}
        \dip{Definizione}{def:naturalmente-ordinato}
        \dip{Definizione}{def:crescente}
    \end{prerequisiti}
    \senzadim
    Tutti gli insiemi naturalmente ordinati non superiormente limitati sono isomorfi (come insiemi ordinati).
\end{Teorema}
\end{concetto}

\dfn[naturali]{I numeri naturali}{
    Per l'assioma di Cantor e il Teorema~\ref{teo:infinito-naturale} esiste un insieme naturalmente ordinato non superiormente
    limitato; ne scegliamo uno, lo indichiamo con $(\mbN, \leq)$ e lo chiamiamo \textbf{insieme dei numeri naturali}
    (per il Teorema~\ref{teo:unicita-naturali} la scelta è indifferente). Inoltre:
    \begin{itemize}
        \item $0 := \min(\mbN)$, $1 := \min(\mbN \setminus \{0\})$, e così via;
        \item per $n \in \mbN$, il \textbf{successore} di $n$ è $n + 1 := \min\{k \in \mbN \mid n < k\}$;
        \item per $n \in \mbN \setminus \{0\}$, il \textbf{predecessore} di $n$ è $n - 1 := \max\{k \in \mbN \mid k < n\}$;
        \item per $m \in \mbN$, poniamo $\mbN_m := \{n \in \mbN \mid m \leq n\}$ (i naturali da $m$ in poi).
    \end{itemize}
}

\begin{concetto}
\begin{Proposizione}{Successore e predecessore}{successore}
    \begin{prerequisiti}
        \dip{Definizione}{def:naturali}
        \dip{Definizione}{def:naturalmente-ordinato}
        \dip{Teorema}{teo:ben-ordinato-totale}
    \end{prerequisiti}
    Il successore e il predecessore esistono, e per ogni $n \in \mbN$:
    \begin{enumerate}
        \item $n < n + 1$ e non esiste $k$ con $n < k < n + 1$ (cioè $n + 1$ copre $n$);
        \item se $n \neq 0$: $n - 1 < n$, non esiste $k$ con $n - 1 < k < n$, e $(n - 1) + 1 = n$.
    \end{enumerate}
\end{Proposizione}
\begin{dimostrazione}
    Ricordiamo che $\mbN$ è totalmente ordinato (Teorema~\ref{teo:ben-ordinato-totale}).
    \begin{caso}{Esistenza}
        $\{k \mid n < k\}$ non è vuoto, perché altrimenti $n$ sarebbe un maggiorante di $\mbN$; essendo $\mbN$ ben ordinato, ha minimo.
        Se $n \neq 0$, $\{k \mid k < n\}$ contiene $0$ e ha il maggiorante $n$: per l'ordine naturale ha massimo.
    \end{caso}
    \begin{caso}{Punto 1}
        $n + 1$ sta in $\{k \mid n < k\}$, quindi $n < n + 1$; se fosse $n < k < n + 1$, $k$ sarebbe in quell'insieme e più piccolo del minimo.
    \end{caso}
    \begin{caso}{Punto 2}
        Allo stesso modo $n - 1 < n$ e nessun $k$ sta strettamente tra i due. Infine $n$ sta in $\{k \mid n - 1 < k\}$; se $j$ sta in
        quest'insieme e fosse $j < n$, allora $j \leq n - 1$ (massimo), contro $n - 1 < j$. Quindi $j \geq n$ per ogni tale $j$, e il
        minimo, cioè $(n - 1) + 1$, è $n$.
    \end{caso}
\end{dimostrazione}
\end{concetto}

\begin{spiegazione}
    In $\mbN$ ogni numero ha un \textbf{numero immediatamente dopo} e (tranne $0$) un \textbf{numero immediatamente prima}: è
    proprio ciò che manca, ad esempio, in $(\mbQ, \leq)$, dove tra due numeri ce n'è sempre un altro. Il diagramma di Hasse di
    $(\mbN, \leq)$ è una catena infinita verso l'alto, in cui ogni numero è coperto dal suo successore:
    \begin{center}
    \begin{tikzpicture}[hasse]
        \foreach \n in {0,...,5} \node[hv] (n\n) at (\n*1.3,0) {$\n$};
        \foreach \n [evaluate=\n as \m using int(\n+1)] in {0,...,4} \draw[hlato, -{Stealth[length=2mm]}] (n\n) -- (n\m);
        \node at (7.2,0) {$\cdots$};
        \draw[hlato, -{Stealth[length=2mm]}] (n5) -- (6.8,0);
        \node[font=\small] at (3.25,-0.65) {la freccia va da $n$ al suo successore $n + 1$ (disegniamo la catena in orizzontale)};
    \end{tikzpicture}
    \end{center}
\end{spiegazione}

\begin{esempio}{L'albergo di Hilbert}\label{es:hilbert}
    La funzione $n \in \mbN \mapsto n + 1 \in \mbN \setminus \{0\}$ è biettiva: ogni $m \neq 0$ è immagine del solo $m - 1$. Quindi
    $\mbN$ è equipotente alla sua parte propria $\mbN \setminus \{0\}$, ed è infinito. È il famoso «albergo di Hilbert»: un albergo
    con infinite stanze tutte occupate può accogliere un nuovo ospite, se ogni cliente si sposta dalla stanza $n$ alla stanza $n + 1$
    e lascia libera la stanza $0$.
\end{esempio}

% ──────────────────────────────────────────────────────────────
\section{Il principio di induzione}\label{sec:induzione}

\begin{concetto}
\begin{Teorema}{Principio di induzione di prima forma}{induzione-prima}
    \begin{prerequisiti}
        \dip{Definizione}{def:naturali}
        \dip{Proposizione}{prop:successore}
        \dip{Definizione}{def:ben-ordinato}
    \end{prerequisiti}
    Sia $x$ una parte non vuota di $\mbN$ tale che $\forall n \in \mbN\,(n \in x \rightarrow n + 1 \in x)$. Allora $x = \mbN_{\min(x)}$. In formule:
    \[ \forall x \in \mcP(\mbN) \setminus \{\varnothing\}\ \Big(\forall n \in \mbN\,\big(n \in x \rightarrow n + 1 \in x\big) \rightarrow x = \mbN_{\min(x)}\Big). \]
\end{Teorema}
\begin{dimostrazione}
    Sia $m = \min(x)$ (esiste perché $\mbN$ è ben ordinato). Ogni elemento di $x$ è $\geq m$, quindi $x \subseteq \mbN_m$.
    Per assurdo sia $x \neq \mbN_m$: allora $y = \mbN_m \setminus x$ non è vuoto e ha un minimo $n$.
    \begin{passi}
        \item $n \neq m$, perché $m \in x$ e $n \notin x$; quindi $m < n$, e in particolare $n \neq 0$.
        \item $n - 1 \geq m$ (perché $m < n$ e $n - 1$ è il massimo dei numeri minori di $n$), cioè $n - 1 \in \mbN_m$.
        \item $n - 1 < n = \min(y)$, quindi $n - 1 \notin y$: deve essere $n - 1 \in x$.
        \item Per ipotesi allora $(n - 1) + 1 \in x$, cioè $n \in x$ (Proposizione~\ref{prop:successore}): ma $n \in y$, assurdo.
    \end{passi}
\end{dimostrazione}
\end{concetto}

\begin{spiegazione}
    È l'effetto domino: se cade la prima tessera ($m \in x$) e ogni tessera che cade fa cadere la successiva ($n \in x \Rightarrow n + 1 \in x$),
    cadono tutte le tessere da $m$ in poi.
    \begin{center}
    \begin{tikzpicture}
        \fill[black!12] (-0.5,-0.06) rectangle (8.9,0);
        \draw[fill=orange!35, draw=orange!80!black, rotate around={-58:(0.12,0)}] (-0.12,0) rectangle (0.12,1.3);
        \draw[fill=orange!35, draw=orange!80!black, rotate around={-32:(1.32,0)}] (1.08,0) rectangle (1.32,1.3);
        \foreach \i in {2,...,7} \draw[fill=primary!15, draw=primary!70!black] (\i*1.2-0.12,0) rectangle (\i*1.2+0.12,1.3);
        \foreach \i/\t in {0/m, 1/m+1, 2/m+2, 3/m+3, 4/m+4, 5/m+5, 6/m+6, 7/m+7} \node[font=\scriptsize] at (\i*1.2,-0.32) {$\t$};
        \draw[-{Stealth}, orange!80!black, line width=0.8pt] (0.5,1.45) to[bend left=30] (1.25,1.5);
        \draw[-{Stealth}, orange!80!black, line width=0.8pt] (1.75,1.5) to[bend left=30] (2.35,1.5);
        \node[font=\small, align=left, anchor=west] at (9.2,0.65) {base: $m \in x$\\ passo: $n \in x \Rightarrow n + 1 \in x$};
    \end{tikzpicture}
    \end{center}
    La dimostrazione formalizza un ragionamento per assurdo: se qualche tessera restasse in piedi, ce ne sarebbe una \textbf{prima}
    rimasta in piedi ($n = \min y$, esiste per il buon ordine); ma la tessera subito prima di lei ($n - 1$) è caduta, e cadendo l'avrebbe
    fatta cadere. È qui che servono le proprietà di $\mbN$: il minimo di ogni parte non vuota e il predecessore.
\end{spiegazione}

Nella pratica il principio si usa così.

\begin{concetto}
\begin{Corollario}{Dimostrazioni per induzione}{induzione-uso}
    \begin{prerequisiti}
        \dip{Teorema}{teo:induzione-prima}
    \end{prerequisiti}
    Sia $P(n)$ una proprietà dei numeri naturali e sia $m \in \mbN$. Supponiamo che
    \begin{itemize}
        \item[] \textbf{(passo base)} $P(m)$ sia vera;
        \item[] \textbf{(passo induttivo)} per ogni $n \geq m$, se $P(n)$ è vera (\textbf{ipotesi induttiva}), allora $P(n + 1)$ è vera.
    \end{itemize}
    Allora $P(n)$ è vera per ogni $n \geq m$.
\end{Corollario}
\begin{dimostrazione}
    Sia $x = \{n \in \mbN_m \mid P(n)\}$. Per il passo base $m \in x$, quindi $x \neq \varnothing$, e $\min(x) = m$ (tutti gli elementi
    di $x$ sono $\geq m$). Per il passo induttivo, $n \in x \rightarrow n + 1 \in x$. Per il Teorema~\ref{teo:induzione-prima}
    $x = \mbN_m$: ogni $n \geq m$ soddisfa $P(n)$.
\end{dimostrazione}
\end{concetto}

\begin{spiegazione}
    \begin{passi}
        \item Si dimostra il \textbf{passo base}: di solito un conto diretto per il primo valore.
        \item Si fissa un $n$ qualunque, si \textbf{suppone} vera $P(n)$ e si dimostra $P(n + 1)$, usando $P(n)$.
        \item Si conclude che $P$ vale da $m$ in poi.
    \end{passi}
    Attenzione: nel passo induttivo non si assume che $P$ valga per tutti gli $n$ (sarebbe la tesi!), ma solo per un $n$ fissato.
\end{spiegazione}

\begin{esempio}{La somma dei primi $n$ numeri}
    Dimostriamo che per ogni $n \in \mbN$: \quad $0 + 1 + \dots + (n - 1) = \dfrac{n(n - 1)}{2}$.
    \begin{passi}
        \item \textbf{Base} ($n = 0$): a sinistra la somma è vuota, cioè $0$; a destra $\frac{0 \cdot (-1)}{2} = 0$.
        \item \textbf{Passo}: supponiamo $0 + 1 + \dots + (n - 1) = \frac{n(n - 1)}{2}$. Aggiungendo $n$ a entrambi i membri:
              \[ 0 + 1 + \dots + (n - 1) + n = \frac{n(n - 1)}{2} + n = \frac{n(n - 1) + 2n}{2} = \frac{(n + 1)n}{2}, \]
              che è la formula per $n + 1$.
    \end{passi}
    Il disegno spiega anche \textbf{da dove viene} la formula: due «scale» da $0 + 1 + 2 + 3$ quadretti formano un rettangolo
    $4 \times 3$, quindi una scala ha $\frac{4 \cdot 3}{2} = 6$ quadretti.
    \begin{center}
    \begin{tikzpicture}[scale=0.45]
        \foreach \c in {1,2,3} \foreach \r in {1,...,\c} \draw[fill=primary!25, draw=primary!70!black] (\c,\r-1) rectangle (\c+1,\r);
        \foreach \c in {0,1,2} \foreach \r in {\c,...,2} \draw[fill=orange!30, draw=orange!80!black] (\c,\r) rectangle (\c+1,\r+1);
        \draw[decorate, decoration={brace, amplitude=4pt, mirror}] (0,-0.2) -- node[below=3pt, font=\small] {$4$} (4,-0.2);
        \draw[decorate, decoration={brace, amplitude=4pt}] (-0.2,0) -- node[left=3pt, font=\small] {$3$} (-0.2,3);
    \end{tikzpicture}
    \end{center}
\end{esempio}

A volte, per dimostrare $P(n)$, non basta sapere che vale $P(n - 1)$: serve che valga per \textbf{tutti} i numeri precedenti.

\begin{concetto}
\begin{Teorema}{Principio di induzione di seconda forma (induzione completa)}{induzione-seconda}
    \begin{prerequisiti}
        \dip{Definizione}{def:naturali}
        \dip{Definizione}{def:ben-ordinato}
    \end{prerequisiti}
    Sia $x$ una parte non vuota di $\mbN$ tale che, per ogni $n \in \mbN$, se tutti i $k$ con $\min(x) \leq k < n$ stanno in $x$,
    allora $n \in x$. Allora $x = \mbN_{\min(x)}$. In formule:
    \[ \forall x \in \mcP(\mbN) \setminus \{\varnothing\}\ \Big(\forall n \in \mbN\,\big(\forall k \in \mbN\,(\min(x) \leq k < n \rightarrow k \in x) \rightarrow n \in x\big)
       \rightarrow x = \mbN_{\min(x)}\Big). \]
\end{Teorema}
\begin{dimostrazione}
    Sia $m = \min(x)$; come prima $x \subseteq \mbN_m$. Per assurdo sia $y = \mbN_m \setminus x \neq \varnothing$ e sia $n = \min(y)$.
    Se $k$ è tale che $m \leq k < n$, allora $k \in \mbN_m$ e $k \notin y$ (perché $k < \min(y)$), quindi $k \in x$. Dunque tutti i $k$
    con $m \leq k < n$ stanno in $x$, e per ipotesi $n \in x$: ma $n \in y$, assurdo.
\end{dimostrazione}
\end{concetto}

\begin{spiegazione}
    È ancora l'effetto domino, ma con tessere più pesanti: per far cadere la tessera $n$ servono \textbf{tutte} quelle prima di lei.
    Il ragionamento per assurdo è lo stesso: la prima tessera rimasta in piedi avrebbe tutte le precedenti cadute, quindi cadrebbe
    anche lei. Nella pratica: si suppone $P(k)$ vera per \textbf{ogni} $k$ con $m \leq k < n$ e si dimostra $P(n)$. Quando il ragionamento
    usa uno dei $k$ precedenti ma non funziona per i primi valori, questi si controllano a mano come passi base.
\end{spiegazione}

\begin{esempio}{Francobolli da $4$ e da $5$}
    Dimostriamo che ogni $n \geq 12$ si scrive $n = 4a + 5b$ con $a, b \in \mbN$ (ogni affrancatura da almeno $12$ centesimi si fa
    con francobolli da $4$ e da $5$).
    \begin{passi}
        \item \textbf{Basi}: $12 = 4 \cdot 3$, \ $13 = 4 \cdot 2 + 5$, \ $14 = 4 + 5 \cdot 2$, \ $15 = 5 \cdot 3$.
        \item \textbf{Passo} ($n > 15$): supponiamo l'enunciato vero per ogni $k$ con $12 \leq k < n$. Allora $12 \leq n - 4 < n$, quindi
              $n - 4 = 4a + 5b$ e $n = 4(a + 1) + 5b$.
    \end{passi}
    \begin{center}
    \begin{tikzpicture}[x=0.8cm]
        \foreach \n in {12,...,23} {
            \pgfmathtruncatemacro{\r}{mod(\n,4)}
            \ifcase\r \def\col{classeA}\or \def\col{classeB}\or \def\col{classeC}\or \def\col{classeD}\fi
            \node[circle, fill=\col, draw=black!30, inner sep=0pt, minimum size=6.5mm, font=\small] (k\n) at (\n,0) {$\n$};
        }
        \foreach \n in {12,...,15} \node[font=\scriptsize] at (\n,-0.6) {base};
        \foreach \n in {12,13,14,15,16,17,18,19} {
            \pgfmathtruncatemacro{\m}{\n+4}
            \draw[-{Stealth}, line width=0.6pt, black!60] (k\n.north) to[bend left=40] (k\m.north);
        }
    \end{tikzpicture}
    \end{center}
    Ogni freccia aggiunge un francobollo da $4$: partendo dalle quattro basi si raggiungono tutti i numeri da $12$ in poi. Qui
    serve la seconda forma perché da $n$ si torna a $n - 4$, non a $n - 1$.
\end{esempio}

% ──────────────────────────────────────────────────────────────
\section{Potenze e multipli}\label{sec:potenze}

Il principio di induzione permette anche di \textbf{definire} oggetti «passo per passo»: si dice cosa succede in un primo
valore e come si passa da $n$ a $n + 1$. Si parla di \textbf{definizione ricorsiva}. Così si definiscono le potenze.

\dfn[potenze]{Potenze e multipli}{
    Sia $(s, *)$ un semigruppo e sia $x \in s$. Le \textbf{potenze} di $x$ sono definite ricorsivamente da
    \[ x^1 := x, \qquad x^{n + 1} := x^n * x \quad (n \geq 1). \]
    Se $(s, *, 1_s)$ è un monoide si pone anche $x^0 := 1_s$. Se $(s, *)$ è un gruppo e $n \geq 1$ si pone $x^{-n} := (x^{-1})^n$.
    In notazione additiva si parla di \textbf{multipli} e si scrive $nx$ invece di $x^n$: $1x = x$, $(n + 1)x = nx + x$, $0x = 0$,
    $(-n)x = n(-x)$.
}

\begin{concetto}
\begin{Teorema}{Regole delle potenze}{regole-potenze}
    \begin{prerequisiti}
        \dip{Definizione}{def:potenze}
        \dip{Corollario}{cor:induzione-uso}
        \dip{Definizione}{def:semigruppo}
    \end{prerequisiti}
    Sia $(s, *)$ un semigruppo e sia $x \in s$. Per ogni $m, n \geq 1$:
    \[ \text{(a)}\ \ x^m * x^n = x^{m + n} \qquad\qquad \text{(b)}\ \ (x^m)^n = x^{mn}. \]
    Se $s$ è un monoide, (a) e (b) valgono per ogni $m, n \geq 0$.
\end{Teorema}
\begin{dimostrazione}
    Fissiamo $m \geq 1$ e procediamo per induzione su $n$.
    \begin{caso}{(a) $x^m * x^n = x^{m + n}$}
        \textbf{Base} $n = 1$: $x^m * x^1 = x^m * x = x^{m + 1}$ per definizione. \textbf{Passo}: se $x^m * x^n = x^{m + n}$, allora, usando
        l'associatività,
        \[ x^m * x^{n + 1} = x^m * (x^n * x) = (x^m * x^n) * x = x^{m + n} * x = x^{m + n + 1}. \]
    \end{caso}
    \begin{caso}{(b) $(x^m)^n = x^{mn}$}
        \textbf{Base} $n = 1$: $(x^m)^1 = x^m$. \textbf{Passo}: se $(x^m)^n = x^{mn}$, allora per definizione e per (a)
        \[ (x^m)^{n + 1} = (x^m)^n * x^m = x^{mn} * x^m = x^{mn + m} = x^{m(n + 1)}. \]
    \end{caso}
    \begin{caso}{Esponenti nulli in un monoide}
        $x^m * x^0 = x^m * 1_s = x^m = x^{m + 0}$, \ $x^0 * x^n = x^n$, \ $(x^m)^0 = 1_s = x^0 = x^{m \cdot 0}$, \ $(x^0)^n = 1_s^n = 1_s = x^0$
        (l'ultima perché $1_s^{n + 1} = 1_s^n * 1_s = 1_s^n$, per induzione).
    \end{caso}
\end{dimostrazione}
\end{concetto}

\begin{spiegazione}
    Nel conto $x^2 * x^3 = (x * x) * (x * x * x) = x * x * x * x * x = x^5$ «togliamo le parentesi» grazie all'associatività e contiamo
    le $x$. Ma «e così via» non è una dimostrazione: l'induzione è il modo rigoroso di dirlo. Seguiamo il passo induttivo per $m = 2$:
    \begin{passi}
        \item $x^2 * x^1 = x^2 * x = x^3$: è la definizione di $x^3$ (base).
        \item $x^2 * x^2 = x^2 * (x^1 * x) = (x^2 * x^1) * x = x^3 * x = x^4$: si usa il passo precedente.
        \item $x^2 * x^3 = x^2 * (x^2 * x) = (x^2 * x^2) * x = x^4 * x = x^5$: e così via, ogni volta con il risultato del passo prima.
    \end{passi}
\end{spiegazione}

\begin{concetto}
\begin{Corollario}{Le potenze di uno stesso elemento commutano}{potenze-commutano}
    \begin{prerequisiti}
        \dip{Teorema}{teo:regole-potenze}
    \end{prerequisiti}
    In un semigruppo $(s, *)$, per ogni $x \in s$ e $m, n \geq 1$ (anche $\geq 0$ in un monoide): $x^m * x^n = x^n * x^m$.
\end{Corollario}
\begin{dimostrazione}
    $x^m * x^n = x^{m + n} = x^{n + m} = x^n * x^m$, per il Teorema~\ref{teo:regole-potenze} e la commutatività della somma in $\mbN$.
\end{dimostrazione}
\end{concetto}

\begin{spiegazione}
    Anche se l'operazione non è commutativa, $x^2 * x^3$ e $x^3 * x^2$ sono entrambi «cinque $x$ in fila»: le potenze di uno
    \textbf{stesso} elemento non si accorgono della mancanza di commutatività. Con elementi diversi invece $x * y$ e $y * x$ possono
    essere diversi, come nel monoide delle parole: $ab \neq ba$.
\end{spiegazione}

\begin{concetto}
\begin{Proposizione}{Potenze dell'inverso}{potenze-inverso}
    \begin{prerequisiti}
        \dip{Definizione}{def:potenze}
        \dip{Proposizione}{prop:unita-gruppo}
        \dip{Corollario}{cor:potenze-commutano}
        \dip{Corollario}{cor:induzione-uso}
    \end{prerequisiti}
    Sia $(g, *)$ un gruppo e $x \in g$. Per ogni $n \geq 1$: $(x^{-1})^n = (x^n)^{-1}$, cioè $x^{-n}$ è l'inverso di $x^n$.
\end{Proposizione}
\begin{dimostrazione}
    Per induzione su $n$. \textbf{Base}: $(x^{-1})^1 = x^{-1} = (x^1)^{-1}$. \textbf{Passo}: se $(x^{-1})^n = (x^n)^{-1}$, allora
    \[ (x^{-1})^{n + 1} = (x^{-1})^n * x^{-1} = (x^n)^{-1} * x^{-1} = (x * x^n)^{-1} = (x^{n + 1})^{-1}, \]
    dove abbiamo usato $(a * b)^{-1} = b^{-1} * a^{-1}$ (Proposizione~\ref{prop:unita-gruppo}) e $x * x^n = x^1 * x^n = x^{n + 1}$.
\end{dimostrazione}
\end{concetto}

\begin{spiegazione}
    Per «disfare» $x^3 = x * x * x$ bisogna togliere le tre $x$ una alla volta, partendo dall'ultima: $x^{-1} * x^{-1} * x^{-1}$.
    In $(\mbZ, +)$ questo dice che l'opposto di $3x$ è $3(-x)$; in $(\mbQ \setminus \{0\}, \cdot)$ che l'inverso di $2^3 = 8$ è
    $(\frac{1}{2})^3 = \frac{1}{8}$.
\end{spiegazione}

\info{Potenze con esponente intero}{
    In un gruppo le regole $x^m * x^n = x^{m + n}$ e $(x^m)^n = x^{mn}$ valgono per \textbf{tutti} gli $m, n \in \mbZ$: si dimostra
    distinguendo i casi secondo i segni di $m$ e $n$ e usando la Proposizione~\ref{prop:potenze-inverso}. Ad esempio
    $x^5 * x^{-2} = x * x * x * x * x * x^{-1} * x^{-1} = x^3$.
}

Le potenze descrivono le sottostrutture generate da un solo elemento.

\begin{concetto}
\begin{Proposizione}{Sottostrutture cicliche}{cicliche}
    \begin{prerequisiti}
        \dip{Teorema}{teo:sottomonoidi-generati}
        \dip{Teorema}{teo:sottogruppi-generati}
        \dip{Definizione}{def:potenze}
        \dip{Teorema}{teo:regole-potenze}
    \end{prerequisiti}
    Sia $x$ un elemento di un monoide $(m, *, 1_m)$. Il sottomonoide generato da $x$ è $\langle x \rangle = \{x^n \mid n \in \mbN\}$.
    Se $m$ è un gruppo, il sottogruppo generato da $x$ è $\langle x \rangle = \{x^n \mid n \in \mbZ\}$.
\end{Proposizione}
\begin{dimostrazione}
    Per il Teorema~\ref{teo:sottomonoidi-generati} con $t = \{x\}$, $\langle x \rangle$ è formato da $1_m = x^0$ e dai prodotti
    $x * x * \dots * x$ di $n \geq 1$ fattori, che sono le potenze $x^n$. Nel caso di un gruppo, per il
    Teorema~\ref{teo:sottogruppi-generati} gli elementi di $\langle x \rangle$ sono i prodotti di fattori $x$ e $x^{-1}$; per le regole delle
    potenze con esponente intero ognuno di essi è una potenza $x^k$, con $k$ uguale al numero di fattori $x$ meno il numero di fattori
    $x^{-1}$; e ogni $x^k$ è un prodotto di questo tipo.
\end{dimostrazione}
\end{concetto}

\begin{spiegazione}
    Partendo da un solo elemento e potendo solo moltiplicarlo per sé stesso (e, in un gruppo, per il suo inverso) si ottengono
    solo le sue potenze. Ad esempio in $(\mbZ, +)$, in notazione additiva, $\langle 1 \rangle = \{n \cdot 1 \mid n \in \mbZ\} = \mbZ$ (il
    gruppo è \textbf{ciclico}) e $\langle m \rangle = \{nm \mid n \in \mbZ\} = m\mbZ$, l'insieme dei multipli di $m$; in $(\mbN, +, 0)$
    il sottomonoide generato da $2$ è $\{0, 2, 4, \dots\} = 2\mbN$.
\end{spiegazione}

Negli anelli i multipli permettono di confrontare la moltiplicazione per un intero con la moltiplicazione per un elemento
dell'anello (Definizione~\ref{def:notazioni-anello}).

\begin{concetto}
\begin{Proposizione}{Moltiplicazione per un multiplo dell'unità}{multipli-unita}
    \begin{prerequisiti}
        \dip{Definizione}{def:notazioni-anello}
        \dip{Lemma}{lem:regole-anello}
        \dip{Corollario}{cor:induzione-uso}
    \end{prerequisiti}
    Sia $(a, +, \cdot)$ un anello unitario. Per ogni $x \in a$ e $n \in \mbZ$: \quad $(n 1_a) \cdot x = nx = x \cdot (n 1_a)$.
\end{Proposizione}
\begin{dimostrazione}
    \begin{caso}{Passo 1: $(ky)x = k(yx)$ e $x(ky) = k(xy)$ per ogni $k \geq 1$ e $y \in a$}
        Per induzione su $k$. Base: $(1y)x = yx = 1(yx)$. Passo: per la distributività,
        $((k + 1)y)x = (ky + y)x = (ky)x + yx = k(yx) + yx = (k + 1)(yx)$. A destra è identico.
    \end{caso}
    \begin{caso}{Passo 2: i tre casi di $n$}
        Se $n \geq 1$, per il passo 1 con $y = 1_a$: $(n1_a)x = n(1_a x) = nx$. Se $n = 0$: $(0 \cdot 1_a)x = 0_a x = 0_a = 0x$ (Lemma~\ref{lem:regole-anello}).
        Se $n = -k$ con $k \geq 1$: per definizione $n1_a = k(-1_a)$, quindi per il passo 1 e il Lemma~\ref{lem:regole-anello}
        $(n1_a)x = k((-1_a)x) = k(-x) = nx$. A destra è identico.
    \end{caso}
\end{dimostrazione}
\end{concetto}

\begin{spiegazione}
    Moltiplicare per «$1 + 1 + 1$» equivale a sommare tre volte: $(3 \cdot 1_a)x = (1_a + 1_a + 1_a)x = x + x + x = 3x$. La proposizione
    dice che questo vale per ogni intero, anche negativo, e sia a sinistra sia a destra, anche in un anello non commutativo.
    In $\mbZ$ sembra ovvio; in un anello come $(\mcP(s), \Delta, \cap)$ dice ad esempio che $(2 \cdot s) \cap x = 2x = x \,\Delta\, x = \varnothing$.
\end{spiegazione}


  % Insiemi finiti, unioni e prodotti, inclusione-esclusione, funzioni, permutazioni, parti, monoidi e anelli finiti
  \chapter{Insiemi finiti e calcolo combinatorio}\label{cap:finiti}

Con i numeri naturali possiamo finalmente \textbf{contare}: definire quanti elementi ha un insieme finito e calcolare quante
sono le funzioni, le funzioni iniettive, le biezioni tra due insiemi finiti e le parti di un insieme. Vedremo anche che negli
insiemi finiti iniettività, suriettività e biettività coincidono (per funzioni tra insiemi con lo stesso numero di elementi), e
ne ricaveremo due fatti sulle strutture algebriche finite lasciati in sospeso nella Parte~III.

% ──────────────────────────────────────────────────────────────
\section{Insiemi finiti}\label{sec:finiti}

\dfn[finito]{Segmenti iniziali, insieme finito, ordine}{
    Poniamo $I_0 := \varnothing$ e, per $n \in \mbN \setminus \{0\}$, $I_n := \{1, 2, \dots, n\} = \{k \in \mbN \mid 1 \leq k \leq n\}$: sono i
    \textbf{segmenti iniziali} di $\mbN \setminus \{0\}$. Un insieme $a$ si dice \textbf{finito} se è equipotente a $I_n$ per qualche
    $n \in \mbN$; $n$ si dice \textbf{ordine} (o \textbf{cardinalità}) di $a$ e si scrive $|a| = n$.
}

Contare gli elementi di $a$ vuol dire proprio costruire una biezione $I_n \to a$: «il primo, il secondo, \dots, l'$n$-esimo».
Ad esempio $|a| = 0$ significa che esiste una biezione $a \to I_0 = \varnothing$, il che obbliga $a = \varnothing$ (se $a$ avesse un
elemento, la funzione non saprebbe dove mandarlo). Perché la definizione di ordine abbia senso, bisogna escludere che uno stesso
insieme sia equipotente a due segmenti diversi: segue dal teorema seguente.

\begin{concetto}
\begin{Teorema}{I segmenti iniziali non sono infiniti}{in-non-infinito}
    \begin{prerequisiti}
        \dip{Definizione}{def:finito}
        \dip{Definizione}{def:infinito}
        \dip{Corollario}{cor:induzione-uso}
        \dip{Definizione}{def:funzioni-notevoli}
    \end{prerequisiti}
    Per ogni $n \in \mbN$, $I_n$ non è infinito: non esiste una biezione tra $I_n$ e una sua parte propria.
\end{Teorema}
\begin{dimostrazione}
    Per induzione su $n$. \textbf{Base}: $I_0 = \varnothing$ non ha parti proprie. \textbf{Passo}: supponiamo che $I_n$ non sia infinito e,
    per assurdo, sia $f \colon I_{n + 1} \to x$ biettiva con $x \subset I_{n + 1}$ parte propria.
    \begin{caso}{Caso 1: $n + 1 \notin x$}
        Allora $x \subseteq I_n$ e la restrizione $f_{|I_n}$ è una biezione tra $I_n$ e $x \setminus \{f(n + 1)\}$, che è una parte propria di $I_n$
        (le manca almeno $f(n + 1) \in I_n$). Assurdo per l'ipotesi induttiva.
    \end{caso}
    \begin{caso}{Caso 2: $n + 1 \in x$ e $f(n + 1) = n + 1$}
        Allora $f_{|I_n}$ è una biezione tra $I_n$ e $x \setminus \{n + 1\}$, che è una parte propria di $I_n$ (perché $x \neq I_{n + 1}$). Assurdo.
    \end{caso}
    \begin{caso}{Caso 3: $n + 1 \in x$ e $f(n + 1) = h \neq n + 1$}
        Sia $k \in I_n$ l'elemento con $f(k) = n + 1$. La funzione $g$ che coincide con $f$ su $I_n \setminus \{k\}$ e manda $k$ in $h$ è una
        biezione tra $I_n$ e $x \setminus \{n + 1\}$, parte propria di $I_n$. Assurdo.
    \end{caso}
\end{dimostrazione}
\end{concetto}

\begin{spiegazione}
    L'idea è ridursi da $I_{n + 1}$ a $I_n$ \textbf{togliendo l'ultimo elemento} da entrambe le parti, in modo da usare l'ipotesi
    induttiva. Se $f$ manda $n + 1$ in sé stesso basta tagliarlo via (caso 2); se lo manda altrove, prima si «aggiusta» $f$
    scambiando due frecce, come in figura per $n + 1 = 4$ (caso 3):
    \begin{center}
    \begin{tikzpicture}
        \foreach \i in {1,...,4} { \node[elem] (a\i) at (0,-\i*0.6) {$\i$}; \node[elem] (b\i) at (2.4,-\i*0.6) {$\i$}; }
        \draw[fun] (a1) -- (b2); \draw[fun] (a2) -- (b4); \draw[fun] (a3) -- (b1); \draw[fun2] (a4) -- (b3);
        \node[font=\small] at (1.2,-3.0) {$f \colon I_4 \to x = I_4$};
        \draw[-{Stealth}, line width=0.9pt, primary] (3.3,-1.5) -- node[above, font=\small] {aggiusto} (4.7,-1.5);
        \begin{scope}[xshift=5.6cm]
            \foreach \i in {1,...,3} { \node[elem] (c\i) at (0,-\i*0.6) {$\i$}; \node[elem] (d\i) at (2.4,-\i*0.6) {$\i$}; }
            \node[elem, black!30] at (0,-2.4) {$4$}; \node[elem, black!30] at (2.4,-2.4) {$4$};
            \draw[fun] (c1) -- (d2); \draw[fun3, line width=1pt] (c2) -- (d3); \draw[fun] (c3) -- (d1);
            \node[font=\small, align=center] at (1.2,-3.15) {$g \colon I_3 \to I_3$: $2$ va in $3 = f(4)$\\ al posto di $4 = f(2)$};
        \end{scope}
    \end{tikzpicture}
    \end{center}
    Il disegno mostra il meccanismo di aggiustamento su una biezione di $I_4$ in sé: $f(4) = 3$ e $f(2) = 4$; la $g$ manda $2$ in $3$ e
    dimentica $4$. Nella dimostrazione lo stesso aggiustamento produrrebbe una biezione di $I_n$ su una sua parte propria, che per
    l'ipotesi induttiva non esiste.
\end{spiegazione}

Ne seguono le proprietà che rendono il contare affidabile. Per dimostrarle servono due lemmi su come cambia l'ordine
togliendo o aggiungendo un elemento.

\begin{concetto}
\begin{Lemma}{Togliere o aggiungere un elemento}{togli-elemento}
    \begin{prerequisiti}
        \dip{Definizione}{def:finito}
        \dip{Corollario}{cor:comp-biettive}
    \end{prerequisiti}
    \begin{enumerate}
        \item Se $|a| = n \geq 1$ e $x \in a$, allora $|a \setminus \{x\}| = n - 1$.
        \item Se $|b| = n$ e $x \notin b$, allora $|b \cup \{x\}| = n + 1$.
    \end{enumerate}
\end{Lemma}
\begin{dimostrazione}
    \begin{caso}{Punto 1}
        Sia $f \colon a \to I_n$ biettiva e sia $\tau \colon I_n \to I_n$ la funzione che scambia $f(x)$ e $n$ e lascia fermi gli altri numeri
        ($\tau \circ \tau = \mathrm{id}$, quindi $\tau$ è biettiva). Allora $g = \tau \circ f$ è biettiva e $g(x) = n$. La restrizione di $g$ ad
        $a \setminus \{x\}$ è iniettiva e prende tutti i valori di $I_n$ tranne $n$: è una biezione tra $a \setminus \{x\}$ e $I_{n - 1}$.
    \end{caso}
    \begin{caso}{Punto 2}
        Se $f \colon b \to I_n$ è biettiva, la funzione che coincide con $f$ su $b$ e manda $x$ in $n + 1$ è una biezione tra $b \cup \{x\}$ e $I_{n + 1}$.
    \end{caso}
\end{dimostrazione}
\end{concetto}

\begin{spiegazione}
    Se abbiamo numerato gli elementi di $a$ da $1$ a $n$ e vogliamo togliere $x$, scambiamo il numero di $x$ con quello
    dell'ultimo elemento e poi togliamo l'ultimo: i numeri rimasti sono ancora $1, \dots, n - 1$, senza buchi. Per aggiungere un
    elemento basta dargli il numero $n + 1$.
\end{spiegazione}

\begin{concetto}
\begin{Lemma}{Le parti di un insieme finito sono finite}{parti-finite}
    \begin{prerequisiti}
        \dip{Lemma}{lem:togli-elemento}
        \dip{Corollario}{cor:induzione-uso}
    \end{prerequisiti}
    Se $s$ è finito e $t \subseteq s$, allora $t$ è finito e $|t| \leq |s|$.
\end{Lemma}
\begin{dimostrazione}
    Per induzione su $n = |s|$. \textbf{Base}: se $|s| = 0$, allora $s = \varnothing$ e $t = \varnothing$. \textbf{Passo}: sia $|s| = n + 1$, sia
    $x \in s$ e $s' = s \setminus \{x\}$, che ha ordine $n$ (Lemma~\ref{lem:togli-elemento}).
    Se $x \notin t$, allora $t \subseteq s'$ e per ipotesi induttiva $|t| \leq n \leq n + 1$. Se $x \in t$, allora $t \setminus \{x\} \subseteq s'$, che
    per ipotesi induttiva ha ordine $k \leq n$, e $|t| = k + 1 \leq n + 1$ per il Lemma~\ref{lem:togli-elemento}, punto 2.
\end{dimostrazione}
\end{concetto}

\begin{spiegazione}
    Si guarda un elemento $x$ alla volta: o resta fuori da $t$, o entra in $t$; in entrambi i casi $t$ «guadagna» al più un
    elemento rispetto a una parte di un insieme più piccolo. Sembra ovvio, ma è proprio questo che ci autorizza a dire che una
    parte ha al più tanti elementi quanti l'insieme.
\end{spiegazione}

\begin{concetto}
\begin{Teorema}{L'ordine di un insieme finito}{ordine-finito}
    \begin{prerequisiti}
        \dip{Teorema}{teo:in-non-infinito}
        \dip{Lemma}{lem:parti-finite}
        \dip{Proposizione}{prop:equipotenza}
    \end{prerequisiti}
    \begin{enumerate}
        \item Se $a$ è equipotente sia a $I_m$ sia a $I_n$, allora $m = n$: l'ordine di un insieme finito è ben definito.
        \item Un insieme finito non è infinito.
        \item Se $t$ è una parte \textbf{propria} di un insieme finito $s$, allora $|t| < |s|$.
    \end{enumerate}
\end{Teorema}
\begin{dimostrazione}
    Il punto 2 il programma lo indica senza dimostrazione; la diamo comunque perché serve per il punto 3.
    \begin{caso}{Punto 1}
        $I_m$ e $I_n$ sono equipotenti (Proposizione~\ref{prop:equipotenza}). Se fosse ad esempio $m < n$, $I_m$ sarebbe una parte propria di
        $I_n$ equipotente a $I_n$, e $I_n$ sarebbe infinito, contro il Teorema~\ref{teo:in-non-infinito}.
    \end{caso}
    \begin{caso}{Punto 2}
        Sia $f \colon s \to I_n$ biettiva e, per assurdo, $g \colon s \to t$ biettiva con $t \subset s$. Allora $f \circ g \circ f^{-1}$ è una biezione tra
        $I_n$ e $\overrightarrow{f}(t)$, che è una parte propria di $I_n$ (se $y \in s \setminus t$, $f(y) \notin \overrightarrow{f}(t)$ perché $f$ è
        iniettiva). Assurdo per il Teorema~\ref{teo:in-non-infinito}.
    \end{caso}
    \begin{caso}{Punto 3}
        Per il Lemma~\ref{lem:parti-finite} $|t| \leq |s|$. Se fosse $|t| = |s| = n$, $s$ e $t$ sarebbero entrambi equipotenti a $I_n$, quindi tra loro:
        $s$ sarebbe equipotente alla sua parte propria $t$, cioè infinito, contro il punto 2.
    \end{caso}
\end{dimostrazione}
\end{concetto}

\begin{spiegazione}
    Questi tre fatti sono il «contratto» del contare: contando gli stessi oggetti in due ordini diversi si ottiene lo stesso numero
    (punto 1), e togliendo qualcosa a un insieme finito il numero scende davvero (punto 3). Con gli insiemi infiniti il punto 3 è
    falso: $\mbN \setminus \{0\}$ è una parte propria di $\mbN$ ma i due sono equipotenti (albergo di Hilbert, Esempio~\ref{es:hilbert}).
\end{spiegazione}

% ──────────────────────────────────────────────────────────────
\section{Unioni e prodotti di insiemi finiti}\label{sec:unioni-finite}

\begin{concetto}
\begin{Lemma}{Unione disgiunta e prodotto cartesiano}{somma-prodotto}
    \begin{prerequisiti}
        \dip{Lemma}{lem:togli-elemento}
        \dip{Corollario}{cor:induzione-uso}
        \dip{Definizione}{def:prodotto-cartesiano}
    \end{prerequisiti}
    Siano $a$ e $b$ insiemi finiti.
    \begin{enumerate}
        \item Se $a \cap b = \varnothing$, allora $|a \cup b| = |a| + |b|$.
        \item $|a \times b| = |a| \cdot |b|$.
        \item Se $p = \{w_1, \dots, w_r\}$ è una partizione di un insieme finito $c$ in $r$ blocchi, allora $|c| = |w_1| + \dots + |w_r|$; in particolare,
              se tutti i blocchi hanno $k$ elementi, $|c| = rk$.
    \end{enumerate}
\end{Lemma}
\begin{dimostrazione}
    \begin{caso}{Punto 1, per induzione su $|b|$}
        Se $|b| = 0$, $a \cup \varnothing = a$. Se $|b| = k + 1$, sia $y \in b$ e $b' = b \setminus \{y\}$: per ipotesi induttiva $|a \cup b'| = |a| + k$;
        inoltre $y \notin a \cup b'$ (perché $a \cap b = \varnothing$), quindi $|a \cup b| = |(a \cup b') \cup \{y\}| = |a| + k + 1$.
    \end{caso}
    \begin{caso}{Punto 2, per induzione su $|b|$}
        Se $b = \varnothing$, $a \times b = \varnothing$. Altrimenti, con $y$ e $b'$ come sopra, $a \times b = (a \times b') \cup (a \times \{y\})$, unione
        disgiunta; $x \mapsto (x, y)$ è una biezione tra $a$ e $a \times \{y\}$, quindi per il punto 1 $|a \times b| = |a| \cdot |b'| + |a| = |a|(|b'| + 1)$.
    \end{caso}
    \begin{caso}{Punto 3, per induzione su $r$}
        I blocchi sono a due a due disgiunti, quindi si aggiunge un blocco alla volta e si applica il punto 1.
    \end{caso}
\end{dimostrazione}
\end{concetto}

\begin{spiegazione}
    Sono le due regole base del contare. \textbf{Somma}: se scelgo una cosa tra $3$ magliette oppure tra $2$ felpe (insiemi
    disgiunti), ho $3 + 2 = 5$ possibilità. \textbf{Prodotto}: se scelgo una maglietta \textbf{e} un paio di pantaloni tra $4$, ho
    $3 \cdot 4 = 12$ coppie: la tabella delle coppie ha $3$ righe e $4$ colonne. La partizione è la somma ripetuta: un'aula con
    $r$ file da $k$ posti ha $rk$ posti.
\end{spiegazione}

\dfn[unione-finita]{Unione di un numero finito di insiemi}{
    Se $a_1, \dots, a_n$ sono insiemi, scriviamo $\displaystyle\bigcup_{i = 1}^{n} a_i := a_1 \cup a_2 \cup \dots \cup a_n = \bigcup\{a_1, \dots, a_n\}$.
}

Se gli insiemi non sono disgiunti la somma degli ordini conta due volte gli elementi in comune. Il principio di
inclusione-esclusione corregge l'errore.

\begin{concetto}
\begin{Teorema}{Principio di inclusione-esclusione}{inclusione-esclusione}
    \begin{prerequisiti}
        \dip{Lemma}{lem:somma-prodotto}
        \dip{Lemma}{lem:parti-finite}
        \dip{Proposizione}{prop:distributivita-ins}
        \dip{Proposizione}{prop:terzo-escluso-ins}
    \end{prerequisiti}
    Siano $a$, $b$, $c$ insiemi finiti. Allora
    \begin{enumerate}
        \item $|a \cup b| = |a| + |b| - |a \cap b|$;
        \item $|a \cup b \cup c| = |a| + |b| + |c| - |a \cap b| - |a \cap c| - |b \cap c| + |a \cap b \cap c|$.
    \end{enumerate}
    Più in generale (\textbf{senza dimostrazione}), per $a_1, \dots, a_n$ finiti:
    \[ \Big|\bigcup_{i = 1}^{n} a_i\Big| = \sum_{i} |a_i| - \sum_{i < j} |a_i \cap a_j| + \sum_{i < j < k} |a_i \cap a_j \cap a_k| - \dots + (-1)^{n - 1} |a_1 \cap \dots \cap a_n|. \]
\end{Teorema}
\begin{dimostrazione}
    Tutti gli insiemi coinvolti sono finiti, perché parti di $a \cup b \cup c$ (Lemma~\ref{lem:parti-finite}).
    \begin{caso}{Punto 1}
        $a \cup b = (a \setminus b) \cup b$ e $a = (a \setminus b) \cup (a \cap b)$, entrambe unioni disgiunte. Per il Lemma~\ref{lem:somma-prodotto}
        $|a \cup b| = |a \setminus b| + |b|$ e $|a| = |a \setminus b| + |a \cap b|$. Sottraendo membro a membro: $|a \cup b| - |a| = |b| - |a \cap b|$.
    \end{caso}
    \begin{caso}{Punto 2}
        Per il punto 1 applicato a $a \cup b$ e $c$: $|a \cup b \cup c| = |a \cup b| + |c| - |(a \cup b) \cap c|$. Per la distributività
        $(a \cup b) \cap c = (a \cap c) \cup (b \cap c)$, e per il punto 1 ha ordine $|a \cap c| + |b \cap c| - |a \cap b \cap c|$. Sostituendo anche
        $|a \cup b| = |a| + |b| - |a \cap b|$ si ottiene la formula.
    \end{caso}
\end{dimostrazione}
\end{concetto}

\begin{spiegazione}
    Con i diagrammi di Venn: nella somma $|a| + |b| + |c|$ ogni zona viene contata tante volte quanti sono gli insiemi che la
    contengono. Togliendo le intersezioni a due, la zona centrale (contata $3$ volte) viene tolta $3$ volte: va rimessa una volta.
    \begin{center}
    \begin{tikzpicture}[scale=0.85]
        \begin{scope}
            \clip (0,0.55) circle (1.1); \clip (-0.6,-0.45) circle (1.1); \fill[orange!45] (0.6,-0.45) circle (1.1);
        \end{scope}
        \begin{scope}
            \clip (0,0.55) circle (1.1); \fill[primary!18] (-0.6,-0.45) circle (1.1); \fill[primary!18] (0.6,-0.45) circle (1.1);
        \end{scope}
        \begin{scope}
            \clip (-0.6,-0.45) circle (1.1); \fill[primary!18] (0.6,-0.45) circle (1.1);
        \end{scope}
        \begin{scope}
            \clip (0,0.55) circle (1.1); \clip (-0.6,-0.45) circle (1.1); \fill[orange!45] (0.6,-0.45) circle (1.1);
        \end{scope}
        \draw (0,0.55) circle (1.1); \draw (-0.6,-0.45) circle (1.1); \draw (0.6,-0.45) circle (1.1);
        \node at (0,1.2) {$1$}; \node at (-1.15,-0.8) {$1$}; \node at (1.15,-0.8) {$1$};
        \node at (-0.62,0.25) {$2$}; \node at (0.62,0.25) {$2$}; \node at (0,-0.85) {$2$}; \node at (0,-0.05) {$3$};
        \node[font=\small] at (0,1.95) {$a$}; \node[font=\small] at (-1.9,-1.1) {$b$}; \node[font=\small] at (1.9,-1.1) {$c$};
        \node[font=\small, align=left, anchor=west] at (2.5,0.2)
            {quante volte ogni zona è contata in $|a| + |b| + |c|$;\\
             le intersezioni a due tolgono $1$ dalle zone «$2$» e $3$ dalla zona «$3$»;\\
             aggiungendo $|a \cap b \cap c|$ ogni zona è contata esattamente una volta.};
    \end{tikzpicture}
    \end{center}
    Ad esempio, se $|a| = 5$, $|b| = 9$ e $|a \cap b| = 3$, allora $|a \cup b| = 5 + 9 - 3 = 11$.
\end{spiegazione}

% ──────────────────────────────────────────────────────────────
\section{Contare le funzioni}\label{sec:contare-funzioni}

\dfn[fattoriale]{Fattoriale}{
    Il \textbf{fattoriale} di un numero naturale è definito ricorsivamente da $0! := 1$ e $(n + 1)! := (n + 1) \cdot n!$.
    Quindi $1! = 1$, $2! = 2$, $3! = 6$, $4! = 24$ e, per $n \geq 1$, $n! = 1 \cdot 2 \cdot \ldots \cdot n$.
}

\begin{concetto}
\begin{Teorema}{Numero delle funzioni}{conta-map}
    \begin{prerequisiti}
        \dip{Definizione}{def:map-in-su-bi}
        \dip{Lemma}{lem:togli-elemento}
        \dip{Lemma}{lem:somma-prodotto}
        \dip{Corollario}{cor:induzione-uso}
    \end{prerequisiti}
    Se $|a| = m$ e $|b| = n$, allora $|\mathrm{Map}(a, b)| = n^m$.
\end{Teorema}
\begin{dimostrazione}
    Per induzione su $m$. \textbf{Base}: se $m = 0$, $a = \varnothing$ e l'unica funzione da $\varnothing$ a $b$ è $(\varnothing, b, \varnothing)$: $1 = n^0$.
    \textbf{Passo}: sia $|a| = m + 1$, sia $x \in a$ e $a' = a \setminus \{x\}$, di ordine $m$. La funzione
    \[ \Psi \colon f \in \mathrm{Map}(a, b) \mapsto (f_{|a'}, f(x)) \in \mathrm{Map}(a', b) \times b \]
    è biettiva: la sua inversa associa alla coppia $(g, y)$ la funzione che coincide con $g$ su $a'$ e manda $x$ in $y$. Per il
    Lemma~\ref{lem:somma-prodotto} e l'ipotesi induttiva, $|\mathrm{Map}(a, b)| = |\mathrm{Map}(a', b)| \cdot |b| = n^m \cdot n = n^{m + 1}$.
\end{dimostrazione}
\end{concetto}

\begin{spiegazione}
    Una funzione da $a$ a $b$ è una sequenza di \textbf{scelte indipendenti}: per ognuno degli $m$ elementi di $a$ scegliamo
    un'immagine tra $n$. Con $a = \{1, 2\}$ e $b = \{p, q, r\}$ l'albero delle scelte ha $3 \cdot 3 = 3^2 = 9$ foglie:
    \begin{center}
    \begin{tikzpicture}[level distance=1.1cm, every node/.style={font=\small}, edge from parent/.style={draw=black!50},
                        level 1/.style={sibling distance=3.2cm}, level 2/.style={sibling distance=0.85cm}]
        \node {scelta di $f(1)$}
            child {node {$p$} child {node {$p$}} child {node {$q$}} child {node {$r$}}}
            child {node {$q$} child {node {$p$}} child {node {$q$}} child {node {$r$}}}
            child {node {$r$} child {node {$p$}} child {node {$q$}} child {node {$r$}}};
        \node[anchor=west, font=\small] at (5.0,-1.1) {$f(1)$: $3$ scelte};
        \node[anchor=west, font=\small] at (5.0,-2.2) {$f(2)$: $3$ scelte per ognuna, $9$ funzioni in tutto};
    \end{tikzpicture}
    \end{center}
    La biezione $\Psi$ della dimostrazione dice proprio questo: scegliere $f$ è come scegliere cosa fa sugli altri elementi
    ($f_{|a'}$) e, indipendentemente, dove manda $x$.
\end{spiegazione}

\begin{concetto}
\begin{Teorema}{Numero delle funzioni iniettive}{conta-iniettive}
    \begin{prerequisiti}
        \dip{Definizione}{def:fattoriale}
        \dip{Lemma}{lem:parti-finite}
        \dip{Lemma}{lem:somma-prodotto}
        \dip{Proposizione}{prop:ridotta}
        \dip{Teorema}{teo:comp-iniettive}
    \end{prerequisiti}
    Siano $|a| = m$ e $|b| = n$. Esistono funzioni iniettive da $a$ a $b$ se e solo se $m \leq n$, e in tal caso
    \[ |\mathrm{In}(a, b)| = \frac{n!}{(n - m)!} = n(n - 1) \cdots (n - m + 1). \]
\end{Teorema}
\begin{dimostrazione}
    \begin{caso}{$(\rightarrow)$ Se $f \in \mathrm{In}(a, b)$, allora $m \leq n$}
        La ridotta $a \to \mathrm{Im} f$ è biettiva, quindi $|\mathrm{Im} f| = m$; e $\mathrm{Im} f \subseteq b$, quindi $m \leq n$ (Lemma~\ref{lem:parti-finite}).
    \end{caso}
    \begin{caso}{$(\leftarrow)$ Se $m \leq n$, esistono funzioni iniettive}
        Siano $\alpha \colon a \to I_m$ e $\beta \colon b \to I_n$ biettive. Poiché $I_m \subseteq I_n$, l'immersione $\iota \colon I_m \to I_n$ è iniettiva,
        e $\beta^{-1} \circ \iota \circ \alpha \colon a \to b$ è iniettiva (composta di iniettive).
    \end{caso}
    \begin{caso}{Il conteggio, per induzione su $m \leq n$}
        Se $m = 0$, c'è solo la funzione vuota: $1 = n!/n!$. Se $m \geq 1$, sia $x \in a$ e $a' = a \setminus \{x\}$. Come nel teorema
        precedente, $f \mapsto (f_{|a'}, f(x))$ è una biezione tra $\mathrm{In}(a, b)$ e l'insieme delle coppie $(g, y)$ con $g \in \mathrm{In}(a', b)$ e
        $y \in b \setminus \mathrm{Im}\, g$ (perché $f$ resti iniettiva, $x$ deve andare in un elemento non ancora usato). Per ogni $g$ ci sono
        $|b \setminus \mathrm{Im}\, g| = n - (m - 1)$ scelte di $y$, quindi (Lemma~\ref{lem:somma-prodotto}, punto 3) e per ipotesi induttiva
        \[ |\mathrm{In}(a, b)| = \frac{n!}{(n - m + 1)!} \cdot (n - m + 1) = \frac{n!}{(n - m)!}. \]
    \end{caso}
\end{dimostrazione}
\end{concetto}

\begin{spiegazione}
    Ora le scelte \textbf{non} sono indipendenti: ogni elemento di $a$ deve andare in un elemento di $b$ non ancora usato. Il primo
    ha $n$ scelte, il secondo $n - 1$, \dots, l'$m$-esimo $n - m + 1$. Con $a = \{1, 2\}$ e $b = \{p, q, r\}$: $f(1)$ ha $3$ scelte, $f(2)$
    solo $2$, per un totale di $3 \cdot 2 = 6 = \frac{3!}{1!}$ funzioni iniettive. È il numero di modi di \textbf{ordinare} $m$ oggetti
    scelti tra $n$: ad esempio di assegnare oro e argento a $2$ dei $3$ concorrenti di una gara.
\end{spiegazione}

\begin{concetto}
\begin{Proposizione}{Esistenza di funzioni suriettive}{esistono-suriettive}
    \begin{prerequisiti}
        \dip{Teorema}{teo:suriettiva-sezione}
        \dip{Teorema}{teo:composte-caratteristiche}
        \dip{Teorema}{teo:conta-iniettive}
    \end{prerequisiti}
    Siano $|a| = m$ e $|b| = n$. Esistono funzioni suriettive da $a$ a $b$ se e solo se $a = b = \varnothing$ oppure $0 < n \leq m$.
\end{Proposizione}
\begin{dimostrazione}
    \begin{caso}{$(\rightarrow)$}
        Sia $f \colon a \to b$ suriettiva. Se $b = \varnothing$, anche $a = \varnothing$ (non ci sono funzioni da un insieme non vuoto al vuoto).
        Se $b \neq \varnothing$, $n > 0$; $f$ ha una sezione $g \colon b \to a$ (Teorema~\ref{teo:suriettiva-sezione}), e $f \circ g = \mathrm{id}_b$ è iniettiva,
        quindi $g$ è iniettiva (Teorema~\ref{teo:composte-caratteristiche}): per il Teorema~\ref{teo:conta-iniettive} $n \leq m$.
    \end{caso}
    \begin{caso}{$(\leftarrow)$}
        La funzione vuota da $\varnothing$ a $\varnothing$ è suriettiva. Se $0 < n \leq m$, la funzione $h \colon I_m \to I_n$ con $h(k) = k$ per $k \leq n$ e
        $h(k) = n$ per $k > n$ è suriettiva; componendo con biezioni $a \to I_m$ e $I_n \to b$ si ottiene una suriettiva da $a$ a $b$.
    \end{caso}
\end{dimostrazione}
\end{concetto}

\begin{spiegazione}
    Per «coprire» tutti gli $n$ elementi di $b$ servono almeno $n$ frecce che partono da elementi diversi, quindi almeno $n$
    elementi in $a$. Il caso $n = 0$ va trattato a parte: verso l'insieme vuoto parte solo la funzione vuota.
\end{spiegazione}

Quando $a$ e $b$ hanno lo stesso numero di elementi succede qualcosa di notevole.

\begin{concetto}
\begin{Teorema}{Iniettive, suriettive e biettive tra insiemi con lo stesso ordine}{stesso-ordine}
    \begin{prerequisiti}
        \dip{Teorema}{teo:ordine-finito}
        \dip{Proposizione}{prop:ridotta}
        \dip{Teorema}{teo:suriettiva-sezione}
        \dip{Teorema}{teo:composte-caratteristiche}
        \dip{Teorema}{teo:biettiva-inversa}
    \end{prerequisiti}
    Siano $a$ e $b$ finiti con $|a| = |b| = n$ e sia $f \colon a \to b$. Allora
    \[ f \text{ è iniettiva} \ \leftrightarrow\ f \text{ è biettiva} \ \leftrightarrow\ f \text{ è suriettiva}. \]
    Cioè $\mathrm{In}(a, b) = \mathrm{Bi}(a, b) = \mathrm{Su}(a, b)$.
\end{Teorema}
\begin{dimostrazione}
    Basta mostrare che ogni iniettiva e ogni suriettiva sono biettive.
    \begin{caso}{Iniettiva $\Rightarrow$ biettiva}
        La ridotta $a \to \mathrm{Im} f$ è biettiva, quindi $|\mathrm{Im} f| = n = |b|$. Se $\mathrm{Im} f$ fosse una parte propria di $b$, avrebbe ordine
        minore di $n$ (Teorema~\ref{teo:ordine-finito}, punto 3). Quindi $\mathrm{Im} f = b$: $f$ è suriettiva.
    \end{caso}
    \begin{caso}{Suriettiva $\Rightarrow$ biettiva}
        $f$ ha una sezione $g \colon b \to a$, con $f \circ g = \mathrm{id}_b$; $g$ è iniettiva (Teorema~\ref{teo:composte-caratteristiche}), quindi biettiva per il
        caso precedente (con i ruoli di $a$ e $b$ scambiati). Allora $f = f \circ g \circ g^{-1} = g^{-1}$ è biettiva.
    \end{caso}
\end{dimostrazione}
\end{concetto}

\begin{spiegazione}
    È il \textbf{principio dei cassetti}: se metto $n$ oggetti in $n$ cassetti senza mai metterne due nello stesso (iniettiva), li
    riempio tutti (suriettiva); e se li riempio tutti, non posso averne messi due nello stesso. Con gli insiemi infiniti è falso:
    \begin{itemize}
        \item $n \in \mbN \mapsto n + 1 \in \mbN$ è iniettiva ma non suriettiva ($0$ resta scoperto);
        \item $n \in \mbN \mapsto n - 1 \in \mbN$ per $n \geq 1$, con $0 \mapsto 0$, è suriettiva ma non iniettiva ($0$ e $1$ vanno entrambi in $0$).
    \end{itemize}
\end{spiegazione}

\begin{concetto}
\begin{Corollario}{Numero delle funzioni biettive}{conta-biettive}
    \begin{prerequisiti}
        \dip{Teorema}{teo:conta-iniettive}
        \dip{Teorema}{teo:stesso-ordine}
        \dip{Teorema}{teo:biettiva-inversa}
    \end{prerequisiti}
    Siano $|a| = m$ e $|b| = n$. Esistono funzioni biettive da $a$ a $b$ se e solo se $m = n$, e in tal caso $|\mathrm{Bi}(a, b)| = n!$.
    In particolare l'insieme $\mathrm{Bi}(a) := \mathrm{Bi}(a, a)$ delle biezioni di $a$ in sé (le \textbf{permutazioni} di $a$) ha $m!$ elementi.
\end{Corollario}
\begin{dimostrazione}
    Se $f \colon a \to b$ è biettiva, $f$ e $f^{-1}$ sono iniettive, quindi $m \leq n$ e $n \leq m$ (Teorema~\ref{teo:conta-iniettive}). Se $m = n$,
    per il Teorema~\ref{teo:stesso-ordine} $\mathrm{Bi}(a, b) = \mathrm{In}(a, b)$, che ha $\frac{n!}{(n - n)!} = \frac{n!}{0!} = n!$ elementi.
\end{dimostrazione}
\end{concetto}

\begin{spiegazione}
    Le permutazioni di $\{1, 2, 3\}$ sono $3! = 6$: il primo posto si sceglie in $3$ modi, il secondo in $2$, il terzo in $1$. Eccole,
    scritte come la sequenza $f(1)f(2)f(3)$: \quad $123$, \ $132$, \ $213$, \ $231$, \ $312$, \ $321$.
\end{spiegazione}

% ──────────────────────────────────────────────────────────────
\section{Insiemi infiniti a confronto}\label{sec:confronto-infiniti}

Anche tra insiemi infiniti si possono confrontare le «grandezze» con le biezioni. Qui diamo solo qualche cenno.

\begin{concetto}
\begin{Proposizione}{$\mbN$ e $\mbZ$ sono equipotenti}{n-z}
    \begin{prerequisiti}
        \dip{Definizione}{def:equipotenti}
        \dip{Teorema}{teo:biettiva-inversa}
    \end{prerequisiti}
    La funzione $f \colon \mbN \to \mbZ$ definita da $f(n) = -\frac{n}{2}$ se $n$ è pari e $f(n) = \frac{n + 1}{2}$ se $n$ è dispari è biettiva.
\end{Proposizione}
\begin{dimostrazione}
    Sia $g \colon \mbZ \to \mbN$ con $g(z) = -2z$ se $z \leq 0$ e $g(z) = 2z - 1$ se $z > 0$. Se $n$ è pari, $f(n) = -\frac{n}{2} \leq 0$ e
    $g(f(n)) = n$; se $n$ è dispari, $f(n) = \frac{n + 1}{2} > 0$ e $g(f(n)) = n$. Se $z \leq 0$, $g(z) = -2z$ è pari e $f(g(z)) = z$; se $z > 0$,
    $g(z) = 2z - 1$ è dispari e $f(g(z)) = z$. Quindi $g$ è l'inversa di $f$, e $f$ è biettiva (Teorema~\ref{teo:biettiva-inversa}).
\end{dimostrazione}
\end{concetto}

\begin{spiegazione}
    Sembra che $\mbZ$ abbia «il doppio» degli elementi di $\mbN$, eppure si possono elencare tutti gli interi in una sola fila,
    \textbf{a zig-zag}: $0, 1, -1, 2, -2, 3, -3, \dots$ È quello che fa $f$:
    \begin{center}
    \renewcommand{\arraystretch}{1.2}
    \begin{tabular}{c|ccccccccc}
        $n$ & $0$ & $1$ & $2$ & $3$ & $4$ & $5$ & $6$ & $7$ & $\cdots$ \\ \hline
        $f(n)$ & $0$ & $1$ & $-1$ & $2$ & $-2$ & $3$ & $-3$ & $4$ & $\cdots$ \\
    \end{tabular}

    \medskip
    \begin{tikzpicture}[x=1.0cm]
        \foreach \z in {-3,...,3} \node[circle, draw=black!40, fill=primary!7, inner sep=0pt, minimum size=6.5mm, font=\small] (z\z) at (\z,0) {$\z$};
        \foreach \a/\b in {0/1, 1/-1, -1/2, 2/-2, -2/3, 3/-3}
            \draw[-{Stealth}, orange!80!black, line width=0.7pt] (z\a) to[bend left=50] (z\b);
        \node[font=\small, anchor=west] at (3.8,0) {la spirale tocca ogni intero una volta sola};
    \end{tikzpicture}
    \end{center}
\end{spiegazione}

\begin{concetto}
\begin{Teorema}{Teorema di Cantor}{cantor}
    \begin{prerequisiti}
        \dip{Definizione}{def:suriettiva}
        \dip{Assioma}{ass:separazione}
    \end{prerequisiti}
    Per nessun insieme $a$ esiste una funzione suriettiva da $a$ a $\mcP(a)$. In particolare $\mbN$ e $\mcP(\mbN)$ non sono equipotenti.
\end{Teorema}
\begin{dimostrazione}
    \facoltativa
    Per assurdo sia $f \colon a \to \mcP(a)$ suriettiva e sia $d = \{x \in a \mid x \notin f(x)\}$. Poiché $d \in \mcP(a)$, esiste $y \in a$ con
    $f(y) = d$. Allora $y \in d \leftrightarrow y \notin f(y) \leftrightarrow y \notin d$: assurdo.
\end{dimostrazione}
\end{concetto}

\begin{spiegazione}
    Ogni elemento $x$ «propone» una parte $f(x)$; la parte $d$ è costruita per essere \textbf{diversa da tutte le proposte}: per
    ogni $x$, $d$ e $f(x)$ non sono d'accordo proprio su $x$. Quindi nessuno propone $d$. È lo stesso trucco del paradosso di
    Russell, e dice che esistono infiniti «più grandi» di altri: $\mcP(\mbN)$ non si può elencare in una fila come $\mbN$.
\end{spiegazione}

% ──────────────────────────────────────────────────────────────
\section{Funzioni caratteristiche e insieme delle parti}\label{sec:caratteristiche}

\dfn[funzione-caratteristica]{Funzione caratteristica}{
    Sia $s$ un insieme e $t \subseteq s$. La funzione
    \[ \chi_{t, s} \colon x \in s \mapsto \begin{cases} 1 & \text{se } x \in t \\ 0 & \text{se } x \notin t \end{cases} \in \{0, 1\} \]
    si chiama \textbf{funzione caratteristica} di $t$ in $s$.
}

\begin{concetto}
\begin{Teorema}{Parti e funzioni caratteristiche}{parti-caratteristiche}
    \begin{prerequisiti}
        \dip{Definizione}{def:funzione-caratteristica}
        \dip{Definizione}{def:biettiva}
        \dip{Corollario}{cor:terna}
    \end{prerequisiti}
    Sia $s$ un insieme. La funzione $\varphi \colon t \in \mcP(s) \mapsto \chi_{t, s} \in \mathrm{Map}(s, \{0, 1\})$ è biettiva.
\end{Teorema}
\begin{dimostrazione}
    \begin{caso}{Suriettiva}
        Sia $f \in \mathrm{Map}(s, \{0, 1\})$ e sia $t = \{x \in s \mid f(x) = 1\}$. Per $x \in t$, $\chi_{t, s}(x) = 1 = f(x)$; per $x \in s \setminus t$,
        $\chi_{t, s}(x) = 0 = f(x)$. Le due funzioni hanno lo stesso grafico, quindi $\varphi(t) = \chi_{t, s} = f$.
    \end{caso}
    \begin{caso}{Iniettiva}
        Siano $t \neq u$ parti di $s$: senza perdere di generalità esiste $x \in t \setminus u$. Allora $\chi_{t, s}(x) = 1 \neq 0 = \chi_{u, s}(x)$, quindi
        $\chi_{t, s} \neq \chi_{u, s}$.
    \end{caso}
\end{dimostrazione}
\end{concetto}

\begin{spiegazione}
    Una parte di $s$ è una lista di risposte «sì/no» alla domanda «$x$ sta in $t$?», una per ogni elemento di $s$. Se gli elementi
    di $s$ sono messi in fila, è una \textbf{stringa di bit}. Per $s = \{a, b, c\}$:
    \begin{center}
    \renewcommand{\arraystretch}{1.2}
    \begin{tabular}{c|cccccccc}
        $t$ & $\varnothing$ & $\{c\}$ & $\{b\}$ & $\{b, c\}$ & $\{a\}$ & $\{a, c\}$ & $\{a, b\}$ & $\{a, b, c\}$ \\ \hline
        $\chi_{t, s}(a)\,\chi_{t, s}(b)\,\chi_{t, s}(c)$ & $000$ & $001$ & $010$ & $011$ & $100$ & $101$ & $110$ & $111$ \\
    \end{tabular}
    \end{center}
    Le parti e le stringhe si corrispondono una a una: è l'idea con cui un computer memorizza un sottoinsieme.
\end{spiegazione}

\begin{concetto}
\begin{Corollario}{Le parti di un insieme finito}{cardinalita-parti}
    \begin{prerequisiti}
        \dip{Teorema}{teo:parti-caratteristiche}
        \dip{Teorema}{teo:conta-map}
    \end{prerequisiti}
    Se $|s| = n$, allora $|\mcP(s)| = 2^n$.
\end{Corollario}
\begin{dimostrazione}
    Per il Teorema~\ref{teo:parti-caratteristiche} $\mcP(s)$ è equipotente a $\mathrm{Map}(s, \{0, 1\})$, che ha $2^n$ elementi per il
    Teorema~\ref{teo:conta-map}.
\end{dimostrazione}
\end{concetto}

\begin{spiegazione}
    Si può vedere anche per induzione: le parti di $s \cup \{x\}$ (con $x \notin s$) sono quelle che \textbf{non} contengono $x$, cioè le
    parti di $s$, e quelle che lo contengono, ottenute aggiungendo $x$ a una parte di $s$. Ogni elemento in più \textbf{raddoppia} il
    numero delle parti: $2^n + 2^n = 2^{n + 1}$.
    \begin{center}
    \begin{tikzpicture}[every node/.style={font=\small}]
        \node[hv] at (0,0) {$\varnothing$}; \node[hv] at (1.0,0) {$\{a\}$}; \node[hv] at (2.0,0) {$\{b\}$}; \node[hv] at (3.2,0) {$\{a, b\}$};
        \node[hv evid] at (0,-0.9) {$\{c\}$}; \node[hv evid] at (1.15,-0.9) {$\{a, c\}$}; \node[hv evid] at (2.4,-0.9) {$\{b, c\}$};
        \node[hv evid] at (3.75,-0.9) {$\{a, b, c\}$};
        \node[anchor=west] at (4.8,0) {parti di $\{a, b\}$: $4$};
        \node[anchor=west] at (4.8,-0.9) {le stesse con $c$ aggiunto: altre $4$};
    \end{tikzpicture}
    \end{center}
\end{spiegazione}

% ──────────────────────────────────────────────────────────────
\section{Monoidi e anelli finiti}\label{sec:monoidi-finiti}

Nel Capitolo~\ref{cap:gruppi} abbiamo visto che in un monoide ogni elemento invertibile è cancellabile, ma che il viceversa è
falso in generale (in $(\mbZ, \cdot, 1)$ il $2$ è cancellabile ma non invertibile). Negli insiemi finiti il
Teorema~\ref{teo:stesso-ordine} ribalta la situazione.

\begin{concetto}
\begin{Teorema}{Monoide finito: invertibile se e solo se cancellabile}{monoide-finito}
    \begin{prerequisiti}
        \dip{Teorema}{teo:invertibile-cancellabile}
        \dip{Proposizione}{prop:traslazioni}
        \dip{Teorema}{teo:stesso-ordine}
        \dip{Teorema}{teo:unicita-inverso}
    \end{prerequisiti}
    Sia $(m, *, 1_m)$ un monoide finito. Allora $a \in m$ è invertibile se e solo se $a$ è cancellabile.
\end{Teorema}
\begin{dimostrazione}
    \begin{caso}{$(\rightarrow)$ Se $a$ è invertibile, allora è cancellabile}
        È il Teorema~\ref{teo:invertibile-cancellabile}.
    \end{caso}
    \begin{caso}{$(\leftarrow)$ Se $a$ è cancellabile, allora è invertibile}
        Le traslazioni $\sigma_a$ e $\delta_a$ sono iniettive (Proposizione~\ref{prop:traslazioni}). Sono funzioni dall'insieme finito $m$ in sé,
        quindi sono anche suriettive (Teorema~\ref{teo:stesso-ordine}). Allora esistono $b, c \in m$ con
        \[ \sigma_a(b) = a * b = 1_m \qquad\qquad \delta_a(c) = c * a = 1_m: \]
        $a$ ha un inverso destro e uno sinistro, quindi è invertibile (e $b = c$ per il Teorema~\ref{teo:unicita-inverso}).
    \end{caso}
\end{dimostrazione}
\end{concetto}

\begin{spiegazione}
    Sulla tavola di Cayley: $a$ cancellabile a sinistra vuol dire che nella riga di $a$ non ci sono ripetizioni
    (Proposizione~\ref{prop:lettura-cayley}). Se la tavola è finita, una riga di $n$ caselle senza ripetizioni contiene
    \textbf{tutti} gli $n$ elementi, quindi anche $1_m$: e la colonna in cui compare $1_m$ è l'inverso destro di $a$. In $(\mbZ, \cdot)$ la
    riga del $2$ ($\dots, -2, 0, 2, 4, \dots$) non ha ripetizioni ma, essendo infinita, può «saltare» l'$1$.
\end{spiegazione}

\begin{concetto}
\begin{Teorema}{Anelli integri finiti}{integri-finiti}
    \begin{prerequisiti}
        \dip{Teorema}{teo:monoide-finito}
        \dip{Teorema}{teo:divisori-cancellabili}
        \dip{Definizione}{def:corpo-campo}
        \dip{Definizione}{def:lap}
        \dip{Lemma}{lem:regole-anello}
    \end{prerequisiti}
    Sia $(a, +, \cdot)$ un anello unitario finito, integro, con almeno due elementi. Allora $a$ è un corpo.
    In particolare, ogni dominio di integrità finito è un campo.
\end{Teorema}
\begin{dimostrazione}
    \begin{caso}{$(a \setminus \{0_a\}, \cdot, 1_a)$ è un monoide}
        $1_a \neq 0_a$ (Lemma~\ref{lem:regole-anello}, punto 4), quindi $1_a \in a \setminus \{0_a\}$. Per la LAP, se $x, y \neq 0_a$ allora $xy \neq 0_a$:
        $a \setminus \{0_a\}$ è una parte chiusa di $(a, \cdot)$.
    \end{caso}
    \begin{caso}{Ogni suo elemento è invertibile}
        Sia $x \neq 0_a$. Per la LAP $x$ non è divisore dello zero, quindi è cancellabile (Teorema~\ref{teo:divisori-cancellabili}), anche nel
        monoide finito $a \setminus \{0_a\}$. Per il Teorema~\ref{teo:monoide-finito} $x$ è invertibile. Quindi $(a \setminus \{0_a\}, \cdot)$ è un gruppo:
        $a$ è un corpo. Se $a$ è anche commutativo (dominio di integrità), è un campo.
    \end{caso}
\end{dimostrazione}
\end{concetto}

\begin{spiegazione}
    In un anello integro moltiplicare per $x \neq 0$ non «schiaccia» mai due elementi diversi nello stesso risultato
    ($xy = xz \Rightarrow y = z$). Se gli elementi sono finiti, moltiplicare per $x$ è quindi una permutazione degli elementi non nulli,
    e prima o poi si ottiene $1$: lì c'è l'inverso di $x$. Vedremo nel Capitolo~\ref{cap:congruenze} gli esempi più importanti:
    gli anelli $\mbZ_p$ con $p$ primo.
\end{spiegazione}


  % k-parti, somma, simmetria, formula ricorsiva, triangolo di Tartaglia, formula chiusa, teorema binomiale
  \chapter{Coefficienti binomiali}\label{cap:binomiali}

Quante sono le parti con esattamente $k$ elementi di un insieme con $n$ elementi? Ad esempio, in quanti modi si scelgono
$3$ persone su $10$ per formare una commissione? La risposta è un numero che compare ovunque in matematica: il
\textbf{coefficiente binomiale} $\binom{n}{k}$.

% ──────────────────────────────────────────────────────────────
\section{Definizione e prime proprietà}\label{sec:def-binomiali}

\dfn[binomiale]{$k$-parti e coefficiente binomiale}{
    Sia $s$ un insieme e $k \in \mbN$. Indichiamo con $\mcP_k(s) := \{t \in \mcP(s) \mid |t| = k\}$ l'insieme delle parti di $s$ con
    esattamente $k$ elementi (le \textbf{$k$-parti} di $s$). Per $n, k \in \mbN$ il \textbf{coefficiente binomiale} «$n$ su $k$» è
    \[ \binom{n}{k} := |\mcP_k(I_n)|, \]
    il numero di parti con $k$ elementi di $\{1, \dots, n\}$.
}

\begin{esempio}{Primi valori}
    Per ogni $n$: $\binom{n}{0} = 1$ (l'unica $0$-parte è $\varnothing$), $\binom{n}{n} = 1$ (l'unica $n$-parte di $I_n$ è $I_n$), $\binom{n}{1} = n$
    (i singleton). Se $k > n$, $\binom{n}{k} = 0$, perché una parte di $I_n$ ha al più $n$ elementi. Le $2$-parti di $I_4$ sono
    $\{1, 2\}, \{1, 3\}, \{1, 4\}, \{2, 3\}, \{2, 4\}, \{3, 4\}$, quindi $\binom{4}{2} = 6$.
\end{esempio}

La definizione usa $I_n$, ma vale per qualunque insieme con $n$ elementi: contano solo quanti sono, non come si chiamano.

\begin{concetto}
\begin{Proposizione}{Le $k$-parti di un insieme qualunque}{k-parti}
    \begin{prerequisiti}
        \dip{Definizione}{def:binomiale}
        \dip{Definizione}{def:immagine-antiimmagine}
        \dip{Proposizione}{prop:ridotta}
    \end{prerequisiti}
    Se $|s| = n$, allora $|\mcP_k(s)| = \binom{n}{k}$ per ogni $k \in \mbN$.
\end{Proposizione}
\begin{dimostrazione}
    Sia $f \colon s \to I_n$ biettiva. Per ogni $t \subseteq s$, la restrizione di $f$ a $t$ è una biezione tra $t$ e $\overrightarrow{f}(t)$, quindi
    $|\overrightarrow{f}(t)| = |t|$. Allora $t \mapsto \overrightarrow{f}(t)$ manda $\mcP_k(s)$ in $\mcP_k(I_n)$, ed è biettiva: la sua inversa è
    $u \mapsto \overrightarrow{f^{-1}}(u)$. Quindi $|\mcP_k(s)| = |\mcP_k(I_n)| = \binom{n}{k}$.
\end{dimostrazione}
\end{concetto}

\begin{spiegazione}
    Una biezione è un cambio di nomi: rinominando gli elementi di $s$ come $1, \dots, n$, le $k$-parti di $s$ diventano le $k$-parti di
    $I_n$, una per una. Ad esempio le $2$-parti di $\{a, b, c, d\}$ sono $6$ come quelle di $I_4$.
\end{spiegazione}

\begin{concetto}
\begin{Teorema}{La somma dei coefficienti binomiali}{somma-binomiali}
    \begin{prerequisiti}
        \dip{Definizione}{def:binomiale}
        \dip{Definizione}{def:partizione}
        \dip{Lemma}{lem:somma-prodotto}
        \dip{Corollario}{cor:cardinalita-parti}
    \end{prerequisiti}
    Per ogni $n \in \mbN$: \quad $\displaystyle\sum_{k = 0}^{n} \binom{n}{k} = \binom{n}{0} + \binom{n}{1} + \dots + \binom{n}{n} = 2^n$.
\end{Teorema}
\begin{dimostrazione}
    Ogni parte di $I_n$ ha un numero di elementi $k$ compreso tra $0$ e $n$, e uno solo; inoltre $\mcP_k(I_n) \neq \varnothing$ per $0 \leq k \leq n$.
    Quindi $\{\mcP_k(I_n) \mid 0 \leq k \leq n\}$ è una partizione di $\mcP(I_n)$ in $n + 1$ blocchi, e per il Lemma~\ref{lem:somma-prodotto}
    \[ \sum_{k = 0}^{n} |\mcP_k(I_n)| = |\mcP(I_n)| = 2^n \]
    (Corollario~\ref{cor:cardinalita-parti}).
\end{dimostrazione}
\end{concetto}

\begin{spiegazione}
    Si contano le parti di $I_n$ \textbf{raggruppandole per numero di elementi}. Per $n = 3$ i gruppi sono le righe del diagramma
    di Hasse di $\mcP(\{a, b, c\})$:
    \begin{center}
    \begin{tikzpicture}[hasse]
        \node[hv] (v1) at (0,0) {$\varnothing$};
        \node[hv evid] (a) at (-1.5,1.1) {$\{a\}$}; \node[hv evid] (b) at (0,1.1) {$\{b\}$}; \node[hv evid] (c) at (1.5,1.1) {$\{c\}$};
        \node[hv evid2] (ab) at (-1.5,2.2) {$\{a, b\}$}; \node[hv evid2] (ac) at (0,2.2) {$\{a, c\}$}; \node[hv evid2] (bc) at (1.5,2.2) {$\{b, c\}$};
        \node[hv] (abc) at (0,3.3) {$\{a, b, c\}$};
        \draw[hlato] (v1) -- (a); \draw[hlato] (v1) -- (b); \draw[hlato] (v1) -- (c);
        \draw[hlato] (a) -- (ab); \draw[hlato] (a) -- (ac); \draw[hlato] (b) -- (ab); \draw[hlato] (b) -- (bc);
        \draw[hlato] (c) -- (ac); \draw[hlato] (c) -- (bc);
        \draw[hlato] (ab) -- (abc); \draw[hlato] (ac) -- (abc); \draw[hlato] (bc) -- (abc);
        \foreach \y/\k/\v in {0/0/1, 1.1/1/3, 2.2/2/3, 3.3/3/1}
            \node[font=\small, anchor=west] at (2.6,\y) {$\binom{3}{\k} = \v$};
        \node[font=\small, anchor=west] at (4.6,1.65) {$1 + 3 + 3 + 1 = 8 = 2^3$};
    \end{tikzpicture}
    \end{center}
\end{spiegazione}

\begin{concetto}
\begin{Teorema}{Simmetria dei coefficienti binomiali}{simmetria-binomiali}
    \begin{prerequisiti}
        \dip{Definizione}{def:binomiale}
        \dip{Teorema}{teo:composte-caratteristiche}
        \dip{Teorema}{teo:inclusione-esclusione}
        \dip{Proposizione}{prop:doppia-neg-ins}
    \end{prerequisiti}
    Per ogni $n, k \in \mbN$ con $k \leq n$: \quad $\displaystyle\binom{n}{k} = \binom{n}{n - k}$.
\end{Teorema}
\begin{dimostrazione}
    La funzione $x \in \mcP(I_n) \mapsto I_n \setminus x \in \mcP(I_n)$ è biettiva: applicata due volte dà $I_n \setminus (I_n \setminus x) = x$ (doppia
    negazione insiemistica), quindi è l'inversa di sé stessa. Se $|x| = k$, allora $|I_n \setminus x| = n - k$, perché $I_n$ è unione disgiunta di
    $x$ e $I_n \setminus x$. Quindi la funzione manda $\mcP_k(I_n)$ in $\mcP_{n - k}(I_n)$, e la sua restrizione è una biezione tra i due
    (con inversa la restrizione $\mcP_{n - k}(I_n) \to \mcP_k(I_n)$ della stessa funzione). Dunque $|\mcP_k(I_n)| = |\mcP_{n - k}(I_n)|$.
\end{dimostrazione}
\end{concetto}

\begin{spiegazione}
    Scegliere i $k$ elementi \textbf{da prendere} equivale a scegliere gli $n - k$ elementi \textbf{da lasciare}. Scegliere $3$ persone su
    $10$ per una commissione è come scegliere le $7$ che restano fuori: $\binom{10}{3} = \binom{10}{7}$. Nel diagramma di Hasse qui sopra,
    il complementare scambia la riga $k$ con la riga $n - k$: è il diagramma capovolto.
\end{spiegazione}

% ──────────────────────────────────────────────────────────────
\section{La formula ricorsiva e il triangolo di Tartaglia}\label{sec:tartaglia}

\begin{concetto}
\begin{Teorema}{Formula ricorsiva dei coefficienti binomiali}{stifel}
    \begin{prerequisiti}
        \dip{Definizione}{def:binomiale}
        \dip{Lemma}{lem:somma-prodotto}
        \dip{Lemma}{lem:togli-elemento}
    \end{prerequisiti}
    Per ogni $n, k \in \mbN$ con $k \leq n$: \quad $\displaystyle\binom{n + 1}{k + 1} = \binom{n}{k + 1} + \binom{n}{k}$.
\end{Teorema}
\begin{dimostrazione}
    Dividiamo le $(k + 1)$-parti di $I_{n + 1}$ in due gruppi, secondo se contengono o no l'ultimo elemento $n + 1$:
    \[ a = \{x \in \mcP_{k + 1}(I_{n + 1}) \mid n + 1 \notin x\}, \qquad b = \{x \in \mcP_{k + 1}(I_{n + 1}) \mid n + 1 \in x\}. \]
    $a$ e $b$ sono disgiunti e la loro unione è $\mcP_{k + 1}(I_{n + 1})$, quindi $\binom{n + 1}{k + 1} = |a| + |b|$ (Lemma~\ref{lem:somma-prodotto}).
    \begin{caso}{$|a| = \binom{n}{k + 1}$}
        Le parti di $I_{n + 1}$ che non contengono $n + 1$ sono le parti di $I_n$: quindi $a = \mcP_{k + 1}(I_n)$.
    \end{caso}
    \begin{caso}{$|b| = \binom{n}{k}$}
        $x \in b \mapsto x \setminus \{n + 1\} \in \mcP_k(I_n)$ è biettiva (Lemma~\ref{lem:togli-elemento}): la sua inversa aggiunge $n + 1$, cioè
        $y \mapsto y \cup \{n + 1\}$.
    \end{caso}
\end{dimostrazione}
\end{concetto}

\begin{spiegazione}
    Per scegliere $k + 1$ persone tra $n + 1$, guardiamo una persona particolare, l'ultima: o \textbf{non} viene scelta, e allora le
    $k + 1$ si scelgono tra le altre $n$; oppure viene scelta, e allora ne mancano $k$ da scegliere tra le altre $n$. Ad esempio
    $\binom{4}{2} = \binom{3}{2} + \binom{3}{1} = 3 + 3 = 6$: le $2$-parti di $I_4$ senza il $4$ sono $\{1,2\}, \{1,3\}, \{2,3\}$; quelle con il $4$ sono
    $\{1,4\}, \{2,4\}, \{3,4\}$.
\end{spiegazione}

La formula ricorsiva permette di calcolare tutti i coefficienti binomiali uno dopo l'altro, disponendoli nel
\textbf{triangolo di Tartaglia} (o di Pascal): la riga $n$ contiene $\binom{n}{0}, \binom{n}{1}, \dots, \binom{n}{n}$, i bordi valgono $1$ e
\textbf{ogni numero interno è la somma dei due numeri subito sopra di lui}.

\begin{concetto}
\begin{center}
\begin{tikzpicture}[x=0.62cm, y=0.62cm]
    \foreach \n in {0,...,6} {
        \foreach \k in {0,...,\n} {
            \pgfmathtruncatemacro{\v}{round(factorial(\n)/(factorial(\k)*factorial(\n-\k)))}
            \node[circle, minimum size=7mm, inner sep=0pt, font=\small, fill=primary!7, draw=primary!40] (t\n-\k) at (2*\k-\n, -1.6*\n) {$\v$};
        }
        \node[font=\small, anchor=east] at (-8.2, -1.6*\n) {$n = \n$};
    }
    \node[circle, minimum size=7mm, inner sep=0pt, font=\small, fill=orange!30, draw=orange!80!black] at (t3-1) {$3$};
    \node[circle, minimum size=7mm, inner sep=0pt, font=\small, fill=orange!30, draw=orange!80!black] at (t3-2) {$3$};
    \node[circle, minimum size=7mm, inner sep=0pt, font=\small, fill=seagreen!30, draw=seagreen!70!black] at (t4-2) {$6$};
    \draw[-{Stealth}, seagreen!70!black, line width=0.8pt] (t3-1) -- (t4-2);
    \draw[-{Stealth}, seagreen!70!black, line width=0.8pt] (t3-2) -- (t4-2);
    \node[font=\small, anchor=west] at (7.4,-6.4) {$\binom{4}{2} = \binom{3}{1} + \binom{3}{2} = 3 + 3$};
\end{tikzpicture}
\end{center}
Nella riga $n$ si leggono anche la simmetria $\binom{n}{k} = \binom{n}{n - k}$ (la riga è palindroma) e la somma $2^n$
(ad esempio $1 + 4 + 6 + 4 + 1 = 16 = 2^4$).
\end{concetto}

% ──────────────────────────────────────────────────────────────
\section{La formula chiusa}\label{sec:formula-chiusa}

\begin{concetto}
\begin{Teorema}{Formula chiusa dei coefficienti binomiali}{formula-binomiali}
    \begin{prerequisiti}
        \dip{Teorema}{teo:stifel}
        \dip{Definizione}{def:fattoriale}
        \dip{Corollario}{cor:induzione-uso}
    \end{prerequisiti}
    Per ogni $n, k \in \mbN$ con $k \leq n$: \quad $\displaystyle\binom{n}{k} = \frac{n!}{k!\,(n - k)!}$.
\end{Teorema}
\begin{dimostrazione}
    Per induzione su $n$, dimostrando la formula per tutti i $k \leq n$ contemporaneamente.
    \textbf{Base} $n = 0$: $\binom{0}{0} = 1 = \frac{0!}{0!\,0!}$. \textbf{Passo}: supponiamo la formula vera per $n$ e ogni $k \leq n$, e sia $k \leq n + 1$.
    \begin{caso}{$k = 0$ oppure $k = n + 1$}
        $\binom{n + 1}{0} = 1 = \frac{(n + 1)!}{0!\,(n + 1)!}$ e $\binom{n + 1}{n + 1} = 1 = \frac{(n + 1)!}{(n + 1)!\,0!}$.
    \end{caso}
    \begin{caso}{$1 \leq k \leq n$}
        Scriviamo $k = j + 1$ con $0 \leq j \leq n - 1$. Per la formula ricorsiva e l'ipotesi induttiva (applicata a $j + 1 \leq n$ e a $j \leq n$):
        \[ \binom{n + 1}{j + 1} = \binom{n}{j + 1} + \binom{n}{j} = \frac{n!}{(j + 1)!\,(n - j - 1)!} + \frac{n!}{j!\,(n - j)!}
           = \frac{n!\,\big((n - j) + (j + 1)\big)}{(j + 1)!\,(n - j)!} = \frac{(n + 1)!}{(j + 1)!\,(n - j)!}, \]
        dove abbiamo portato le due frazioni allo stesso denominatore $(j + 1)!\,(n - j)!$, moltiplicando la prima per $n - j$ e la
        seconda per $j + 1$. Poiché $j + 1 = k$ e $n - j = (n + 1) - k$, è la formula per $n + 1$ e $k$.
    \end{caso}
\end{dimostrazione}
\end{concetto}

\begin{spiegazione}
    Ecco un altro modo di capire la formula, che usa le funzioni iniettive (Teorema~\ref{teo:conta-iniettive}). Le funzioni iniettive
    da $I_k$ in $I_n$ sono $\frac{n!}{(n - k)!}$: sono i modi di \textbf{mettere in fila} $k$ elementi scelti tra $n$. Ma mettere in fila $k$
    elementi si fa in due passi:
    \begin{passi}
        \item si sceglie \textbf{quali} elementi: una $k$-parte, in $\binom{n}{k}$ modi;
        \item si sceglie \textbf{in che ordine} metterli: una permutazione dei $k$ elementi, in $k!$ modi (Corollario~\ref{cor:conta-biettive}).
    \end{passi}
    Quindi $\binom{n}{k} \cdot k! = \frac{n!}{(n - k)!}$, cioè $\binom{n}{k} = \frac{n!}{k!\,(n - k)!}$. Ad esempio, i podi (oro, argento, bronzo) con
    $10$ atleti sono $10 \cdot 9 \cdot 8 = 720$; ogni terna di atleti compare in $3! = 6$ podi diversi, quindi le terne sono $\binom{10}{3} = 720 / 6 = 120$.
\end{spiegazione}

\begin{esempio}{Contare sottoinsiemi}
    Sia $a = \{0, 1, \dots, 9\}$, con $10$ elementi.
    \begin{itemize}
        \item Le parti di $a$ con $3$ elementi sono $\binom{10}{3} = \frac{10!}{3!\,7!} = \frac{10 \cdot 9 \cdot 8}{6} = 120$; quelle con $11$ elementi sono $0$.
        \item Le parti di $a$ con $7$ elementi sono $\binom{10}{7} = \binom{10}{3} = 120$ (simmetria).
        \item Le parti formate da $0$ e da altri tre elementi distinti si ottengono scegliendo $3$ elementi tra i $9$ diversi da $0$:
              $\binom{9}{3} = \frac{9 \cdot 8 \cdot 7}{6} = 84$.
    \end{itemize}
\end{esempio}

% ──────────────────────────────────────────────────────────────
\section{Il teorema binomiale}\label{sec:newton}

I coefficienti binomiali devono il loro nome alla formula per la potenza di un \textbf{binomio} $a + b$.

\begin{concetto}
\begin{Teorema}{Teorema binomiale (formula di Newton)}{newton}
    \begin{prerequisiti}
        \dip{Teorema}{teo:stifel}
        \dip{Teorema}{teo:regole-potenze}
        \dip{Proposizione}{prop:multipli-unita}
        \dip{Corollario}{cor:induzione-uso}
    \end{prerequisiti}
    Sia $(s, +, \cdot)$ un anello unitario e siano $a, b \in s$ con $ab = ba$. Per ogni $n \in \mbN$:
    \[ (a + b)^n = \sum_{k = 0}^{n} \binom{n}{k} a^k b^{n - k} = b^n + \binom{n}{1} a b^{n - 1} + \dots + \binom{n}{n - 1} a^{n - 1} b + a^n. \]
\end{Teorema}
\begin{dimostrazione}
    \facoltativa
    Per induzione su $n$. \textbf{Base}: $(a + b)^0 = 1 = \binom{0}{0} a^0 b^0$. \textbf{Passo}: usando l'ipotesi induttiva, la distributività e
    $ab = ba$ (che permette di spostare le $a$ prima delle $b$),
    \[ (a + b)^{n + 1} = (a + b)\sum_{k = 0}^{n} \binom{n}{k} a^k b^{n - k} = \sum_{k = 0}^{n} \binom{n}{k} a^{k + 1} b^{n - k} + \sum_{k = 0}^{n} \binom{n}{k} a^k b^{n + 1 - k}. \]
    Nella prima somma il termine $a^j b^{n + 1 - j}$ compare con coefficiente $\binom{n}{j - 1}$, nella seconda con $\binom{n}{j}$: per la formula
    ricorsiva il coefficiente totale è $\binom{n}{j - 1} + \binom{n}{j} = \binom{n + 1}{j}$ (per $j = 0$ e $j = n + 1$ c'è un solo termine, con coefficiente
    $1 = \binom{n + 1}{0} = \binom{n + 1}{n + 1}$).
\end{dimostrazione}
\end{concetto}

\begin{spiegazione}
    Sviluppando $(a + b)^3 = (a + b)(a + b)(a + b)$ si sceglie, da ogni fattore, $a$ oppure $b$, e si moltiplicano le scelte. Il
    termine $a^2 b$ nasce da tutte le scelte con \textbf{due} $a$ e una $b$: tante quante le $2$-parti dei tre fattori, cioè $\binom{3}{2} = 3$
    ($aab$, $aba$, $baa$). Quindi
    \[ (a + b)^3 = b^3 + 3ab^2 + 3a^2 b + a^3, \]
    con i coefficienti $1, 3, 3, 1$ della riga $3$ del triangolo di Tartaglia. L'ipotesi $ab = ba$ serve a riordinare $aba$ e $baa$ in
    $a^2 b$: in un anello non commutativo la formula è falsa.
\end{spiegazione}


  % Parte 6: Aritmetica (divisibilità, congruenze, equazioni)

  \part{Aritmetica}

  % Divisori, MCD, divisione euclidea, Euclide, Bézout, coprimi, primi, teorema fondamentale, sottogruppi di Z
  \chapter{Divisibilità nei numeri interi}\label{cap:divisibilita}

In questa parte studiamo l'aritmetica dei numeri interi con gli strumenti costruiti finora. In questo capitolo: divisori e
multipli, la divisione con resto, il massimo comun divisore e l'algoritmo per calcolarlo, il teorema di Bézout, i numeri
primi e il \textbf{teorema fondamentale dell'aritmetica}, che dice che ogni intero si scompone in modo essenzialmente unico in
fattori primi.

% ──────────────────────────────────────────────────────────────
\section{Divisori, multipli, elementi associati}\label{sec:divisori}

\dfn[divide]{Divisibilità}{
    Sia $s = \mbN$ oppure $s = \mbZ$ e siano $a, b \in s$. Diciamo che \textbf{$a$ divide $b$} (o che $a$ è un \textbf{divisore} di $b$, o che
    $b$ è un \textbf{multiplo} di $a$) e scriviamo $a \mid b$ se $\exists k \in s\,(b = ak)$. Poniamo
    \[ \mathrm{Div}(a) := \{b \in s \mid b \mid a\}, \qquad \mathrm{Mult}(a) := \{b \in s \mid a \mid b\}. \]
    Il \textbf{valore assoluto} di $n \in \mbZ$ è $|n| := \max\{n, -n\}$. I divisori $\pm 1$ e $\pm n$ di un intero $n$ si dicono \textbf{divisori banali} di $n$.
}

\dfn[associati]{Elementi associati}{
    Due interi $a$ e $b$ si dicono \textbf{associati} se $a \mid b$ e $b \mid a$.
}

\begin{concetto}
\begin{Proposizione}{Proprietà della divisibilità in $\mbZ$}{prop-divisibilita}
    \begin{prerequisiti}
        \dip{Definizione}{def:divide}
        \dip{Definizione}{def:associati}
        \dip{Definizione}{def:unita}
    \end{prerequisiti}
    Siano $a, b, c \in \mbZ$.
    \begin{enumerate}
        \item $a \mid a$; se $a \mid b$ e $b \mid c$, allora $a \mid c$; inoltre $1 \mid a$ e $a \mid 0$.
        \item Se $a \mid b$ e $a \mid c$, allora $a \mid bu + cv$ per ogni $u, v \in \mbZ$: un divisore comune divide ogni \textbf{combinazione}.
        \item $a$ e $b$ sono associati se e solo se $b = a$ oppure $b = -a$.
        \item $a$ e $b$ sono associati se e solo se $\mathrm{Div}(a) = \mathrm{Div}(b)$, se e solo se $\mathrm{Mult}(a) = \mathrm{Mult}(b)$.
    \end{enumerate}
\end{Proposizione}
\begin{dimostrazione}
    \begin{caso}{Punti 1 e 2}
        $a = a \cdot 1$; se $b = ah$ e $c = bk$ allora $c = a(hk)$; $a = 1 \cdot a$; $0 = a \cdot 0$. Se $b = ah$ e $c = ak$, allora
        $bu + cv = a(hu + kv)$.
    \end{caso}
    \begin{caso}{Punto 3}
        $(\leftarrow)$ $a = (-a)(-1)$ e $-a = a(-1)$. $(\rightarrow)$ Siano $b = ah$ e $a = bk$. Se $a = 0$, allora $b = 0 = a$. Se $a \neq 0$, da
        $a = a(hk)$ e dalla cancellabilità di $a$ in $(\mbZ, \cdot)$ si ha $hk = 1$: $h$ è invertibile in $\mbZ$, cioè $h = \pm 1$, e $b = \pm a$.
    \end{caso}
    \begin{caso}{Punto 4}
        $(\rightarrow)$ Se $a \mid b$ e $b \mid a$, per la transitività un divisore di $a$ divide $b$ e viceversa; allo stesso modo per i multipli.
        $(\leftarrow)$ Se $\mathrm{Div}(a) = \mathrm{Div}(b)$, da $a \in \mathrm{Div}(a)$ segue $a \mid b$ e da $b \in \mathrm{Div}(b)$ segue $b \mid a$. Se
        $\mathrm{Mult}(a) = \mathrm{Mult}(b)$, da $b \in \mathrm{Mult}(b) = \mathrm{Mult}(a)$ segue $a \mid b$, e simmetricamente $b \mid a$.
    \end{caso}
\end{dimostrazione}
\end{concetto}

\begin{spiegazione}
    In $\mbZ$ la divisibilità non distingue un numero dal suo opposto: $6$ e $-6$ hanno esattamente gli stessi divisori
    ($\pm 1, \pm 2, \pm 3, \pm 6$) e gli stessi multipli. Per questo $\mid$ su $\mbZ$ non è antisimmetrica (lo abbiamo notato
    negli esempi della Sezione~\ref{sec:proprieta-relazioni}) e, quando cercheremo «il» massimo comun divisore, ne troveremo due, uno opposto dell'altro. Il punto 2 è quello
    che useremo più spesso: ad esempio, siccome $3 \mid 12$ e $3 \mid 21$, allora $3$ divide anche $12 \cdot 5 - 21 \cdot 2 = 18$.
\end{spiegazione}

\dfn[mcd]{Massimo comun divisore, minimo comune multiplo}{
    Sia $s = \mbN$ oppure $\mbZ$ e siano $a, b \in s$. Definiamo
    \[ \mathrm{MCD}(a, b) := \{d \in \mathrm{Div}(a) \cap \mathrm{Div}(b) \mid \forall z \in \mathrm{Div}(a) \cap \mathrm{Div}(b)\,(z \mid d)\}, \]
    \[ \mathrm{mcm}(a, b) := \{d \in \mathrm{Mult}(a) \cap \mathrm{Mult}(b) \mid \forall z \in \mathrm{Mult}(a) \cap \mathrm{Mult}(b)\,(d \mid z)\}: \]
    un massimo comun divisore è un divisore comune \textbf{diviso da tutti} i divisori comuni; un minimo comune multiplo è un
    multiplo comune che \textbf{divide tutti} i multipli comuni.
}

\begin{concetto}
\begin{Proposizione}{I massimi comuni divisori sono associati}{mcd-associati}
    \begin{prerequisiti}
        \dip{Definizione}{def:mcd}
        \dip{Proposizione}{prop:prop-divisibilita}
    \end{prerequisiti}
    Siano $a, b \in \mbZ$. Due elementi di $\mathrm{MCD}(a, b)$ sono associati, e un associato di un elemento di $\mathrm{MCD}(a, b)$ sta in
    $\mathrm{MCD}(a, b)$. Quindi, se $d \in \mathrm{MCD}(a, b)$, allora $\mathrm{MCD}(a, b) = \{d, -d\}$. Lo stesso vale per $\mathrm{mcm}(a, b)$.
\end{Proposizione}
\begin{dimostrazione}
    Se $d, d' \in \mathrm{MCD}(a, b)$, sono entrambi divisori comuni; $d'$ divide $d$ perché $d$ è diviso da tutti i divisori comuni, e
    simmetricamente $d \mid d'$. Se invece $d \in \mathrm{MCD}(a, b)$ e $d'$ è associato a $d$, $d'$ ha gli stessi divisori e multipli di $d$
    (Proposizione~\ref{prop:prop-divisibilita}, punto 4): è ancora un divisore comune ed è diviso da tutti i divisori comuni. L'ultima
    affermazione segue dal punto 3 della stessa proposizione. Per $\mathrm{mcm}$ il ragionamento è identico.
\end{dimostrazione}
\end{concetto}

\begin{spiegazione}
    «Massimo» va inteso rispetto alla divisibilità, non rispetto a $\leq$: si guarda il diagramma di Hasse di $(\mbN, \mid)$.
    Per $a = 12$ e $b = 18$, in $\mbN$ i divisori comuni sono $1, 2, 3, 6$ e tutti dividono $6$: quindi $\mathrm{MCD}_{\mbN}(12, 18) = \{6\}$, e in
    $\mbZ$ $\mathrm{MCD}(12, 18) = \{6, -6\}$.
    \begin{center}
    \begin{tikzpicture}[hasse]
        \node[hv evid2] (1) at (0,0) {$1$}; \node[hv evid2] (2) at (-0.9,1.0) {$2$}; \node[hv evid2] (3) at (0.9,1.0) {$3$};
        \node[hv] (4) at (-1.8,2.0) {$4$}; \node[hv evid2] (6) at (0,2.0) {$6$}; \node[hv] (9) at (1.8,2.0) {$9$};
        \node[hv evid] (12) at (-0.9,3.0) {$12$}; \node[hv evid] (18) at (0.9,3.0) {$18$};
        \draw[hlato] (1) -- (2) -- (4) -- (12); \draw[hlato] (1) -- (3) -- (6) -- (12);
        \draw[hlato] (2) -- (6) -- (18); \draw[hlato] (3) -- (9) -- (18);
        \node[font=\small, align=left, anchor=west] at (2.8,1.5)
            {divisori di $12$ e di $18$ in $(\mbN, \mid)$;\\ in verde i divisori comuni:\\ il «più alto» è $6$, e tutti lo dividono.};
    \end{tikzpicture}
    \end{center}
    Un caso curioso: ogni intero divide $0$, quindi $\mathrm{Div}(0) = \mbZ$ e $\mathrm{MCD}(0, 0) = \{0\}$: in $(\mbN, \mid)$ lo $0$ è il massimo.
\end{spiegazione}

% ──────────────────────────────────────────────────────────────
\section{La divisione euclidea}\label{sec:divisione-euclidea}

\begin{concetto}
\begin{Teorema}{Teorema della divisione euclidea (con resto)}{divisione-euclidea}
    \begin{prerequisiti}
        \dip{Definizione}{def:divide}
        \dip{Teorema}{teo:induzione-seconda}
    \end{prerequisiti}
    Per ogni $n, m \in \mbZ$ con $m \neq 0$ esiste una e una sola coppia $(q, r) \in \mbZ \times \mbN$ tale che
    \[ n = mq + r \qquad\text{e}\qquad 0 \leq r < |m|. \]
    $q$ si chiama \textbf{quoziente} e $r$ \textbf{resto} della divisione di $n$ per $m$.
\end{Teorema}
\begin{dimostrazione}
    \begin{caso}{Esistenza per $n \geq 0$, per induzione di seconda forma su $n$}
        Se $n < |m|$ (in particolare se $n = 0$), basta $q = 0$, $r = n$. Se $n \geq |m|$, allora $0 \leq n - |m| < n$ e per ipotesi induttiva
        $n - |m| = mq_1 + r_1$ con $0 \leq r_1 < |m|$. Quindi $n = mq_1 + |m| + r_1$: se $m > 0$, $n = m(q_1 + 1) + r_1$; se $m < 0$, $|m| = -m$ e
        $n = m(q_1 - 1) + r_1$.
    \end{caso}
    \begin{caso}{Esistenza per $n < 0$}
        Per il caso precedente $-n = mq' + r'$ con $0 \leq r' < |m|$. Se $r' = 0$, $n = m(-q')$. Se $r' > 0$, scriviamo
        $n = -mq' - |m| + (|m| - r')$, con $0 < |m| - r' < |m|$; poi $-mq' - |m| = m(-q' - 1)$ se $m > 0$ e $= m(-q' + 1)$ se $m < 0$.
    \end{caso}
    \begin{caso}{Unicità}
        Se $n = mq_1 + r_1 = mq_2 + r_2$ con $0 \leq r_1, r_2 < |m|$, allora $m(q_1 - q_2) = r_2 - r_1$ e quindi $|m|\,|q_1 - q_2| = |r_2 - r_1| < |m|$.
        Poiché $|m| \neq 0$, si ha $|q_1 - q_2| < 1$, cioè $q_1 = q_2$, e di conseguenza $r_1 = r_2$.
    \end{caso}
\end{dimostrazione}
\end{concetto}

\begin{spiegazione}
    Sulla retta dei numeri i multipli di $m$ formano una «scala» con gradini lunghi $|m|$: ogni intero $n$ sta su un gradino o tra due
    gradini consecutivi. $mq$ è il gradino \textbf{subito a sinistra} di $n$ (o $n$ stesso), e $r$ è quanto manca per arrivare a $n$.
    Per $n = -21$ e $m = 4$:
    \begin{center}
    \begin{tikzpicture}[x=0.55cm]
        \draw[-{Stealth}, black!50] (-27,0) -- (2,0);
        \foreach \k [evaluate=\k as \v using int(4*\k)] in {-6,...,0} { \draw[primary!80!black, line width=1pt] (4*\k,-0.15) -- (4*\k,0.15); \node[font=\scriptsize, primary!80!black] at (4*\k,0.45) {$\v$}; }
        \node[circle, fill=orange!80!black, inner sep=1.5pt] at (-21,0) {};
        \node[font=\scriptsize, orange!80!black] at (-21,-0.45) {$n = -21$};
        \draw[decorate, decoration={brace, amplitude=3pt, mirror}] (-24,-0.75) -- node[below=3pt, font=\scriptsize] {$r = 3$} (-21,-0.75);
        \node[font=\scriptsize, align=left, anchor=west] at (-15,-0.8) {$-21 = 4 \cdot (-6) + 3$: \ $q = -6$, $r = 3$\\ (il resto è sempre $\geq 0$!)};
    \end{tikzpicture}
    \end{center}
    Attenzione al segno: $-21 = 4 \cdot (-5) - 1$ non va bene, perché il resto deve essere tra $0$ e $|m| - 1$. Altri esempi:
    $21 = 4 \cdot 5 + 1$, \ $21 = (-4)(-5) + 1$, \ $-21 = (-4) \cdot 6 + 3$, \ $11 = 2 \cdot 5 + 1$, \ $10 = 5 \cdot 2 + 0$.
\end{spiegazione}

% ──────────────────────────────────────────────────────────────
\section{L'algoritmo delle divisioni successive}\label{sec:euclide}

Poiché $\mathrm{MCD}(a, b) = \{d, -d\}$, per trovarlo basta trovare $d \geq 0$; e poiché $a$ e $-a$ hanno gli stessi divisori, possiamo
supporre $a, b \in \mbN$. L'idea per calcolarlo è che dividere con resto \textbf{non cambia i divisori comuni}.

\begin{concetto}
\begin{Lemma}{La divisione non cambia i divisori comuni}{mcd-resto}
    \begin{prerequisiti}
        \dip{Proposizione}{prop:prop-divisibilita}
        \dip{Definizione}{def:mcd}
    \end{prerequisiti}
    Se $a = bq + r$ con $a, b, q, r \in \mbZ$, allora $\mathrm{Div}(a) \cap \mathrm{Div}(b) = \mathrm{Div}(b) \cap \mathrm{Div}(r)$, e quindi
    $\mathrm{MCD}(a, b) = \mathrm{MCD}(b, r)$. Inoltre $\mathrm{MCD}(a, 0) = \{a, -a\}$.
\end{Lemma}
\begin{dimostrazione}
    Se $z \mid a$ e $z \mid b$, allora $z \mid a - bq = r$; se $z \mid b$ e $z \mid r$, allora $z \mid bq + r = a$ (Proposizione~\ref{prop:prop-divisibilita},
    punto 2). Avendo gli stessi divisori comuni, le due coppie hanno gli stessi massimi comuni divisori. Infine, ogni intero divide $0$,
    quindi i divisori comuni di $a$ e $0$ sono i divisori di $a$, e quelli divisi da tutti sono $a$ e $-a$.
\end{dimostrazione}
\end{concetto}

\begin{spiegazione}
    Ad esempio $111 = 17 \cdot 6 + 9$: un numero che divide $111$ e $17$ divide anche $9 = 111 - 6 \cdot 17$, e uno che divide $17$ e $9$
    divide anche $111$. Quindi cercare l'MCD di $111$ e $17$ è come cercare quello di $17$ e $9$: numeri più piccoli.
\end{spiegazione}

\begin{concetto}
\begin{Teorema}{Algoritmo delle divisioni successive (di Euclide)}{euclide}
    \begin{prerequisiti}
        \dip{Teorema}{teo:divisione-euclidea}
        \dip{Lemma}{lem:mcd-resto}
        \dip{Definizione}{def:ben-ordinato}
    \end{prerequisiti}
    Siano $a, b \in \mbN$. Si ripete: se $b = 0$, ci si ferma; altrimenti si divide $a$ per $b$, $a = bq + r$, e si sostituisce la coppia $(a, b)$
    con $(b, r)$. Il procedimento termina, e quando si ferma, con la coppia $(a', 0)$, si ha $a' \in \mathrm{MCD}(a, b)$.
\end{Teorema}
\begin{dimostrazione}
    Posto $r_0 = b$, l'algoritmo produce i resti $r_1 > r_2 > \dots \geq 0$: ogni resto è minore del divisore usato, che è il resto precedente.
    Una successione strettamente decrescente di numeri naturali non può essere infinita (l'insieme dei suoi termini avrebbe un minimo, ma
    il termine successivo sarebbe più piccolo), quindi a un certo passo il resto è $0$. Per il Lemma~\ref{lem:mcd-resto}, ad ogni passo
    $\mathrm{MCD}(a, b)$ resta lo stesso, e alla fine $\mathrm{MCD}(a, b) = \mathrm{MCD}(a', 0) = \{a', -a'\}$.
\end{dimostrazione}
\end{concetto}

\begin{spiegazione}
    Ecco l'algoritmo come diagramma di flusso (per $a, b \in \mbN$ non entrambi nulli) e un esempio passo passo con $a = 111$, $b = 17$:
    \begin{center}
    \begin{tikzpicture}[node distance=0.9cm, every node/.style={font=\small},
        blocco/.style={rectangle, draw=primary!70!black, fill=primary!7, rounded corners=2pt, minimum height=6mm, inner xsep=6pt},
        scelta/.style={diamond, aspect=2.2, draw=orange!80!black, fill=orange!12, inner sep=1pt},
        estremo/.style={ellipse, draw=black!60, fill=black!5, minimum height=6mm},
        flusso/.style={-{Stealth[length=2mm]}, line width=0.7pt, draw=black!70}]
        \node[estremo] (inizio) at (0,0) {Inizio};
        \node[scelta] (a0) at (0,-1.2) {$a = 0$?};
        \node[scelta] (b0) at (0,-2.7) {$b = 0$?};
        \node[blocco] (de) at (0,-4.1) {$(q, r) :=$ divisione di $a$ per $b$};
        \node[blocco] (sc) at (0,-5.2) {$(a, b) := (b, r)$};
        \node[blocco] (ra) at (3.6,-1.2) {$\mathrm{MCD} := \{b\}$};
        \node[blocco] (rb) at (3.6,-2.7) {$\mathrm{MCD} := \{a\}$};
        \node[estremo] (fine) at (5.9,-1.95) {Fine};
        \draw[flusso] (inizio) -- (a0); \draw[flusso] (a0) -- node[right] {no} (b0); \draw[flusso] (b0) -- node[right] {no} (de);
        \draw[flusso] (de) -- (sc); \draw[flusso] (sc.west) -- ++(-1.9,0) |- (b0.west);
        \draw[flusso] (a0) -- node[above] {sì} (ra); \draw[flusso] (b0) -- node[above] {sì} (rb);
        \draw[flusso] (ra) -| (fine); \draw[flusso] (rb) -| (fine);
        \node[anchor=west, align=left] at (7.0,-2.6) {
            \renewcommand{\arraystretch}{1.15}
            \begin{tabular}{c|c|l}
                passo & $(a, b)$ & divisione \\ \hline
                $1$ & $(111, 17)$ & $111 = 17 \cdot 6 + 9$ \\
                $2$ & $(17, 9)$ & $17 = 9 \cdot 1 + 8$ \\
                $3$ & $(9, 8)$ & $9 = 8 \cdot 1 + 1$ \\
                $4$ & $(8, 1)$ & $8 = 1 \cdot 8 + 0$ \\
                fine & $(1, 0)$ & $\mathrm{MCD} = \{1\}$ in $\mbN$ \\
            \end{tabular}};
    \end{tikzpicture}
    \end{center}
    I resti $9, 8, 1, 0$ scendono sempre, quindi l'algoritmo non può andare avanti all'infinito. L'ultimo resto non nullo, $1$, è il
    massimo comun divisore: $\mathrm{MCD}(111, 17) = \{1, -1\}$ in $\mbZ$.
\end{spiegazione}

% ──────────────────────────────────────────────────────────────
\section{Il teorema di Bézout}\label{sec:bezout}

\begin{concetto}
\begin{Teorema}{Teorema di Bézout}{bezout}
    \begin{prerequisiti}
        \dip{Teorema}{teo:euclide}
        \dip{Proposizione}{prop:mcd-associati}
        \dip{Teorema}{teo:induzione-seconda}
    \end{prerequisiti}
    Siano $a, b \in \mbZ$ non entrambi nulli e sia $d \in \mathrm{MCD}(a, b)$. Allora esistono $u, v \in \mbZ$ tali che $d = au + bv$.
\end{Teorema}
\begin{dimostrazione}
    Se vale per $d$ vale per $-d = a(-u) + b(-v)$; e poiché $a = \pm |a|$, $b = \pm |b|$, una combinazione di $|a|$ e $|b|$ è anche una
    combinazione di $a$ e $b$. Possiamo quindi supporre $a, b \in \mbN$ e $d \geq 0$. Se $b = 0$, $d = a = a \cdot 1 + b \cdot 0$. Altrimenti
    consideriamo i resti dell'algoritmo di Euclide, $r_{-1} = a$, $r_0 = b$, $r_1, \dots, r_{t - 1}$, $r_t = 0$, con $r_{i - 2} = r_{i - 1}q_i + r_i$ e
    $d = r_{t - 1}$. Dimostriamo di più: \textbf{ogni resto è una combinazione di $a$ e $b$}. Per assurdo sia $k$ il minimo indice per cui
    $r_k$ non lo è. Non è $k = -1$ né $k = 0$ ($a = a \cdot 1 + b \cdot 0$, $b = a \cdot 0 + b \cdot 1$). Allora $r_{k - 2} = au' + bv'$ e
    $r_{k - 1} = au'' + bv''$, e
    \[ r_k = r_{k - 2} - r_{k - 1}q_k = a(u' - u''q_k) + b(v' - v''q_k), \]
    che è una combinazione di $a$ e $b$: assurdo. In particolare $d = r_{t - 1}$ è una combinazione di $a$ e $b$.
\end{dimostrazione}
\end{concetto}

\begin{spiegazione}
    La dimostrazione è anche un \textbf{algoritmo} (algoritmo delle divisioni successive \textbf{esteso}): si scrive accanto a ogni resto
    come combinazione di $a$ e $b$, usando le due righe precedenti. Per $a = 111$, $b = 17$:
    \begin{center}
    \renewcommand{\arraystretch}{1.25}
    \begin{tabular}{l|c|cc|l}
        \rowcolor{primary!12} riga & resto & $u$ & $v$ & come si ottiene \\
        $r_{-1}$ & $111$ & $1$ & $0$ & $a$ \\
        $r_0$ & $17$ & $0$ & $1$ & $b$ \\
        $r_1$ & $9$ & $1$ & $-6$ & $r_{-1} - 6\,r_0$ \quad ($111 = 17 \cdot 6 + 9$) \\
        $r_2$ & $8$ & $-1$ & $7$ & $r_0 - 1\,r_1$ \quad ($17 = 9 \cdot 1 + 8$) \\
        $r_3$ & $1$ & $2$ & $-13$ & $r_1 - 1\,r_2$ \quad ($9 = 8 \cdot 1 + 1$) \\
    \end{tabular}
    \end{center}
    Ogni riga si ottiene sottraendo dalla penultima l'ultima moltiplicata per il quoziente, \textbf{anche nelle colonne $u$ e $v$}.
    Risultato: $1 = 111 \cdot 2 + 17 \cdot (-13)$. Verifica: $222 - 221 = 1$.
\end{spiegazione}

% ──────────────────────────────────────────────────────────────
\section{Numeri coprimi, lemma di Euclide, numeri primi}\label{sec:primi}

\dfn[coprimi]{Numeri coprimi}{
    Due interi $a$ e $b$ si dicono \textbf{coprimi} se $1 \in \mathrm{MCD}(a, b)$, cioè se i loro unici divisori comuni sono $1$ e $-1$.
}

\begin{concetto}
\begin{Teorema}{Lemma di Euclide}{lemma-euclide}
    \begin{prerequisiti}
        \dip{Definizione}{def:coprimi}
        \dip{Teorema}{teo:bezout}
        \dip{Proposizione}{prop:prop-divisibilita}
    \end{prerequisiti}
    Siano $a, b, c \in \mbZ$. Se $a$ e $b$ sono coprimi e $a \mid bc$, allora $a \mid c$.
\end{Teorema}
\begin{dimostrazione}
    Per ipotesi $bc = ah$ per qualche $h$. Per il teorema di Bézout esistono $u, v$ con $1 = au + bv$. Moltiplicando per $c$:
    \[ c = acu + bcv = acu + ahv = a(cu + hv), \]
    quindi $a \mid c$.
\end{dimostrazione}
\end{concetto}

\begin{spiegazione}
    Se $a$ divide un prodotto ma non ha niente in comune con il primo fattore, deve «stare tutto» nel secondo. Ad esempio $4 \mid 3 \cdot 8$
    e $\mathrm{MCD}(4, 3) = \{\pm 1\}$, quindi $4 \mid 8$. Senza l'ipotesi di coprimalità è falso: $6 \mid 4 \cdot 9 = 36$, ma $6 \nmid 4$ e $6 \nmid 9$
    ($6$ si «divide» tra i due fattori: il $2$ va nel $4$, il $3$ nel $9$).
\end{spiegazione}

\dfn[primo]{Numero primo, elemento irriducibile}{
    Un intero $p \notin \{-1, 0, 1\}$ si dice
    \begin{itemize}
        \item \textbf{primo} se $\forall a, b \in \mbZ\,(p \mid ab \rightarrow p \mid a \vee p \mid b)$;
        \item \textbf{irriducibile} se ha solo i divisori banali $\pm 1$, $\pm p$.
    \end{itemize}
}

\begin{concetto}
\begin{Teorema}{In $\mbZ$ primi e irriducibili coincidono}{primi-irriducibili}
    \begin{prerequisiti}
        \dip{Definizione}{def:primo}
        \dip{Teorema}{teo:lemma-euclide}
        \dip{Proposizione}{prop:prop-divisibilita}
    \end{prerequisiti}
    Sia $p \in \mbZ \setminus \{-1, 0, 1\}$. Allora $p$ è primo se e solo se è irriducibile.
\end{Teorema}
\begin{dimostrazione}
    \begin{caso}{$(\rightarrow)$ Primo $\Rightarrow$ irriducibile}
        Sia $a$ un divisore di $p$: $p = ab$. Allora $p \mid ab$ e, essendo primo, $p \mid a$ oppure $p \mid b$. Se $p \mid a$, $a = pk$ e
        $p = pkb$, da cui $kb = 1$ ($p \neq 0$ è cancellabile): $b = \pm 1$ e $a = \pm p$. Se $p \mid b$, allo stesso modo $a = \pm 1$.
        Quindi i divisori di $p$ sono solo quelli banali.
    \end{caso}
    \begin{caso}{$(\leftarrow)$ Irriducibile $\Rightarrow$ primo}
        Sia $p \mid ab$ e supponiamo $p \nmid a$. I divisori comuni di $p$ e $a$ sono divisori di $p$, quindi tra $\pm 1, \pm p$; ma $\pm p$ non
        dividono $a$. Quindi $p$ e $a$ sono coprimi e, per il lemma di Euclide, $p \mid b$.
    \end{caso}
\end{dimostrazione}
\end{concetto}

\begin{spiegazione}
    «Irriducibile» dice che $p$ non si spezza in fattori più piccoli; «primo» dice che quando divide un prodotto divide uno dei fattori.
    In $\mbZ$ sono la stessa cosa. Ad esempio $7$ è irriducibile, e infatti se $7 \mid ab$ allora $7$ divide $a$ o $b$; $6$ non è irriducibile
    ($6 = 2 \cdot 3$) e infatti non è primo: $6 \mid 2 \cdot 3$ ma $6 \nmid 2$ e $6 \nmid 3$. I primi positivi sono $2, 3, 5, 7, 11, 13, \dots$; anche $-2, -3, \dots$
    sono primi.
\end{spiegazione}

% ──────────────────────────────────────────────────────────────
\section{Il teorema fondamentale dell'aritmetica}\label{sec:tfa}

\begin{concetto}
\begin{Lemma}{Un primo che divide un prodotto divide un fattore}{primo-prodotto}
    \begin{prerequisiti}
        \dip{Definizione}{def:primo}
        \dip{Corollario}{cor:induzione-uso}
    \end{prerequisiti}
    Se $p$ è primo e $p \mid a_1 a_2 \cdots a_s$ con $s \geq 1$, allora $p \mid a_j$ per qualche $j \in \{1, \dots, s\}$.
\end{Lemma}
\begin{dimostrazione}
    Per induzione su $s$. Per $s = 1$ è ovvio. Se $p \mid (a_1 \cdots a_s) a_{s + 1}$, per definizione di primo $p \mid a_1 \cdots a_s$ (e si conclude con
    l'ipotesi induttiva) oppure $p \mid a_{s + 1}$.
\end{dimostrazione}
\end{concetto}

\begin{concetto}
\begin{Teorema}{Teorema fondamentale dell'aritmetica}{tfa}
    \begin{prerequisiti}
        \dip{Teorema}{teo:primi-irriducibili}
        \dip{Lemma}{lem:primo-prodotto}
        \dip{Teorema}{teo:induzione-seconda}
        \dip{Corollario}{cor:induzione-uso}
    \end{prerequisiti}
    Sia $m \in \mbZ \setminus \{-1, 0, 1\}$. Allora:
    \begin{enumerate}
        \item (\textbf{esistenza}) esistono primi $p_1, \dots, p_r$ tali che $m = p_1 p_2 \cdots p_r$;
        \item (\textbf{essenziale unicità}) se anche $m = q_1 \cdots q_s$ con i $q_j$ primi, allora $r = s$ ed esiste una funzione biettiva
              $\sigma \colon I_r \to I_r$ tale che $p_i = \pm q_{\sigma(i)}$ per ogni $i \in I_r$.
    \end{enumerate}
\end{Teorema}
\begin{dimostrazione}
    \begin{caso}{Esistenza per $m \geq 2$, per induzione di seconda forma}
        $2$ è primo, perché i suoi unici divisori positivi sono $1$ e $2$. Sia $m > 2$ e supponiamo la tesi per ogni $k$ con $2 \leq k < m$.
        Se $m$ è primo abbiamo finito ($r = 1$). Altrimenti $m$ non è irriducibile e ha un divisore non banale: $m = ab$ con $a, b \in \mbN$ e
        $1 < a, b < m$. Per ipotesi induttiva $a$ e $b$ sono prodotti di primi, quindi lo è anche $m = ab$.
    \end{caso}
    \begin{caso}{Esistenza per $m \leq -2$}
        $-m = p_1 \cdots p_r$, quindi $m = (-p_1) p_2 \cdots p_r$, e $-p_1$ è primo (ha gli stessi divisori di $p_1$).
    \end{caso}
    \begin{caso}{Essenziale unicità, per induzione su $r$}
        Sia $p_1 \cdots p_r = q_1 \cdots q_s$. Poiché $p_r$ divide il secondo membro, per il Lemma~\ref{lem:primo-prodotto} divide qualche $q_j$;
        rinumerando i $q$ possiamo supporre $p_r \mid q_s$. Poiché $q_s$ è irriducibile e $p_r \neq \pm 1$, si ha $q_s = \pm p_r$. Cancellando
        $p_r$ ($\neq 0$) si ottiene $p_1 \cdots p_{r - 1} = \pm q_1 \cdots q_{s - 1}$. Se $r = 1$, il primo membro è $1$, quindi anche $s = 1$ (nessun primo
        divide $\pm 1$). Se $r > 1$, allo stesso modo $s > 1$ (altrimenti $p_1$ dividerebbe $\pm 1$), e si applica l'ipotesi induttiva (assorbendo il segno in $q_1$): $r - 1 = s - 1$ e c'è una biezione
        $\tau$ di $I_{r - 1}$ con $p_i = \pm q_{\tau(i)}$; estendendo $\tau$ con $\sigma(r) = r$ (dopo la rinumerazione) si ottiene $\sigma$.
    \end{caso}
\end{dimostrazione}
\end{concetto}

\begin{spiegazione}
    \textbf{Esistenza}: si spezza il numero finché tutti i pezzi sono primi; l'induzione garantisce che lo spezzare finisca, perché i
    pezzi sono sempre più piccoli. Ecco l'«albero» di $360$:
    \begin{center}
    \begin{tikzpicture}[level distance=0.85cm, every node/.style={font=\small},
                        level 1/.style={sibling distance=2.6cm}, level 2/.style={sibling distance=1.3cm}, level 3/.style={sibling distance=0.8cm},
                        edge from parent/.style={draw=black!50}]
        \node {$360$}
            child {node {$12$} child {node {$4$} child {node[hv evid] {$2$}} child {node[hv evid] {$2$}}} child {node[hv evid] {$3$}}}
            child {node {$30$} child {node[hv evid] {$2$}} child {node {$15$} child {node[hv evid] {$3$}} child {node[hv evid] {$5$}}}};
        \node[anchor=west, font=\small] at (3.2,-1.4) {$360 = 2 \cdot 2 \cdot 3 \cdot 2 \cdot 3 \cdot 5 = 2^3 \cdot 3^2 \cdot 5$};
    \end{tikzpicture}
    \end{center}
    \textbf{Unicità}: spezzando in modo diverso ($360 = 10 \cdot 36$, ad esempio) si ottengono gli stessi primi, al più in un altro ordine e
    con segni diversi, come in $12 = 2 \cdot 2 \cdot 3 = 3 \cdot (-2) \cdot (-2)$. La dimostrazione «stacca» un primo alla volta: $p_r$ deve comparire
    anche nell'altra scomposizione (è qui che serve che sia \textbf{primo}), lo si cancella da entrambe e si continua.
\end{spiegazione}

\begin{concetto}
\begin{Corollario}{Esistono infiniti numeri primi (Euclide)}{infiniti-primi}
    \begin{prerequisiti}
        \dip{Teorema}{teo:tfa}
        \dip{Proposizione}{prop:prop-divisibilita}
    \end{prerequisiti}
    L'insieme dei numeri primi non è finito.
\end{Corollario}
\begin{dimostrazione}
    Per assurdo siano $p_1, \dots, p_t$ tutti i primi positivi e sia $m = p_1 \cdots p_t + 1 \geq 2$. Per il Teorema~\ref{teo:tfa} qualche primo
    positivo $p$ divide $m$, e $p = p_i$ per qualche $i$. Ma $p_i$ divide anche $p_1 \cdots p_t$, quindi divide la differenza
    $m - p_1 \cdots p_t = 1$: assurdo.
\end{dimostrazione}
\end{concetto}

\begin{spiegazione}
    Data una lista finita di primi, il numero $m = (\text{prodotto della lista}) + 1$ lascia resto $1$ nella divisione per ciascuno di essi:
    quindi i suoi fattori primi non sono nella lista. Con la lista $2, 3, 5$: $m = 31$, che è un primo nuovo. (Attenzione: $m$ non è sempre
    primo: con $2, 3, 5, 7, 11, 13$ si ha $m = 30031 = 59 \cdot 509$; ma i suoi fattori primi sono comunque nuovi.)
\end{spiegazione}

Dalla scomposizione in fattori primi si leggono tutti i divisori di un numero, e quindi anche massimo comun divisore e minimo
comune multiplo.

\begin{concetto}
\begin{Proposizione}{Divisori, MCD e mcm dalla scomposizione}{divisori-scomposizione}
    \begin{prerequisiti}
        \dip{Teorema}{teo:tfa}
        \dip{Lemma}{lem:somma-prodotto}
        \dip{Definizione}{def:mcd}
    \end{prerequisiti}
    Sia $n = p_1^{e_1} \cdots p_r^{e_r}$ con $p_1, \dots, p_r$ primi positivi distinti ed $e_i \geq 1$. Allora:
    \begin{enumerate}
        \item i divisori positivi di $n$ sono esattamente i numeri $p_1^{f_1} \cdots p_r^{f_r}$ con $0 \leq f_i \leq e_i$; sono $(e_1 + 1) \cdots (e_r + 1)$, e
              i divisori in $\mbZ$ sono il doppio;
        \item se $a = p_1^{e_1} \cdots p_r^{e_r}$ e $b = p_1^{f_1} \cdots p_r^{f_r}$ (con esponenti $\geq 0$), allora
              $d = \prod p_i^{\min(e_i, f_i)} \in \mathrm{MCD}(a, b)$, $\ell = \prod p_i^{\max(e_i, f_i)} \in \mathrm{mcm}(a, b)$ e $d\ell = ab$.
    \end{enumerate}
\end{Proposizione}
\begin{dimostrazione}
    \begin{caso}{Punto 1}
        Ogni $p_1^{f_1} \cdots p_r^{f_r}$ con $f_i \leq e_i$ divide $n$ (il quoziente è $\prod p_i^{e_i - f_i}$). Viceversa, se $c > 0$ e $n = ck$, scomponendo
        $c$ e $k$ in primi positivi si ottiene una scomposizione di $n$; per l'unicità i primi di $c$ sono tra i $p_i$, ognuno con esponente
        al più $e_i$. Per l'unicità, esponenti diversi danno numeri diversi; quindi i divisori positivi corrispondono alle $r$-uple
        $(f_1, \dots, f_r)$, che sono $(e_1 + 1) \cdots (e_r + 1)$ (Lemma~\ref{lem:somma-prodotto}). In $\mbZ$ si aggiungono gli opposti.
    \end{caso}
    \begin{caso}{Punto 2}
        Per il punto 1, i divisori comuni positivi di $a$ e $b$ sono i $\prod p_i^{g_i}$ con $g_i \leq e_i$ e $g_i \leq f_i$, cioè $g_i \leq \min(e_i, f_i)$:
        tutti dividono $d$, che è uno di loro. Allo stesso modo i multipli comuni hanno esponenti $\geq \max(e_i, f_i)$ e sono multipli di $\ell$.
        Infine $\min(e_i, f_i) + \max(e_i, f_i) = e_i + f_i$, quindi $d\ell = ab$.
    \end{caso}
\end{dimostrazione}
\end{concetto}

\begin{spiegazione}
    Un divisore di $12 = 2^2 \cdot 3$ si costruisce scegliendo quante volte usare il $2$ ($0$, $1$ o $2$) e quante il $3$ ($0$ o $1$): $3 \cdot 2 = 6$
    divisori, disposti in una griglia che è proprio il diagramma di Hasse dei divisori di $12$:
    \begin{center}
    \begin{tikzpicture}[hasse]
        \foreach \f/\g/\v in {0/0/1, 1/0/2, 2/0/4, 0/1/3, 1/1/6, 2/1/12} \node[hv] (d\f\g) at (1.5*\f-0.8*\g, 0.9*\f+1.0*\g) {$\v$};
        \draw[hlato] (d00) -- (d10) -- (d20); \draw[hlato] (d01) -- (d11) -- (d21);
        \draw[hlato] (d00) -- (d01); \draw[hlato] (d10) -- (d11); \draw[hlato] (d20) -- (d21);
        \node[font=\small, align=left, anchor=west] at (3.8,1.3)
            {verso destra: un $2$ in più;\\ verso l'alto a sinistra: un $3$ in più.\\[1mm]
             $60 = 2^2 \cdot 3 \cdot 5$ ha $3 \cdot 2 \cdot 2 = 12$ divisori positivi;\\
             $\mathrm{MCD}(12, 18) \ni 2^{\min(2,1)} \cdot 3^{\min(1,2)} = 6$,\\ $\mathrm{mcm}(12, 18) \ni 2^2 \cdot 3^2 = 36$, \ $6 \cdot 36 = 12 \cdot 18$.};
    \end{tikzpicture}
    \end{center}
\end{spiegazione}

% ──────────────────────────────────────────────────────────────
\section{I sottogruppi di $(\mbZ, +)$}\label{sec:sottogruppi-z}

Per ogni $m \in \mbZ$, l'insieme dei multipli $m\mbZ = \langle m \rangle$ è un sottogruppo di $(\mbZ, +)$ (Proposizione~\ref{prop:cicliche}). Non
ce ne sono altri.

\begin{concetto}
\begin{Proposizione}{I sottogruppi di $\mbZ$}{sottogruppi-z}
    \begin{prerequisiti}
        \dip{Definizione}{def:sottostruttura}
        \dip{Teorema}{teo:divisione-euclidea}
        \dip{Definizione}{def:ben-ordinato}
    \end{prerequisiti}
    I sottogruppi di $(\mbZ, +)$ sono tutti e soli gli insiemi $m\mbZ$ con $m \in \mbN$.
\end{Proposizione}
\begin{dimostrazione}
    \facoltativa
    Sia $h$ un sottogruppo. Se $h = \{0\}$, $h = 0\mbZ$. Altrimenti $h$ contiene un elemento non nullo e il suo opposto, quindi un elemento
    positivo; sia $m$ il minimo degli elementi positivi di $h$. I multipli di $m$ stanno in $h$ (somme di $m$ e di $-m$), quindi $m\mbZ \subseteq h$.
    Se $x \in h$, dividiamo: $x = mq + r$ con $0 \leq r < m$; allora $r = x - mq \in h$ e, per la minimalità di $m$, $r = 0$: $x \in m\mbZ$.
\end{dimostrazione}
\end{concetto}

\begin{spiegazione}
    Un sottogruppo di $\mbZ$ è chiuso per somme e differenze. Se contiene il suo più piccolo elemento positivo $m$, contiene tutti i
    multipli di $m$; e se contenesse un numero che non è multiplo di $m$, sottraendo multipli di $m$ si otterrebbe un resto positivo più
    piccolo di $m$, impossibile. Ad esempio il sottogruppo generato da $\{4, 6\}$ contiene $6 - 4 = 2$, quindi è $2\mbZ$.
\end{spiegazione}


  % Classi di resto, mod, congruenze nelle strutture, strutture quoziente, anelli Z_m, Z_m campo
  \chapter{Congruenze e anelli $\mbZ_m$}\label{cap:congruenze}

Nel Capitolo~\ref{cap:equivalenze} abbiamo visto la congruenza modulo $m$ come relazione di equivalenza. Qui ne studiamo le classi,
le \textbf{classi di resto}, e scopriamo che con esse si possono fare somme e prodotti: si ottengono gli anelli $\mbZ_m$, l'aritmetica
«dell'orologio». Per farlo introduciamo un'idea generale: le \textbf{congruenze} di una struttura algebrica, cioè le relazioni di
equivalenza che rispettano le operazioni, e le \textbf{strutture quoziente}.

% ──────────────────────────────────────────────────────────────
\section{Classi di resto}\label{sec:classi-resto}

Ricordiamo (Definizione~\ref{def:congruenza-modulo}) che $a \equiv_m b$ se $b - a$ è un multiplo di $m$, cioè se $m \mid b - a$.

\dfn[zm]{Classi di resto, l'insieme $\mbZ_m$}{
    Sia $m \in \mbZ$. Poniamo $\mbZ_m := \mbZ/{\equiv_m}$ e, per $a \in \mbZ$, indichiamo la classe di $a$ con
    \[ [a]_m := [a]_{\equiv_m} = \{a + mk \mid k \in \mbZ\} = a + m\mbZ. \]
    Le classi $[a]_m$ si chiamano \textbf{classi di resto} modulo $m$.
}

\dfn[mod]{L'operazione $\bmod$}{
    Per $a \in \mbZ$ e $m \in \mbZ \setminus \{0\}$ poniamo $a \bmod m := \min([a]_m \cap \mbN)$ (si scrive anche $a \,\%\, m$): il più piccolo numero
    naturale congruo ad $a$ modulo $m$.
}

\begin{concetto}
\begin{Proposizione}{$a \bmod m$ è il resto della divisione}{mod-resto}
    \begin{prerequisiti}
        \dip{Definizione}{def:mod}
        \dip{Teorema}{teo:divisione-euclidea}
    \end{prerequisiti}
    Siano $a \in \mbZ$, $m \neq 0$ e $a = mq + r$ con $0 \leq r < |m|$. Allora $a \bmod m = r$. In particolare $0 \leq a \bmod m < |m|$.
\end{Proposizione}
\begin{dimostrazione}
    $r = a - mq \in [a]_m \cap \mbN$. Se fosse $x \in [a]_m \cap \mbN$ con $x < r$, allora $x \equiv_m r$ e $0 < r - x < |m|$: $r - x$ sarebbe un multiplo di $m$
    non nullo di valore assoluto minore di $|m|$, impossibile. Quindi $r$ è il minimo.
\end{dimostrazione}
\end{concetto}

\begin{spiegazione}
    Per trovare il rappresentante «canonico» di una classe basta dividere: $101 \bmod 10 = 1$, \ $30093 \bmod 3 = 0$,
    \ $-7 \bmod 3 = 2$ (perché $-7 = 3 \cdot (-3) + 2$). Su una calcolatrice o in un linguaggio di programmazione l'operatore $\%$ può
    dare risultati negativi per $a < 0$: con la nostra definizione il risultato è sempre tra $0$ e $|m| - 1$.
\end{spiegazione}

\begin{concetto}
\begin{Teorema}{Descrizione esplicita di $\mbZ_m$}{zm-esplicito}
    \begin{prerequisiti}
        \dip{Definizione}{def:zm}
        \dip{Teorema}{teo:divisione-euclidea}
        \dip{Teorema}{teo:proprieta-classi}
    \end{prerequisiti}
    Sia $m \in \mbZ \setminus \{0\}$. Allora $\mbZ_m = \{[0]_m, [1]_m, \dots, [|m| - 1]_m\}$ e queste $|m|$ classi sono a due a due distinte; in particolare
    $|\mbZ_m| = |m|$.
\end{Teorema}
\begin{dimostrazione}
    Poiché $\equiv_m$ e $\equiv_{-m}$ coincidono, possiamo supporre $m > 0$.
    \begin{caso}{Ogni classe è una delle $[r]_m$}
        Sia $a \in \mbZ$ e sia $a = mq + r$ con $0 \leq r < m$. Allora $a - r = mq$, cioè $a \equiv_m r$, e quindi $[a]_m = [r]_m$.
    \end{caso}
    \begin{caso}{Le $[r]_m$ sono distinte}
        Siano $0 \leq i \leq j \leq m - 1$ con $[i]_m = [j]_m$. Allora $m \mid j - i$ con $0 \leq j - i < m$, quindi $j - i = 0$: $i = j$.
    \end{caso}
\end{dimostrazione}
\end{concetto}

\begin{spiegazione}
    Gli interi si «avvolgono» attorno a un orologio con $m$ posizioni: ogni posizione è una classe di resto, e i numeri che ci
    finiscono sopra differiscono per giri completi, cioè per multipli di $m$. Per $m = 5$:
    \begin{center}
    \begin{tikzpicture}
        \draw[black!40, line width=0.8pt] (0,0) circle (1.3);
        \foreach \r/\ang/\col in {0/90/classeA, 1/18/classeB, 2/-54/classeC, 3/-126/classeD, 4/162/classeE} {
            \node[circle, fill=\col, draw=black!40, minimum size=8mm, inner sep=0pt, font=\small] at (\ang:1.3) {$[\r]_5$};
        }
        \node[font=\scriptsize, anchor=south] at (90:1.75) {$\dots, -5, 0, 5, 10, \dots$};
        \node[font=\scriptsize, anchor=west] at (18:1.75) {$\dots, -4, 1, 6, 11, \dots$};
        \node[font=\scriptsize, anchor=west] at (-54:1.75) {$\dots, -3, 2, 7, 12, \dots$};
        \node[font=\scriptsize, anchor=east] at (-126:1.75) {$\dots, -2, 3, 8, 13, \dots$};
        \node[font=\scriptsize, anchor=east] at (162:1.75) {$\dots, -1, 4, 9, 14, \dots$};
    \end{tikzpicture}
    \end{center}
    La dimostrazione dice due cose: ogni intero cade in una delle $5$ posizioni (quella del suo resto), e due posizioni diverse
    non coincidono (la distanza tra $0, 1, \dots, 4$ è sempre minore di un giro intero).
\end{spiegazione}

% ──────────────────────────────────────────────────────────────
\section{Congruenze in una struttura algebrica}\label{sec:congruenze-strutture}

Vorremmo sommare e moltiplicare le classi di resto ponendo $[a]_m + [b]_m := [a + b]_m$. Ma una classe ha tanti rappresentanti: la
definizione ha senso solo se il risultato \textbf{non dipende da quali rappresentanti scegliamo}. Questo succede quando la
relazione di equivalenza «va d'accordo» con l'operazione.

\dfn[congruenza]{Compatibilità, congruenza}{
    Sia $s \neq \varnothing$ e sia $*$ un'operazione binaria interna su $s$. Una relazione di equivalenza $\sim$ su $s$ si dice
    \begin{itemize}
        \item \textbf{compatibile a sinistra} con $*$ se $\forall a, b, c \in s\,(a \sim b \rightarrow c * a \sim c * b)$;
        \item \textbf{compatibile a destra} con $*$ se $\forall a, b, c \in s\,(a \sim b \rightarrow a * c \sim b * c)$.
    \end{itemize}
    Se $*_1, \dots, *_n$ sono operazioni binarie interne su $s$, $\sim$ si dice una \textbf{congruenza} in $(s, *_1, \dots, *_n)$ se, per ogni $i$,
    \[ \forall a, b, c, d \in s\,\big((a \sim b \wedge c \sim d) \rightarrow a *_i c \sim b *_i d\big). \]
}

\begin{concetto}
\begin{Teorema}{Congruenze e compatibilità}{congruenza-compatibile}
    \begin{prerequisiti}
        \dip{Definizione}{def:congruenza}
        \dip{Definizione}{def:tipi-relazioni}
    \end{prerequisiti}
    Una relazione di equivalenza $\sim$ su $s$ è una congruenza in $(s, *_1, \dots, *_n)$ se e solo se è compatibile a destra e a sinistra
    con ogni operazione $*_i$.
\end{Teorema}
\begin{dimostrazione}
    Basta considerare una sola operazione $*$.
    \begin{caso}{$(\rightarrow)$}
        Sia $a \sim b$. Applicando la definizione di congruenza con $c \sim c$ (riflessività) si ottiene $a * c \sim b * c$: compatibilità a destra.
        Applicandola alle coppie $c \sim c$ e $a \sim b$ si ottiene $c * a \sim c * b$: compatibilità a sinistra.
    \end{caso}
    \begin{caso}{$(\leftarrow)$}
        Siano $a \sim b$ e $c \sim d$. Per la compatibilità a destra $a * c \sim b * c$; per quella a sinistra $b * c \sim b * d$. Per la
        transitività $a * c \sim b * d$.
    \end{caso}
\end{dimostrazione}
\end{concetto}

\begin{spiegazione}
    Per sostituire \textbf{entrambi} i fattori con elementi equivalenti, li si sostituisce \textbf{uno alla volta}: prima quello a sinistra
    (compatibilità a destra), poi quello a destra (compatibilità a sinistra), e si collegano i due passi con la transitività:
    \[ a * c \ \overset{\text{cambio } a \to b}{\sim}\ b * c \ \overset{\text{cambio } c \to d}{\sim}\ b * d. \]
\end{spiegazione}

\begin{concetto}
\begin{Proposizione}{Le operazioni quoziente}{operazioni-quoziente}
    \begin{prerequisiti}
        \dip{Definizione}{def:congruenza}
        \dip{Teorema}{teo:proprieta-classi}
        \dip{Definizione}{def:funzione}
    \end{prerequisiti}
    Sia $\sim$ una congruenza in $(s, *_1, \dots, *_n)$. Allora per ogni $i$ è ben definita l'operazione binaria interna su $s/{\sim}$
    \[ (*_i)_\sim \colon ([x]_\sim, [y]_\sim) \in s/{\sim} \times s/{\sim} \mapsto [x *_i y]_\sim \in s/{\sim}, \]
    detta \textbf{operazione quoziente}. La struttura $(s/{\sim}, (*_1)_\sim, \dots, (*_n)_\sim)$ si chiama \textbf{struttura quoziente}.
\end{Proposizione}
\begin{dimostrazione}
    Bisogna verificare che il risultato non dipenda dai rappresentanti. Se $[x] = [x']$ e $[y] = [y']$, allora $x \sim x'$ e $y \sim y'$
    (Teorema~\ref{teo:proprieta-classi}); essendo $\sim$ una congruenza, $x *_i y \sim x' *_i y'$, cioè $[x *_i y] = [x' *_i y']$.
\end{dimostrazione}
\end{concetto}

\begin{spiegazione}
    Si opera sulle classi \textbf{scegliendo un rappresentante qualunque per ciascuna}, facendo il conto con i rappresentanti e
    prendendo la classe del risultato. Ad esempio in $\mbZ_5$: $[3]_5 + [4]_5 = [7]_5 = [2]_5$; scegliendo altri rappresentanti,
    $[8]_5 + [-1]_5 = [7]_5$ ancora, e $[13]_5 + [9]_5 = [22]_5 = [2]_5$. La proposizione garantisce che vada sempre così.
\end{spiegazione}

\begin{esempio}{Una relazione di equivalenza che non è una congruenza}
    Su $\mbZ$ consideriamo l'operazione $a * b = a + b$ se $b$ è dispari e $a * b = a + \frac{b}{2}$ se $b$ è pari. La congruenza $\equiv_2$
    \textbf{non} è compatibile con $*$: $0 \equiv_2 2$, ma $1 * 0 = 1$ e $1 * 2 = 2$, che non sono congrui modulo $2$. Quindi «$[1]_2 * [0]_2$» non si può
    definire: dipenderebbe dal rappresentante scelto per $[0]_2$. Invece con l'operazione $a * b = 2ab$ la relazione $\equiv_2$ è una congruenza:
    il risultato è sempre pari, quindi tutti i risultati sono congrui tra loro.
\end{esempio}

\begin{concetto}
\begin{Teorema}{L'epimorfismo canonico}{epimorfismo-canonico}
    \begin{prerequisiti}
        \dip{Proposizione}{prop:operazioni-quoziente}
        \dip{Definizione}{def:omomorfismo}
        \dip{Teorema}{teo:epimorfismi}
        \dip{Proposizione}{prop:proiezione}
    \end{prerequisiti}
    Sia $\sim$ una congruenza in $(s, *_1, \dots, *_n)$. La proiezione canonica $\pi \colon x \in s \mapsto [x]_\sim \in s/{\sim}$ è un epimorfismo (per
    ogni operazione), detto \textbf{epimorfismo canonico}. Di conseguenza la struttura quoziente eredita associatività, commutatività,
    elementi neutri e simmetrici, e la distributività di un'operazione rispetto a un'altra. In particolare i quozienti di semigruppi,
    monoidi, gruppi (abeliani) e anelli (commutativi, unitari) sono strutture dello stesso tipo.
\end{Teorema}
\begin{dimostrazione}
    $\pi$ è suriettiva (Proposizione~\ref{prop:proiezione}) e, per definizione di operazione quoziente, $\pi(x *_i y) = [x *_i y] = [x] (*_i)_\sim [y]
    = \pi(x) (*_i)_\sim \pi(y)$: è un epimorfismo. Il Teorema~\ref{teo:epimorfismi} dà allora associatività, commutatività, neutri e simmetrici.
    Per la distributività: $[x]([y] + [z]) = [x(y + z)] = [xy + xz] = [x][y] + [x][z]$ (scrivendo $+$ e $\cdot$ per due delle operazioni).
\end{dimostrazione}
\end{concetto}

\begin{spiegazione}
    Ogni proprietà espressa da un'uguaglianza tra conti «passa al quoziente»: per verificarla sulle classi la si verifica sui
    rappresentanti e poi si prende la classe. Ad esempio la commutatività in $\mbZ_m$: $[a][b] = [ab] = [ba] = [b][a]$. Il neutro di $\mbZ_m$
    per la somma è $[0]$, l'opposto di $[a]$ è $[-a]$, il neutro per il prodotto è $[1]$.
\end{spiegazione}

% ──────────────────────────────────────────────────────────────
\section{Gli anelli $\mbZ_m$}\label{sec:anelli-zm}

\begin{concetto}
\begin{Proposizione}{La congruenza modulo $m$ è una congruenza di $(\mbZ, +, \cdot)$}{congruenza-z}
    \begin{prerequisiti}
        \dip{Teorema}{teo:congruenza-compatibile}
        \dip{Teorema}{teo:epimorfismo-canonico}
        \dip{Proposizione}{prop:prop-divisibilita}
    \end{prerequisiti}
    Per ogni $m \in \mbZ$, $\equiv_m$ è una congruenza in $(\mbZ, +, \cdot)$. Quindi $(\mbZ_m, +, \cdot)$, con
    \[ [a]_m + [b]_m := [a + b]_m, \qquad [a]_m \cdot [b]_m := [ab]_m, \]
    è un anello commutativo unitario, con zero $[0]_m$ e unità $[1]_m$.
\end{Proposizione}
\begin{dimostrazione}
    Poiché somma e prodotto in $\mbZ$ sono commutativi, basta la compatibilità a destra (Teorema~\ref{teo:congruenza-compatibile}). Sia
    $a \equiv_m b$, cioè $m \mid b - a$. Allora $(b + c) - (a + c) = b - a$ e $bc - ac = (b - a)c$ sono multipli di $m$: $a + c \equiv_m b + c$ e
    $ac \equiv_m bc$. La seconda parte segue dal Teorema~\ref{teo:epimorfismo-canonico}, perché $(\mbZ, +, \cdot)$ è un anello commutativo unitario.
\end{dimostrazione}
\end{concetto}

\begin{spiegazione}
    È l'aritmetica dell'orologio: in $\mbZ_{12}$, $[5] + [9] = [14] = [2]$ dice che se sono le $5$, tra $9$ ore saranno le $2$. Ecco le tavole
    di Cayley di $\mbZ_4$ (rappresentanti $0, 1, 2, 3$; scriviamo $r$ per $[r]_4$):
    \begin{center}
    \renewcommand{\arraystretch}{1.2}
    \begin{tabular}{c|cccc}
        $+$ & $0$ & $1$ & $2$ & $3$ \\ \hline
        $0$ & $0$ & $1$ & $2$ & $3$ \\
        $1$ & $1$ & $2$ & $3$ & $0$ \\
        $2$ & $2$ & $3$ & $0$ & $1$ \\
        $3$ & $3$ & $0$ & $1$ & $2$ \\
    \end{tabular}
    \qquad\qquad
    \begin{tabular}{c|cccc}
        $\cdot$ & $0$ & $1$ & $2$ & $3$ \\ \hline
        $0$ & $0$ & $0$ & $0$ & $0$ \\
        $1$ & $0$ & $1$ & $2$ & $3$ \\
        $2$ & $0$ & $2$ & \cellcolor{cherryred!20}$0$ & $2$ \\
        $3$ & $0$ & $3$ & $2$ & \cellcolor{seagreen!25}$1$ \\
    \end{tabular}
    \end{center}
    Si legge che $[2]_4 \cdot [2]_4 = [0]_4$: in $\mbZ_4$ c'è un \textbf{divisore dello zero}, quindi $\mbZ_4$ non è un dominio di integrità; invece
    $[3]_4 \cdot [3]_4 = [1]_4$: $[3]_4$ è invertibile.
\end{spiegazione}

\begin{concetto}
\begin{Teorema}{Quando $\mbZ_m$ è un campo}{zm-campo}
    \begin{prerequisiti}
        \dip{Proposizione}{prop:congruenza-z}
        \dip{Proposizione}{prop:campo-dominio}
        \dip{Definizione}{def:primo}
        \dip{Teorema}{teo:primi-irriducibili}
        \dip{Teorema}{teo:bezout}
        \dip{Teorema}{teo:zm-esplicito}
    \end{prerequisiti}
    Sia $m \in \mbZ \setminus \{0\}$. Sono equivalenti:
    \begin{enumerate}
        \item $(\mbZ_m, +, \cdot)$ è un campo;
        \item $(\mbZ_m, +, \cdot)$ è un dominio di integrità;
        \item $m$ è primo.
    \end{enumerate}
\end{Teorema}
\begin{dimostrazione}
    \begin{caso}{(1) $\rightarrow$ (2)}
        Ogni campo è un dominio di integrità (Proposizione~\ref{prop:campo-dominio}).
    \end{caso}
    \begin{caso}{(2) $\rightarrow$ (3)}
        In un dominio di integrità $[0] \neq [1]$, quindi $|m| = |\mbZ_m| \geq 2$, cioè $m \notin \{-1, 0, 1\}$. Sia $m \mid ab$: allora
        $[a][b] = [ab] = [0]$ e, per la legge di annullamento del prodotto, $[a] = [0]$ oppure $[b] = [0]$, cioè $m \mid a$ oppure $m \mid b$. Quindi $m$ è primo.
    \end{caso}
    \begin{caso}{(3) $\rightarrow$ (1)}
        $\mbZ_m$ è commutativo e unitario, con $[0] \neq [1]$ perché $|m| \geq 2$. Sia $[n] \neq [0]$, con $1 \leq n < |m|$ (Teorema~\ref{teo:zm-esplicito}). Poiché $m$
        è primo, è irriducibile: i divisori comuni di $m$ e $n$ sono tra $\pm 1, \pm m$, e $\pm m$ non dividono $n$. Quindi $m$ e $n$ sono coprimi e per
        Bézout $1 = nu + mv$ per certi $u, v$. Allora $[n][u] = [nu] = [1 - mv] = [1]$: ogni classe non nulla è invertibile.
    \end{caso}
\end{dimostrazione}
\end{concetto}

\begin{spiegazione}
    Confrontiamo le righe del prodotto (senza lo zero) in $\mbZ_6$ e in $\mbZ_7$:
    \begin{center}
    \renewcommand{\arraystretch}{1.15}
    \begin{tabular}{c|ccccc}
        $\mbZ_6$ & $1$ & $2$ & $3$ & $4$ & $5$ \\ \hline
        $1$ & $1$ & $2$ & $3$ & $4$ & $5$ \\
        $2$ & $2$ & $4$ & \cellcolor{cherryred!20}$0$ & $2$ & $4$ \\
        $3$ & $3$ & \cellcolor{cherryred!20}$0$ & $3$ & \cellcolor{cherryred!20}$0$ & $3$ \\
        $4$ & $4$ & $2$ & \cellcolor{cherryred!20}$0$ & $4$ & $2$ \\
        $5$ & $5$ & $4$ & $3$ & $2$ & \cellcolor{seagreen!25}$1$ \\
    \end{tabular}
    \qquad\qquad
    \begin{tabular}{c|cccccc}
        $\mbZ_7$ & $1$ & $2$ & $3$ & $4$ & $5$ & $6$ \\ \hline
        $1$ & \cellcolor{seagreen!25}$1$ & $2$ & $3$ & $4$ & $5$ & $6$ \\
        $2$ & $2$ & $4$ & $6$ & \cellcolor{seagreen!25}$1$ & $3$ & $5$ \\
        $3$ & $3$ & $6$ & $2$ & $5$ & \cellcolor{seagreen!25}$1$ & $4$ \\
        $4$ & $4$ & \cellcolor{seagreen!25}$1$ & $5$ & $2$ & $6$ & $3$ \\
        $5$ & $5$ & $3$ & \cellcolor{seagreen!25}$1$ & $6$ & $4$ & $2$ \\
        $6$ & $6$ & $5$ & $4$ & $3$ & $2$ & \cellcolor{seagreen!25}$1$ \\
    \end{tabular}
    \end{center}
    In $\mbZ_6$, con $6 = 2 \cdot 3$ non primo, $[2][3] = [0]$: la fattorizzazione di $6$ produce divisori dello zero, e le righe di $2, 3, 4$ hanno
    ripetizioni e non contengono $1$. In $\mbZ_7$ ogni riga è una permutazione di $1, \dots, 6$ e contiene l'$1$: ogni elemento non nullo è
    invertibile ($[3]^{-1} = [5]$, perché $3 \cdot 5 = 15 = 2 \cdot 7 + 1$).
\end{spiegazione}


\end{document}
